\documentclass[10pt,a4paper]{article}
\usepackage[T1]{fontenc}
\usepackage{lmodern}
\usepackage[utf8]{inputenc}
\usepackage[margin=18mm]{geometry}
\usepackage{amsmath,amssymb,booktabs,makecell,xcolor,graphicx,textcomp}
\usepackage[hidelinks,pdfauthor={Przemyslaw Bielecki and Tomasz Hachaj},pdftitle={Supplementary Material: Imputation of Missing Traffic Sensor Data}]{hyperref}
\setlength{\parindent}{0pt}
\setlength{\parskip}{6pt}
\setcounter{topnumber}{3}
\setcounter{bottomnumber}{3}
\setcounter{totalnumber}{4}
\renewcommand{\topfraction}{0.95}
\renewcommand{\bottomfraction}{0.95}
\renewcommand{\textfraction}{0.05}
\renewcommand{\floatpagefraction}{0.75}
\begin{document}
\begin{center}
{\LARGE\bfseries Supplementary Material\par}
\vspace{8mm}
{\Large Imputation of Missing Traffic Sensor Data:\\
A Comparative Study of Classical, Graph-Based,\\
and Deep Spatio-Temporal Methods\par}
\vspace{5mm}
Przemys{\l}aw Bielecki and Tomasz Hachaj
\end{center}
\vspace{8mm}
This supplement contains the complete tabulated results accompanying the revised manuscript, using the completed road graphs and the five prescribed seeds. It reports test results from January 2023 and October--December 2022 in Darmstadt. Graph Laplacian, Graph Temporal and GRIN use graphs with 254 and 252 edges, respectively, with no isolated sensors after training-based qualification and connection through excluded sensors.

\textbf{Contents}
\begin{itemize}
\item \hyperref[tab:supp_p01_1]{Tables S1--S24}: complete primary results for all 13 masking scenarios, all 12 methods and MAE, RMSE and MRE.
\item \hyperref[tab:supp_p01_25]{Tables S25--S96}: corresponding results for sensors A023, A027 and A142.
\item \hyperref[tab:p03_A1]{Tables A1--A31}: persistent stream loss, graph and scaling controls, traffic-count and graph-degree strata, daily aggregation and the separate forecasting check.
\end{itemize}

Machine-readable summaries and individual-run results are provided in the accompanying \texttt{data/} directory. The \texttt{README.md} file maps the CSV files to the tables and describes how to compile this document. Training records and rounding comparisons are included as CSV files.

Every reported interval uses complete seed runs as the replication unit. The five-run intervals are conditional on fixed chronological splits, not estimates of between-city uncertainty. Natural-missingness realism has not been quantitatively validated. The two periods come from one city.

This document is self-contained and does not require the main manuscript's LaTeX class or bibliography.
\clearpage
\renewcommand{\theHtable}{S\arabic{table}}
% P01 full results: LaTeX supplement fragment, generated from final CSV.

% Include after the manuscript or in a supplement wrapper using xcolor, booktabs and makecell.

\clearpage

\begingroup

\hypersetup{linkcolor=black}

\renewcommand{\thetable}{S\arabic{table}}

\setcounter{table}{0}

\section*{{Supplementary results: complete primary and selected-sensor comparisons}}

{Tables S1--S24 report all 936 dataset--scenario--method--metric combinations in the main benchmark. Tables S25--S96 report all 2808 corresponding combinations for sensors A023, A027 and A142. Every estimate uses five seed runs on the same chronological split. Cells give mean $\pm$ sample SD and an unadjusted run-level 95\% Student-$t$ confidence interval. Negative lower bounds are retained when produced by this interval calculation; they do not imply negative error values. MAE and RMSE have units of vehicles per ten-minute interval; MRE is a dimensionless ratio evaluated on nonzero targets. Blocks use independent starts by sensor unless labelled synchronous. Bold identifies the lowest unrounded mean in a column, not a significant difference. The two KNN variants have the different scaling and $k$ selection policies specified in Methods.}

\begin{table}[htbp]
\centering
\caption{{MAE for January 2023 (vehicles per ten-minute interval). Each cell gives mean $\pm$ sample SD, followed by the 95\% confidence interval; $n=5$ seed runs. Bold means identify the lowest unrounded mean within each column, not statistical significance.}}
\label{tab:supp_p01_1}
\begingroup
\footnotesize
\setlength{\tabcolsep}{3pt}
\renewcommand{\arraystretch}{1.15}
\begin{tabular}{lcccc}
\toprule
Method & Block $B=1$ & Block $B=2$ & Block $B=4$ & Block $B=8$ \\
\midrule
Linear & \makecell[c]{$\mathbf{23.69}\,\pm\,3.80$\\$[18.97,\ 28.41]$} & \makecell[c]{$23.65\,\pm\,5.01$\\$[17.43,\ 29.87]$} & \makecell[c]{$26.68\,\pm\,1.89$\\$[24.34,\ 29.02]$} & \makecell[c]{$28.62\,\pm\,5.57$\\$[21.69,\ 35.54]$} \\
Median & \makecell[c]{$116.05\,\pm\,6.82$\\$[107.58,\ 124.52]$} & \makecell[c]{$111.70\,\pm\,15.42$\\$[92.55,\ 130.85]$} & \makecell[c]{$109.70\,\pm\,10.91$\\$[96.15,\ 123.24]$} & \makecell[c]{$107.90\,\pm\,13.95$\\$[90.58,\ 125.22]$} \\
FFill + BFill & \makecell[c]{$31.74\,\pm\,5.34$\\$[25.12,\ 38.37]$} & \makecell[c]{$29.11\,\pm\,7.49$\\$[19.81,\ 38.41]$} & \makecell[c]{$34.51\,\pm\,2.36$\\$[31.58,\ 37.44]$} & \makecell[c]{$41.97\,\pm\,8.25$\\$[31.73,\ 52.21]$} \\
Rolling mean & \makecell[c]{$31.37\,\pm\,6.07$\\$[23.82,\ 38.91]$} & \makecell[c]{$29.46\,\pm\,3.12$\\$[25.58,\ 33.33]$} & \makecell[c]{$34.20\,\pm\,1.86$\\$[31.89,\ 36.50]$} & \makecell[c]{$33.21\,\pm\,5.58$\\$[26.27,\ 40.14]$} \\
KNN (scaled) & \makecell[c]{$126.02\,\pm\,7.26$\\$[117.00,\ 135.03]$} & \makecell[c]{$115.63\,\pm\,18.70$\\$[92.41,\ 138.84]$} & \makecell[c]{$113.30\,\pm\,11.34$\\$[99.22,\ 127.38]$} & \makecell[c]{$109.32\,\pm\,13.93$\\$[92.02,\ 126.62]$} \\
KNN (unscaled) & \makecell[c]{$48.64\,\pm\,10.50$\\$[35.61,\ 61.68]$} & \makecell[c]{$36.93\,\pm\,6.76$\\$[28.54,\ 45.32]$} & \makecell[c]{$37.65\,\pm\,4.03$\\$[32.64,\ 42.65]$} & \makecell[c]{$37.88\,\pm\,10.21$\\$[25.21,\ 50.56]$} \\
Historical & \makecell[c]{$66.44\,\pm\,11.50$\\$[52.17,\ 80.72]$} & \makecell[c]{$57.65\,\pm\,12.70$\\$[41.88,\ 73.42]$} & \makecell[c]{$55.95\,\pm\,5.15$\\$[49.56,\ 62.35]$} & \makecell[c]{$54.15\,\pm\,8.55$\\$[43.53,\ 64.77]$} \\
Graph Laplacian & \makecell[c]{$67.48\,\pm\,9.66$\\$[55.49,\ 79.47]$} & \makecell[c]{$67.64\,\pm\,9.10$\\$[56.34,\ 78.94]$} & \makecell[c]{$62.54\,\pm\,5.24$\\$[56.03,\ 69.05]$} & \makecell[c]{$66.61\,\pm\,10.78$\\$[53.23,\ 80.00]$} \\
Graph Temporal & \makecell[c]{$52.98\,\pm\,11.73$\\$[38.41,\ 67.54]$} & \makecell[c]{$61.59\,\pm\,5.65$\\$[54.57,\ 68.60]$} & \makecell[c]{$60.43\,\pm\,4.51$\\$[54.84,\ 66.03]$} & \makecell[c]{$66.70\,\pm\,13.06$\\$[50.48,\ 82.91]$} \\
Sage-Voronoi & \makecell[c]{$33.38\,\pm\,7.14$\\$[24.51,\ 42.25]$} & \makecell[c]{$28.06\,\pm\,3.36$\\$[23.89,\ 32.23]$} & \makecell[c]{$30.61\,\pm\,3.76$\\$[25.93,\ 35.28]$} & \makecell[c]{$36.01\,\pm\,6.81$\\$[27.55,\ 44.48]$} \\
GAIN & \makecell[c]{$46.37\,\pm\,9.23$\\$[34.91,\ 57.84]$} & \makecell[c]{$38.36\,\pm\,6.85$\\$[29.86,\ 46.86]$} & \makecell[c]{$39.24\,\pm\,5.42$\\$[32.51,\ 45.98]$} & \makecell[c]{$37.05\,\pm\,8.69$\\$[26.26,\ 47.84]$} \\
GRIN & \makecell[c]{$25.69\,\pm\,5.38$\\$[19.00,\ 32.37]$} & \makecell[c]{$\mathbf{20.15}\,\pm\,5.01$\\$[13.93,\ 26.37]$} & \makecell[c]{$\mathbf{23.43}\,\pm\,3.63$\\$[18.92,\ 27.94]$} & \makecell[c]{$\mathbf{24.62}\,\pm\,8.01$\\$[14.68,\ 34.56]$} \\
\bottomrule
\end{tabular}
\endgroup
\end{table}

\begin{table}[htbp]
\centering
\caption{{MAE for January 2023 (vehicles per ten-minute interval). Each cell gives mean $\pm$ sample SD, followed by the 95\% confidence interval; $n=5$ seed runs. Bold means identify the lowest unrounded mean within each column, not statistical significance.}}
\label{tab:supp_p01_2}
\begingroup
\footnotesize
\setlength{\tabcolsep}{3pt}
\renewcommand{\arraystretch}{1.15}
\begin{tabular}{lccc}
\toprule
Method & Block $B=16$ & Block $B=32$ & Block $B=64$ \\
\midrule
Linear & \makecell[c]{$34.94\,\pm\,7.85$\\$[25.19,\ 44.69]$} & \makecell[c]{$49.84\,\pm\,3.07$\\$[46.02,\ 53.65]$} & \makecell[c]{$75.75\,\pm\,5.45$\\$[68.97,\ 82.52]$} \\
Median & \makecell[c]{$105.25\,\pm\,6.61$\\$[97.05,\ 113.46]$} & \makecell[c]{$113.51\,\pm\,8.80$\\$[102.58,\ 124.44]$} & \makecell[c]{$114.91\,\pm\,7.02$\\$[106.20,\ 123.62]$} \\
FFill + BFill & \makecell[c]{$62.45\,\pm\,11.08$\\$[48.69,\ 76.20]$} & \makecell[c]{$90.90\,\pm\,7.37$\\$[81.75,\ 100.05]$} & \makecell[c]{$129.79\,\pm\,15.96$\\$[109.97,\ 149.60]$} \\
Rolling mean & \makecell[c]{$38.48\,\pm\,8.16$\\$[28.35,\ 48.61]$} & \makecell[c]{$50.79\,\pm\,3.27$\\$[46.72,\ 54.85]$} & \makecell[c]{$75.82\,\pm\,5.58$\\$[68.89,\ 82.74]$} \\
KNN (scaled) & \makecell[c]{$112.97\,\pm\,7.51$\\$[103.64,\ 122.30]$} & \makecell[c]{$116.55\,\pm\,8.55$\\$[105.94,\ 127.16]$} & \makecell[c]{$116.04\,\pm\,11.38$\\$[101.90,\ 130.17]$} \\
KNN (unscaled) & \makecell[c]{$39.42\,\pm\,6.23$\\$[31.69,\ 47.15]$} & \makecell[c]{$46.60\,\pm\,12.23$\\$[31.41,\ 61.79]$} & \makecell[c]{$43.92\,\pm\,6.69$\\$[35.61,\ 52.23]$} \\
Historical & \makecell[c]{$58.55\,\pm\,4.66$\\$[52.76,\ 64.33]$} & \makecell[c]{$63.87\,\pm\,11.37$\\$[49.76,\ 77.98]$} & \makecell[c]{$62.10\,\pm\,7.54$\\$[52.74,\ 71.47]$} \\
Graph Laplacian & \makecell[c]{$57.27\,\pm\,2.12$\\$[54.63,\ 59.90]$} & \makecell[c]{$67.90\,\pm\,9.54$\\$[56.06,\ 79.74]$} & \makecell[c]{$69.89\,\pm\,8.47$\\$[59.37,\ 80.41]$} \\
Graph Temporal & \makecell[c]{$61.16\,\pm\,1.87$\\$[58.83,\ 63.48]$} & \makecell[c]{$70.88\,\pm\,11.54$\\$[56.55,\ 85.21]$} & \makecell[c]{$76.18\,\pm\,10.80$\\$[62.77,\ 89.59]$} \\
Sage-Voronoi & \makecell[c]{$45.26\,\pm\,11.65$\\$[30.80,\ 59.73]$} & \makecell[c]{$52.78\,\pm\,9.63$\\$[40.83,\ 64.73]$} & \makecell[c]{$58.89\,\pm\,7.01$\\$[50.19,\ 67.60]$} \\
GAIN & \makecell[c]{$39.43\,\pm\,6.95$\\$[30.80,\ 48.06]$} & \makecell[c]{$45.51\,\pm\,10.92$\\$[31.95,\ 59.08]$} & \makecell[c]{$42.61\,\pm\,6.41$\\$[34.65,\ 50.57]$} \\
GRIN & \makecell[c]{$\mathbf{29.20}\,\pm\,4.49$\\$[23.62,\ 34.78]$} & \makecell[c]{$\mathbf{40.71}\,\pm\,11.45$\\$[26.49,\ 54.93]$} & \makecell[c]{$\mathbf{41.20}\,\pm\,6.65$\\$[32.95,\ 49.46]$} \\
\bottomrule
\end{tabular}
\endgroup
\end{table}

\clearpage

\begin{table}[htbp]
\centering
\caption{{MAE for January 2023 (vehicles per ten-minute interval). Each cell gives mean $\pm$ sample SD, followed by the 95\% confidence interval; $n=5$ seed runs. Bold means identify the lowest unrounded mean within each column, not statistical significance.}}
\label{tab:supp_p01_3}
\begingroup
\footnotesize
\setlength{\tabcolsep}{3pt}
\renewcommand{\arraystretch}{1.15}
\begin{tabular}{lcc}
\toprule
Method & Block $B=128$ & Block $B=144$ \\
\midrule
Linear & \makecell[c]{$125.97\,\pm\,6.87$\\$[117.43,\ 134.51]$} & \makecell[c]{$134.40\,\pm\,5.56$\\$[127.49,\ 141.30]$} \\
Median & \makecell[c]{$111.59\,\pm\,5.48$\\$[104.79,\ 118.39]$} & \makecell[c]{$114.06\,\pm\,4.63$\\$[108.31,\ 119.80]$} \\
FFill + BFill & \makecell[c]{$142.98\,\pm\,11.03$\\$[129.29,\ 156.67]$} & \makecell[c]{$141.63\,\pm\,7.80$\\$[131.95,\ 151.31]$} \\
Rolling mean & \makecell[c]{$125.75\,\pm\,6.98$\\$[117.09,\ 134.41]$} & \makecell[c]{$134.52\,\pm\,5.57$\\$[127.60,\ 141.44]$} \\
KNN (scaled) & \makecell[c]{$110.90\,\pm\,5.14$\\$[104.53,\ 117.28]$} & \makecell[c]{$113.60\,\pm\,5.43$\\$[106.85,\ 120.34]$} \\
KNN (unscaled) & \makecell[c]{$41.29\,\pm\,3.91$\\$[36.44,\ 46.14]$} & \makecell[c]{$42.70\,\pm\,3.37$\\$[38.52,\ 46.89]$} \\
Historical & \makecell[c]{$59.00\,\pm\,5.72$\\$[51.89,\ 66.11]$} & \makecell[c]{$59.26\,\pm\,3.51$\\$[54.91,\ 63.62]$} \\
Graph Laplacian & \makecell[c]{$64.59\,\pm\,4.94$\\$[58.46,\ 70.72]$} & \makecell[c]{$66.43\,\pm\,5.07$\\$[60.14,\ 72.72]$} \\
Graph Temporal & \makecell[c]{$70.54\,\pm\,5.96$\\$[63.14,\ 77.93]$} & \makecell[c]{$72.02\,\pm\,4.78$\\$[66.09,\ 77.96]$} \\
Sage-Voronoi & \makecell[c]{$62.91\,\pm\,10.29$\\$[50.13,\ 75.69]$} & \makecell[c]{$60.19\,\pm\,3.87$\\$[55.39,\ 64.99]$} \\
GAIN & \makecell[c]{$38.32\,\pm\,4.66$\\$[32.54,\ 44.10]$} & \makecell[c]{$40.35\,\pm\,4.79$\\$[34.40,\ 46.29]$} \\
GRIN & \makecell[c]{$\mathbf{38.19}\,\pm\,4.66$\\$[32.40,\ 43.98]$} & \makecell[c]{$\mathbf{40.27}\,\pm\,4.54$\\$[34.63,\ 45.92]$} \\
\bottomrule
\end{tabular}
\endgroup
\end{table}

\begin{table}[htbp]
\centering
\caption{{MAE for January 2023 (vehicles per ten-minute interval). Each cell gives mean $\pm$ sample SD, followed by the 95\% confidence interval; $n=5$ seed runs. Bold means identify the lowest unrounded mean within each column, not statistical significance.}}
\label{tab:supp_p01_4}
\begingroup
\footnotesize
\setlength{\tabcolsep}{3pt}
\renewcommand{\arraystretch}{1.15}
\begin{tabular}{lcccc}
\toprule
Method & MCAR 10\% & MCAR 30\% & MCAR 50\% & Synchronous $B=144$ \\
\midrule
Linear & \makecell[c]{$22.49\,\pm\,0.38$\\$[22.02,\ 22.96]$} & \makecell[c]{$\mathbf{23.45}\,\pm\,0.19$\\$[23.21,\ 23.69]$} & \makecell[c]{$\mathbf{24.68}\,\pm\,0.19$\\$[24.45,\ 24.91]$} & \makecell[c]{$129.11\,\pm\,28.07$\\$[94.26,\ 163.96]$} \\
Median & \makecell[c]{$111.81\,\pm\,1.10$\\$[110.44,\ 113.18]$} & \makecell[c]{$112.11\,\pm\,0.51$\\$[111.48,\ 112.75]$} & \makecell[c]{$112.04\,\pm\,0.46$\\$[111.47,\ 112.62]$} & \makecell[c]{$114.94\,\pm\,18.87$\\$[91.51,\ 138.37]$} \\
FFill + BFill & \makecell[c]{$27.85\,\pm\,0.65$\\$[27.04,\ 28.65]$} & \makecell[c]{$29.78\,\pm\,0.30$\\$[29.41,\ 30.15]$} & \makecell[c]{$32.15\,\pm\,0.15$\\$[31.96,\ 32.34]$} & \makecell[c]{$128.06\,\pm\,26.80$\\$[94.78,\ 161.33]$} \\
Rolling mean & \makecell[c]{$29.87\,\pm\,0.53$\\$[29.22,\ 30.52]$} & \makecell[c]{$30.83\,\pm\,0.29$\\$[30.47,\ 31.19]$} & \makecell[c]{$31.79\,\pm\,0.10$\\$[31.67,\ 31.92]$} & \makecell[c]{$129.27\,\pm\,27.80$\\$[94.75,\ 163.79]$} \\
KNN (scaled) & \makecell[c]{$111.50\,\pm\,1.25$\\$[109.95,\ 113.05]$} & \makecell[c]{$104.12\,\pm\,1.46$\\$[102.30,\ 105.93]$} & \makecell[c]{$94.51\,\pm\,1.29$\\$[92.91,\ 96.11]$} & \makecell[c]{$117.73\,\pm\,17.44$\\$[96.08,\ 139.38]$} \\
KNN (unscaled) & \makecell[c]{$41.89\,\pm\,1.19$\\$[40.41,\ 43.36]$} & \makecell[c]{$43.13\,\pm\,0.10$\\$[43.00,\ 43.26]$} & \makecell[c]{$44.49\,\pm\,0.18$\\$[44.27,\ 44.71]$} & \makecell[c]{$117.73\,\pm\,17.44$\\$[96.08,\ 139.38]$} \\
Historical & \makecell[c]{$58.60\,\pm\,1.29$\\$[57.00,\ 60.20]$} & \makecell[c]{$59.35\,\pm\,0.27$\\$[59.02,\ 59.68]$} & \makecell[c]{$59.69\,\pm\,0.45$\\$[59.13,\ 60.25]$} & \makecell[c]{$\mathbf{61.54}\,\pm\,5.85$\\$[54.28,\ 68.80]$} \\
Graph Laplacian & \makecell[c]{$65.70\,\pm\,1.90$\\$[63.34,\ 68.05]$} & \makecell[c]{$66.42\,\pm\,0.46$\\$[65.84,\ 66.99]$} & \makecell[c]{$67.64\,\pm\,0.24$\\$[67.34,\ 67.95]$} & \makecell[c]{$117.73\,\pm\,17.44$\\$[96.08,\ 139.38]$} \\
Graph Temporal & \makecell[c]{$51.95\,\pm\,1.56$\\$[50.01,\ 53.89]$} & \makecell[c]{$53.79\,\pm\,0.45$\\$[53.23,\ 54.34]$} & \makecell[c]{$56.96\,\pm\,0.30$\\$[56.59,\ 57.33]$} & \makecell[c]{$161.98\,\pm\,53.62$\\$[95.40,\ 228.56]$} \\
Sage-Voronoi & \makecell[c]{$29.08\,\pm\,0.89$\\$[27.98,\ 30.18]$} & \makecell[c]{$31.08\,\pm\,0.60$\\$[30.33,\ 31.82]$} & \makecell[c]{$34.28\,\pm\,0.91$\\$[33.15,\ 35.40]$} & \makecell[c]{$126.29\,\pm\,29.88$\\$[89.19,\ 163.38]$} \\
GAIN & \makecell[c]{$39.46\,\pm\,0.66$\\$[38.64,\ 40.27]$} & \makecell[c]{$42.68\,\pm\,1.05$\\$[41.37,\ 43.99]$} & \makecell[c]{$56.56\,\pm\,1.85$\\$[54.27,\ 58.86]$} & \makecell[c]{$152.17\,\pm\,10.18$\\$[139.53,\ 164.81]$} \\
GRIN & \makecell[c]{$\mathbf{21.89}\,\pm\,0.87$\\$[20.81,\ 22.97]$} & \makecell[c]{$23.80\,\pm\,0.44$\\$[23.25,\ 24.36]$} & \makecell[c]{$26.56\,\pm\,0.34$\\$[26.14,\ 26.97]$} & \makecell[c]{$118.37\,\pm\,21.18$\\$[92.07,\ 144.67]$} \\
\bottomrule
\end{tabular}
\endgroup
\end{table}

\clearpage

\begin{table}[htbp]
\centering
\caption{{RMSE for January 2023 (vehicles per ten-minute interval). Each cell gives mean $\pm$ sample SD, followed by the 95\% confidence interval; $n=5$ seed runs. Bold means identify the lowest unrounded mean within each column, not statistical significance.}}
\label{tab:supp_p01_5}
\begingroup
\footnotesize
\setlength{\tabcolsep}{3pt}
\renewcommand{\arraystretch}{1.15}
\begin{tabular}{lcccc}
\toprule
Method & Block $B=1$ & Block $B=2$ & Block $B=4$ & Block $B=8$ \\
\midrule
Linear & \makecell[c]{$\mathbf{70.47}\,\pm\,20.29$\\$[45.28,\ 95.66]$} & \makecell[c]{$\mathbf{73.00}\,\pm\,32.00$\\$[33.27,\ 112.73]$} & \makecell[c]{$\mathbf{90.18}\,\pm\,15.90$\\$[70.43,\ 109.93]$} & \makecell[c]{$102.16\,\pm\,43.66$\\$[47.94,\ 156.37]$} \\
Median & \makecell[c]{$237.40\,\pm\,63.05$\\$[159.11,\ 315.68]$} & \makecell[c]{$197.12\,\pm\,44.82$\\$[141.46,\ 252.77]$} & \makecell[c]{$194.44\,\pm\,30.41$\\$[156.68,\ 232.20]$} & \makecell[c]{$187.63\,\pm\,63.93$\\$[108.25,\ 267.01]$} \\
FFill + BFill & \makecell[c]{$107.72\,\pm\,45.95$\\$[50.67,\ 164.78]$} & \makecell[c]{$91.46\,\pm\,52.07$\\$[26.81,\ 156.11]$} & \makecell[c]{$111.16\,\pm\,33.06$\\$[70.10,\ 152.21]$} & \makecell[c]{$110.70\,\pm\,36.55$\\$[65.31,\ 156.08]$} \\
Rolling mean & \makecell[c]{$97.51\,\pm\,43.90$\\$[43.00,\ 152.01]$} & \makecell[c]{$91.19\,\pm\,32.58$\\$[50.73,\ 131.65]$} & \makecell[c]{$108.42\,\pm\,30.19$\\$[70.94,\ 145.91]$} & \makecell[c]{$102.65\,\pm\,33.45$\\$[61.12,\ 144.18]$} \\
KNN (scaled) & \makecell[c]{$264.65\,\pm\,58.09$\\$[192.51,\ 336.78]$} & \makecell[c]{$224.98\,\pm\,53.22$\\$[158.90,\ 291.06]$} & \makecell[c]{$217.79\,\pm\,33.66$\\$[175.99,\ 259.58]$} & \makecell[c]{$207.70\,\pm\,63.05$\\$[129.41,\ 285.99]$} \\
KNN (unscaled) & \makecell[c]{$195.02\,\pm\,83.22$\\$[91.69,\ 298.35]$} & \makecell[c]{$128.02\,\pm\,45.15$\\$[71.96,\ 184.08]$} & \makecell[c]{$131.66\,\pm\,35.71$\\$[87.32,\ 176.00]$} & \makecell[c]{$125.07\,\pm\,74.00$\\$[33.18,\ 216.95]$} \\
Historical & \makecell[c]{$203.84\,\pm\,84.28$\\$[99.19,\ 308.50]$} & \makecell[c]{$152.95\,\pm\,54.08$\\$[85.80,\ 220.10]$} & \makecell[c]{$154.36\,\pm\,42.95$\\$[101.02,\ 207.69]$} & \makecell[c]{$130.99\,\pm\,68.66$\\$[45.74,\ 216.25]$} \\
Graph Laplacian & \makecell[c]{$213.11\,\pm\,79.85$\\$[113.97,\ 312.25]$} & \makecell[c]{$235.15\,\pm\,66.73$\\$[152.30,\ 318.01]$} & \makecell[c]{$187.75\,\pm\,42.60$\\$[134.85,\ 240.65]$} & \makecell[c]{$233.29\,\pm\,72.04$\\$[143.85,\ 322.73]$} \\
Graph Temporal & \makecell[c]{$163.17\,\pm\,66.48$\\$[80.62,\ 245.72]$} & \makecell[c]{$193.68\,\pm\,25.20$\\$[162.40,\ 224.96]$} & \makecell[c]{$187.32\,\pm\,21.22$\\$[160.97,\ 213.66]$} & \makecell[c]{$228.60\,\pm\,52.96$\\$[162.84,\ 294.35]$} \\
Sage-Voronoi & \makecell[c]{$129.20\,\pm\,72.98$\\$[38.58,\ 219.82]$} & \makecell[c]{$89.69\,\pm\,23.27$\\$[60.80,\ 118.58]$} & \makecell[c]{$94.52\,\pm\,29.23$\\$[58.23,\ 130.81]$} & \makecell[c]{$120.80\,\pm\,48.17$\\$[60.99,\ 180.61]$} \\
GAIN & \makecell[c]{$185.71\,\pm\,90.17$\\$[73.76,\ 297.67]$} & \makecell[c]{$124.79\,\pm\,46.37$\\$[67.21,\ 182.36]$} & \makecell[c]{$134.96\,\pm\,42.75$\\$[81.88,\ 188.05]$} & \makecell[c]{$119.36\,\pm\,70.36$\\$[31.99,\ 206.72]$} \\
GRIN & \makecell[c]{$106.52\,\pm\,62.14$\\$[29.36,\ 183.68]$} & \makecell[c]{$76.15\,\pm\,47.05$\\$[17.72,\ 134.57]$} & \makecell[c]{$96.26\,\pm\,31.96$\\$[56.57,\ 135.94]$} & \makecell[c]{$\mathbf{85.09}\,\pm\,72.27$\\$[-4.65,\ 174.83]$} \\
\bottomrule
\end{tabular}
\endgroup
\end{table}

\begin{table}[htbp]
\centering
\caption{{RMSE for January 2023 (vehicles per ten-minute interval). Each cell gives mean $\pm$ sample SD, followed by the 95\% confidence interval; $n=5$ seed runs. Bold means identify the lowest unrounded mean within each column, not statistical significance.}}
\label{tab:supp_p01_6}
\begingroup
\footnotesize
\setlength{\tabcolsep}{3pt}
\renewcommand{\arraystretch}{1.15}
\begin{tabular}{lccc}
\toprule
Method & Block $B=16$ & Block $B=32$ & Block $B=64$ \\
\midrule
Linear & \makecell[c]{$117.45\,\pm\,56.62$\\$[47.15,\ 187.76]$} & \makecell[c]{$\mathbf{140.47}\,\pm\,17.59$\\$[118.62,\ 162.31]$} & \makecell[c]{$\mathbf{160.97}\,\pm\,25.02$\\$[129.91,\ 192.04]$} \\
Median & \makecell[c]{$186.19\,\pm\,27.29$\\$[152.31,\ 220.07]$} & \makecell[c]{$232.15\,\pm\,52.86$\\$[166.51,\ 297.78]$} & \makecell[c]{$227.39\,\pm\,42.78$\\$[174.27,\ 280.51]$} \\
FFill + BFill & \makecell[c]{$185.62\,\pm\,91.30$\\$[72.26,\ 298.98]$} & \makecell[c]{$220.66\,\pm\,30.07$\\$[183.32,\ 257.99]$} & \makecell[c]{$253.82\,\pm\,55.99$\\$[184.29,\ 323.35]$} \\
Rolling mean & \makecell[c]{$128.77\,\pm\,65.04$\\$[48.02,\ 209.53]$} & \makecell[c]{$143.87\,\pm\,16.75$\\$[123.07,\ 164.67]$} & \makecell[c]{$162.22\,\pm\,24.63$\\$[131.63,\ 192.81]$} \\
KNN (scaled) & \makecell[c]{$212.30\,\pm\,26.08$\\$[179.92,\ 244.68]$} & \makecell[c]{$250.05\,\pm\,48.76$\\$[189.50,\ 310.60]$} & \makecell[c]{$244.24\,\pm\,47.35$\\$[185.45,\ 303.03]$} \\
KNN (unscaled) & \makecell[c]{$132.30\,\pm\,46.42$\\$[74.65,\ 189.94]$} & \makecell[c]{$187.95\,\pm\,75.65$\\$[94.02,\ 281.88]$} & \makecell[c]{$175.93\,\pm\,53.25$\\$[109.82,\ 242.05]$} \\
Historical & \makecell[c]{$161.28\,\pm\,38.73$\\$[113.19,\ 209.37]$} & \makecell[c]{$204.91\,\pm\,70.19$\\$[117.76,\ 292.06]$} & \makecell[c]{$187.97\,\pm\,54.68$\\$[120.08,\ 255.87]$} \\
Graph Laplacian & \makecell[c]{$156.35\,\pm\,26.22$\\$[123.79,\ 188.91]$} & \makecell[c]{$233.91\,\pm\,46.78$\\$[175.83,\ 291.99]$} & \makecell[c]{$230.50\,\pm\,51.56$\\$[166.48,\ 294.52]$} \\
Graph Temporal & \makecell[c]{$182.17\,\pm\,20.91$\\$[156.20,\ 208.14]$} & \makecell[c]{$242.25\,\pm\,53.27$\\$[176.11,\ 308.40]$} & \makecell[c]{$253.29\,\pm\,51.05$\\$[189.90,\ 316.68]$} \\
Sage-Voronoi & \makecell[c]{$183.51\,\pm\,116.16$\\$[39.27,\ 327.75]$} & \makecell[c]{$191.09\,\pm\,37.95$\\$[143.97,\ 238.21]$} & \makecell[c]{$210.21\,\pm\,44.80$\\$[154.58,\ 265.84]$} \\
GAIN & \makecell[c]{$137.08\,\pm\,42.82$\\$[83.92,\ 190.24]$} & \makecell[c]{$191.34\,\pm\,71.70$\\$[102.31,\ 280.36]$} & \makecell[c]{$173.80\,\pm\,55.91$\\$[104.38,\ 243.21]$} \\
GRIN & \makecell[c]{$\mathbf{105.59}\,\pm\,38.09$\\$[58.29,\ 152.88]$} & \makecell[c]{$176.13\,\pm\,73.30$\\$[85.11,\ 267.15]$} & \makecell[c]{$170.73\,\pm\,54.86$\\$[102.60,\ 238.85]$} \\
\bottomrule
\end{tabular}
\endgroup
\end{table}

\clearpage

\begin{table}[htbp]
\centering
\caption{{RMSE for January 2023 (vehicles per ten-minute interval). Each cell gives mean $\pm$ sample SD, followed by the 95\% confidence interval; $n=5$ seed runs. Bold means identify the lowest unrounded mean within each column, not statistical significance.}}
\label{tab:supp_p01_7}
\begingroup
\footnotesize
\setlength{\tabcolsep}{3pt}
\renewcommand{\arraystretch}{1.15}
\begin{tabular}{lcc}
\toprule
Method & Block $B=128$ & Block $B=144$ \\
\midrule
Linear & \makecell[c]{$236.86\,\pm\,25.90$\\$[204.70,\ 269.01]$} & \makecell[c]{$266.46\,\pm\,21.85$\\$[239.33,\ 293.59]$} \\
Median & \makecell[c]{$209.61\,\pm\,36.21$\\$[164.65,\ 254.56]$} & \makecell[c]{$218.97\,\pm\,25.18$\\$[187.71,\ 250.24]$} \\
FFill + BFill & \makecell[c]{$265.15\,\pm\,36.01$\\$[220.44,\ 309.86]$} & \makecell[c]{$297.06\,\pm\,38.22$\\$[249.60,\ 344.52]$} \\
Rolling mean & \makecell[c]{$236.51\,\pm\,25.89$\\$[204.35,\ 268.66]$} & \makecell[c]{$266.98\,\pm\,22.11$\\$[239.53,\ 294.43]$} \\
KNN (scaled) & \makecell[c]{$227.42\,\pm\,32.08$\\$[187.59,\ 267.25]$} & \makecell[c]{$237.08\,\pm\,24.33$\\$[206.87,\ 267.29]$} \\
KNN (unscaled) & \makecell[c]{$161.62\,\pm\,40.44$\\$[111.41,\ 211.83]$} & \makecell[c]{$163.75\,\pm\,31.19$\\$[125.02,\ 202.48]$} \\
Historical & \makecell[c]{$177.85\,\pm\,40.97$\\$[126.98,\ 228.72]$} & \makecell[c]{$180.08\,\pm\,28.24$\\$[145.01,\ 215.15]$} \\
Graph Laplacian & \makecell[c]{$212.51\,\pm\,39.39$\\$[163.60,\ 261.41]$} & \makecell[c]{$217.62\,\pm\,34.26$\\$[175.09,\ 260.16]$} \\
Graph Temporal & \makecell[c]{$232.74\,\pm\,34.71$\\$[189.65,\ 275.83]$} & \makecell[c]{$240.03\,\pm\,26.62$\\$[206.97,\ 273.08]$} \\
Sage-Voronoi & \makecell[c]{$249.63\,\pm\,122.27$\\$[97.81,\ 401.45]$} & \makecell[c]{$199.49\,\pm\,31.27$\\$[160.66,\ 238.33]$} \\
GAIN & \makecell[c]{$154.00\,\pm\,46.06$\\$[96.81,\ 211.18]$} & \makecell[c]{$162.37\,\pm\,34.50$\\$[119.54,\ 205.20]$} \\
GRIN & \makecell[c]{$\mathbf{153.85}\,\pm\,46.64$\\$[95.93,\ 211.77]$} & \makecell[c]{$\mathbf{160.53}\,\pm\,34.04$\\$[118.27,\ 202.80]$} \\
\bottomrule
\end{tabular}
\endgroup
\end{table}

\begin{table}[htbp]
\centering
\caption{{RMSE for January 2023 (vehicles per ten-minute interval). Each cell gives mean $\pm$ sample SD, followed by the 95\% confidence interval; $n=5$ seed runs. Bold means identify the lowest unrounded mean within each column, not statistical significance.}}
\label{tab:supp_p01_8}
\begingroup
\footnotesize
\setlength{\tabcolsep}{3pt}
\renewcommand{\arraystretch}{1.15}
\begin{tabular}{lcccc}
\toprule
Method & MCAR 10\% & MCAR 30\% & MCAR 50\% & Synchronous $B=144$ \\
\midrule
Linear & \makecell[c]{$\mathbf{73.08}\,\pm\,2.79$\\$[69.62,\ 76.55]$} & \makecell[c]{$\mathbf{77.99}\,\pm\,1.36$\\$[76.30,\ 79.67]$} & \makecell[c]{$\mathbf{85.18}\,\pm\,1.47$\\$[83.36,\ 87.01]$} & \makecell[c]{$260.20\,\pm\,69.44$\\$[173.98,\ 346.42]$} \\
Median & \makecell[c]{$214.16\,\pm\,6.83$\\$[205.68,\ 222.65]$} & \makecell[c]{$217.15\,\pm\,1.89$\\$[214.80,\ 219.50]$} & \makecell[c]{$217.70\,\pm\,2.21$\\$[214.95,\ 220.45]$} & \makecell[c]{$227.38\,\pm\,50.65$\\$[164.50,\ 290.27]$} \\
FFill + BFill & \makecell[c]{$93.34\,\pm\,4.77$\\$[87.41,\ 99.27]$} & \makecell[c]{$101.05\,\pm\,3.34$\\$[96.89,\ 105.20]$} & \makecell[c]{$109.34\,\pm\,2.06$\\$[106.78,\ 111.90]$} & \makecell[c]{$257.26\,\pm\,63.29$\\$[178.67,\ 335.85]$} \\
Rolling mean & \makecell[c]{$97.62\,\pm\,5.34$\\$[90.99,\ 104.24]$} & \makecell[c]{$102.53\,\pm\,3.71$\\$[97.93,\ 107.13]$} & \makecell[c]{$106.85\,\pm\,1.51$\\$[104.98,\ 108.73]$} & \makecell[c]{$260.17\,\pm\,69.00$\\$[174.50,\ 345.85]$} \\
KNN (scaled) & \makecell[c]{$231.20\,\pm\,6.22$\\$[223.48,\ 238.93]$} & \makecell[c]{$227.19\,\pm\,1.50$\\$[225.33,\ 229.06]$} & \makecell[c]{$220.11\,\pm\,1.07$\\$[218.78,\ 221.44]$} & \makecell[c]{$231.67\,\pm\,47.12$\\$[173.17,\ 290.18]$} \\
KNN (unscaled) & \makecell[c]{$164.99\,\pm\,9.55$\\$[153.13,\ 176.86]$} & \makecell[c]{$170.92\,\pm\,2.42$\\$[167.92,\ 173.92]$} & \makecell[c]{$173.74\,\pm\,2.88$\\$[170.17,\ 177.32]$} & \makecell[c]{$231.67\,\pm\,47.12$\\$[173.17,\ 290.18]$} \\
Historical & \makecell[c]{$178.95\,\pm\,8.36$\\$[168.57,\ 189.33]$} & \makecell[c]{$184.79\,\pm\,1.76$\\$[182.60,\ 186.97]$} & \makecell[c]{$186.38\,\pm\,2.39$\\$[183.41,\ 189.34]$} & \makecell[c]{$\mathbf{196.38}\,\pm\,45.90$\\$[139.38,\ 253.37]$} \\
Graph Laplacian & \makecell[c]{$219.24\,\pm\,12.47$\\$[203.76,\ 234.72]$} & \makecell[c]{$222.58\,\pm\,1.63$\\$[220.56,\ 224.60]$} & \makecell[c]{$223.94\,\pm\,1.58$\\$[221.98,\ 225.90]$} & \makecell[c]{$231.67\,\pm\,47.12$\\$[173.17,\ 290.18]$} \\
Graph Temporal & \makecell[c]{$164.59\,\pm\,6.30$\\$[156.77,\ 172.41]$} & \makecell[c]{$173.76\,\pm\,1.50$\\$[171.90,\ 175.62]$} & \makecell[c]{$185.57\,\pm\,2.11$\\$[182.95,\ 188.19]$} & \makecell[c]{$318.46\,\pm\,96.51$\\$[198.62,\ 438.29]$} \\
Sage-Voronoi & \makecell[c]{$106.03\,\pm\,5.00$\\$[99.83,\ 112.23]$} & \makecell[c]{$115.38\,\pm\,6.97$\\$[106.72,\ 124.03]$} & \makecell[c]{$126.29\,\pm\,6.11$\\$[118.71,\ 133.88]$} & \makecell[c]{$267.60\,\pm\,61.56$\\$[191.16,\ 344.03]$} \\
GAIN & \makecell[c]{$161.64\,\pm\,9.57$\\$[149.76,\ 173.52]$} & \makecell[c]{$163.12\,\pm\,3.90$\\$[158.28,\ 167.96]$} & \makecell[c]{$169.15\,\pm\,3.01$\\$[165.41,\ 172.89]$} & \makecell[c]{$287.15\,\pm\,29.25$\\$[250.83,\ 323.47]$} \\
GRIN & \makecell[c]{$94.67\,\pm\,8.69$\\$[83.87,\ 105.46]$} & \makecell[c]{$106.46\,\pm\,7.05$\\$[97.71,\ 115.21]$} & \makecell[c]{$118.22\,\pm\,6.31$\\$[110.38,\ 126.06]$} & \makecell[c]{$241.27\,\pm\,50.79$\\$[178.21,\ 304.34]$} \\
\bottomrule
\end{tabular}
\endgroup
\end{table}

\clearpage

\begin{table}[htbp]
\centering
\caption{{MRE for January 2023 (dimensionless ratio). Each cell gives mean $\pm$ sample SD, followed by the 95\% confidence interval; $n=5$ seed runs. Bold means identify the lowest unrounded mean within each column, not statistical significance.}}
\label{tab:supp_p01_9}
\begingroup
\footnotesize
\setlength{\tabcolsep}{3pt}
\renewcommand{\arraystretch}{1.15}
\begin{tabular}{lcccc}
\toprule
Method & Block $B=1$ & Block $B=2$ & Block $B=4$ & Block $B=8$ \\
\midrule
Linear & \makecell[c]{$0.232\,\pm\,0.035$\\$[0.188,\ 0.276]$} & \makecell[c]{$0.292\,\pm\,0.058$\\$[0.220,\ 0.364]$} & \makecell[c]{$0.302\,\pm\,0.033$\\$[0.261,\ 0.342]$} & \makecell[c]{$0.332\,\pm\,0.040$\\$[0.282,\ 0.383]$} \\
Median & \makecell[c]{$1.781\,\pm\,0.317$\\$[1.387,\ 2.175]$} & \makecell[c]{$2.332\,\pm\,0.465$\\$[1.754,\ 2.910]$} & \makecell[c]{$2.576\,\pm\,0.364$\\$[2.124,\ 3.028]$} & \makecell[c]{$2.237\,\pm\,0.424$\\$[1.710,\ 2.764]$} \\
FFill + BFill & \makecell[c]{$0.271\,\pm\,0.063$\\$[0.192,\ 0.350]$} & \makecell[c]{$0.344\,\pm\,0.050$\\$[0.282,\ 0.405]$} & \makecell[c]{$0.348\,\pm\,0.049$\\$[0.287,\ 0.409]$} & \makecell[c]{$0.449\,\pm\,0.058$\\$[0.377,\ 0.521]$} \\
Rolling mean & \makecell[c]{$0.287\,\pm\,0.041$\\$[0.236,\ 0.337]$} & \makecell[c]{$0.351\,\pm\,0.086$\\$[0.245,\ 0.458]$} & \makecell[c]{$0.351\,\pm\,0.048$\\$[0.292,\ 0.410]$} & \makecell[c]{$0.360\,\pm\,0.051$\\$[0.297,\ 0.424]$} \\
KNN (scaled) & \makecell[c]{$0.847\,\pm\,0.123$\\$[0.693,\ 1.000]$} & \makecell[c]{$1.044\,\pm\,0.159$\\$[0.847,\ 1.242]$} & \makecell[c]{$1.166\,\pm\,0.146$\\$[0.984,\ 1.348]$} & \makecell[c]{$0.982\,\pm\,0.173$\\$[0.767,\ 1.197]$} \\
KNN (unscaled) & \makecell[c]{$0.405\,\pm\,0.093$\\$[0.290,\ 0.520]$} & \makecell[c]{$0.423\,\pm\,0.048$\\$[0.363,\ 0.483]$} & \makecell[c]{$0.450\,\pm\,0.069$\\$[0.364,\ 0.536]$} & \makecell[c]{$0.413\,\pm\,0.069$\\$[0.327,\ 0.499]$} \\
Historical & \makecell[c]{$0.514\,\pm\,0.089$\\$[0.403,\ 0.625]$} & \makecell[c]{$0.541\,\pm\,0.097$\\$[0.420,\ 0.661]$} & \makecell[c]{$0.512\,\pm\,0.080$\\$[0.412,\ 0.611]$} & \makecell[c]{$0.540\,\pm\,0.048$\\$[0.480,\ 0.600]$} \\
Graph Laplacian & \makecell[c]{$0.727\,\pm\,0.089$\\$[0.616,\ 0.838]$} & \makecell[c]{$0.838\,\pm\,0.153$\\$[0.648,\ 1.028]$} & \makecell[c]{$1.078\,\pm\,0.147$\\$[0.895,\ 1.260]$} & \makecell[c]{$0.949\,\pm\,0.292$\\$[0.588,\ 1.311]$} \\
Graph Temporal & \makecell[c]{$0.675\,\pm\,0.207$\\$[0.418,\ 0.933]$} & \makecell[c]{$0.805\,\pm\,0.298$\\$[0.435,\ 1.176]$} & \makecell[c]{$1.362\,\pm\,0.350$\\$[0.927,\ 1.797]$} & \makecell[c]{$1.087\,\pm\,0.268$\\$[0.754,\ 1.420]$} \\
Sage-Voronoi & \makecell[c]{$0.255\,\pm\,0.038$\\$[0.208,\ 0.303]$} & \makecell[c]{$0.311\,\pm\,0.053$\\$[0.245,\ 0.376]$} & \makecell[c]{$0.364\,\pm\,0.067$\\$[0.281,\ 0.447]$} & \makecell[c]{$0.377\,\pm\,0.067$\\$[0.293,\ 0.461]$} \\
GAIN & \makecell[c]{$0.311\,\pm\,0.073$\\$[0.220,\ 0.402]$} & \makecell[c]{$0.305\,\pm\,0.031$\\$[0.267,\ 0.343]$} & \makecell[c]{$0.320\,\pm\,0.032$\\$[0.280,\ 0.359]$} & \makecell[c]{$0.317\,\pm\,0.030$\\$[0.281,\ 0.354]$} \\
GRIN & \makecell[c]{$\mathbf{0.200}\,\pm\,0.013$\\$[0.184,\ 0.217]$} & \makecell[c]{$\mathbf{0.214}\,\pm\,0.026$\\$[0.181,\ 0.246]$} & \makecell[c]{$\mathbf{0.261}\,\pm\,0.052$\\$[0.196,\ 0.326]$} & \makecell[c]{$\mathbf{0.292}\,\pm\,0.104$\\$[0.163,\ 0.420]$} \\
\bottomrule
\end{tabular}
\endgroup
\end{table}

\begin{table}[htbp]
\centering
\caption{{MRE for January 2023 (dimensionless ratio). Each cell gives mean $\pm$ sample SD, followed by the 95\% confidence interval; $n=5$ seed runs. Bold means identify the lowest unrounded mean within each column, not statistical significance.}}
\label{tab:supp_p01_10}
\begingroup
\footnotesize
\setlength{\tabcolsep}{3pt}
\renewcommand{\arraystretch}{1.15}
\begin{tabular}{lccc}
\toprule
Method & Block $B=16$ & Block $B=32$ & Block $B=64$ \\
\midrule
Linear & \makecell[c]{$0.376\,\pm\,0.060$\\$[0.302,\ 0.450]$} & \makecell[c]{$0.886\,\pm\,0.094$\\$[0.770,\ 1.003]$} & \makecell[c]{$1.628\,\pm\,0.127$\\$[1.471,\ 1.786]$} \\
Median & \makecell[c]{$2.015\,\pm\,0.438$\\$[1.472,\ 2.559]$} & \makecell[c]{$2.383\,\pm\,0.339$\\$[1.962,\ 2.804]$} & \makecell[c]{$2.272\,\pm\,0.231$\\$[1.985,\ 2.559]$} \\
FFill + BFill & \makecell[c]{$0.619\,\pm\,0.065$\\$[0.538,\ 0.699]$} & \makecell[c]{$1.136\,\pm\,0.173$\\$[0.920,\ 1.351]$} & \makecell[c]{$2.281\,\pm\,0.351$\\$[1.845,\ 2.716]$} \\
Rolling mean & \makecell[c]{$0.380\,\pm\,0.049$\\$[0.319,\ 0.440]$} & \makecell[c]{$0.809\,\pm\,0.095$\\$[0.691,\ 0.928]$} & \makecell[c]{$1.578\,\pm\,0.124$\\$[1.424,\ 1.733]$} \\
KNN (scaled) & \makecell[c]{$0.940\,\pm\,0.153$\\$[0.749,\ 1.131]$} & \makecell[c]{$1.034\,\pm\,0.063$\\$[0.955,\ 1.112]$} & \makecell[c]{$1.034\,\pm\,0.087$\\$[0.926,\ 1.141]$} \\
KNN (unscaled) & \makecell[c]{$0.420\,\pm\,0.096$\\$[0.300,\ 0.539]$} & \makecell[c]{$0.445\,\pm\,0.039$\\$[0.396,\ 0.494]$} & \makecell[c]{$0.430\,\pm\,0.033$\\$[0.389,\ 0.471]$} \\
Historical & \makecell[c]{$0.506\,\pm\,0.078$\\$[0.410,\ 0.603]$} & \makecell[c]{$0.547\,\pm\,0.018$\\$[0.525,\ 0.569]$} & \makecell[c]{$0.568\,\pm\,0.079$\\$[0.470,\ 0.665]$} \\
Graph Laplacian & \makecell[c]{$0.868\,\pm\,0.316$\\$[0.476,\ 1.261]$} & \makecell[c]{$0.953\,\pm\,0.157$\\$[0.758,\ 1.148]$} & \makecell[c]{$0.922\,\pm\,0.159$\\$[0.725,\ 1.119]$} \\
Graph Temporal & \makecell[c]{$0.879\,\pm\,0.495$\\$[0.265,\ 1.493]$} & \makecell[c]{$1.342\,\pm\,0.638$\\$[0.550,\ 2.134]$} & \makecell[c]{$1.008\,\pm\,0.423$\\$[0.482,\ 1.533]$} \\
Sage-Voronoi & \makecell[c]{$0.441\,\pm\,0.100$\\$[0.316,\ 0.565]$} & \makecell[c]{$0.576\,\pm\,0.095$\\$[0.458,\ 0.694]$} & \makecell[c]{$0.650\,\pm\,0.106$\\$[0.517,\ 0.782]$} \\
GAIN & \makecell[c]{$0.305\,\pm\,0.057$\\$[0.234,\ 0.376]$} & \makecell[c]{$\mathbf{0.314}\,\pm\,0.036$\\$[0.269,\ 0.359]$} & \makecell[c]{$\mathbf{0.311}\,\pm\,0.026$\\$[0.279,\ 0.343]$} \\
GRIN & \makecell[c]{$\mathbf{0.290}\,\pm\,0.071$\\$[0.202,\ 0.377]$} & \makecell[c]{$0.463\,\pm\,0.239$\\$[0.166,\ 0.760]$} & \makecell[c]{$0.416\,\pm\,0.168$\\$[0.206,\ 0.625]$} \\
\bottomrule
\end{tabular}
\endgroup
\end{table}

\clearpage

\begin{table}[htbp]
\centering
\caption{{MRE for January 2023 (dimensionless ratio). Each cell gives mean $\pm$ sample SD, followed by the 95\% confidence interval; $n=5$ seed runs. Bold means identify the lowest unrounded mean within each column, not statistical significance.}}
\label{tab:supp_p01_11}
\begingroup
\footnotesize
\setlength{\tabcolsep}{3pt}
\renewcommand{\arraystretch}{1.15}
\begin{tabular}{lcc}
\toprule
Method & Block $B=128$ & Block $B=144$ \\
\midrule
Linear & \makecell[c]{$2.858\,\pm\,0.141$\\$[2.683,\ 3.034]$} & \makecell[c]{$2.936\,\pm\,0.296$\\$[2.568,\ 3.303]$} \\
Median & \makecell[c]{$2.292\,\pm\,0.093$\\$[2.176,\ 2.408]$} & \makecell[c]{$2.269\,\pm\,0.066$\\$[2.188,\ 2.351]$} \\
FFill + BFill & \makecell[c]{$3.035\,\pm\,0.088$\\$[2.926,\ 3.144]$} & \makecell[c]{$2.918\,\pm\,0.233$\\$[2.629,\ 3.206]$} \\
Rolling mean & \makecell[c]{$2.847\,\pm\,0.139$\\$[2.674,\ 3.020]$} & \makecell[c]{$2.935\,\pm\,0.295$\\$[2.568,\ 3.301]$} \\
KNN (scaled) & \makecell[c]{$1.031\,\pm\,0.053$\\$[0.965,\ 1.098]$} & \makecell[c]{$1.051\,\pm\,0.068$\\$[0.966,\ 1.136]$} \\
KNN (unscaled) & \makecell[c]{$0.429\,\pm\,0.024$\\$[0.400,\ 0.459]$} & \makecell[c]{$0.457\,\pm\,0.035$\\$[0.414,\ 0.501]$} \\
Historical & \makecell[c]{$0.545\,\pm\,0.036$\\$[0.501,\ 0.590]$} & \makecell[c]{$0.552\,\pm\,0.018$\\$[0.530,\ 0.574]$} \\
Graph Laplacian & \makecell[c]{$0.901\,\pm\,0.030$\\$[0.864,\ 0.938]$} & \makecell[c]{$0.900\,\pm\,0.061$\\$[0.824,\ 0.976]$} \\
Graph Temporal & \makecell[c]{$1.222\,\pm\,0.194$\\$[0.981,\ 1.464]$} & \makecell[c]{$1.050\,\pm\,0.158$\\$[0.853,\ 1.247]$} \\
Sage-Voronoi & \makecell[c]{$0.787\,\pm\,0.154$\\$[0.596,\ 0.978]$} & \makecell[c]{$0.796\,\pm\,0.132$\\$[0.632,\ 0.959]$} \\
GAIN & \makecell[c]{$\mathbf{0.325}\,\pm\,0.030$\\$[0.287,\ 0.362]$} & \makecell[c]{$\mathbf{0.357}\,\pm\,0.038$\\$[0.309,\ 0.404]$} \\
GRIN & \makecell[c]{$0.397\,\pm\,0.032$\\$[0.357,\ 0.436]$} & \makecell[c]{$0.440\,\pm\,0.025$\\$[0.410,\ 0.471]$} \\
\bottomrule
\end{tabular}
\endgroup
\end{table}

\begin{table}[htbp]
\centering
\caption{{MRE for January 2023 (dimensionless ratio). Each cell gives mean $\pm$ sample SD, followed by the 95\% confidence interval; $n=5$ seed runs. Bold means identify the lowest unrounded mean within each column, not statistical significance.}}
\label{tab:supp_p01_12}
\begingroup
\footnotesize
\setlength{\tabcolsep}{3pt}
\renewcommand{\arraystretch}{1.15}
\begin{tabular}{lcccc}
\toprule
Method & MCAR 10\% & MCAR 30\% & MCAR 50\% & Synchronous $B=144$ \\
\midrule
Linear & \makecell[c]{$0.258\,\pm\,0.003$\\$[0.254,\ 0.263]$} & \makecell[c]{$0.266\,\pm\,0.003$\\$[0.262,\ 0.269]$} & \makecell[c]{$0.277\,\pm\,0.004$\\$[0.273,\ 0.282]$} & \makecell[c]{$4.287\,\pm\,0.995$\\$[3.051,\ 5.524]$} \\
Median & \makecell[c]{$2.268\,\pm\,0.038$\\$[2.221,\ 2.315]$} & \makecell[c]{$2.269\,\pm\,0.020$\\$[2.244,\ 2.294]$} & \makecell[c]{$2.270\,\pm\,0.010$\\$[2.258,\ 2.283]$} & \makecell[c]{$2.500\,\pm\,0.369$\\$[2.042,\ 2.958]$} \\
FFill + BFill & \makecell[c]{$0.306\,\pm\,0.004$\\$[0.301,\ 0.311]$} & \makecell[c]{$0.318\,\pm\,0.004$\\$[0.313,\ 0.323]$} & \makecell[c]{$0.338\,\pm\,0.003$\\$[0.333,\ 0.342]$} & \makecell[c]{$3.846\,\pm\,1.724$\\$[1.706,\ 5.987]$} \\
Rolling mean & \makecell[c]{$0.305\,\pm\,0.004$\\$[0.300,\ 0.310]$} & \makecell[c]{$0.316\,\pm\,0.004$\\$[0.311,\ 0.321]$} & \makecell[c]{$0.328\,\pm\,0.005$\\$[0.322,\ 0.334]$} & \makecell[c]{$4.284\,\pm\,0.999$\\$[3.044,\ 5.524]$} \\
KNN (scaled) & \makecell[c]{$1.028\,\pm\,0.027$\\$[0.994,\ 1.062]$} & \makecell[c]{$1.045\,\pm\,0.024$\\$[1.015,\ 1.074]$} & \makecell[c]{$1.008\,\pm\,0.047$\\$[0.950,\ 1.067]$} & \makecell[c]{$2.599\,\pm\,0.414$\\$[2.085,\ 3.113]$} \\
KNN (unscaled) & \makecell[c]{$0.433\,\pm\,0.005$\\$[0.427,\ 0.439]$} & \makecell[c]{$0.469\,\pm\,0.014$\\$[0.452,\ 0.486]$} & \makecell[c]{$0.530\,\pm\,0.016$\\$[0.510,\ 0.550]$} & \makecell[c]{$2.599\,\pm\,0.414$\\$[2.085,\ 3.113]$} \\
Historical & \makecell[c]{$0.534\,\pm\,0.007$\\$[0.526,\ 0.542]$} & \makecell[c]{$0.539\,\pm\,0.006$\\$[0.532,\ 0.546]$} & \makecell[c]{$0.541\,\pm\,0.002$\\$[0.539,\ 0.544]$} & \makecell[c]{$\mathbf{0.590}\,\pm\,0.326$\\$[0.186,\ 0.995]$} \\
Graph Laplacian & \makecell[c]{$0.888\,\pm\,0.016$\\$[0.869,\ 0.907]$} & \makecell[c]{$0.884\,\pm\,0.017$\\$[0.863,\ 0.906]$} & \makecell[c]{$0.885\,\pm\,0.029$\\$[0.849,\ 0.921]$} & \makecell[c]{$2.599\,\pm\,0.414$\\$[2.085,\ 3.113]$} \\
Graph Temporal & \makecell[c]{$0.712\,\pm\,0.051$\\$[0.648,\ 0.776]$} & \makecell[c]{$0.750\,\pm\,0.031$\\$[0.711,\ 0.788]$} & \makecell[c]{$0.786\,\pm\,0.029$\\$[0.750,\ 0.822]$} & \makecell[c]{$5.008\,\pm\,1.311$\\$[3.380,\ 6.635]$} \\
Sage-Voronoi & \makecell[c]{$0.292\,\pm\,0.011$\\$[0.278,\ 0.306]$} & \makecell[c]{$0.306\,\pm\,0.011$\\$[0.292,\ 0.320]$} & \makecell[c]{$0.331\,\pm\,0.014$\\$[0.313,\ 0.348]$} & \makecell[c]{$3.198\,\pm\,1.499$\\$[1.336,\ 5.059]$} \\
GAIN & \makecell[c]{$0.333\,\pm\,0.019$\\$[0.309,\ 0.357]$} & \makecell[c]{$0.528\,\pm\,0.037$\\$[0.481,\ 0.574]$} & \makecell[c]{$1.000\,\pm\,0.061$\\$[0.925,\ 1.076]$} & \makecell[c]{$4.976\,\pm\,0.899$\\$[3.860,\ 6.092]$} \\
GRIN & \makecell[c]{$\mathbf{0.227}\,\pm\,0.006$\\$[0.220,\ 0.234]$} & \makecell[c]{$\mathbf{0.242}\,\pm\,0.007$\\$[0.234,\ 0.251]$} & \makecell[c]{$\mathbf{0.274}\,\pm\,0.012$\\$[0.259,\ 0.290]$} & \makecell[c]{$1.902\,\pm\,0.452$\\$[1.341,\ 2.463]$} \\
\bottomrule
\end{tabular}
\endgroup
\end{table}

\clearpage

\begin{table}[htbp]
\centering
\caption{{MAE for October--December 2022 (vehicles per ten-minute interval). Each cell gives mean $\pm$ sample SD, followed by the 95\% confidence interval; $n=5$ seed runs. Bold means identify the lowest unrounded mean within each column, not statistical significance.}}
\label{tab:supp_p01_13}
\begingroup
\footnotesize
\setlength{\tabcolsep}{3pt}
\renewcommand{\arraystretch}{1.15}
\begin{tabular}{lcccc}
\toprule
Method & Block $B=1$ & Block $B=2$ & Block $B=4$ & Block $B=8$ \\
\midrule
Linear & \makecell[c]{$19.81\,\pm\,2.44$\\$[16.78,\ 22.84]$} & \makecell[c]{$20.40\,\pm\,3.17$\\$[16.47,\ 24.34]$} & \makecell[c]{$19.29\,\pm\,1.95$\\$[16.87,\ 21.71]$} & \makecell[c]{$24.73\,\pm\,2.04$\\$[22.20,\ 27.27]$} \\
Median & \makecell[c]{$102.68\,\pm\,8.81$\\$[91.74,\ 113.62]$} & \makecell[c]{$98.69\,\pm\,6.38$\\$[90.76,\ 106.62]$} & \makecell[c]{$99.62\,\pm\,4.84$\\$[93.60,\ 105.63]$} & \makecell[c]{$101.06\,\pm\,2.28$\\$[98.23,\ 103.89]$} \\
FFill + BFill & \makecell[c]{$24.12\,\pm\,1.84$\\$[21.84,\ 26.41]$} & \makecell[c]{$24.05\,\pm\,3.15$\\$[20.14,\ 27.97]$} & \makecell[c]{$25.98\,\pm\,2.06$\\$[23.42,\ 28.53]$} & \makecell[c]{$33.80\,\pm\,2.86$\\$[30.25,\ 37.34]$} \\
Rolling mean & \makecell[c]{$25.79\,\pm\,2.25$\\$[23.01,\ 28.58]$} & \makecell[c]{$28.69\,\pm\,2.92$\\$[25.07,\ 32.32]$} & \makecell[c]{$26.10\,\pm\,2.46$\\$[23.05,\ 29.16]$} & \makecell[c]{$27.79\,\pm\,2.29$\\$[24.94,\ 30.63]$} \\
KNN (scaled) & \makecell[c]{$41.94\,\pm\,6.22$\\$[34.22,\ 49.66]$} & \makecell[c]{$37.53\,\pm\,3.09$\\$[33.70,\ 41.37]$} & \makecell[c]{$35.73\,\pm\,7.47$\\$[26.45,\ 45.01]$} & \makecell[c]{$34.38\,\pm\,1.87$\\$[32.05,\ 36.71]$} \\
KNN (unscaled) & \makecell[c]{$44.49\,\pm\,8.13$\\$[34.40,\ 54.58]$} & \makecell[c]{$42.32\,\pm\,3.99$\\$[37.37,\ 47.27]$} & \makecell[c]{$39.84\,\pm\,6.45$\\$[31.84,\ 47.85]$} & \makecell[c]{$38.07\,\pm\,2.23$\\$[35.31,\ 40.84]$} \\
Historical & \makecell[c]{$53.64\,\pm\,8.97$\\$[42.49,\ 64.78]$} & \makecell[c]{$56.73\,\pm\,4.86$\\$[50.69,\ 62.76]$} & \makecell[c]{$53.69\,\pm\,6.27$\\$[45.91,\ 61.48]$} & \makecell[c]{$53.52\,\pm\,3.51$\\$[49.16,\ 57.88]$} \\
Graph Laplacian & \makecell[c]{$47.62\,\pm\,8.24$\\$[37.39,\ 57.85]$} & \makecell[c]{$40.26\,\pm\,5.14$\\$[33.88,\ 46.64]$} & \makecell[c]{$40.84\,\pm\,4.58$\\$[35.15,\ 46.53]$} & \makecell[c]{$41.82\,\pm\,4.33$\\$[36.44,\ 47.19]$} \\
Graph Temporal & \makecell[c]{$36.23\,\pm\,4.40$\\$[30.77,\ 41.70]$} & \makecell[c]{$34.03\,\pm\,5.69$\\$[26.96,\ 41.10]$} & \makecell[c]{$38.15\,\pm\,3.88$\\$[33.33,\ 42.97]$} & \makecell[c]{$40.25\,\pm\,2.94$\\$[36.60,\ 43.91]$} \\
Sage-Voronoi & \makecell[c]{$22.43\,\pm\,3.53$\\$[18.05,\ 26.82]$} & \makecell[c]{$22.24\,\pm\,2.59$\\$[19.03,\ 25.45]$} & \makecell[c]{$21.08\,\pm\,2.36$\\$[18.14,\ 24.01]$} & \makecell[c]{$24.73\,\pm\,0.95$\\$[23.56,\ 25.91]$} \\
GAIN & \makecell[c]{$44.62\,\pm\,8.77$\\$[33.72,\ 55.51]$} & \makecell[c]{$39.57\,\pm\,4.64$\\$[33.81,\ 45.33]$} & \makecell[c]{$38.02\,\pm\,5.42$\\$[31.29,\ 44.75]$} & \makecell[c]{$38.20\,\pm\,1.73$\\$[36.05,\ 40.35]$} \\
GRIN & \makecell[c]{$\mathbf{18.30}\,\pm\,3.21$\\$[14.31,\ 22.29]$} & \makecell[c]{$\mathbf{17.70}\,\pm\,2.86$\\$[14.15,\ 21.26]$} & \makecell[c]{$\mathbf{15.90}\,\pm\,2.27$\\$[13.09,\ 18.72]$} & \makecell[c]{$\mathbf{18.45}\,\pm\,0.99$\\$[17.22,\ 19.68]$} \\
\bottomrule
\end{tabular}
\endgroup
\end{table}

\begin{table}[htbp]
\centering
\caption{{MAE for October--December 2022 (vehicles per ten-minute interval). Each cell gives mean $\pm$ sample SD, followed by the 95\% confidence interval; $n=5$ seed runs. Bold means identify the lowest unrounded mean within each column, not statistical significance.}}
\label{tab:supp_p01_14}
\begingroup
\footnotesize
\setlength{\tabcolsep}{3pt}
\renewcommand{\arraystretch}{1.15}
\begin{tabular}{lccc}
\toprule
Method & Block $B=16$ & Block $B=32$ & Block $B=64$ \\
\midrule
Linear & \makecell[c]{$29.55\,\pm\,2.19$\\$[26.83,\ 32.28]$} & \makecell[c]{$38.63\,\pm\,4.86$\\$[32.61,\ 44.66]$} & \makecell[c]{$64.95\,\pm\,3.00$\\$[61.23,\ 68.68]$} \\
Median & \makecell[c]{$110.37\,\pm\,6.77$\\$[101.96,\ 118.78]$} & \makecell[c]{$102.83\,\pm\,4.48$\\$[97.27,\ 108.39]$} & \makecell[c]{$104.40\,\pm\,6.13$\\$[96.79,\ 112.01]$} \\
FFill + BFill & \makecell[c]{$49.46\,\pm\,2.29$\\$[46.62,\ 52.30]$} & \makecell[c]{$73.11\,\pm\,6.86$\\$[64.60,\ 81.62]$} & \makecell[c]{$108.15\,\pm\,5.41$\\$[101.43,\ 114.88]$} \\
Rolling mean & \makecell[c]{$31.48\,\pm\,1.67$\\$[29.40,\ 33.55]$} & \makecell[c]{$39.65\,\pm\,4.86$\\$[33.61,\ 45.68]$} & \makecell[c]{$64.99\,\pm\,3.27$\\$[60.93,\ 69.05]$} \\
KNN (scaled) & \makecell[c]{$40.64\,\pm\,4.06$\\$[35.60,\ 45.68]$} & \makecell[c]{$34.90\,\pm\,3.22$\\$[30.90,\ 38.90]$} & \makecell[c]{$39.09\,\pm\,6.33$\\$[31.24,\ 46.95]$} \\
KNN (unscaled) & \makecell[c]{$43.65\,\pm\,4.47$\\$[38.10,\ 49.20]$} & \makecell[c]{$38.40\,\pm\,3.75$\\$[33.74,\ 43.05]$} & \makecell[c]{$42.98\,\pm\,6.74$\\$[34.62,\ 51.35]$} \\
Historical & \makecell[c]{$59.68\,\pm\,3.64$\\$[55.17,\ 64.19]$} & \makecell[c]{$50.95\,\pm\,4.80$\\$[44.99,\ 56.92]$} & \makecell[c]{$55.25\,\pm\,6.06$\\$[47.72,\ 62.78]$} \\
Graph Laplacian & \makecell[c]{$48.49\,\pm\,5.12$\\$[42.13,\ 54.85]$} & \makecell[c]{$42.23\,\pm\,4.01$\\$[37.25,\ 47.20]$} & \makecell[c]{$46.34\,\pm\,8.24$\\$[36.11,\ 56.57]$} \\
Graph Temporal & \makecell[c]{$48.14\,\pm\,5.03$\\$[41.90,\ 54.38]$} & \makecell[c]{$42.41\,\pm\,5.12$\\$[36.06,\ 48.76]$} & \makecell[c]{$46.77\,\pm\,6.80$\\$[38.33,\ 55.21]$} \\
Sage-Voronoi & \makecell[c]{$29.14\,\pm\,1.16$\\$[27.70,\ 30.58]$} & \makecell[c]{$33.67\,\pm\,3.43$\\$[29.42,\ 37.93]$} & \makecell[c]{$42.28\,\pm\,7.28$\\$[33.24,\ 51.32]$} \\
GAIN & \makecell[c]{$46.75\,\pm\,5.10$\\$[40.42,\ 53.09]$} & \makecell[c]{$38.15\,\pm\,3.65$\\$[33.62,\ 42.68]$} & \makecell[c]{$42.50\,\pm\,5.96$\\$[35.10,\ 49.90]$} \\
GRIN & \makecell[c]{$\mathbf{23.92}\,\pm\,1.32$\\$[22.28,\ 25.55]$} & \makecell[c]{$\mathbf{25.18}\,\pm\,2.95$\\$[21.52,\ 28.84]$} & \makecell[c]{$\mathbf{32.65}\,\pm\,6.35$\\$[24.77,\ 40.53]$} \\
\bottomrule
\end{tabular}
\endgroup
\end{table}

\clearpage

\begin{table}[htbp]
\centering
\caption{{MAE for October--December 2022 (vehicles per ten-minute interval). Each cell gives mean $\pm$ sample SD, followed by the 95\% confidence interval; $n=5$ seed runs. Bold means identify the lowest unrounded mean within each column, not statistical significance.}}
\label{tab:supp_p01_15}
\begingroup
\footnotesize
\setlength{\tabcolsep}{3pt}
\renewcommand{\arraystretch}{1.15}
\begin{tabular}{lcc}
\toprule
Method & Block $B=128$ & Block $B=144$ \\
\midrule
Linear & \makecell[c]{$114.24\,\pm\,4.65$\\$[108.46,\ 120.01]$} & \makecell[c]{$117.12\,\pm\,6.64$\\$[108.88,\ 125.37]$} \\
Median & \makecell[c]{$101.95\,\pm\,2.80$\\$[98.47,\ 105.43]$} & \makecell[c]{$103.00\,\pm\,3.04$\\$[99.23,\ 106.78]$} \\
FFill + BFill & \makecell[c]{$129.04\,\pm\,4.19$\\$[123.83,\ 134.24]$} & \makecell[c]{$118.93\,\pm\,7.12$\\$[110.08,\ 127.77]$} \\
Rolling mean & \makecell[c]{$114.13\,\pm\,4.70$\\$[108.30,\ 119.97]$} & \makecell[c]{$117.13\,\pm\,6.61$\\$[108.93,\ 125.34]$} \\
KNN (scaled) & \makecell[c]{$35.83\,\pm\,0.94$\\$[34.67,\ 37.00]$} & \makecell[c]{$35.87\,\pm\,1.81$\\$[33.63,\ 38.11]$} \\
KNN (unscaled) & \makecell[c]{$39.48\,\pm\,1.41$\\$[37.73,\ 41.23]$} & \makecell[c]{$39.83\,\pm\,2.04$\\$[37.30,\ 42.36]$} \\
Historical & \makecell[c]{$52.27\,\pm\,2.41$\\$[49.27,\ 55.26]$} & \makecell[c]{$52.31\,\pm\,2.66$\\$[49.00,\ 55.61]$} \\
Graph Laplacian & \makecell[c]{$43.96\,\pm\,2.46$\\$[40.91,\ 47.01]$} & \makecell[c]{$44.07\,\pm\,2.86$\\$[40.52,\ 47.63]$} \\
Graph Temporal & \makecell[c]{$44.68\,\pm\,1.77$\\$[42.47,\ 46.88]$} & \makecell[c]{$44.67\,\pm\,3.18$\\$[40.72,\ 48.62]$} \\
Sage-Voronoi & \makecell[c]{$42.32\,\pm\,2.54$\\$[39.16,\ 45.47]$} & \makecell[c]{$42.62\,\pm\,3.06$\\$[38.82,\ 46.42]$} \\
GAIN & \makecell[c]{$39.46\,\pm\,2.29$\\$[36.61,\ 42.30]$} & \makecell[c]{$39.37\,\pm\,3.60$\\$[34.90,\ 43.85]$} \\
GRIN & \makecell[c]{$\mathbf{32.24}\,\pm\,1.78$\\$[30.03,\ 34.45]$} & \makecell[c]{$\mathbf{31.50}\,\pm\,2.96$\\$[27.83,\ 35.18]$} \\
\bottomrule
\end{tabular}
\endgroup
\end{table}

\begin{table}[htbp]
\centering
\caption{{MAE for October--December 2022 (vehicles per ten-minute interval). Each cell gives mean $\pm$ sample SD, followed by the 95\% confidence interval; $n=5$ seed runs. Bold means identify the lowest unrounded mean within each column, not statistical significance.}}
\label{tab:supp_p01_16}
\begingroup
\footnotesize
\setlength{\tabcolsep}{3pt}
\renewcommand{\arraystretch}{1.15}
\begin{tabular}{lcccc}
\toprule
Method & MCAR 10\% & MCAR 30\% & MCAR 50\% & Synchronous $B=144$ \\
\midrule
Linear & \makecell[c]{$20.26\,\pm\,0.22$\\$[19.98,\ 20.54]$} & \makecell[c]{$20.88\,\pm\,0.14$\\$[20.71,\ 21.05]$} & \makecell[c]{$21.88\,\pm\,0.10$\\$[21.75,\ 22.00]$} & \makecell[c]{$107.07\,\pm\,29.19$\\$[70.82,\ 143.31]$} \\
Median & \makecell[c]{$102.31\,\pm\,0.40$\\$[101.81,\ 102.81]$} & \makecell[c]{$101.91\,\pm\,0.34$\\$[101.49,\ 102.34]$} & \makecell[c]{$101.92\,\pm\,0.13$\\$[101.76,\ 102.08]$} & \makecell[c]{$97.61\,\pm\,10.68$\\$[84.35,\ 110.87]$} \\
FFill + BFill & \makecell[c]{$24.71\,\pm\,0.31$\\$[24.32,\ 25.09]$} & \makecell[c]{$25.86\,\pm\,0.19$\\$[25.63,\ 26.09]$} & \makecell[c]{$27.92\,\pm\,0.17$\\$[27.71,\ 28.13]$} & \makecell[c]{$111.84\,\pm\,35.75$\\$[67.46,\ 156.23]$} \\
Rolling mean & \makecell[c]{$25.58\,\pm\,0.28$\\$[25.23,\ 25.92]$} & \makecell[c]{$26.34\,\pm\,0.14$\\$[26.16,\ 26.51]$} & \makecell[c]{$27.28\,\pm\,0.10$\\$[27.15,\ 27.41]$} & \makecell[c]{$107.24\,\pm\,29.28$\\$[70.88,\ 143.60]$} \\
KNN (scaled) & \makecell[c]{$36.81\,\pm\,0.23$\\$[36.52,\ 37.09]$} & \makecell[c]{$36.78\,\pm\,0.19$\\$[36.55,\ 37.02]$} & \makecell[c]{$36.87\,\pm\,0.10$\\$[36.75,\ 36.99]$} & \makecell[c]{$98.48\,\pm\,11.67$\\$[83.99,\ 112.97]$} \\
KNN (unscaled) & \makecell[c]{$40.56\,\pm\,0.35$\\$[40.12,\ 40.99]$} & \makecell[c]{$39.99\,\pm\,0.11$\\$[39.85,\ 40.13]$} & \makecell[c]{$39.49\,\pm\,0.06$\\$[39.41,\ 39.56]$} & \makecell[c]{$98.48\,\pm\,11.67$\\$[83.99,\ 112.97]$} \\
Historical & \makecell[c]{$54.69\,\pm\,0.15$\\$[54.50,\ 54.87]$} & \makecell[c]{$54.56\,\pm\,0.23$\\$[54.27,\ 54.85]$} & \makecell[c]{$54.47\,\pm\,0.10$\\$[54.34,\ 54.59]$} & \makecell[c]{$\mathbf{58.25}\,\pm\,15.81$\\$[38.63,\ 77.88]$} \\
Graph Laplacian & \makecell[c]{$43.81\,\pm\,0.57$\\$[43.11,\ 44.52]$} & \makecell[c]{$43.96\,\pm\,0.20$\\$[43.72,\ 44.21]$} & \makecell[c]{$44.34\,\pm\,0.20$\\$[44.09,\ 44.58]$} & \makecell[c]{$98.48\,\pm\,11.67$\\$[83.99,\ 112.97]$} \\
Graph Temporal & \makecell[c]{$34.05\,\pm\,0.39$\\$[33.56,\ 34.54]$} & \makecell[c]{$35.37\,\pm\,0.16$\\$[35.17,\ 35.57]$} & \makecell[c]{$37.00\,\pm\,0.05$\\$[36.93,\ 37.06]$} & \makecell[c]{$107.96\,\pm\,29.90$\\$[70.83,\ 145.09]$} \\
Sage-Voronoi & \makecell[c]{$22.04\,\pm\,0.14$\\$[21.88,\ 22.21]$} & \makecell[c]{$23.06\,\pm\,0.12$\\$[22.91,\ 23.21]$} & \makecell[c]{$24.70\,\pm\,0.23$\\$[24.42,\ 24.98]$} & \makecell[c]{$121.20\,\pm\,37.88$\\$[74.16,\ 168.23]$} \\
GAIN & \makecell[c]{$37.68\,\pm\,0.47$\\$[37.10,\ 38.27]$} & \makecell[c]{$41.58\,\pm\,0.27$\\$[41.24,\ 41.91]$} & \makecell[c]{$62.44\,\pm\,0.82$\\$[61.42,\ 63.46]$} & \makecell[c]{$205.60\,\pm\,16.50$\\$[185.11,\ 226.10]$} \\
GRIN & \makecell[c]{$\mathbf{16.92}\,\pm\,0.21$\\$[16.65,\ 17.18]$} & \makecell[c]{$\mathbf{17.86}\,\pm\,0.12$\\$[17.72,\ 18.01]$} & \makecell[c]{$\mathbf{19.50}\,\pm\,0.13$\\$[19.34,\ 19.66]$} & \makecell[c]{$86.52\,\pm\,18.84$\\$[63.12,\ 109.91]$} \\
\bottomrule
\end{tabular}
\endgroup
\end{table}

\clearpage

\begin{table}[htbp]
\centering
\caption{{RMSE for October--December 2022 (vehicles per ten-minute interval). Each cell gives mean $\pm$ sample SD, followed by the 95\% confidence interval; $n=5$ seed runs. Bold means identify the lowest unrounded mean within each column, not statistical significance.}}
\label{tab:supp_p01_17}
\begingroup
\footnotesize
\setlength{\tabcolsep}{3pt}
\renewcommand{\arraystretch}{1.15}
\begin{tabular}{lcccc}
\toprule
Method & Block $B=1$ & Block $B=2$ & Block $B=4$ & Block $B=8$ \\
\midrule
Linear & \makecell[c]{$\mathbf{48.70}\,\pm\,19.58$\\$[24.39,\ 73.02]$} & \makecell[c]{$64.12\,\pm\,25.26$\\$[32.75,\ 95.48]$} & \makecell[c]{$43.97\,\pm\,10.67$\\$[30.72,\ 57.23]$} & \makecell[c]{$81.48\,\pm\,15.10$\\$[62.73,\ 100.23]$} \\
Median & \makecell[c]{$190.11\,\pm\,33.16$\\$[148.93,\ 231.28]$} & \makecell[c]{$171.44\,\pm\,6.46$\\$[163.41,\ 179.46]$} & \makecell[c]{$163.38\,\pm\,10.42$\\$[150.44,\ 176.33]$} & \makecell[c]{$170.33\,\pm\,13.77$\\$[153.23,\ 187.42]$} \\
FFill + BFill & \makecell[c]{$57.25\,\pm\,8.04$\\$[47.27,\ 67.23]$} & \makecell[c]{$65.09\,\pm\,21.81$\\$[38.00,\ 92.17]$} & \makecell[c]{$59.25\,\pm\,14.69$\\$[41.00,\ 77.49]$} & \makecell[c]{$85.02\,\pm\,12.83$\\$[69.09,\ 100.95]$} \\
Rolling mean & \makecell[c]{$62.85\,\pm\,9.04$\\$[51.62,\ 74.07]$} & \makecell[c]{$90.44\,\pm\,18.07$\\$[68.00,\ 112.88]$} & \makecell[c]{$58.83\,\pm\,13.81$\\$[41.69,\ 75.97]$} & \makecell[c]{$80.48\,\pm\,9.52$\\$[68.66,\ 92.31]$} \\
KNN (scaled) & \makecell[c]{$135.85\,\pm\,30.64$\\$[97.80,\ 173.90]$} & \makecell[c]{$119.99\,\pm\,23.49$\\$[90.82,\ 149.16]$} & \makecell[c]{$111.06\,\pm\,29.17$\\$[74.84,\ 147.29]$} & \makecell[c]{$99.38\,\pm\,5.94$\\$[92.01,\ 106.76]$} \\
KNN (unscaled) & \makecell[c]{$147.28\,\pm\,44.29$\\$[92.29,\ 202.27]$} & \makecell[c]{$127.35\,\pm\,24.35$\\$[97.12,\ 157.58]$} & \makecell[c]{$119.27\,\pm\,29.31$\\$[82.87,\ 155.67]$} & \makecell[c]{$106.21\,\pm\,7.92$\\$[96.38,\ 116.05]$} \\
Historical & \makecell[c]{$123.40\,\pm\,30.73$\\$[85.25,\ 161.55]$} & \makecell[c]{$134.09\,\pm\,20.77$\\$[108.30,\ 159.88]$} & \makecell[c]{$127.28\,\pm\,28.96$\\$[91.31,\ 163.24]$} & \makecell[c]{$115.94\,\pm\,14.26$\\$[98.23,\ 133.65]$} \\
Graph Laplacian & \makecell[c]{$152.68\,\pm\,69.65$\\$[66.20,\ 239.16]$} & \makecell[c]{$113.05\,\pm\,33.33$\\$[71.67,\ 154.44]$} & \makecell[c]{$105.54\,\pm\,22.11$\\$[78.09,\ 132.99]$} & \makecell[c]{$104.37\,\pm\,12.71$\\$[88.58,\ 120.15]$} \\
Graph Temporal & \makecell[c]{$98.75\,\pm\,25.65$\\$[66.90,\ 130.60]$} & \makecell[c]{$97.39\,\pm\,32.22$\\$[57.38,\ 137.40]$} & \makecell[c]{$94.64\,\pm\,16.66$\\$[73.95,\ 115.32]$} & \makecell[c]{$96.85\,\pm\,5.39$\\$[90.16,\ 103.54]$} \\
Sage-Voronoi & \makecell[c]{$61.17\,\pm\,21.19$\\$[34.86,\ 87.48]$} & \makecell[c]{$70.71\,\pm\,22.51$\\$[42.77,\ 98.66]$} & \makecell[c]{$46.06\,\pm\,11.98$\\$[31.19,\ 60.94]$} & \makecell[c]{$64.42\,\pm\,10.05$\\$[51.94,\ 76.90]$} \\
GAIN & \makecell[c]{$149.38\,\pm\,58.47$\\$[76.77,\ 221.98]$} & \makecell[c]{$118.97\,\pm\,33.12$\\$[77.84,\ 160.10]$} & \makecell[c]{$99.89\,\pm\,23.60$\\$[70.58,\ 129.19]$} & \makecell[c]{$103.36\,\pm\,12.19$\\$[88.23,\ 118.50]$} \\
GRIN & \makecell[c]{$57.08\,\pm\,22.33$\\$[29.35,\ 84.81]$} & \makecell[c]{$\mathbf{61.47}\,\pm\,26.31$\\$[28.79,\ 94.14]$} & \makecell[c]{$\mathbf{39.72}\,\pm\,11.46$\\$[25.49,\ 53.94]$} & \makecell[c]{$\mathbf{51.10}\,\pm\,11.92$\\$[36.30,\ 65.91]$} \\
\bottomrule
\end{tabular}
\endgroup
\end{table}

\begin{table}[htbp]
\centering
\caption{{RMSE for October--December 2022 (vehicles per ten-minute interval). Each cell gives mean $\pm$ sample SD, followed by the 95\% confidence interval; $n=5$ seed runs. Bold means identify the lowest unrounded mean within each column, not statistical significance.}}
\label{tab:supp_p01_18}
\begingroup
\footnotesize
\setlength{\tabcolsep}{3pt}
\renewcommand{\arraystretch}{1.15}
\begin{tabular}{lccc}
\toprule
Method & Block $B=16$ & Block $B=32$ & Block $B=64$ \\
\midrule
Linear & \makecell[c]{$88.28\,\pm\,17.81$\\$[66.16,\ 110.39]$} & \makecell[c]{$103.15\,\pm\,30.10$\\$[65.78,\ 140.52]$} & \makecell[c]{$123.56\,\pm\,5.17$\\$[117.14,\ 129.99]$} \\
Median & \makecell[c]{$214.65\,\pm\,37.38$\\$[168.23,\ 261.06]$} & \makecell[c]{$171.87\,\pm\,9.89$\\$[159.59,\ 184.15]$} & \makecell[c]{$193.37\,\pm\,48.32$\\$[133.38,\ 253.37]$} \\
FFill + BFill & \makecell[c]{$115.93\,\pm\,17.78$\\$[93.86,\ 138.00]$} & \makecell[c]{$148.75\,\pm\,25.30$\\$[117.34,\ 180.17]$} & \makecell[c]{$192.64\,\pm\,16.65$\\$[171.97,\ 213.31]$} \\
Rolling mean & \makecell[c]{$90.55\,\pm\,15.04$\\$[71.88,\ 109.22]$} & \makecell[c]{$104.19\,\pm\,29.10$\\$[68.07,\ 140.32]$} & \makecell[c]{$123.50\,\pm\,5.63$\\$[116.50,\ 130.49]$} \\
KNN (scaled) & \makecell[c]{$148.81\,\pm\,36.38$\\$[103.64,\ 193.97]$} & \makecell[c]{$102.49\,\pm\,14.64$\\$[84.31,\ 120.67]$} & \makecell[c]{$141.63\,\pm\,61.66$\\$[65.07,\ 218.19]$} \\
KNN (unscaled) & \makecell[c]{$152.10\,\pm\,41.16$\\$[100.99,\ 203.20]$} & \makecell[c]{$107.12\,\pm\,17.46$\\$[85.43,\ 128.80]$} & \makecell[c]{$145.63\,\pm\,61.89$\\$[68.78,\ 222.48]$} \\
Historical & \makecell[c]{$147.75\,\pm\,16.78$\\$[126.91,\ 168.58]$} & \makecell[c]{$112.92\,\pm\,17.03$\\$[91.78,\ 134.07]$} & \makecell[c]{$141.88\,\pm\,55.98$\\$[72.37,\ 211.39]$} \\
Graph Laplacian & \makecell[c]{$184.46\,\pm\,60.45$\\$[109.40,\ 259.52]$} & \makecell[c]{$119.64\,\pm\,30.13$\\$[82.23,\ 157.05]$} & \makecell[c]{$148.28\,\pm\,59.60$\\$[74.27,\ 222.29]$} \\
Graph Temporal & \makecell[c]{$177.38\,\pm\,55.64$\\$[108.29,\ 246.47]$} & \makecell[c]{$117.54\,\pm\,31.61$\\$[78.29,\ 156.79]$} & \makecell[c]{$148.82\,\pm\,57.78$\\$[77.08,\ 220.56]$} \\
Sage-Voronoi & \makecell[c]{$\mathbf{80.56}\,\pm\,7.93$\\$[70.71,\ 90.40]$} & \makecell[c]{$98.89\,\pm\,35.02$\\$[55.41,\ 142.37]$} & \makecell[c]{$139.53\,\pm\,62.33$\\$[62.14,\ 216.92]$} \\
GAIN & \makecell[c]{$176.64\,\pm\,49.86$\\$[114.73,\ 238.55]$} & \makecell[c]{$108.84\,\pm\,20.77$\\$[83.05,\ 134.62]$} & \makecell[c]{$144.39\,\pm\,63.17$\\$[65.95,\ 222.83]$} \\
GRIN & \makecell[c]{$81.76\,\pm\,13.09$\\$[65.50,\ 98.02]$} & \makecell[c]{$\mathbf{80.37}\,\pm\,25.40$\\$[48.83,\ 111.91]$} & \makecell[c]{$\mathbf{119.40}\,\pm\,62.00$\\$[42.41,\ 196.38]$} \\
\bottomrule
\end{tabular}
\endgroup
\end{table}

\clearpage

\begin{table}[htbp]
\centering
\caption{{RMSE for October--December 2022 (vehicles per ten-minute interval). Each cell gives mean $\pm$ sample SD, followed by the 95\% confidence interval; $n=5$ seed runs. Bold means identify the lowest unrounded mean within each column, not statistical significance.}}
\label{tab:supp_p01_19}
\begingroup
\footnotesize
\setlength{\tabcolsep}{3pt}
\renewcommand{\arraystretch}{1.15}
\begin{tabular}{lcc}
\toprule
Method & Block $B=128$ & Block $B=144$ \\
\midrule
Linear & \makecell[c]{$205.33\,\pm\,14.53$\\$[187.29,\ 223.37]$} & \makecell[c]{$211.45\,\pm\,23.51$\\$[182.26,\ 240.64]$} \\
Median & \makecell[c]{$182.06\,\pm\,9.05$\\$[170.83,\ 193.30]$} & \makecell[c]{$178.53\,\pm\,15.52$\\$[159.26,\ 197.79]$} \\
FFill + BFill & \makecell[c]{$229.57\,\pm\,13.27$\\$[213.09,\ 246.05]$} & \makecell[c]{$213.10\,\pm\,22.83$\\$[184.74,\ 241.45]$} \\
Rolling mean & \makecell[c]{$205.44\,\pm\,14.63$\\$[187.28,\ 223.60]$} & \makecell[c]{$211.41\,\pm\,23.40$\\$[182.36,\ 240.46]$} \\
KNN (scaled) & \makecell[c]{$\mathbf{114.72}\,\pm\,4.15$\\$[109.58,\ 119.87]$} & \makecell[c]{$112.74\,\pm\,6.58$\\$[104.57,\ 120.92]$} \\
KNN (unscaled) & \makecell[c]{$121.32\,\pm\,6.58$\\$[113.16,\ 129.48]$} & \makecell[c]{$119.51\,\pm\,8.30$\\$[109.21,\ 129.81]$} \\
Historical & \makecell[c]{$119.85\,\pm\,6.95$\\$[111.22,\ 128.49]$} & \makecell[c]{$121.60\,\pm\,10.58$\\$[108.46,\ 134.74]$} \\
Graph Laplacian & \makecell[c]{$143.95\,\pm\,10.98$\\$[130.32,\ 157.58]$} & \makecell[c]{$134.49\,\pm\,22.30$\\$[106.80,\ 162.17]$} \\
Graph Temporal & \makecell[c]{$144.26\,\pm\,9.70$\\$[132.22,\ 156.31]$} & \makecell[c]{$133.66\,\pm\,21.55$\\$[106.90,\ 160.42]$} \\
Sage-Voronoi & \makecell[c]{$139.61\,\pm\,14.10$\\$[122.11,\ 157.11]$} & \makecell[c]{$129.11\,\pm\,18.89$\\$[105.66,\ 152.56]$} \\
GAIN & \makecell[c]{$131.87\,\pm\,13.45$\\$[115.16,\ 148.57]$} & \makecell[c]{$125.05\,\pm\,22.19$\\$[97.50,\ 152.60]$} \\
GRIN & \makecell[c]{$115.58\,\pm\,10.78$\\$[102.19,\ 128.97]$} & \makecell[c]{$\mathbf{103.73}\,\pm\,18.67$\\$[80.55,\ 126.91]$} \\
\bottomrule
\end{tabular}
\endgroup
\end{table}

\begin{table}[htbp]
\centering
\caption{{RMSE for October--December 2022 (vehicles per ten-minute interval). Each cell gives mean $\pm$ sample SD, followed by the 95\% confidence interval; $n=5$ seed runs. Bold means identify the lowest unrounded mean within each column, not statistical significance.}}
\label{tab:supp_p01_20}
\begingroup
\footnotesize
\setlength{\tabcolsep}{3pt}
\renewcommand{\arraystretch}{1.15}
\begin{tabular}{lcccc}
\toprule
Method & MCAR 10\% & MCAR 30\% & MCAR 50\% & Synchronous $B=144$ \\
\midrule
Linear & \makecell[c]{$57.36\,\pm\,1.52$\\$[55.47,\ 59.25]$} & \makecell[c]{$60.94\,\pm\,1.12$\\$[59.55,\ 62.34]$} & \makecell[c]{$66.26\,\pm\,0.47$\\$[65.68,\ 66.84]$} & \makecell[c]{$182.74\,\pm\,49.56$\\$[121.20,\ 244.29]$} \\
Median & \makecell[c]{$185.32\,\pm\,1.82$\\$[183.06,\ 187.59]$} & \makecell[c]{$183.28\,\pm\,0.95$\\$[182.10,\ 184.46]$} & \makecell[c]{$182.76\,\pm\,0.41$\\$[182.25,\ 183.27]$} & \makecell[c]{$160.36\,\pm\,15.33$\\$[141.33,\ 179.40]$} \\
FFill + BFill & \makecell[c]{$72.44\,\pm\,1.71$\\$[70.31,\ 74.57]$} & \makecell[c]{$77.74\,\pm\,1.59$\\$[75.76,\ 79.71]$} & \makecell[c]{$83.62\,\pm\,0.93$\\$[82.46,\ 84.78]$} & \makecell[c]{$194.65\,\pm\,55.75$\\$[125.43,\ 263.87]$} \\
Rolling mean & \makecell[c]{$74.51\,\pm\,1.05$\\$[73.20,\ 75.82]$} & \makecell[c]{$78.12\,\pm\,1.33$\\$[76.47,\ 79.76]$} & \makecell[c]{$81.68\,\pm\,0.72$\\$[80.79,\ 82.58]$} & \makecell[c]{$182.98\,\pm\,49.65$\\$[121.33,\ 244.63]$} \\
KNN (scaled) & \makecell[c]{$125.22\,\pm\,3.16$\\$[121.29,\ 129.15]$} & \makecell[c]{$123.63\,\pm\,1.28$\\$[122.05,\ 125.22]$} & \makecell[c]{$121.96\,\pm\,0.88$\\$[120.86,\ 123.05]$} & \makecell[c]{$159.65\,\pm\,15.04$\\$[140.98,\ 178.32]$} \\
KNN (unscaled) & \makecell[c]{$130.28\,\pm\,2.95$\\$[126.61,\ 133.94]$} & \makecell[c]{$128.46\,\pm\,0.92$\\$[127.32,\ 129.61]$} & \makecell[c]{$126.81\,\pm\,0.74$\\$[125.88,\ 127.73]$} & \makecell[c]{$159.65\,\pm\,15.04$\\$[140.98,\ 178.32]$} \\
Historical & \makecell[c]{$133.35\,\pm\,2.19$\\$[130.63,\ 136.06]$} & \makecell[c]{$132.80\,\pm\,1.33$\\$[131.15,\ 134.45]$} & \makecell[c]{$131.89\,\pm\,0.62$\\$[131.11,\ 132.66]$} & \makecell[c]{$\mathbf{124.80}\,\pm\,16.44$\\$[104.40,\ 145.21]$} \\
Graph Laplacian & \makecell[c]{$143.40\,\pm\,4.69$\\$[137.57,\ 149.22]$} & \makecell[c]{$141.51\,\pm\,0.84$\\$[140.47,\ 142.55]$} & \makecell[c]{$141.47\,\pm\,0.37$\\$[141.01,\ 141.92]$} & \makecell[c]{$159.65\,\pm\,15.04$\\$[140.98,\ 178.32]$} \\
Graph Temporal & \makecell[c]{$99.21\,\pm\,2.01$\\$[96.71,\ 101.71]$} & \makecell[c]{$104.23\,\pm\,1.62$\\$[102.22,\ 106.23]$} & \makecell[c]{$110.78\,\pm\,0.76$\\$[109.84,\ 111.72]$} & \makecell[c]{$189.22\,\pm\,49.38$\\$[127.90,\ 250.54]$} \\
Sage-Voronoi & \makecell[c]{$66.76\,\pm\,0.62$\\$[66.00,\ 67.53]$} & \makecell[c]{$71.12\,\pm\,1.10$\\$[69.76,\ 72.48]$} & \makecell[c]{$75.69\,\pm\,0.97$\\$[74.48,\ 76.90]$} & \makecell[c]{$214.66\,\pm\,56.69$\\$[144.27,\ 285.05]$} \\
GAIN & \makecell[c]{$132.67\,\pm\,5.47$\\$[125.88,\ 139.46]$} & \makecell[c]{$127.18\,\pm\,1.01$\\$[125.92,\ 128.43]$} & \makecell[c]{$141.50\,\pm\,0.62$\\$[140.73,\ 142.27]$} & \makecell[c]{$360.95\,\pm\,17.67$\\$[339.01,\ 382.90]$} \\
GRIN & \makecell[c]{$\mathbf{54.40}\,\pm\,1.41$\\$[52.65,\ 56.15]$} & \makecell[c]{$\mathbf{58.21}\,\pm\,1.23$\\$[56.68,\ 59.74]$} & \makecell[c]{$\mathbf{62.90}\,\pm\,0.76$\\$[61.96,\ 63.85]$} & \makecell[c]{$147.33\,\pm\,26.73$\\$[114.13,\ 180.52]$} \\
\bottomrule
\end{tabular}
\endgroup
\end{table}

\clearpage

\begin{table}[htbp]
\centering
\caption{{MRE for October--December 2022 (dimensionless ratio). Each cell gives mean $\pm$ sample SD, followed by the 95\% confidence interval; $n=5$ seed runs. Bold means identify the lowest unrounded mean within each column, not statistical significance.}}
\label{tab:supp_p01_21}
\begingroup
\footnotesize
\setlength{\tabcolsep}{3pt}
\renewcommand{\arraystretch}{1.15}
\begin{tabular}{lcccc}
\toprule
Method & Block $B=1$ & Block $B=2$ & Block $B=4$ & Block $B=8$ \\
\midrule
Linear & \makecell[c]{$0.226\,\pm\,0.037$\\$[0.181,\ 0.271]$} & \makecell[c]{$0.240\,\pm\,0.028$\\$[0.206,\ 0.274]$} & \makecell[c]{$0.244\,\pm\,0.031$\\$[0.206,\ 0.283]$} & \makecell[c]{$0.298\,\pm\,0.037$\\$[0.252,\ 0.343]$} \\
Median & \makecell[c]{$2.383\,\pm\,0.463$\\$[1.807,\ 2.958]$} & \makecell[c]{$2.405\,\pm\,0.422$\\$[1.881,\ 2.929]$} & \makecell[c]{$2.405\,\pm\,0.471$\\$[1.820,\ 2.990]$} & \makecell[c]{$2.243\,\pm\,0.256$\\$[1.926,\ 2.561]$} \\
FFill + BFill & \makecell[c]{$0.272\,\pm\,0.048$\\$[0.213,\ 0.332]$} & \makecell[c]{$0.286\,\pm\,0.029$\\$[0.250,\ 0.322]$} & \makecell[c]{$0.309\,\pm\,0.051$\\$[0.245,\ 0.373]$} & \makecell[c]{$0.382\,\pm\,0.044$\\$[0.327,\ 0.436]$} \\
Rolling mean & \makecell[c]{$0.258\,\pm\,0.022$\\$[0.230,\ 0.285]$} & \makecell[c]{$0.272\,\pm\,0.014$\\$[0.254,\ 0.290]$} & \makecell[c]{$0.303\,\pm\,0.045$\\$[0.247,\ 0.359]$} & \makecell[c]{$0.310\,\pm\,0.026$\\$[0.277,\ 0.342]$} \\
KNN (scaled) & \makecell[c]{$0.453\,\pm\,0.303$\\$[0.076,\ 0.829]$} & \makecell[c]{$0.493\,\pm\,0.222$\\$[0.216,\ 0.769]$} & \makecell[c]{$0.429\,\pm\,0.225$\\$[0.150,\ 0.709]$} & \makecell[c]{$0.341\,\pm\,0.083$\\$[0.238,\ 0.443]$} \\
KNN (unscaled) & \makecell[c]{$0.529\,\pm\,0.309$\\$[0.146,\ 0.912]$} & \makecell[c]{$0.558\,\pm\,0.193$\\$[0.318,\ 0.798]$} & \makecell[c]{$0.483\,\pm\,0.189$\\$[0.249,\ 0.718]$} & \makecell[c]{$0.416\,\pm\,0.073$\\$[0.326,\ 0.507]$} \\
Historical & \makecell[c]{$0.621\,\pm\,0.101$\\$[0.496,\ 0.746]$} & \makecell[c]{$0.746\,\pm\,0.133$\\$[0.582,\ 0.911]$} & \makecell[c]{$0.598\,\pm\,0.132$\\$[0.435,\ 0.762]$} & \makecell[c]{$0.645\,\pm\,0.050$\\$[0.583,\ 0.707]$} \\
Graph Laplacian & \makecell[c]{$0.578\,\pm\,0.129$\\$[0.418,\ 0.737]$} & \makecell[c]{$0.480\,\pm\,0.078$\\$[0.384,\ 0.576]$} & \makecell[c]{$0.524\,\pm\,0.065$\\$[0.444,\ 0.605]$} & \makecell[c]{$0.572\,\pm\,0.092$\\$[0.457,\ 0.686]$} \\
Graph Temporal & \makecell[c]{$0.471\,\pm\,0.088$\\$[0.361,\ 0.580]$} & \makecell[c]{$0.419\,\pm\,0.063$\\$[0.342,\ 0.497]$} & \makecell[c]{$0.488\,\pm\,0.034$\\$[0.446,\ 0.530]$} & \makecell[c]{$0.581\,\pm\,0.097$\\$[0.460,\ 0.701]$} \\
Sage-Voronoi & \makecell[c]{$0.226\,\pm\,0.033$\\$[0.186,\ 0.267]$} & \makecell[c]{$0.225\,\pm\,0.022$\\$[0.197,\ 0.253]$} & \makecell[c]{$0.255\,\pm\,0.017$\\$[0.234,\ 0.276]$} & \makecell[c]{$0.277\,\pm\,0.022$\\$[0.249,\ 0.305]$} \\
GAIN & \makecell[c]{$0.332\,\pm\,0.061$\\$[0.257,\ 0.408]$} & \makecell[c]{$0.319\,\pm\,0.034$\\$[0.277,\ 0.360]$} & \makecell[c]{$0.324\,\pm\,0.050$\\$[0.262,\ 0.385]$} & \makecell[c]{$0.337\,\pm\,0.066$\\$[0.256,\ 0.419]$} \\
GRIN & \makecell[c]{$\mathbf{0.182}\,\pm\,0.015$\\$[0.163,\ 0.201]$} & \makecell[c]{$\mathbf{0.184}\,\pm\,0.017$\\$[0.163,\ 0.204]$} & \makecell[c]{$\mathbf{0.190}\,\pm\,0.030$\\$[0.152,\ 0.227]$} & \makecell[c]{$\mathbf{0.232}\,\pm\,0.061$\\$[0.157,\ 0.308]$} \\
\bottomrule
\end{tabular}
\endgroup
\end{table}

\begin{table}[htbp]
\centering
\caption{{MRE for October--December 2022 (dimensionless ratio). Each cell gives mean $\pm$ sample SD, followed by the 95\% confidence interval; $n=5$ seed runs. Bold means identify the lowest unrounded mean within each column, not statistical significance.}}
\label{tab:supp_p01_22}
\begingroup
\footnotesize
\setlength{\tabcolsep}{3pt}
\renewcommand{\arraystretch}{1.15}
\begin{tabular}{lccc}
\toprule
Method & Block $B=16$ & Block $B=32$ & Block $B=64$ \\
\midrule
Linear & \makecell[c]{$0.373\,\pm\,0.052$\\$[0.308,\ 0.438]$} & \makecell[c]{$0.700\,\pm\,0.129$\\$[0.540,\ 0.860]$} & \makecell[c]{$1.282\,\pm\,0.121$\\$[1.131,\ 1.432]$} \\
Median & \makecell[c]{$2.531\,\pm\,0.525$\\$[1.880,\ 3.183]$} & \makecell[c]{$2.695\,\pm\,0.487$\\$[2.090,\ 3.299]$} & \makecell[c]{$2.311\,\pm\,0.110$\\$[2.175,\ 2.448]$} \\
FFill + BFill & \makecell[c]{$0.560\,\pm\,0.059$\\$[0.487,\ 0.633]$} & \makecell[c]{$1.068\,\pm\,0.102$\\$[0.941,\ 1.194]$} & \makecell[c]{$1.840\,\pm\,0.187$\\$[1.607,\ 2.073]$} \\
Rolling mean & \makecell[c]{$0.363\,\pm\,0.047$\\$[0.304,\ 0.421]$} & \makecell[c]{$0.685\,\pm\,0.108$\\$[0.551,\ 0.819]$} & \makecell[c]{$1.248\,\pm\,0.126$\\$[1.092,\ 1.404]$} \\
KNN (scaled) & \makecell[c]{$0.636\,\pm\,0.509$\\$[0.004,\ 1.268]$} & \makecell[c]{$0.628\,\pm\,0.348$\\$[0.197,\ 1.060]$} & \makecell[c]{$0.398\,\pm\,0.097$\\$[0.278,\ 0.518]$} \\
KNN (unscaled) & \makecell[c]{$0.796\,\pm\,0.663$\\$[-0.027,\ 1.618]$} & \makecell[c]{$0.692\,\pm\,0.369$\\$[0.233,\ 1.150]$} & \makecell[c]{$0.473\,\pm\,0.100$\\$[0.349,\ 0.598]$} \\
Historical & \makecell[c]{$0.852\,\pm\,0.318$\\$[0.457,\ 1.246]$} & \makecell[c]{$0.809\,\pm\,0.171$\\$[0.597,\ 1.021]$} & \makecell[c]{$0.674\,\pm\,0.119$\\$[0.526,\ 0.822]$} \\
Graph Laplacian & \makecell[c]{$0.544\,\pm\,0.165$\\$[0.340,\ 0.749]$} & \makecell[c]{$0.661\,\pm\,0.165$\\$[0.456,\ 0.865]$} & \makecell[c]{$0.583\,\pm\,0.108$\\$[0.448,\ 0.717]$} \\
Graph Temporal & \makecell[c]{$0.552\,\pm\,0.189$\\$[0.317,\ 0.787]$} & \makecell[c]{$0.684\,\pm\,0.187$\\$[0.451,\ 0.916]$} & \makecell[c]{$0.593\,\pm\,0.104$\\$[0.464,\ 0.722]$} \\
Sage-Voronoi & \makecell[c]{$0.294\,\pm\,0.026$\\$[0.261,\ 0.326]$} & \makecell[c]{$0.410\,\pm\,0.046$\\$[0.354,\ 0.467]$} & \makecell[c]{$0.416\,\pm\,0.080$\\$[0.318,\ 0.515]$} \\
GAIN & \makecell[c]{$0.415\,\pm\,0.166$\\$[0.210,\ 0.621]$} & \makecell[c]{$0.439\,\pm\,0.121$\\$[0.288,\ 0.589]$} & \makecell[c]{$0.362\,\pm\,0.097$\\$[0.241,\ 0.483]$} \\
GRIN & \makecell[c]{$\mathbf{0.237}\,\pm\,0.074$\\$[0.145,\ 0.329]$} & \makecell[c]{$\mathbf{0.350}\,\pm\,0.129$\\$[0.189,\ 0.510]$} & \makecell[c]{$\mathbf{0.335}\,\pm\,0.098$\\$[0.213,\ 0.457]$} \\
\bottomrule
\end{tabular}
\endgroup
\end{table}

\clearpage

\begin{table}[htbp]
\centering
\caption{{MRE for October--December 2022 (dimensionless ratio). Each cell gives mean $\pm$ sample SD, followed by the 95\% confidence interval; $n=5$ seed runs. Bold means identify the lowest unrounded mean within each column, not statistical significance.}}
\label{tab:supp_p01_23}
\begingroup
\footnotesize
\setlength{\tabcolsep}{3pt}
\renewcommand{\arraystretch}{1.15}
\begin{tabular}{lcc}
\toprule
Method & Block $B=128$ & Block $B=144$ \\
\midrule
Linear & \makecell[c]{$2.356\,\pm\,0.113$\\$[2.216,\ 2.496]$} & \makecell[c]{$2.552\,\pm\,0.112$\\$[2.413,\ 2.691]$} \\
Median & \makecell[c]{$2.399\,\pm\,0.103$\\$[2.270,\ 2.527]$} & \makecell[c]{$2.449\,\pm\,0.054$\\$[2.382,\ 2.516]$} \\
FFill + BFill & \makecell[c]{$2.561\,\pm\,0.050$\\$[2.499,\ 2.623]$} & \makecell[c]{$2.499\,\pm\,0.118$\\$[2.351,\ 2.646]$} \\
Rolling mean & \makecell[c]{$2.349\,\pm\,0.110$\\$[2.212,\ 2.486]$} & \makecell[c]{$2.551\,\pm\,0.112$\\$[2.411,\ 2.690]$} \\
KNN (scaled) & \makecell[c]{$0.516\,\pm\,0.158$\\$[0.319,\ 0.713]$} & \makecell[c]{$0.491\,\pm\,0.125$\\$[0.335,\ 0.646]$} \\
KNN (unscaled) & \makecell[c]{$0.591\,\pm\,0.174$\\$[0.376,\ 0.806]$} & \makecell[c]{$0.588\,\pm\,0.166$\\$[0.382,\ 0.793]$} \\
Historical & \makecell[c]{$0.735\,\pm\,0.077$\\$[0.640,\ 0.830]$} & \makecell[c]{$0.749\,\pm\,0.056$\\$[0.680,\ 0.818]$} \\
Graph Laplacian & \makecell[c]{$0.531\,\pm\,0.059$\\$[0.457,\ 0.605]$} & \makecell[c]{$0.588\,\pm\,0.069$\\$[0.503,\ 0.674]$} \\
Graph Temporal & \makecell[c]{$0.543\,\pm\,0.052$\\$[0.478,\ 0.608]$} & \makecell[c]{$0.620\,\pm\,0.089$\\$[0.509,\ 0.731]$} \\
Sage-Voronoi & \makecell[c]{$0.402\,\pm\,0.048$\\$[0.342,\ 0.462]$} & \makecell[c]{$0.440\,\pm\,0.040$\\$[0.390,\ 0.491]$} \\
GAIN & \makecell[c]{$0.353\,\pm\,0.062$\\$[0.276,\ 0.430]$} & \makecell[c]{$0.441\,\pm\,0.189$\\$[0.206,\ 0.676]$} \\
GRIN & \makecell[c]{$\mathbf{0.299}\,\pm\,0.057$\\$[0.228,\ 0.371]$} & \makecell[c]{$\mathbf{0.339}\,\pm\,0.058$\\$[0.267,\ 0.411]$} \\
\bottomrule
\end{tabular}
\endgroup
\end{table}

\begin{table}[htbp]
\centering
\caption{{MRE for October--December 2022 (dimensionless ratio). Each cell gives mean $\pm$ sample SD, followed by the 95\% confidence interval; $n=5$ seed runs. Bold means identify the lowest unrounded mean within each column, not statistical significance.}}
\label{tab:supp_p01_24}
\begingroup
\footnotesize
\setlength{\tabcolsep}{3pt}
\renewcommand{\arraystretch}{1.15}
\begin{tabular}{lcccc}
\toprule
Method & MCAR 10\% & MCAR 30\% & MCAR 50\% & Synchronous $B=144$ \\
\midrule
Linear & \makecell[c]{$0.252\,\pm\,0.003$\\$[0.249,\ 0.256]$} & \makecell[c]{$0.254\,\pm\,0.001$\\$[0.253,\ 0.256]$} & \makecell[c]{$0.263\,\pm\,0.002$\\$[0.261,\ 0.265]$} & \makecell[c]{$1.794\,\pm\,1.318$\\$[0.158,\ 3.430]$} \\
Median & \makecell[c]{$2.408\,\pm\,0.022$\\$[2.381,\ 2.435]$} & \makecell[c]{$2.386\,\pm\,0.012$\\$[2.370,\ 2.401]$} & \makecell[c]{$2.387\,\pm\,0.012$\\$[2.372,\ 2.402]$} & \makecell[c]{$2.003\,\pm\,0.772$\\$[1.044,\ 2.962]$} \\
FFill + BFill & \makecell[c]{$0.298\,\pm\,0.005$\\$[0.292,\ 0.304]$} & \makecell[c]{$0.301\,\pm\,0.002$\\$[0.298,\ 0.303]$} & \makecell[c]{$0.318\,\pm\,0.003$\\$[0.315,\ 0.322]$} & \makecell[c]{$1.857\,\pm\,1.390$\\$[0.131,\ 3.583]$} \\
Rolling mean & \makecell[c]{$0.290\,\pm\,0.004$\\$[0.285,\ 0.295]$} & \makecell[c]{$0.294\,\pm\,0.002$\\$[0.291,\ 0.297]$} & \makecell[c]{$0.306\,\pm\,0.002$\\$[0.303,\ 0.308]$} & \makecell[c]{$1.797\,\pm\,1.319$\\$[0.159,\ 3.434]$} \\
KNN (scaled) & \makecell[c]{$0.460\,\pm\,0.022$\\$[0.433,\ 0.487]$} & \makecell[c]{$0.457\,\pm\,0.008$\\$[0.447,\ 0.467]$} & \makecell[c]{$0.444\,\pm\,0.003$\\$[0.440,\ 0.447]$} & \makecell[c]{$2.032\,\pm\,0.780$\\$[1.063,\ 3.001]$} \\
KNN (unscaled) & \makecell[c]{$0.527\,\pm\,0.022$\\$[0.500,\ 0.554]$} & \makecell[c]{$0.519\,\pm\,0.014$\\$[0.501,\ 0.536]$} & \makecell[c]{$0.492\,\pm\,0.008$\\$[0.482,\ 0.503]$} & \makecell[c]{$2.032\,\pm\,0.780$\\$[1.063,\ 3.001]$} \\
Historical & \makecell[c]{$0.714\,\pm\,0.013$\\$[0.698,\ 0.730]$} & \makecell[c]{$0.714\,\pm\,0.006$\\$[0.707,\ 0.722]$} & \makecell[c]{$0.714\,\pm\,0.005$\\$[0.708,\ 0.719]$} & \makecell[c]{$\mathbf{0.825}\,\pm\,0.480$\\$[0.230,\ 1.421]$} \\
Graph Laplacian & \makecell[c]{$0.551\,\pm\,0.008$\\$[0.542,\ 0.561]$} & \makecell[c]{$0.549\,\pm\,0.005$\\$[0.543,\ 0.555]$} & \makecell[c]{$0.551\,\pm\,0.004$\\$[0.546,\ 0.556]$} & \makecell[c]{$2.032\,\pm\,0.780$\\$[1.063,\ 3.001]$} \\
Graph Temporal & \makecell[c]{$0.448\,\pm\,0.008$\\$[0.438,\ 0.458]$} & \makecell[c]{$0.457\,\pm\,0.003$\\$[0.454,\ 0.461]$} & \makecell[c]{$0.473\,\pm\,0.003$\\$[0.469,\ 0.477]$} & \makecell[c]{$1.852\,\pm\,1.321$\\$[0.212,\ 3.492]$} \\
Sage-Voronoi & \makecell[c]{$0.252\,\pm\,0.004$\\$[0.247,\ 0.257]$} & \makecell[c]{$0.255\,\pm\,0.005$\\$[0.249,\ 0.261]$} & \makecell[c]{$0.270\,\pm\,0.005$\\$[0.263,\ 0.276]$} & \makecell[c]{$1.477\,\pm\,0.982$\\$[0.257,\ 2.697]$} \\
GAIN & \makecell[c]{$0.353\,\pm\,0.011$\\$[0.339,\ 0.367]$} & \makecell[c]{$0.547\,\pm\,0.019$\\$[0.523,\ 0.571]$} & \makecell[c]{$1.109\,\pm\,0.033$\\$[1.068,\ 1.150]$} & \makecell[c]{$5.137\,\pm\,1.664$\\$[3.071,\ 7.203]$} \\
GRIN & \makecell[c]{$\mathbf{0.202}\,\pm\,0.002$\\$[0.199,\ 0.204]$} & \makecell[c]{$\mathbf{0.212}\,\pm\,0.002$\\$[0.209,\ 0.215]$} & \makecell[c]{$\mathbf{0.238}\,\pm\,0.005$\\$[0.232,\ 0.245]$} & \makecell[c]{$1.631\,\pm\,0.801$\\$[0.636,\ 2.625]$} \\
\bottomrule
\end{tabular}
\endgroup
\end{table}

\clearpage

\begin{table}[htbp]
\centering
\caption{{MAE for sensor A023, January 2023 (vehicles per ten-minute interval). Each cell gives mean $\pm$ sample SD, followed by the 95\% confidence interval; $n=5$ seed runs. Bold means identify the lowest unrounded mean within each column, not statistical significance.}}
\label{tab:supp_p01_25}
\begingroup
\footnotesize
\setlength{\tabcolsep}{3pt}
\renewcommand{\arraystretch}{1.15}
\begin{tabular}{lcccc}
\toprule
Method & Block $B=1$ & Block $B=2$ & Block $B=4$ & Block $B=8$ \\
\midrule
Linear & \makecell[c]{$\mathbf{9.50}\,\pm\,12.81$\\$[-6.41,\ 25.41]$} & \makecell[c]{$29.70\,\pm\,18.44$\\$[6.80,\ 52.60]$} & \makecell[c]{$29.85\,\pm\,17.04$\\$[8.69,\ 51.01]$} & \makecell[c]{$27.06\,\pm\,16.68$\\$[6.35,\ 47.76]$} \\
Median & \makecell[c]{$86.70\,\pm\,59.96$\\$[12.25,\ 161.15]$} & \makecell[c]{$129.10\,\pm\,71.94$\\$[39.77,\ 218.43]$} & \makecell[c]{$122.30\,\pm\,69.94$\\$[35.45,\ 209.15]$} & \makecell[c]{$129.53\,\pm\,54.40$\\$[61.97,\ 197.08]$} \\
FFill + BFill & \makecell[c]{$17.60\,\pm\,18.24$\\$[-5.05,\ 40.25]$} & \makecell[c]{$35.90\,\pm\,28.29$\\$[0.78,\ 71.02]$} & \makecell[c]{$42.45\,\pm\,25.54$\\$[10.74,\ 74.16]$} & \makecell[c]{$59.08\,\pm\,35.84$\\$[14.58,\ 103.57]$} \\
Rolling mean & \makecell[c]{$28.95\,\pm\,21.35$\\$[2.44,\ 55.46]$} & \makecell[c]{$23.75\,\pm\,20.62$\\$[-1.85,\ 49.35]$} & \makecell[c]{$37.80\,\pm\,20.39$\\$[12.48,\ 63.12]$} & \makecell[c]{$41.88\,\pm\,24.47$\\$[11.49,\ 72.27]$} \\
KNN (scaled) & \makecell[c]{$122.14\,\pm\,70.14$\\$[35.05,\ 209.23]$} & \makecell[c]{$126.73\,\pm\,74.56$\\$[34.15,\ 219.31]$} & \makecell[c]{$60.38\,\pm\,33.68$\\$[18.56,\ 102.20]$} & \makecell[c]{$116.92\,\pm\,73.41$\\$[25.76,\ 208.07]$} \\
KNN (unscaled) & \makecell[c]{$23.12\,\pm\,6.43$\\$[15.13,\ 31.11]$} & \makecell[c]{$21.52\,\pm\,6.98$\\$[12.86,\ 30.18]$} & \makecell[c]{$17.41\,\pm\,4.65$\\$[11.63,\ 23.19]$} & \makecell[c]{$\mathbf{21.20}\,\pm\,5.65$\\$[14.18,\ 28.21]$} \\
Historical & \makecell[c]{$45.87\,\pm\,20.58$\\$[20.33,\ 71.42]$} & \makecell[c]{$29.36\,\pm\,43.43$\\$[-24.57,\ 83.28]$} & \makecell[c]{$53.68\,\pm\,52.04$\\$[-10.94,\ 118.29]$} & \makecell[c]{$59.90\,\pm\,41.33$\\$[8.58,\ 111.22]$} \\
Graph Laplacian & \makecell[c]{$45.59\,\pm\,32.46$\\$[5.29,\ 85.89]$} & \makecell[c]{$58.69\,\pm\,34.09$\\$[16.36,\ 101.02]$} & \makecell[c]{$27.78\,\pm\,14.99$\\$[9.16,\ 46.40]$} & \makecell[c]{$48.25\,\pm\,25.53$\\$[16.55,\ 79.94]$} \\
Graph Temporal & \makecell[c]{$19.84\,\pm\,21.43$\\$[-6.77,\ 46.44]$} & \makecell[c]{$30.41\,\pm\,16.66$\\$[9.73,\ 51.10]$} & \makecell[c]{$18.04\,\pm\,9.95$\\$[5.68,\ 30.39]$} & \makecell[c]{$32.46\,\pm\,18.74$\\$[9.19,\ 55.73]$} \\
Sage-Voronoi & \makecell[c]{$14.30\,\pm\,8.44$\\$[3.82,\ 24.78]$} & \makecell[c]{$18.91\,\pm\,11.52$\\$[4.61,\ 33.22]$} & \makecell[c]{$19.98\,\pm\,11.08$\\$[6.23,\ 33.73]$} & \makecell[c]{$30.16\,\pm\,16.90$\\$[9.17,\ 51.14]$} \\
GAIN & \makecell[c]{$15.09\,\pm\,16.95$\\$[-5.96,\ 36.13]$} & \makecell[c]{$24.97\,\pm\,15.22$\\$[6.07,\ 43.87]$} & \makecell[c]{$11.88\,\pm\,2.73$\\$[8.49,\ 15.27]$} & \makecell[c]{$26.08\,\pm\,11.69$\\$[11.56,\ 40.60]$} \\
GRIN & \makecell[c]{$10.05\,\pm\,8.73$\\$[-0.79,\ 20.89]$} & \makecell[c]{$\mathbf{18.90}\,\pm\,5.05$\\$[12.63,\ 25.17]$} & \makecell[c]{$\mathbf{11.58}\,\pm\,3.81$\\$[6.85,\ 16.31]$} & \makecell[c]{$22.07\,\pm\,8.18$\\$[11.90,\ 32.23]$} \\
\bottomrule
\end{tabular}
\endgroup
\end{table}

\begin{table}[htbp]
\centering
\caption{{MAE for sensor A023, January 2023 (vehicles per ten-minute interval). Each cell gives mean $\pm$ sample SD, followed by the 95\% confidence interval; $n=5$ seed runs. Bold means identify the lowest unrounded mean within each column, not statistical significance.}}
\label{tab:supp_p01_26}
\begingroup
\footnotesize
\setlength{\tabcolsep}{3pt}
\renewcommand{\arraystretch}{1.15}
\begin{tabular}{lccc}
\toprule
Method & Block $B=16$ & Block $B=32$ & Block $B=64$ \\
\midrule
Linear & \makecell[c]{$29.37\,\pm\,16.24$\\$[9.21,\ 49.53]$} & \makecell[c]{$27.57\,\pm\,12.94$\\$[11.50,\ 43.63]$} & \makecell[c]{$92.53\,\pm\,68.66$\\$[7.27,\ 177.78]$} \\
Median & \makecell[c]{$126.67\,\pm\,36.86$\\$[80.91,\ 172.44]$} & \makecell[c]{$122.70\,\pm\,6.91$\\$[114.12,\ 131.28]$} & \makecell[c]{$109.46\,\pm\,48.15$\\$[49.67,\ 169.25]$} \\
FFill + BFill & \makecell[c]{$58.96\,\pm\,42.35$\\$[6.38,\ 111.55]$} & \makecell[c]{$50.11\,\pm\,31.77$\\$[10.66,\ 89.56]$} & \makecell[c]{$102.43\,\pm\,56.51$\\$[32.27,\ 172.60]$} \\
Rolling mean & \makecell[c]{$19.89\,\pm\,7.61$\\$[10.44,\ 29.34]$} & \makecell[c]{$28.08\,\pm\,11.50$\\$[13.81,\ 42.36]$} & \makecell[c]{$90.26\,\pm\,66.87$\\$[7.24,\ 173.29]$} \\
KNN (scaled) & \makecell[c]{$79.97\,\pm\,49.21$\\$[18.87,\ 141.07]$} & \makecell[c]{$160.32\,\pm\,62.91$\\$[82.20,\ 238.44]$} & \makecell[c]{$121.84\,\pm\,45.88$\\$[64.87,\ 178.81]$} \\
KNN (unscaled) & \makecell[c]{$23.73\,\pm\,8.99$\\$[12.57,\ 34.89]$} & \makecell[c]{$\mathbf{24.12}\,\pm\,3.80$\\$[19.40,\ 28.83]$} & \makecell[c]{$24.36\,\pm\,2.65$\\$[21.07,\ 27.65]$} \\
Historical & \makecell[c]{$31.11\,\pm\,11.29$\\$[17.09,\ 45.13]$} & \makecell[c]{$62.08\,\pm\,23.73$\\$[32.61,\ 91.55]$} & \makecell[c]{$48.88\,\pm\,16.23$\\$[28.73,\ 69.02]$} \\
Graph Laplacian & \makecell[c]{$46.51\,\pm\,19.99$\\$[21.69,\ 71.33]$} & \makecell[c]{$43.12\,\pm\,12.15$\\$[28.03,\ 58.22]$} & \makecell[c]{$43.92\,\pm\,11.97$\\$[29.06,\ 58.78]$} \\
Graph Temporal & \makecell[c]{$30.40\,\pm\,9.28$\\$[18.88,\ 41.92]$} & \makecell[c]{$24.39\,\pm\,3.36$\\$[20.21,\ 28.56]$} & \makecell[c]{$30.10\,\pm\,3.38$\\$[25.90,\ 34.29]$} \\
Sage-Voronoi & \makecell[c]{$28.91\,\pm\,10.83$\\$[15.46,\ 42.36]$} & \makecell[c]{$31.02\,\pm\,4.74$\\$[25.13,\ 36.91]$} & \makecell[c]{$29.29\,\pm\,8.87$\\$[18.28,\ 40.31]$} \\
GAIN & \makecell[c]{$20.05\,\pm\,9.32$\\$[8.48,\ 31.62]$} & \makecell[c]{$24.24\,\pm\,6.35$\\$[16.35,\ 32.13]$} & \makecell[c]{$24.32\,\pm\,8.90$\\$[13.27,\ 35.38]$} \\
GRIN & \makecell[c]{$\mathbf{16.68}\,\pm\,5.55$\\$[9.79,\ 23.57]$} & \makecell[c]{$24.81\,\pm\,6.74$\\$[16.44,\ 33.17]$} & \makecell[c]{$\mathbf{20.18}\,\pm\,5.75$\\$[13.05,\ 27.32]$} \\
\bottomrule
\end{tabular}
\endgroup
\end{table}

\clearpage

\begin{table}[htbp]
\centering
\caption{{MAE for sensor A023, January 2023 (vehicles per ten-minute interval). Each cell gives mean $\pm$ sample SD, followed by the 95\% confidence interval; $n=5$ seed runs. Bold means identify the lowest unrounded mean within each column, not statistical significance.}}
\label{tab:supp_p01_27}
\begingroup
\footnotesize
\setlength{\tabcolsep}{3pt}
\renewcommand{\arraystretch}{1.15}
\begin{tabular}{lcc}
\toprule
Method & Block $B=128$ & Block $B=144$ \\
\midrule
Linear & \makecell[c]{$119.08\,\pm\,41.15$\\$[67.99,\ 170.17]$} & \makecell[c]{$132.96\,\pm\,19.36$\\$[108.92,\ 157.00]$} \\
Median & \makecell[c]{$108.20\,\pm\,18.21$\\$[85.59,\ 130.81]$} & \makecell[c]{$120.96\,\pm\,3.80$\\$[116.24,\ 125.67]$} \\
FFill + BFill & \makecell[c]{$141.16\,\pm\,59.97$\\$[66.69,\ 215.63]$} & \makecell[c]{$130.77\,\pm\,14.73$\\$[112.48,\ 149.06]$} \\
Rolling mean & \makecell[c]{$118.22\,\pm\,42.29$\\$[65.71,\ 170.74]$} & \makecell[c]{$132.90\,\pm\,19.10$\\$[109.19,\ 156.61]$} \\
KNN (scaled) & \makecell[c]{$100.37\,\pm\,21.52$\\$[73.65,\ 127.09]$} & \makecell[c]{$110.16\,\pm\,27.75$\\$[75.71,\ 144.61]$} \\
KNN (unscaled) & \makecell[c]{$24.21\,\pm\,2.86$\\$[20.66,\ 27.76]$} & \makecell[c]{$25.92\,\pm\,3.03$\\$[22.16,\ 29.68]$} \\
Historical & \makecell[c]{$46.06\,\pm\,3.89$\\$[41.24,\ 50.89]$} & \makecell[c]{$45.33\,\pm\,2.81$\\$[41.84,\ 48.82]$} \\
Graph Laplacian & \makecell[c]{$34.98\,\pm\,6.38$\\$[27.06,\ 42.90]$} & \makecell[c]{$43.96\,\pm\,7.67$\\$[34.43,\ 53.49]$} \\
Graph Temporal & \makecell[c]{$24.43\,\pm\,3.90$\\$[19.59,\ 29.26]$} & \makecell[c]{$30.14\,\pm\,1.15$\\$[28.70,\ 31.57]$} \\
Sage-Voronoi & \makecell[c]{$26.14\,\pm\,5.25$\\$[19.62,\ 32.65]$} & \makecell[c]{$37.67\,\pm\,5.62$\\$[30.69,\ 44.64]$} \\
GAIN & \makecell[c]{$22.36\,\pm\,3.47$\\$[18.05,\ 26.67]$} & \makecell[c]{$23.04\,\pm\,2.55$\\$[19.88,\ 26.21]$} \\
GRIN & \makecell[c]{$\mathbf{21.06}\,\pm\,1.92$\\$[18.68,\ 23.44]$} & \makecell[c]{$\mathbf{22.10}\,\pm\,1.35$\\$[20.43,\ 23.77]$} \\
\bottomrule
\end{tabular}
\endgroup
\end{table}

\begin{table}[htbp]
\centering
\caption{{MAE for sensor A023, January 2023 (vehicles per ten-minute interval). Each cell gives mean $\pm$ sample SD, followed by the 95\% confidence interval; $n=5$ seed runs. Bold means identify the lowest unrounded mean within each column, not statistical significance.}}
\label{tab:supp_p01_28}
\begingroup
\footnotesize
\setlength{\tabcolsep}{3pt}
\renewcommand{\arraystretch}{1.15}
\begin{tabular}{lcccc}
\toprule
Method & MCAR 10\% & MCAR 30\% & MCAR 50\% & Synchronous $B=144$ \\
\midrule
Linear & \makecell[c]{$21.43\,\pm\,1.58$\\$[19.46,\ 23.39]$} & \makecell[c]{$21.50\,\pm\,1.41$\\$[19.76,\ 23.25]$} & \makecell[c]{$21.74\,\pm\,0.70$\\$[20.87,\ 22.60]$} & \makecell[c]{$126.44\,\pm\,10.34$\\$[113.61,\ 139.27]$} \\
Median & \makecell[c]{$114.44\,\pm\,3.81$\\$[109.71,\ 119.18]$} & \makecell[c]{$112.57\,\pm\,2.85$\\$[109.03,\ 116.12]$} & \makecell[c]{$111.77\,\pm\,1.64$\\$[109.73,\ 113.81]$} & \makecell[c]{$117.44\,\pm\,7.52$\\$[108.10,\ 126.78]$} \\
FFill + BFill & \makecell[c]{$25.13\,\pm\,2.26$\\$[22.33,\ 27.93]$} & \makecell[c]{$25.97\,\pm\,0.76$\\$[25.02,\ 26.92]$} & \makecell[c]{$28.74\,\pm\,0.63$\\$[27.96,\ 29.53]$} & \makecell[c]{$135.15\,\pm\,21.97$\\$[107.88,\ 162.43]$} \\
Rolling mean & \makecell[c]{$24.86\,\pm\,2.27$\\$[22.04,\ 27.68]$} & \makecell[c]{$25.86\,\pm\,1.21$\\$[24.36,\ 27.36]$} & \makecell[c]{$27.38\,\pm\,1.00$\\$[26.14,\ 28.62]$} & \makecell[c]{$126.66\,\pm\,10.65$\\$[113.43,\ 139.88]$} \\
KNN (scaled) & \makecell[c]{$108.96\,\pm\,7.18$\\$[100.04,\ 117.87]$} & \makecell[c]{$101.72\,\pm\,2.60$\\$[98.49,\ 104.95]$} & \makecell[c]{$90.09\,\pm\,3.46$\\$[85.79,\ 94.38]$} & \makecell[c]{$121.75\,\pm\,13.61$\\$[104.86,\ 138.65]$} \\
KNN (unscaled) & \makecell[c]{$23.81\,\pm\,2.66$\\$[20.51,\ 27.11]$} & \makecell[c]{$26.88\,\pm\,0.61$\\$[26.13,\ 27.64]$} & \makecell[c]{$27.55\,\pm\,1.15$\\$[26.13,\ 28.98]$} & \makecell[c]{$121.75\,\pm\,13.61$\\$[104.86,\ 138.65]$} \\
Historical & \makecell[c]{$45.48\,\pm\,2.90$\\$[41.88,\ 49.08]$} & \makecell[c]{$49.09\,\pm\,1.54$\\$[47.18,\ 51.00]$} & \makecell[c]{$47.91\,\pm\,0.76$\\$[46.97,\ 48.85]$} & \makecell[c]{$\mathbf{48.22}\,\pm\,8.25$\\$[37.97,\ 58.46]$} \\
Graph Laplacian & \makecell[c]{$43.48\,\pm\,2.65$\\$[40.18,\ 46.77]$} & \makecell[c]{$45.79\,\pm\,1.43$\\$[44.01,\ 47.56]$} & \makecell[c]{$47.59\,\pm\,1.38$\\$[45.87,\ 49.31]$} & \makecell[c]{$121.75\,\pm\,13.61$\\$[104.86,\ 138.65]$} \\
Graph Temporal & \makecell[c]{$25.64\,\pm\,0.75$\\$[24.72,\ 26.57]$} & \makecell[c]{$27.46\,\pm\,1.02$\\$[26.18,\ 28.73]$} & \makecell[c]{$27.59\,\pm\,0.62$\\$[26.82,\ 28.36]$} & \makecell[c]{$159.28\,\pm\,53.93$\\$[92.31,\ 226.24]$} \\
Sage-Voronoi & \makecell[c]{$20.95\,\pm\,1.87$\\$[18.63,\ 23.27]$} & \makecell[c]{$22.74\,\pm\,2.00$\\$[20.25,\ 25.22]$} & \makecell[c]{$23.89\,\pm\,1.47$\\$[22.06,\ 25.72]$} & \makecell[c]{$119.28\,\pm\,31.27$\\$[80.45,\ 158.11]$} \\
GAIN & \makecell[c]{$22.61\,\pm\,1.89$\\$[20.25,\ 24.96]$} & \makecell[c]{$27.96\,\pm\,2.40$\\$[24.98,\ 30.94]$} & \makecell[c]{$41.97\,\pm\,3.84$\\$[37.20,\ 46.73]$} & \makecell[c]{$121.41\,\pm\,12.95$\\$[105.33,\ 137.48]$} \\
GRIN & \makecell[c]{$\mathbf{17.45}\,\pm\,0.28$\\$[17.10,\ 17.80]$} & \makecell[c]{$\mathbf{18.54}\,\pm\,1.12$\\$[17.14,\ 19.93]$} & \makecell[c]{$\mathbf{19.98}\,\pm\,1.18$\\$[18.51,\ 21.44]$} & \makecell[c]{$132.16\,\pm\,25.48$\\$[100.52,\ 163.80]$} \\
\bottomrule
\end{tabular}
\endgroup
\end{table}

\clearpage

\begin{table}[htbp]
\centering
\caption{{RMSE for sensor A023, January 2023 (vehicles per ten-minute interval). Each cell gives mean $\pm$ sample SD, followed by the 95\% confidence interval; $n=5$ seed runs. Bold means identify the lowest unrounded mean within each column, not statistical significance.}}
\label{tab:supp_p01_29}
\begingroup
\footnotesize
\setlength{\tabcolsep}{3pt}
\renewcommand{\arraystretch}{1.15}
\begin{tabular}{lcccc}
\toprule
Method & Block $B=1$ & Block $B=2$ & Block $B=4$ & Block $B=8$ \\
\midrule
Linear & \makecell[c]{$\mathbf{9.50}\,\pm\,12.81$\\$[-6.41,\ 25.41]$} & \makecell[c]{$32.82\,\pm\,19.60$\\$[8.48,\ 57.16]$} & \makecell[c]{$33.28\,\pm\,16.60$\\$[12.67,\ 53.89]$} & \makecell[c]{$32.30\,\pm\,17.74$\\$[10.26,\ 54.33]$} \\
Median & \makecell[c]{$86.70\,\pm\,59.96$\\$[12.25,\ 161.15]$} & \makecell[c]{$129.65\,\pm\,71.95$\\$[40.31,\ 218.99]$} & \makecell[c]{$123.41\,\pm\,69.61$\\$[36.98,\ 209.84]$} & \makecell[c]{$133.26\,\pm\,51.40$\\$[69.43,\ 197.09]$} \\
FFill + BFill & \makecell[c]{$17.60\,\pm\,18.24$\\$[-5.05,\ 40.25]$} & \makecell[c]{$38.47\,\pm\,28.23$\\$[3.42,\ 73.53]$} & \makecell[c]{$45.21\,\pm\,26.12$\\$[12.78,\ 77.65]$} & \makecell[c]{$65.28\,\pm\,37.60$\\$[18.60,\ 111.97]$} \\
Rolling mean & \makecell[c]{$28.95\,\pm\,21.35$\\$[2.44,\ 55.46]$} & \makecell[c]{$26.02\,\pm\,22.16$\\$[-1.50,\ 53.54]$} & \makecell[c]{$41.46\,\pm\,20.96$\\$[15.44,\ 67.48]$} & \makecell[c]{$48.04\,\pm\,24.43$\\$[17.70,\ 78.38]$} \\
KNN (scaled) & \makecell[c]{$122.14\,\pm\,70.14$\\$[35.05,\ 209.23]$} & \makecell[c]{$128.99\,\pm\,72.94$\\$[38.42,\ 219.56]$} & \makecell[c]{$66.36\,\pm\,33.09$\\$[25.27,\ 107.45]$} & \makecell[c]{$121.72\,\pm\,73.05$\\$[31.02,\ 212.42]$} \\
KNN (unscaled) & \makecell[c]{$23.12\,\pm\,6.43$\\$[15.13,\ 31.11]$} & \makecell[c]{$23.95\,\pm\,7.51$\\$[14.62,\ 33.27]$} & \makecell[c]{$19.89\,\pm\,5.89$\\$[12.57,\ 27.20]$} & \makecell[c]{$\mathbf{25.36}\,\pm\,4.92$\\$[19.25,\ 31.47]$} \\
Historical & \makecell[c]{$45.87\,\pm\,20.58$\\$[20.33,\ 71.42]$} & \makecell[c]{$30.97\,\pm\,43.05$\\$[-22.48,\ 84.43]$} & \makecell[c]{$56.52\,\pm\,50.84$\\$[-6.60,\ 119.65]$} & \makecell[c]{$64.46\,\pm\,39.68$\\$[15.19,\ 113.74]$} \\
Graph Laplacian & \makecell[c]{$45.59\,\pm\,32.46$\\$[5.29,\ 85.89]$} & \makecell[c]{$61.62\,\pm\,32.29$\\$[21.52,\ 101.72]$} & \makecell[c]{$31.48\,\pm\,13.78$\\$[14.37,\ 48.58]$} & \makecell[c]{$54.21\,\pm\,26.82$\\$[20.91,\ 87.50]$} \\
Graph Temporal & \makecell[c]{$19.84\,\pm\,21.43$\\$[-6.77,\ 46.44]$} & \makecell[c]{$33.39\,\pm\,16.48$\\$[12.92,\ 53.85]$} & \makecell[c]{$20.12\,\pm\,10.54$\\$[7.04,\ 33.20]$} & \makecell[c]{$35.35\,\pm\,18.70$\\$[12.13,\ 58.57]$} \\
Sage-Voronoi & \makecell[c]{$14.30\,\pm\,8.44$\\$[3.82,\ 24.78]$} & \makecell[c]{$\mathbf{20.47}\,\pm\,12.33$\\$[5.16,\ 35.78]$} & \makecell[c]{$23.11\,\pm\,11.09$\\$[9.35,\ 36.88]$} & \makecell[c]{$35.86\,\pm\,20.11$\\$[10.89,\ 60.84]$} \\
GAIN & \makecell[c]{$15.09\,\pm\,16.95$\\$[-5.96,\ 36.13]$} & \makecell[c]{$26.68\,\pm\,14.55$\\$[8.61,\ 44.75]$} & \makecell[c]{$15.02\,\pm\,3.82$\\$[10.28,\ 19.76]$} & \makecell[c]{$29.89\,\pm\,11.95$\\$[15.05,\ 44.73]$} \\
GRIN & \makecell[c]{$10.05\,\pm\,8.73$\\$[-0.79,\ 20.89]$} & \makecell[c]{$21.45\,\pm\,7.70$\\$[11.90,\ 31.01]$} & \makecell[c]{$\mathbf{14.53}\,\pm\,4.00$\\$[9.56,\ 19.50]$} & \makecell[c]{$25.74\,\pm\,8.16$\\$[15.61,\ 35.88]$} \\
\bottomrule
\end{tabular}
\endgroup
\end{table}

\begin{table}[htbp]
\centering
\caption{{RMSE for sensor A023, January 2023 (vehicles per ten-minute interval). Each cell gives mean $\pm$ sample SD, followed by the 95\% confidence interval; $n=5$ seed runs. Bold means identify the lowest unrounded mean within each column, not statistical significance.}}
\label{tab:supp_p01_30}
\begingroup
\footnotesize
\setlength{\tabcolsep}{3pt}
\renewcommand{\arraystretch}{1.15}
\begin{tabular}{lccc}
\toprule
Method & Block $B=16$ & Block $B=32$ & Block $B=64$ \\
\midrule
Linear & \makecell[c]{$35.74\,\pm\,18.49$\\$[12.79,\ 58.69]$} & \makecell[c]{$33.39\,\pm\,14.14$\\$[15.83,\ 50.94]$} & \makecell[c]{$106.03\,\pm\,77.63$\\$[9.64,\ 202.42]$} \\
Median & \makecell[c]{$138.86\,\pm\,36.31$\\$[93.77,\ 183.94]$} & \makecell[c]{$128.38\,\pm\,11.46$\\$[114.14,\ 142.61]$} & \makecell[c]{$119.93\,\pm\,49.56$\\$[58.39,\ 181.47]$} \\
FFill + BFill & \makecell[c]{$77.02\,\pm\,56.40$\\$[6.99,\ 147.06]$} & \makecell[c]{$58.83\,\pm\,37.00$\\$[12.89,\ 104.78]$} & \makecell[c]{$123.09\,\pm\,65.66$\\$[41.57,\ 204.61]$} \\
Rolling mean & \makecell[c]{$23.81\,\pm\,8.74$\\$[12.96,\ 34.66]$} & \makecell[c]{$34.07\,\pm\,12.85$\\$[18.11,\ 50.03]$} & \makecell[c]{$103.73\,\pm\,76.05$\\$[9.30,\ 198.16]$} \\
KNN (scaled) & \makecell[c]{$88.08\,\pm\,46.40$\\$[30.47,\ 145.70]$} & \makecell[c]{$163.49\,\pm\,61.56$\\$[87.06,\ 239.92]$} & \makecell[c]{$135.28\,\pm\,40.58$\\$[84.90,\ 185.67]$} \\
KNN (unscaled) & \makecell[c]{$29.47\,\pm\,12.66$\\$[13.74,\ 45.19]$} & \makecell[c]{$\mathbf{30.04}\,\pm\,5.57$\\$[23.12,\ 36.96]$} & \makecell[c]{$31.02\,\pm\,3.15$\\$[27.11,\ 34.93]$} \\
Historical & \makecell[c]{$38.85\,\pm\,14.07$\\$[21.38,\ 56.32]$} & \makecell[c]{$68.20\,\pm\,24.29$\\$[38.04,\ 98.36]$} & \makecell[c]{$59.70\,\pm\,18.03$\\$[37.32,\ 82.09]$} \\
Graph Laplacian & \makecell[c]{$51.74\,\pm\,17.71$\\$[29.75,\ 73.72]$} & \makecell[c]{$53.04\,\pm\,13.70$\\$[36.02,\ 70.05]$} & \makecell[c]{$53.39\,\pm\,14.74$\\$[35.08,\ 71.69]$} \\
Graph Temporal & \makecell[c]{$35.74\,\pm\,10.94$\\$[22.16,\ 49.32]$} & \makecell[c]{$30.26\,\pm\,4.66$\\$[24.48,\ 36.04]$} & \makecell[c]{$37.12\,\pm\,3.15$\\$[33.21,\ 41.03]$} \\
Sage-Voronoi & \makecell[c]{$35.55\,\pm\,13.82$\\$[18.39,\ 52.71]$} & \makecell[c]{$37.55\,\pm\,5.52$\\$[30.70,\ 44.40]$} & \makecell[c]{$36.87\,\pm\,9.67$\\$[24.87,\ 48.87]$} \\
GAIN & \makecell[c]{$25.41\,\pm\,12.86$\\$[9.43,\ 41.38]$} & \makecell[c]{$30.12\,\pm\,7.63$\\$[20.65,\ 39.59]$} & \makecell[c]{$31.20\,\pm\,11.07$\\$[17.45,\ 44.95]$} \\
GRIN & \makecell[c]{$\mathbf{20.74}\,\pm\,8.16$\\$[10.60,\ 30.87]$} & \makecell[c]{$30.68\,\pm\,7.92$\\$[20.85,\ 40.51]$} & \makecell[c]{$\mathbf{26.24}\,\pm\,6.69$\\$[17.93,\ 34.54]$} \\
\bottomrule
\end{tabular}
\endgroup
\end{table}

\clearpage

\begin{table}[htbp]
\centering
\caption{{RMSE for sensor A023, January 2023 (vehicles per ten-minute interval). Each cell gives mean $\pm$ sample SD, followed by the 95\% confidence interval; $n=5$ seed runs. Bold means identify the lowest unrounded mean within each column, not statistical significance.}}
\label{tab:supp_p01_31}
\begingroup
\footnotesize
\setlength{\tabcolsep}{3pt}
\renewcommand{\arraystretch}{1.15}
\begin{tabular}{lcc}
\toprule
Method & Block $B=128$ & Block $B=144$ \\
\midrule
Linear & \makecell[c]{$135.78\,\pm\,45.41$\\$[79.39,\ 192.16]$} & \makecell[c]{$160.31\,\pm\,28.63$\\$[124.76,\ 195.86]$} \\
Median & \makecell[c]{$121.13\,\pm\,17.71$\\$[99.14,\ 143.12]$} & \makecell[c]{$135.35\,\pm\,3.97$\\$[130.42,\ 140.29]$} \\
FFill + BFill & \makecell[c]{$160.24\,\pm\,65.48$\\$[78.94,\ 241.55]$} & \makecell[c]{$160.80\,\pm\,22.56$\\$[132.78,\ 188.81]$} \\
Rolling mean & \makecell[c]{$134.86\,\pm\,46.98$\\$[76.53,\ 193.18]$} & \makecell[c]{$160.20\,\pm\,28.36$\\$[124.98,\ 195.41]$} \\
KNN (scaled) & \makecell[c]{$116.58\,\pm\,20.32$\\$[91.34,\ 141.81]$} & \makecell[c]{$129.74\,\pm\,24.98$\\$[98.72,\ 160.76]$} \\
KNN (unscaled) & \makecell[c]{$30.31\,\pm\,4.30$\\$[24.97,\ 35.65]$} & \makecell[c]{$32.51\,\pm\,4.26$\\$[27.22,\ 37.80]$} \\
Historical & \makecell[c]{$59.26\,\pm\,7.01$\\$[50.56,\ 67.96]$} & \makecell[c]{$57.99\,\pm\,2.09$\\$[55.39,\ 60.58]$} \\
Graph Laplacian & \makecell[c]{$42.80\,\pm\,6.94$\\$[34.19,\ 51.42]$} & \makecell[c]{$53.33\,\pm\,8.69$\\$[42.54,\ 64.12]$} \\
Graph Temporal & \makecell[c]{$30.90\,\pm\,4.81$\\$[24.93,\ 36.87]$} & \makecell[c]{$37.05\,\pm\,1.34$\\$[35.39,\ 38.71]$} \\
Sage-Voronoi & \makecell[c]{$33.37\,\pm\,7.24$\\$[24.38,\ 42.36]$} & \makecell[c]{$45.83\,\pm\,6.96$\\$[37.19,\ 54.47]$} \\
GAIN & \makecell[c]{$28.32\,\pm\,3.74$\\$[23.68,\ 32.96]$} & \makecell[c]{$29.77\,\pm\,2.73$\\$[26.38,\ 33.17]$} \\
GRIN & \makecell[c]{$\mathbf{27.27}\,\pm\,2.43$\\$[24.25,\ 30.30]$} & \makecell[c]{$\mathbf{28.67}\,\pm\,2.43$\\$[25.65,\ 31.69]$} \\
\bottomrule
\end{tabular}
\endgroup
\end{table}

\begin{table}[htbp]
\centering
\caption{{RMSE for sensor A023, January 2023 (vehicles per ten-minute interval). Each cell gives mean $\pm$ sample SD, followed by the 95\% confidence interval; $n=5$ seed runs. Bold means identify the lowest unrounded mean within each column, not statistical significance.}}
\label{tab:supp_p01_32}
\begingroup
\footnotesize
\setlength{\tabcolsep}{3pt}
\renewcommand{\arraystretch}{1.15}
\begin{tabular}{lcccc}
\toprule
Method & MCAR 10\% & MCAR 30\% & MCAR 50\% & Synchronous $B=144$ \\
\midrule
Linear & \makecell[c]{$28.13\,\pm\,2.28$\\$[25.29,\ 30.96]$} & \makecell[c]{$28.15\,\pm\,1.80$\\$[25.92,\ 30.38]$} & \makecell[c]{$28.30\,\pm\,0.78$\\$[27.33,\ 29.27]$} & \makecell[c]{$168.03\,\pm\,23.27$\\$[139.14,\ 196.93]$} \\
Median & \makecell[c]{$129.95\,\pm\,3.22$\\$[125.95,\ 133.95]$} & \makecell[c]{$128.63\,\pm\,2.88$\\$[125.06,\ 132.20]$} & \makecell[c]{$127.51\,\pm\,1.83$\\$[125.23,\ 129.79]$} & \makecell[c]{$133.91\,\pm\,3.79$\\$[129.21,\ 138.61]$} \\
FFill + BFill & \makecell[c]{$32.17\,\pm\,2.78$\\$[28.71,\ 35.63]$} & \makecell[c]{$34.07\,\pm\,1.07$\\$[32.74,\ 35.40]$} & \makecell[c]{$37.80\,\pm\,1.50$\\$[35.94,\ 39.66]$} & \makecell[c]{$179.50\,\pm\,24.33$\\$[149.29,\ 209.71]$} \\
Rolling mean & \makecell[c]{$31.12\,\pm\,2.66$\\$[27.82,\ 34.42]$} & \makecell[c]{$33.88\,\pm\,1.58$\\$[31.92,\ 35.83]$} & \makecell[c]{$36.11\,\pm\,1.50$\\$[34.24,\ 37.97]$} & \makecell[c]{$168.46\,\pm\,23.12$\\$[139.76,\ 197.16]$} \\
KNN (scaled) & \makecell[c]{$129.01\,\pm\,5.71$\\$[121.92,\ 136.10]$} & \makecell[c]{$124.30\,\pm\,2.25$\\$[121.50,\ 127.10]$} & \makecell[c]{$116.01\,\pm\,3.32$\\$[111.89,\ 120.13]$} & \makecell[c]{$134.87\,\pm\,10.18$\\$[122.22,\ 147.51]$} \\
KNN (unscaled) & \makecell[c]{$31.01\,\pm\,4.00$\\$[26.05,\ 35.98]$} & \makecell[c]{$34.52\,\pm\,0.78$\\$[33.55,\ 35.48]$} & \makecell[c]{$37.21\,\pm\,1.66$\\$[35.14,\ 39.27]$} & \makecell[c]{$134.87\,\pm\,10.18$\\$[122.22,\ 147.51]$} \\
Historical & \makecell[c]{$60.59\,\pm\,3.47$\\$[56.29,\ 64.90]$} & \makecell[c]{$63.32\,\pm\,1.94$\\$[60.91,\ 65.73]$} & \makecell[c]{$61.47\,\pm\,1.19$\\$[59.99,\ 62.94]$} & \makecell[c]{$\mathbf{62.48}\,\pm\,12.94$\\$[46.41,\ 78.55]$} \\
Graph Laplacian & \makecell[c]{$53.51\,\pm\,2.28$\\$[50.68,\ 56.34]$} & \makecell[c]{$57.43\,\pm\,1.79$\\$[55.22,\ 59.65]$} & \makecell[c]{$60.84\,\pm\,2.08$\\$[58.26,\ 63.42]$} & \makecell[c]{$134.87\,\pm\,10.18$\\$[122.22,\ 147.51]$} \\
Graph Temporal & \makecell[c]{$32.24\,\pm\,0.86$\\$[31.17,\ 33.30]$} & \makecell[c]{$34.33\,\pm\,1.69$\\$[32.24,\ 36.42]$} & \makecell[c]{$34.78\,\pm\,0.64$\\$[33.99,\ 35.58]$} & \makecell[c]{$200.07\,\pm\,59.36$\\$[126.37,\ 273.77]$} \\
Sage-Voronoi & \makecell[c]{$27.28\,\pm\,3.04$\\$[23.50,\ 31.06]$} & \makecell[c]{$30.12\,\pm\,2.56$\\$[26.93,\ 33.30]$} & \makecell[c]{$31.55\,\pm\,1.89$\\$[29.20,\ 33.91]$} & \makecell[c]{$157.21\,\pm\,38.71$\\$[109.15,\ 205.27]$} \\
GAIN & \makecell[c]{$29.31\,\pm\,2.09$\\$[26.72,\ 31.90]$} & \makecell[c]{$34.73\,\pm\,2.77$\\$[31.28,\ 38.17]$} & \makecell[c]{$49.99\,\pm\,4.26$\\$[44.70,\ 55.28]$} & \makecell[c]{$134.91\,\pm\,9.11$\\$[123.59,\ 146.23]$} \\
GRIN & \makecell[c]{$\mathbf{22.67}\,\pm\,1.08$\\$[21.33,\ 24.01]$} & \makecell[c]{$\mathbf{24.27}\,\pm\,1.56$\\$[22.34,\ 26.21]$} & \makecell[c]{$\mathbf{26.02}\,\pm\,1.40$\\$[24.28,\ 27.76]$} & \makecell[c]{$148.79\,\pm\,26.15$\\$[116.32,\ 181.26]$} \\
\bottomrule
\end{tabular}
\endgroup
\end{table}

\clearpage

\begin{table}[htbp]
\centering
\caption{{MRE for sensor A023, January 2023 (dimensionless ratio). Each cell gives mean $\pm$ sample SD, followed by the 95\% confidence interval; $n=5$ seed runs. Bold means identify the lowest unrounded mean within each column, not statistical significance.}}
\label{tab:supp_p01_33}
\begingroup
\footnotesize
\setlength{\tabcolsep}{3pt}
\renewcommand{\arraystretch}{1.15}
\begin{tabular}{lcccc}
\toprule
Method & Block $B=1$ & Block $B=2$ & Block $B=4$ & Block $B=8$ \\
\midrule
Linear & \makecell[c]{$\mathbf{0.031}\,\pm\,0.031$\\$[-0.008,\ 0.070]$} & \makecell[c]{$0.163\,\pm\,0.089$\\$[0.052,\ 0.274]$} & \makecell[c]{$0.275\,\pm\,0.159$\\$[0.079,\ 0.472]$} & \makecell[c]{$0.135\,\pm\,0.056$\\$[0.065,\ 0.204]$} \\
Median & \makecell[c]{$0.385\,\pm\,0.427$\\$[-0.146,\ 0.915]$} & \makecell[c]{$1.362\,\pm\,1.660$\\$[-0.699,\ 3.423]$} & \makecell[c]{$1.552\,\pm\,1.663$\\$[-0.512,\ 3.617]$} & \makecell[c]{$1.101\,\pm\,1.372$\\$[-0.602,\ 2.804]$} \\
FFill + BFill & \makecell[c]{$0.056\,\pm\,0.044$\\$[0.002,\ 0.111]$} & \makecell[c]{$0.170\,\pm\,0.061$\\$[0.095,\ 0.246]$} & \makecell[c]{$0.339\,\pm\,0.153$\\$[0.149,\ 0.529]$} & \makecell[c]{$0.339\,\pm\,0.208$\\$[0.081,\ 0.598]$} \\
Rolling mean & \makecell[c]{$0.143\,\pm\,0.161$\\$[-0.057,\ 0.344]$} & \makecell[c]{$0.169\,\pm\,0.171$\\$[-0.044,\ 0.381]$} & \makecell[c]{$0.316\,\pm\,0.115$\\$[0.173,\ 0.459]$} & \makecell[c]{$0.232\,\pm\,0.126$\\$[0.076,\ 0.388]$} \\
KNN (scaled) & \makecell[c]{$0.400\,\pm\,0.107$\\$[0.268,\ 0.533]$} & \makecell[c]{$0.755\,\pm\,0.459$\\$[0.185,\ 1.325]$} & \makecell[c]{$0.733\,\pm\,0.873$\\$[-0.351,\ 1.817]$} & \makecell[c]{$0.650\,\pm\,0.576$\\$[-0.066,\ 1.365]$} \\
KNN (unscaled) & \makecell[c]{$0.092\,\pm\,0.046$\\$[0.034,\ 0.149]$} & \makecell[c]{$0.217\,\pm\,0.254$\\$[-0.099,\ 0.533]$} & \makecell[c]{$0.155\,\pm\,0.081$\\$[0.055,\ 0.256]$} & \makecell[c]{$0.135\,\pm\,0.133$\\$[-0.030,\ 0.300]$} \\
Historical & \makecell[c]{$0.166\,\pm\,0.069$\\$[0.080,\ 0.251]$} & \makecell[c]{$\mathbf{0.128}\,\pm\,0.096$\\$[0.009,\ 0.248]$} & \makecell[c]{$0.514\,\pm\,0.452$\\$[-0.048,\ 1.075]$} & \makecell[c]{$0.283\,\pm\,0.138$\\$[0.112,\ 0.454]$} \\
Graph Laplacian & \makecell[c]{$0.179\,\pm\,0.139$\\$[0.006,\ 0.352]$} & \makecell[c]{$0.420\,\pm\,0.434$\\$[-0.120,\ 0.959]$} & \makecell[c]{$0.364\,\pm\,0.427$\\$[-0.166,\ 0.893]$} & \makecell[c]{$0.286\,\pm\,0.264$\\$[-0.041,\ 0.613]$} \\
Graph Temporal & \makecell[c]{$0.074\,\pm\,0.076$\\$[-0.020,\ 0.168]$} & \makecell[c]{$0.211\,\pm\,0.225$\\$[-0.068,\ 0.491]$} & \makecell[c]{$0.190\,\pm\,0.189$\\$[-0.045,\ 0.424]$} & \makecell[c]{$0.182\,\pm\,0.118$\\$[0.035,\ 0.329]$} \\
Sage-Voronoi & \makecell[c]{$0.057\,\pm\,0.040$\\$[0.007,\ 0.107]$} & \makecell[c]{$0.135\,\pm\,0.128$\\$[-0.024,\ 0.294]$} & \makecell[c]{$0.192\,\pm\,0.101$\\$[0.066,\ 0.318]$} & \makecell[c]{$0.144\,\pm\,0.038$\\$[0.096,\ 0.191]$} \\
GAIN & \makecell[c]{$0.053\,\pm\,0.045$\\$[-0.003,\ 0.109]$} & \makecell[c]{$0.150\,\pm\,0.105$\\$[0.019,\ 0.281]$} & \makecell[c]{$\mathbf{0.099}\,\pm\,0.027$\\$[0.066,\ 0.132]$} & \makecell[c]{$0.153\,\pm\,0.115$\\$[0.011,\ 0.296]$} \\
GRIN & \makecell[c]{$0.036\,\pm\,0.026$\\$[0.004,\ 0.068]$} & \makecell[c]{$0.159\,\pm\,0.157$\\$[-0.036,\ 0.354]$} & \makecell[c]{$0.103\,\pm\,0.039$\\$[0.054,\ 0.151]$} & \makecell[c]{$\mathbf{0.102}\,\pm\,0.038$\\$[0.055,\ 0.149]$} \\
\bottomrule
\end{tabular}
\endgroup
\end{table}

\begin{table}[htbp]
\centering
\caption{{MRE for sensor A023, January 2023 (dimensionless ratio). Each cell gives mean $\pm$ sample SD, followed by the 95\% confidence interval; $n=5$ seed runs. Bold means identify the lowest unrounded mean within each column, not statistical significance.}}
\label{tab:supp_p01_34}
\begingroup
\footnotesize
\setlength{\tabcolsep}{3pt}
\renewcommand{\arraystretch}{1.15}
\begin{tabular}{lccc}
\toprule
Method & Block $B=16$ & Block $B=32$ & Block $B=64$ \\
\midrule
Linear & \makecell[c]{$0.375\,\pm\,0.361$\\$[-0.073,\ 0.824]$} & \makecell[c]{$0.094\,\pm\,0.042$\\$[0.042,\ 0.146]$} & \makecell[c]{$1.270\,\pm\,1.697$\\$[-0.837,\ 3.377]$} \\
Median & \makecell[c]{$1.697\,\pm\,1.171$\\$[0.242,\ 3.151]$} & \makecell[c]{$0.642\,\pm\,0.730$\\$[-0.264,\ 1.548]$} & \makecell[c]{$1.260\,\pm\,1.403$\\$[-0.482,\ 3.003]$} \\
FFill + BFill & \makecell[c]{$0.390\,\pm\,0.253$\\$[0.076,\ 0.703]$} & \makecell[c]{$0.396\,\pm\,0.669$\\$[-0.435,\ 1.227]$} & \makecell[c]{$0.997\,\pm\,1.319$\\$[-0.641,\ 2.636]$} \\
Rolling mean & \makecell[c]{$0.192\,\pm\,0.117$\\$[0.047,\ 0.336]$} & \makecell[c]{$0.101\,\pm\,0.056$\\$[0.032,\ 0.170]$} & \makecell[c]{$1.226\,\pm\,1.660$\\$[-0.835,\ 3.288]$} \\
KNN (scaled) & \makecell[c]{$0.697\,\pm\,0.377$\\$[0.228,\ 1.166]$} & \makecell[c]{$0.527\,\pm\,0.074$\\$[0.435,\ 0.620]$} & \makecell[c]{$0.667\,\pm\,0.368$\\$[0.210,\ 1.125]$} \\
KNN (unscaled) & \makecell[c]{$0.228\,\pm\,0.159$\\$[0.031,\ 0.425]$} & \makecell[c]{$0.094\,\pm\,0.059$\\$[0.020,\ 0.167]$} & \makecell[c]{$0.160\,\pm\,0.115$\\$[0.017,\ 0.303]$} \\
Historical & \makecell[c]{$0.288\,\pm\,0.186$\\$[0.057,\ 0.519]$} & \makecell[c]{$0.230\,\pm\,0.122$\\$[0.079,\ 0.382]$} & \makecell[c]{$0.273\,\pm\,0.171$\\$[0.061,\ 0.485]$} \\
Graph Laplacian & \makecell[c]{$0.477\,\pm\,0.268$\\$[0.144,\ 0.810]$} & \makecell[c]{$0.192\,\pm\,0.151$\\$[0.005,\ 0.380]$} & \makecell[c]{$0.368\,\pm\,0.321$\\$[-0.030,\ 0.766]$} \\
Graph Temporal & \makecell[c]{$0.262\,\pm\,0.139$\\$[0.090,\ 0.435]$} & \makecell[c]{$0.108\,\pm\,0.094$\\$[-0.008,\ 0.224]$} & \makecell[c]{$0.228\,\pm\,0.185$\\$[-0.002,\ 0.457]$} \\
Sage-Voronoi & \makecell[c]{$0.248\,\pm\,0.137$\\$[0.078,\ 0.418]$} & \makecell[c]{$0.163\,\pm\,0.188$\\$[-0.071,\ 0.397]$} & \makecell[c]{$0.219\,\pm\,0.224$\\$[-0.059,\ 0.498]$} \\
GAIN & \makecell[c]{$0.149\,\pm\,0.054$\\$[0.081,\ 0.216]$} & \makecell[c]{$0.092\,\pm\,0.055$\\$[0.024,\ 0.160]$} & \makecell[c]{$0.121\,\pm\,0.045$\\$[0.065,\ 0.177]$} \\
GRIN & \makecell[c]{$\mathbf{0.129}\,\pm\,0.048$\\$[0.069,\ 0.189]$} & \makecell[c]{$\mathbf{0.086}\,\pm\,0.033$\\$[0.045,\ 0.126]$} & \makecell[c]{$\mathbf{0.094}\,\pm\,0.030$\\$[0.056,\ 0.132]$} \\
\bottomrule
\end{tabular}
\endgroup
\end{table}

\clearpage

\begin{table}[htbp]
\centering
\caption{{MRE for sensor A023, January 2023 (dimensionless ratio). Each cell gives mean $\pm$ sample SD, followed by the 95\% confidence interval; $n=5$ seed runs. Bold means identify the lowest unrounded mean within each column, not statistical significance.}}
\label{tab:supp_p01_35}
\begingroup
\footnotesize
\setlength{\tabcolsep}{3pt}
\renewcommand{\arraystretch}{1.15}
\begin{tabular}{lcc}
\toprule
Method & Block $B=128$ & Block $B=144$ \\
\midrule
Linear & \makecell[c]{$0.644\,\pm\,0.302$\\$[0.268,\ 1.019]$} & \makecell[c]{$1.642\,\pm\,0.686$\\$[0.791,\ 2.494]$} \\
Median & \makecell[c]{$1.041\,\pm\,0.464$\\$[0.465,\ 1.617]$} & \makecell[c]{$1.377\,\pm\,0.187$\\$[1.144,\ 1.610]$} \\
FFill + BFill & \makecell[c]{$0.816\,\pm\,0.227$\\$[0.534,\ 1.098]$} & \makecell[c]{$1.619\,\pm\,0.616$\\$[0.854,\ 2.384]$} \\
Rolling mean & \makecell[c]{$0.635\,\pm\,0.302$\\$[0.260,\ 1.009]$} & \makecell[c]{$1.640\,\pm\,0.684$\\$[0.791,\ 2.490]$} \\
KNN (scaled) & \makecell[c]{$0.564\,\pm\,0.120$\\$[0.415,\ 0.713]$} & \makecell[c]{$0.626\,\pm\,0.165$\\$[0.421,\ 0.830]$} \\
KNN (unscaled) & \makecell[c]{$0.159\,\pm\,0.026$\\$[0.126,\ 0.191]$} & \makecell[c]{$0.198\,\pm\,0.037$\\$[0.152,\ 0.244]$} \\
Historical & \makecell[c]{$0.305\,\pm\,0.083$\\$[0.202,\ 0.408]$} & \makecell[c]{$0.263\,\pm\,0.045$\\$[0.207,\ 0.319]$} \\
Graph Laplacian & \makecell[c]{$0.293\,\pm\,0.117$\\$[0.148,\ 0.438]$} & \makecell[c]{$0.384\,\pm\,0.046$\\$[0.327,\ 0.441]$} \\
Graph Temporal & \makecell[c]{$0.172\,\pm\,0.071$\\$[0.084,\ 0.261]$} & \makecell[c]{$0.238\,\pm\,0.033$\\$[0.197,\ 0.278]$} \\
Sage-Voronoi & \makecell[c]{$0.178\,\pm\,0.059$\\$[0.105,\ 0.251]$} & \makecell[c]{$0.351\,\pm\,0.144$\\$[0.172,\ 0.530]$} \\
GAIN & \makecell[c]{$0.127\,\pm\,0.014$\\$[0.109,\ 0.145]$} & \makecell[c]{$0.136\,\pm\,0.021$\\$[0.110,\ 0.163]$} \\
GRIN & \makecell[c]{$\mathbf{0.106}\,\pm\,0.007$\\$[0.098,\ 0.115]$} & \makecell[c]{$\mathbf{0.110}\,\pm\,0.007$\\$[0.101,\ 0.119]$} \\
\bottomrule
\end{tabular}
\endgroup
\end{table}

\begin{table}[htbp]
\centering
\caption{{MRE for sensor A023, January 2023 (dimensionless ratio). Each cell gives mean $\pm$ sample SD, followed by the 95\% confidence interval; $n=5$ seed runs. Bold means identify the lowest unrounded mean within each column, not statistical significance.}}
\label{tab:supp_p01_36}
\begingroup
\footnotesize
\setlength{\tabcolsep}{3pt}
\renewcommand{\arraystretch}{1.15}
\begin{tabular}{lcccc}
\toprule
Method & MCAR 10\% & MCAR 30\% & MCAR 50\% & Synchronous $B=144$ \\
\midrule
Linear & \makecell[c]{$0.126\,\pm\,0.008$\\$[0.117,\ 0.135]$} & \makecell[c]{$0.128\,\pm\,0.012$\\$[0.113,\ 0.143]$} & \makecell[c]{$0.122\,\pm\,0.005$\\$[0.115,\ 0.128]$} & \makecell[c]{$1.782\,\pm\,0.351$\\$[1.346,\ 2.217]$} \\
Median & \makecell[c]{$1.278\,\pm\,0.113$\\$[1.137,\ 1.418]$} & \makecell[c]{$1.226\,\pm\,0.142$\\$[1.050,\ 1.401]$} & \makecell[c]{$1.159\,\pm\,0.090$\\$[1.047,\ 1.270]$} & \makecell[c]{$1.333\,\pm\,0.250$\\$[1.023,\ 1.644]$} \\
FFill + BFill & \makecell[c]{$0.150\,\pm\,0.017$\\$[0.129,\ 0.171]$} & \makecell[c]{$0.154\,\pm\,0.012$\\$[0.139,\ 0.169]$} & \makecell[c]{$0.163\,\pm\,0.008$\\$[0.153,\ 0.173]$} & \makecell[c]{$1.645\,\pm\,0.629$\\$[0.864,\ 2.427]$} \\
Rolling mean & \makecell[c]{$0.152\,\pm\,0.020$\\$[0.127,\ 0.177]$} & \makecell[c]{$0.156\,\pm\,0.012$\\$[0.141,\ 0.171]$} & \makecell[c]{$0.156\,\pm\,0.007$\\$[0.147,\ 0.164]$} & \makecell[c]{$1.783\,\pm\,0.350$\\$[1.349,\ 2.218]$} \\
KNN (scaled) & \makecell[c]{$0.657\,\pm\,0.045$\\$[0.601,\ 0.712]$} & \makecell[c]{$0.610\,\pm\,0.059$\\$[0.537,\ 0.683]$} & \makecell[c]{$0.541\,\pm\,0.019$\\$[0.517,\ 0.564]$} & \makecell[c]{$1.228\,\pm\,0.209$\\$[0.969,\ 1.488]$} \\
KNN (unscaled) & \makecell[c]{$0.189\,\pm\,0.030$\\$[0.151,\ 0.226]$} & \makecell[c]{$0.205\,\pm\,0.021$\\$[0.179,\ 0.232]$} & \makecell[c]{$0.217\,\pm\,0.020$\\$[0.192,\ 0.241]$} & \makecell[c]{$1.228\,\pm\,0.209$\\$[0.969,\ 1.488]$} \\
Historical & \makecell[c]{$0.297\,\pm\,0.016$\\$[0.277,\ 0.317]$} & \makecell[c]{$0.310\,\pm\,0.019$\\$[0.286,\ 0.334]$} & \makecell[c]{$0.292\,\pm\,0.020$\\$[0.267,\ 0.317]$} & \makecell[c]{$\mathbf{0.309}\,\pm\,0.190$\\$[0.073,\ 0.544]$} \\
Graph Laplacian & \makecell[c]{$0.378\,\pm\,0.033$\\$[0.337,\ 0.419]$} & \makecell[c]{$0.378\,\pm\,0.043$\\$[0.324,\ 0.432]$} & \makecell[c]{$0.373\,\pm\,0.021$\\$[0.347,\ 0.399]$} & \makecell[c]{$1.228\,\pm\,0.209$\\$[0.969,\ 1.488]$} \\
Graph Temporal & \makecell[c]{$0.195\,\pm\,0.018$\\$[0.173,\ 0.218]$} & \makecell[c]{$0.198\,\pm\,0.020$\\$[0.173,\ 0.224]$} & \makecell[c]{$0.190\,\pm\,0.017$\\$[0.169,\ 0.211]$} & \makecell[c]{$2.092\,\pm\,0.586$\\$[1.365,\ 2.819]$} \\
Sage-Voronoi & \makecell[c]{$0.127\,\pm\,0.009$\\$[0.116,\ 0.138]$} & \makecell[c]{$0.135\,\pm\,0.017$\\$[0.113,\ 0.156]$} & \makecell[c]{$0.135\,\pm\,0.012$\\$[0.120,\ 0.150]$} & \makecell[c]{$1.333\,\pm\,0.453$\\$[0.770,\ 1.896]$} \\
GAIN & \makecell[c]{$0.143\,\pm\,0.023$\\$[0.115,\ 0.171]$} & \makecell[c]{$0.203\,\pm\,0.035$\\$[0.160,\ 0.247]$} & \makecell[c]{$0.322\,\pm\,0.047$\\$[0.264,\ 0.380]$} & \makecell[c]{$1.241\,\pm\,0.219$\\$[0.969,\ 1.512]$} \\
GRIN & \makecell[c]{$\mathbf{0.099}\,\pm\,0.007$\\$[0.090,\ 0.108]$} & \makecell[c]{$\mathbf{0.101}\,\pm\,0.007$\\$[0.092,\ 0.109]$} & \makecell[c]{$\mathbf{0.104}\,\pm\,0.004$\\$[0.098,\ 0.109]$} & \makecell[c]{$0.999\,\pm\,0.286$\\$[0.643,\ 1.354]$} \\
\bottomrule
\end{tabular}
\endgroup
\end{table}

\clearpage

\begin{table}[htbp]
\centering
\caption{{MAE for sensor A023, October--December 2022 (vehicles per ten-minute interval). Each cell gives mean $\pm$ sample SD, followed by the 95\% confidence interval; $n=5$ seed runs. Bold means identify the lowest unrounded mean within each column, not statistical significance.}}
\label{tab:supp_p01_37}
\begingroup
\footnotesize
\setlength{\tabcolsep}{3pt}
\renewcommand{\arraystretch}{1.15}
\begin{tabular}{lcccc}
\toprule
Method & Block $B=1$ & Block $B=2$ & Block $B=4$ & Block $B=8$ \\
\midrule
Linear & \makecell[c]{$21.80\,\pm\,8.63$\\$[11.09,\ 32.51]$} & \makecell[c]{$21.73\,\pm\,8.38$\\$[11.33,\ 32.14]$} & \makecell[c]{$19.57\,\pm\,8.94$\\$[8.47,\ 30.67]$} & \makecell[c]{$22.59\,\pm\,7.28$\\$[13.56,\ 31.63]$} \\
Median & \makecell[c]{$172.00\,\pm\,73.24$\\$[81.06,\ 262.94]$} & \makecell[c]{$148.40\,\pm\,73.41$\\$[57.25,\ 239.55]$} & \makecell[c]{$159.55\,\pm\,86.97$\\$[51.57,\ 267.53]$} & \makecell[c]{$99.80\,\pm\,71.80$\\$[10.65,\ 188.95]$} \\
FFill + BFill & \makecell[c]{$21.60\,\pm\,15.22$\\$[2.70,\ 40.50]$} & \makecell[c]{$23.80\,\pm\,17.28$\\$[2.35,\ 45.25]$} & \makecell[c]{$24.95\,\pm\,17.47$\\$[3.25,\ 46.65]$} & \makecell[c]{$37.80\,\pm\,18.47$\\$[14.86,\ 60.74]$} \\
Rolling mean & \makecell[c]{$25.45\,\pm\,15.39$\\$[6.34,\ 44.56]$} & \makecell[c]{$16.13\,\pm\,9.66$\\$[4.14,\ 28.13]$} & \makecell[c]{$30.97\,\pm\,16.70$\\$[10.23,\ 51.70]$} & \makecell[c]{$23.99\,\pm\,6.73$\\$[15.63,\ 32.35]$} \\
KNN (scaled) & \makecell[c]{$13.86\,\pm\,15.13$\\$[-4.93,\ 32.65]$} & \makecell[c]{$16.78\,\pm\,7.29$\\$[7.73,\ 25.83]$} & \makecell[c]{$16.12\,\pm\,10.38$\\$[3.22,\ 29.01]$} & \makecell[c]{$19.75\,\pm\,8.73$\\$[8.92,\ 30.59]$} \\
KNN (unscaled) & \makecell[c]{$23.64\,\pm\,12.07$\\$[8.65,\ 38.63]$} & \makecell[c]{$26.76\,\pm\,14.31$\\$[8.99,\ 44.53]$} & \makecell[c]{$17.07\,\pm\,7.08$\\$[8.27,\ 25.87]$} & \makecell[c]{$26.84\,\pm\,15.23$\\$[7.94,\ 45.74]$} \\
Historical & \makecell[c]{$27.08\,\pm\,30.77$\\$[-11.12,\ 65.29]$} & \makecell[c]{$75.03\,\pm\,66.96$\\$[-8.10,\ 158.17]$} & \makecell[c]{$29.15\,\pm\,11.90$\\$[14.37,\ 43.92]$} & \makecell[c]{$49.80\,\pm\,53.33$\\$[-16.42,\ 116.02]$} \\
Graph Laplacian & \makecell[c]{$29.83\,\pm\,25.45$\\$[-1.77,\ 61.43]$} & \makecell[c]{$29.65\,\pm\,9.76$\\$[17.53,\ 41.77]$} & \makecell[c]{$28.86\,\pm\,15.60$\\$[9.49,\ 48.22]$} & \makecell[c]{$25.03\,\pm\,8.83$\\$[14.06,\ 36.00]$} \\
Graph Temporal & \makecell[c]{$32.91\,\pm\,16.68$\\$[12.20,\ 53.61]$} & \makecell[c]{$30.57\,\pm\,8.37$\\$[20.17,\ 40.97]$} & \makecell[c]{$26.84\,\pm\,13.27$\\$[10.37,\ 43.32]$} & \makecell[c]{$24.20\,\pm\,7.41$\\$[15.00,\ 33.39]$} \\
Sage-Voronoi & \makecell[c]{$16.41\,\pm\,15.03$\\$[-2.25,\ 35.07]$} & \makecell[c]{$18.46\,\pm\,8.06$\\$[8.46,\ 28.47]$} & \makecell[c]{$19.89\,\pm\,13.76$\\$[2.80,\ 36.98]$} & \makecell[c]{$19.66\,\pm\,5.30$\\$[13.08,\ 26.24]$} \\
GAIN & \makecell[c]{$19.26\,\pm\,18.06$\\$[-3.16,\ 41.68]$} & \makecell[c]{$23.33\,\pm\,21.91$\\$[-3.87,\ 50.53]$} & \makecell[c]{$15.55\,\pm\,5.11$\\$[9.20,\ 21.90]$} & \makecell[c]{$23.58\,\pm\,4.88$\\$[17.51,\ 29.64]$} \\
GRIN & \makecell[c]{$\mathbf{13.05}\,\pm\,9.83$\\$[0.84,\ 25.26]$} & \makecell[c]{$\mathbf{14.42}\,\pm\,9.85$\\$[2.19,\ 26.65]$} & \makecell[c]{$\mathbf{14.43}\,\pm\,8.96$\\$[3.31,\ 25.56]$} & \makecell[c]{$\mathbf{18.16}\,\pm\,2.53$\\$[15.01,\ 21.30]$} \\
\bottomrule
\end{tabular}
\endgroup
\end{table}

\begin{table}[htbp]
\centering
\caption{{MAE for sensor A023, October--December 2022 (vehicles per ten-minute interval). Each cell gives mean $\pm$ sample SD, followed by the 95\% confidence interval; $n=5$ seed runs. Bold means identify the lowest unrounded mean within each column, not statistical significance.}}
\label{tab:supp_p01_38}
\begingroup
\footnotesize
\setlength{\tabcolsep}{3pt}
\renewcommand{\arraystretch}{1.15}
\begin{tabular}{lccc}
\toprule
Method & Block $B=16$ & Block $B=32$ & Block $B=64$ \\
\midrule
Linear & \makecell[c]{$34.75\,\pm\,7.28$\\$[25.71,\ 43.79]$} & \makecell[c]{$46.95\,\pm\,21.33$\\$[20.47,\ 73.43]$} & \makecell[c]{$70.52\,\pm\,45.49$\\$[14.04,\ 127.00]$} \\
Median & \makecell[c]{$72.59\,\pm\,32.96$\\$[31.67,\ 113.51]$} & \makecell[c]{$115.01\,\pm\,86.59$\\$[7.49,\ 222.52]$} & \makecell[c]{$108.52\,\pm\,61.53$\\$[32.12,\ 184.91]$} \\
FFill + BFill & \makecell[c]{$55.36\,\pm\,52.62$\\$[-9.98,\ 120.70]$} & \makecell[c]{$113.36\,\pm\,65.95$\\$[31.47,\ 195.24]$} & \makecell[c]{$102.60\,\pm\,91.58$\\$[-11.11,\ 216.31]$} \\
Rolling mean & \makecell[c]{$33.36\,\pm\,12.09$\\$[18.35,\ 48.37]$} & \makecell[c]{$54.39\,\pm\,25.62$\\$[22.58,\ 86.20]$} & \makecell[c]{$70.09\,\pm\,48.93$\\$[9.34,\ 130.84]$} \\
KNN (scaled) & \makecell[c]{$21.35\,\pm\,4.59$\\$[15.65,\ 27.05]$} & \makecell[c]{$22.79\,\pm\,6.65$\\$[14.53,\ 31.05]$} & \makecell[c]{$28.29\,\pm\,14.88$\\$[9.82,\ 46.76]$} \\
KNN (unscaled) & \makecell[c]{$25.54\,\pm\,3.61$\\$[21.05,\ 30.02]$} & \makecell[c]{$26.78\,\pm\,8.79$\\$[15.87,\ 37.70]$} & \makecell[c]{$34.25\,\pm\,17.62$\\$[12.37,\ 56.12]$} \\
Historical & \makecell[c]{$22.26\,\pm\,6.47$\\$[14.22,\ 30.29]$} & \makecell[c]{$59.16\,\pm\,29.18$\\$[22.93,\ 95.40]$} & \makecell[c]{$67.80\,\pm\,34.01$\\$[25.57,\ 110.02]$} \\
Graph Laplacian & \makecell[c]{$35.36\,\pm\,4.70$\\$[29.53,\ 41.19]$} & \makecell[c]{$34.35\,\pm\,13.96$\\$[17.01,\ 51.69]$} & \makecell[c]{$35.51\,\pm\,9.10$\\$[24.21,\ 46.81]$} \\
Graph Temporal & \makecell[c]{$36.20\,\pm\,11.04$\\$[22.49,\ 49.91]$} & \makecell[c]{$33.16\,\pm\,13.46$\\$[16.45,\ 49.88]$} & \makecell[c]{$33.95\,\pm\,8.60$\\$[23.26,\ 44.63]$} \\
Sage-Voronoi & \makecell[c]{$40.99\,\pm\,17.79$\\$[18.90,\ 63.08]$} & \makecell[c]{$29.81\,\pm\,13.46$\\$[13.10,\ 46.52]$} & \makecell[c]{$36.56\,\pm\,5.21$\\$[30.09,\ 43.02]$} \\
GAIN & \makecell[c]{$26.63\,\pm\,7.18$\\$[17.72,\ 35.55]$} & \makecell[c]{$26.69\,\pm\,11.34$\\$[12.61,\ 40.77]$} & \makecell[c]{$\mathbf{24.13}\,\pm\,10.09$\\$[11.60,\ 36.66]$} \\
GRIN & \makecell[c]{$\mathbf{19.54}\,\pm\,5.21$\\$[13.08,\ 26.01]$} & \makecell[c]{$\mathbf{19.90}\,\pm\,6.74$\\$[11.53,\ 28.27]$} & \makecell[c]{$25.21\,\pm\,13.00$\\$[9.07,\ 41.35]$} \\
\bottomrule
\end{tabular}
\endgroup
\end{table}

\clearpage

\begin{table}[htbp]
\centering
\caption{{MAE for sensor A023, October--December 2022 (vehicles per ten-minute interval). Each cell gives mean $\pm$ sample SD, followed by the 95\% confidence interval; $n=5$ seed runs. Bold means identify the lowest unrounded mean within each column, not statistical significance.}}
\label{tab:supp_p01_39}
\begingroup
\footnotesize
\setlength{\tabcolsep}{3pt}
\renewcommand{\arraystretch}{1.15}
\begin{tabular}{lcc}
\toprule
Method & Block $B=128$ & Block $B=144$ \\
\midrule
Linear & \makecell[c]{$182.17\,\pm\,65.26$\\$[101.14,\ 263.21]$} & \makecell[c]{$126.79\,\pm\,10.31$\\$[113.99,\ 139.59]$} \\
Median & \makecell[c]{$94.85\,\pm\,15.95$\\$[75.05,\ 114.65]$} & \makecell[c]{$113.29\,\pm\,7.10$\\$[104.48,\ 122.10]$} \\
FFill + BFill & \makecell[c]{$178.65\,\pm\,76.79$\\$[83.30,\ 274.00]$} & \makecell[c]{$117.95\,\pm\,8.44$\\$[107.47,\ 128.43]$} \\
Rolling mean & \makecell[c]{$181.80\,\pm\,65.69$\\$[100.23,\ 263.37]$} & \makecell[c]{$126.61\,\pm\,10.13$\\$[114.04,\ 139.19]$} \\
KNN (scaled) & \makecell[c]{$22.07\,\pm\,2.37$\\$[19.13,\ 25.01]$} & \makecell[c]{$23.57\,\pm\,2.25$\\$[20.78,\ 26.37]$} \\
KNN (unscaled) & \makecell[c]{$26.61\,\pm\,0.45$\\$[26.05,\ 27.17]$} & \makecell[c]{$27.82\,\pm\,2.01$\\$[25.33,\ 30.31]$} \\
Historical & \makecell[c]{$29.97\,\pm\,7.52$\\$[20.63,\ 39.31]$} & \makecell[c]{$45.38\,\pm\,20.47$\\$[19.96,\ 70.80]$} \\
Graph Laplacian & \makecell[c]{$34.23\,\pm\,2.08$\\$[31.65,\ 36.82]$} & \makecell[c]{$35.72\,\pm\,6.70$\\$[27.40,\ 44.05]$} \\
Graph Temporal & \makecell[c]{$35.15\,\pm\,2.99$\\$[31.44,\ 38.87]$} & \makecell[c]{$33.87\,\pm\,6.52$\\$[25.78,\ 41.97]$} \\
Sage-Voronoi & \makecell[c]{$34.79\,\pm\,2.29$\\$[31.95,\ 37.63]$} & \makecell[c]{$32.62\,\pm\,7.72$\\$[23.03,\ 42.21]$} \\
GAIN & \makecell[c]{$23.59\,\pm\,3.55$\\$[19.18,\ 27.99]$} & \makecell[c]{$23.67\,\pm\,2.92$\\$[20.05,\ 27.30]$} \\
GRIN & \makecell[c]{$\mathbf{20.65}\,\pm\,0.66$\\$[19.84,\ 21.47]$} & \makecell[c]{$\mathbf{21.72}\,\pm\,3.57$\\$[17.28,\ 26.15]$} \\
\bottomrule
\end{tabular}
\endgroup
\end{table}

\begin{table}[htbp]
\centering
\caption{{MAE for sensor A023, October--December 2022 (vehicles per ten-minute interval). Each cell gives mean $\pm$ sample SD, followed by the 95\% confidence interval; $n=5$ seed runs. Bold means identify the lowest unrounded mean within each column, not statistical significance.}}
\label{tab:supp_p01_40}
\begingroup
\footnotesize
\setlength{\tabcolsep}{3pt}
\renewcommand{\arraystretch}{1.15}
\begin{tabular}{lcccc}
\toprule
Method & MCAR 10\% & MCAR 30\% & MCAR 50\% & Synchronous $B=144$ \\
\midrule
Linear & \makecell[c]{$22.87\,\pm\,1.48$\\$[21.03,\ 24.71]$} & \makecell[c]{$23.53\,\pm\,0.86$\\$[22.46,\ 24.60]$} & \makecell[c]{$23.94\,\pm\,0.40$\\$[23.44,\ 24.43]$} & \makecell[c]{$124.28\,\pm\,30.70$\\$[86.16,\ 162.40]$} \\
Median & \makecell[c]{$108.24\,\pm\,2.26$\\$[105.44,\ 111.04]$} & \makecell[c]{$110.07\,\pm\,1.82$\\$[107.81,\ 112.34]$} & \makecell[c]{$107.39\,\pm\,0.89$\\$[106.29,\ 108.49]$} & \makecell[c]{$107.86\,\pm\,13.79$\\$[90.73,\ 124.98]$} \\
FFill + BFill & \makecell[c]{$27.91\,\pm\,1.29$\\$[26.30,\ 29.51]$} & \makecell[c]{$28.00\,\pm\,0.89$\\$[26.89,\ 29.10]$} & \makecell[c]{$29.78\,\pm\,0.86$\\$[28.71,\ 30.86]$} & \makecell[c]{$121.93\,\pm\,27.77$\\$[87.46,\ 156.41]$} \\
Rolling mean & \makecell[c]{$26.87\,\pm\,1.65$\\$[24.83,\ 28.92]$} & \makecell[c]{$27.40\,\pm\,0.62$\\$[26.63,\ 28.17]$} & \makecell[c]{$28.59\,\pm\,0.62$\\$[27.83,\ 29.36]$} & \makecell[c]{$124.48\,\pm\,30.67$\\$[86.41,\ 162.56]$} \\
KNN (scaled) & \makecell[c]{$24.76\,\pm\,1.20$\\$[23.27,\ 26.26]$} & \makecell[c]{$24.45\,\pm\,0.94$\\$[23.28,\ 25.61]$} & \makecell[c]{$24.80\,\pm\,0.45$\\$[24.24,\ 25.36]$} & \makecell[c]{$103.28\,\pm\,14.90$\\$[84.78,\ 121.77]$} \\
KNN (unscaled) & \makecell[c]{$29.15\,\pm\,2.09$\\$[26.55,\ 31.75]$} & \makecell[c]{$28.27\,\pm\,0.85$\\$[27.21,\ 29.33]$} & \makecell[c]{$27.83\,\pm\,0.29$\\$[27.47,\ 28.19]$} & \makecell[c]{$103.28\,\pm\,14.90$\\$[84.78,\ 121.77]$} \\
Historical & \makecell[c]{$46.54\,\pm\,2.48$\\$[43.46,\ 49.62]$} & \makecell[c]{$46.45\,\pm\,1.89$\\$[44.11,\ 48.80]$} & \makecell[c]{$46.29\,\pm\,0.56$\\$[45.59,\ 46.99]$} & \makecell[c]{$\mathbf{54.86}\,\pm\,19.45$\\$[30.71,\ 79.00]$} \\
Graph Laplacian & \makecell[c]{$35.28\,\pm\,0.74$\\$[34.35,\ 36.20]$} & \makecell[c]{$36.09\,\pm\,1.07$\\$[34.76,\ 37.41]$} & \makecell[c]{$36.27\,\pm\,0.30$\\$[35.90,\ 36.64]$} & \makecell[c]{$103.28\,\pm\,14.90$\\$[84.78,\ 121.77]$} \\
Graph Temporal & \makecell[c]{$30.38\,\pm\,1.12$\\$[28.98,\ 31.77]$} & \makecell[c]{$31.80\,\pm\,0.41$\\$[31.30,\ 32.31]$} & \makecell[c]{$32.73\,\pm\,0.36$\\$[32.28,\ 33.17]$} & \makecell[c]{$118.67\,\pm\,24.66$\\$[88.06,\ 149.29]$} \\
Sage-Voronoi & \makecell[c]{$22.70\,\pm\,1.32$\\$[21.06,\ 24.35]$} & \makecell[c]{$23.36\,\pm\,0.66$\\$[22.55,\ 24.18]$} & \makecell[c]{$25.06\,\pm\,0.52$\\$[24.42,\ 25.71]$} & \makecell[c]{$130.94\,\pm\,42.92$\\$[77.65,\ 184.24]$} \\
GAIN & \makecell[c]{$24.53\,\pm\,1.41$\\$[22.78,\ 26.28]$} & \makecell[c]{$28.33\,\pm\,0.96$\\$[27.13,\ 29.52]$} & \makecell[c]{$39.49\,\pm\,0.89$\\$[38.38,\ 40.60]$} & \makecell[c]{$104.91\,\pm\,11.67$\\$[90.42,\ 119.41]$} \\
GRIN & \makecell[c]{$\mathbf{19.46}\,\pm\,1.17$\\$[18.01,\ 20.91]$} & \makecell[c]{$\mathbf{19.88}\,\pm\,0.79$\\$[18.90,\ 20.87]$} & \makecell[c]{$\mathbf{20.51}\,\pm\,0.46$\\$[19.94,\ 21.09]$} & \makecell[c]{$89.28\,\pm\,16.66$\\$[68.60,\ 109.96]$} \\
\bottomrule
\end{tabular}
\endgroup
\end{table}

\clearpage

\begin{table}[htbp]
\centering
\caption{{RMSE for sensor A023, October--December 2022 (vehicles per ten-minute interval). Each cell gives mean $\pm$ sample SD, followed by the 95\% confidence interval; $n=5$ seed runs. Bold means identify the lowest unrounded mean within each column, not statistical significance.}}
\label{tab:supp_p01_41}
\begingroup
\footnotesize
\setlength{\tabcolsep}{3pt}
\renewcommand{\arraystretch}{1.15}
\begin{tabular}{lcccc}
\toprule
Method & Block $B=1$ & Block $B=2$ & Block $B=4$ & Block $B=8$ \\
\midrule
Linear & \makecell[c]{$21.80\,\pm\,8.63$\\$[11.09,\ 32.51]$} & \makecell[c]{$24.24\,\pm\,7.45$\\$[14.99,\ 33.49]$} & \makecell[c]{$21.25\,\pm\,9.30$\\$[9.69,\ 32.80]$} & \makecell[c]{$28.19\,\pm\,9.50$\\$[16.40,\ 39.99]$} \\
Median & \makecell[c]{$172.00\,\pm\,73.24$\\$[81.06,\ 262.94]$} & \makecell[c]{$149.35\,\pm\,73.46$\\$[58.14,\ 240.55]$} & \makecell[c]{$160.87\,\pm\,85.49$\\$[54.72,\ 267.02]$} & \makecell[c]{$104.64\,\pm\,69.39$\\$[18.48,\ 190.81]$} \\
FFill + BFill & \makecell[c]{$21.60\,\pm\,15.22$\\$[2.70,\ 40.50]$} & \makecell[c]{$24.61\,\pm\,17.91$\\$[2.37,\ 46.85]$} & \makecell[c]{$28.47\,\pm\,17.87$\\$[6.28,\ 50.66]$} & \makecell[c]{$42.99\,\pm\,19.68$\\$[18.56,\ 67.43]$} \\
Rolling mean & \makecell[c]{$25.45\,\pm\,15.39$\\$[6.34,\ 44.56]$} & \makecell[c]{$18.73\,\pm\,9.90$\\$[6.43,\ 31.03]$} & \makecell[c]{$34.95\,\pm\,16.52$\\$[14.44,\ 55.46]$} & \makecell[c]{$28.76\,\pm\,8.02$\\$[18.80,\ 38.73]$} \\
KNN (scaled) & \makecell[c]{$13.86\,\pm\,15.13$\\$[-4.93,\ 32.65]$} & \makecell[c]{$20.44\,\pm\,9.27$\\$[8.93,\ 31.95]$} & \makecell[c]{$18.10\,\pm\,11.05$\\$[4.38,\ 31.82]$} & \makecell[c]{$24.39\,\pm\,8.71$\\$[13.58,\ 35.21]$} \\
KNN (unscaled) & \makecell[c]{$23.64\,\pm\,12.07$\\$[8.65,\ 38.63]$} & \makecell[c]{$30.39\,\pm\,14.95$\\$[11.83,\ 48.96]$} & \makecell[c]{$20.32\,\pm\,8.77$\\$[9.43,\ 31.21]$} & \makecell[c]{$32.85\,\pm\,14.00$\\$[15.47,\ 50.24]$} \\
Historical & \makecell[c]{$27.08\,\pm\,30.77$\\$[-11.12,\ 65.29]$} & \makecell[c]{$78.58\,\pm\,64.75$\\$[-1.82,\ 158.98]$} & \makecell[c]{$31.73\,\pm\,11.85$\\$[17.02,\ 46.44]$} & \makecell[c]{$53.84\,\pm\,51.67$\\$[-10.32,\ 117.99]$} \\
Graph Laplacian & \makecell[c]{$29.83\,\pm\,25.45$\\$[-1.77,\ 61.43]$} & \makecell[c]{$33.48\,\pm\,10.71$\\$[20.18,\ 46.77]$} & \makecell[c]{$31.73\,\pm\,14.02$\\$[14.32,\ 49.13]$} & \makecell[c]{$29.82\,\pm\,9.86$\\$[17.58,\ 42.06]$} \\
Graph Temporal & \makecell[c]{$32.91\,\pm\,16.68$\\$[12.20,\ 53.61]$} & \makecell[c]{$35.46\,\pm\,9.46$\\$[23.72,\ 47.21]$} & \makecell[c]{$29.52\,\pm\,11.33$\\$[15.46,\ 43.58]$} & \makecell[c]{$29.18\,\pm\,8.05$\\$[19.18,\ 39.17]$} \\
Sage-Voronoi & \makecell[c]{$16.41\,\pm\,15.03$\\$[-2.25,\ 35.07]$} & \makecell[c]{$20.18\,\pm\,7.59$\\$[10.75,\ 29.61]$} & \makecell[c]{$22.51\,\pm\,13.76$\\$[5.43,\ 39.59]$} & \makecell[c]{$24.23\,\pm\,5.71$\\$[17.14,\ 31.32]$} \\
GAIN & \makecell[c]{$19.26\,\pm\,18.06$\\$[-3.16,\ 41.68]$} & \makecell[c]{$26.81\,\pm\,23.25$\\$[-2.06,\ 55.67]$} & \makecell[c]{$19.03\,\pm\,5.59$\\$[12.10,\ 25.97]$} & \makecell[c]{$29.65\,\pm\,4.61$\\$[23.92,\ 35.37]$} \\
GRIN & \makecell[c]{$\mathbf{13.05}\,\pm\,9.83$\\$[0.84,\ 25.26]$} & \makecell[c]{$\mathbf{16.63}\,\pm\,10.94$\\$[3.05,\ 30.21]$} & \makecell[c]{$\mathbf{16.34}\,\pm\,9.64$\\$[4.37,\ 28.31]$} & \makecell[c]{$\mathbf{21.76}\,\pm\,3.40$\\$[17.53,\ 25.98]$} \\
\bottomrule
\end{tabular}
\endgroup
\end{table}

\begin{table}[htbp]
\centering
\caption{{RMSE for sensor A023, October--December 2022 (vehicles per ten-minute interval). Each cell gives mean $\pm$ sample SD, followed by the 95\% confidence interval; $n=5$ seed runs. Bold means identify the lowest unrounded mean within each column, not statistical significance.}}
\label{tab:supp_p01_42}
\begingroup
\footnotesize
\setlength{\tabcolsep}{3pt}
\renewcommand{\arraystretch}{1.15}
\begin{tabular}{lccc}
\toprule
Method & Block $B=16$ & Block $B=32$ & Block $B=64$ \\
\midrule
Linear & \makecell[c]{$41.51\,\pm\,5.68$\\$[34.45,\ 48.56]$} & \makecell[c]{$55.84\,\pm\,27.60$\\$[21.57,\ 90.12]$} & \makecell[c]{$82.58\,\pm\,49.65$\\$[20.93,\ 144.23]$} \\
Median & \makecell[c]{$84.20\,\pm\,37.73$\\$[37.34,\ 131.05]$} & \makecell[c]{$124.45\,\pm\,85.45$\\$[18.35,\ 230.56]$} & \makecell[c]{$122.37\,\pm\,64.04$\\$[42.86,\ 201.88]$} \\
FFill + BFill & \makecell[c]{$67.48\,\pm\,60.58$\\$[-7.74,\ 142.70]$} & \makecell[c]{$125.37\,\pm\,69.46$\\$[39.12,\ 211.62]$} & \makecell[c]{$118.85\,\pm\,90.20$\\$[6.85,\ 230.85]$} \\
Rolling mean & \makecell[c]{$39.87\,\pm\,9.86$\\$[27.63,\ 52.11]$} & \makecell[c]{$63.45\,\pm\,33.63$\\$[21.69,\ 105.20]$} & \makecell[c]{$82.26\,\pm\,53.31$\\$[16.07,\ 148.45]$} \\
KNN (scaled) & \makecell[c]{$29.35\,\pm\,6.17$\\$[21.69,\ 37.01]$} & \makecell[c]{$29.08\,\pm\,9.08$\\$[17.80,\ 40.35]$} & \makecell[c]{$36.73\,\pm\,21.40$\\$[10.16,\ 63.30]$} \\
KNN (unscaled) & \makecell[c]{$32.51\,\pm\,4.38$\\$[27.08,\ 37.95]$} & \makecell[c]{$33.13\,\pm\,10.37$\\$[20.25,\ 46.00]$} & \makecell[c]{$44.29\,\pm\,22.86$\\$[15.90,\ 72.68]$} \\
Historical & \makecell[c]{$29.92\,\pm\,7.25$\\$[20.92,\ 38.92]$} & \makecell[c]{$64.77\,\pm\,31.16$\\$[26.09,\ 103.46]$} & \makecell[c]{$81.56\,\pm\,41.90$\\$[29.53,\ 133.59]$} \\
Graph Laplacian & \makecell[c]{$43.76\,\pm\,6.76$\\$[35.37,\ 52.15]$} & \makecell[c]{$40.72\,\pm\,13.13$\\$[24.42,\ 57.02]$} & \makecell[c]{$44.93\,\pm\,11.94$\\$[30.11,\ 59.75]$} \\
Graph Temporal & \makecell[c]{$45.31\,\pm\,14.34$\\$[27.50,\ 63.12]$} & \makecell[c]{$38.56\,\pm\,13.05$\\$[22.35,\ 54.77]$} & \makecell[c]{$42.80\,\pm\,13.24$\\$[26.37,\ 59.24]$} \\
Sage-Voronoi & \makecell[c]{$48.57\,\pm\,18.47$\\$[25.63,\ 71.50]$} & \makecell[c]{$35.33\,\pm\,14.56$\\$[17.26,\ 53.41]$} & \makecell[c]{$46.25\,\pm\,10.71$\\$[32.95,\ 59.55]$} \\
GAIN & \makecell[c]{$36.95\,\pm\,8.83$\\$[25.98,\ 47.92]$} & \makecell[c]{$33.35\,\pm\,13.12$\\$[17.06,\ 49.65]$} & \makecell[c]{$\mathbf{33.07}\,\pm\,17.86$\\$[10.90,\ 55.25]$} \\
GRIN & \makecell[c]{$\mathbf{27.72}\,\pm\,6.93$\\$[19.12,\ 36.33]$} & \makecell[c]{$\mathbf{24.77}\,\pm\,7.43$\\$[15.55,\ 33.99]$} & \makecell[c]{$33.81\,\pm\,20.02$\\$[8.96,\ 58.66]$} \\
\bottomrule
\end{tabular}
\endgroup
\end{table}

\clearpage

\begin{table}[htbp]
\centering
\caption{{RMSE for sensor A023, October--December 2022 (vehicles per ten-minute interval). Each cell gives mean $\pm$ sample SD, followed by the 95\% confidence interval; $n=5$ seed runs. Bold means identify the lowest unrounded mean within each column, not statistical significance.}}
\label{tab:supp_p01_43}
\begingroup
\footnotesize
\setlength{\tabcolsep}{3pt}
\renewcommand{\arraystretch}{1.15}
\begin{tabular}{lcc}
\toprule
Method & Block $B=128$ & Block $B=144$ \\
\midrule
Linear & \makecell[c]{$213.48\,\pm\,55.88$\\$[144.10,\ 282.87]$} & \makecell[c]{$164.90\,\pm\,21.60$\\$[138.07,\ 191.72]$} \\
Median & \makecell[c]{$117.24\,\pm\,17.04$\\$[96.08,\ 138.40]$} & \makecell[c]{$138.75\,\pm\,10.01$\\$[126.32,\ 151.18]$} \\
FFill + BFill & \makecell[c]{$207.25\,\pm\,71.18$\\$[118.88,\ 295.63]$} & \makecell[c]{$155.89\,\pm\,22.97$\\$[127.37,\ 184.42]$} \\
Rolling mean & \makecell[c]{$213.14\,\pm\,56.41$\\$[143.10,\ 283.18]$} & \makecell[c]{$164.72\,\pm\,21.40$\\$[138.15,\ 191.29]$} \\
KNN (scaled) & \makecell[c]{$28.36\,\pm\,2.58$\\$[25.16,\ 31.56]$} & \makecell[c]{$32.36\,\pm\,6.24$\\$[24.61,\ 40.11]$} \\
KNN (unscaled) & \makecell[c]{$33.78\,\pm\,1.07$\\$[32.45,\ 35.11]$} & \makecell[c]{$36.43\,\pm\,5.70$\\$[29.35,\ 43.51]$} \\
Historical & \makecell[c]{$38.69\,\pm\,8.63$\\$[27.97,\ 49.41]$} & \makecell[c]{$60.50\,\pm\,26.27$\\$[27.88,\ 93.12]$} \\
Graph Laplacian & \makecell[c]{$41.93\,\pm\,2.66$\\$[38.62,\ 45.23]$} & \makecell[c]{$43.75\,\pm\,9.81$\\$[31.57,\ 55.93]$} \\
Graph Temporal & \makecell[c]{$42.52\,\pm\,3.43$\\$[38.26,\ 46.78]$} & \makecell[c]{$42.31\,\pm\,10.47$\\$[29.30,\ 55.32]$} \\
Sage-Voronoi & \makecell[c]{$44.24\,\pm\,3.44$\\$[39.96,\ 48.51]$} & \makecell[c]{$42.97\,\pm\,11.76$\\$[28.36,\ 57.57]$} \\
GAIN & \makecell[c]{$30.13\,\pm\,4.93$\\$[24.01,\ 36.24]$} & \makecell[c]{$32.49\,\pm\,7.06$\\$[23.73,\ 41.26]$} \\
GRIN & \makecell[c]{$\mathbf{26.97}\,\pm\,0.46$\\$[26.39,\ 27.54]$} & \makecell[c]{$\mathbf{29.85}\,\pm\,8.66$\\$[19.10,\ 40.60]$} \\
\bottomrule
\end{tabular}
\endgroup
\end{table}

\begin{table}[htbp]
\centering
\caption{{RMSE for sensor A023, October--December 2022 (vehicles per ten-minute interval). Each cell gives mean $\pm$ sample SD, followed by the 95\% confidence interval; $n=5$ seed runs. Bold means identify the lowest unrounded mean within each column, not statistical significance.}}
\label{tab:supp_p01_44}
\begingroup
\footnotesize
\setlength{\tabcolsep}{3pt}
\renewcommand{\arraystretch}{1.15}
\begin{tabular}{lcccc}
\toprule
Method & MCAR 10\% & MCAR 30\% & MCAR 50\% & Synchronous $B=144$ \\
\midrule
Linear & \makecell[c]{$31.01\,\pm\,2.90$\\$[27.41,\ 34.62]$} & \makecell[c]{$32.66\,\pm\,1.96$\\$[30.23,\ 35.10]$} & \makecell[c]{$33.36\,\pm\,0.60$\\$[32.61,\ 34.10]$} & \makecell[c]{$148.03\,\pm\,39.90$\\$[98.49,\ 197.57]$} \\
Median & \makecell[c]{$133.32\,\pm\,2.13$\\$[130.67,\ 135.96]$} & \makecell[c]{$135.44\,\pm\,1.73$\\$[133.29,\ 137.60]$} & \makecell[c]{$132.83\,\pm\,0.83$\\$[131.79,\ 133.87]$} & \makecell[c]{$129.88\,\pm\,15.38$\\$[110.78,\ 148.97]$} \\
FFill + BFill & \makecell[c]{$37.14\,\pm\,3.18$\\$[33.20,\ 41.09]$} & \makecell[c]{$38.61\,\pm\,2.31$\\$[35.74,\ 41.48]$} & \makecell[c]{$41.34\,\pm\,1.03$\\$[40.06,\ 42.61]$} & \makecell[c]{$146.10\,\pm\,38.02$\\$[98.90,\ 193.30]$} \\
Rolling mean & \makecell[c]{$36.84\,\pm\,3.02$\\$[33.10,\ 40.59]$} & \makecell[c]{$37.72\,\pm\,1.30$\\$[36.10,\ 39.34]$} & \makecell[c]{$39.56\,\pm\,0.63$\\$[38.77,\ 40.34]$} & \makecell[c]{$148.02\,\pm\,39.88$\\$[98.50,\ 197.53]$} \\
KNN (scaled) & \makecell[c]{$33.73\,\pm\,2.16$\\$[31.04,\ 36.41]$} & \makecell[c]{$33.83\,\pm\,1.94$\\$[31.42,\ 36.25]$} & \makecell[c]{$34.72\,\pm\,1.20$\\$[33.24,\ 36.21]$} & \makecell[c]{$118.27\,\pm\,14.45$\\$[100.32,\ 136.22]$} \\
KNN (unscaled) & \makecell[c]{$38.36\,\pm\,3.38$\\$[34.17,\ 42.56]$} & \makecell[c]{$38.03\,\pm\,1.77$\\$[35.82,\ 40.23]$} & \makecell[c]{$38.20\,\pm\,0.88$\\$[37.10,\ 39.30]$} & \makecell[c]{$118.27\,\pm\,14.45$\\$[100.32,\ 136.22]$} \\
Historical & \makecell[c]{$65.14\,\pm\,3.60$\\$[60.68,\ 69.61]$} & \makecell[c]{$65.48\,\pm\,2.66$\\$[62.18,\ 68.78]$} & \makecell[c]{$65.08\,\pm\,0.73$\\$[64.17,\ 65.98]$} & \makecell[c]{$\mathbf{71.49}\,\pm\,23.31$\\$[42.55,\ 100.44]$} \\
Graph Laplacian & \makecell[c]{$44.52\,\pm\,1.70$\\$[42.42,\ 46.63]$} & \makecell[c]{$45.34\,\pm\,1.89$\\$[42.99,\ 47.69]$} & \makecell[c]{$45.90\,\pm\,0.76$\\$[44.96,\ 46.84]$} & \makecell[c]{$118.27\,\pm\,14.45$\\$[100.32,\ 136.22]$} \\
Graph Temporal & \makecell[c]{$38.39\,\pm\,2.47$\\$[35.33,\ 41.45]$} & \makecell[c]{$40.48\,\pm\,1.24$\\$[38.94,\ 42.03]$} & \makecell[c]{$41.98\,\pm\,0.91$\\$[40.84,\ 43.11]$} & \makecell[c]{$141.58\,\pm\,34.27$\\$[99.03,\ 184.14]$} \\
Sage-Voronoi & \makecell[c]{$31.35\,\pm\,3.27$\\$[27.30,\ 35.41]$} & \makecell[c]{$32.52\,\pm\,1.56$\\$[30.59,\ 34.46]$} & \makecell[c]{$34.98\,\pm\,1.30$\\$[33.36,\ 36.60]$} & \makecell[c]{$163.35\,\pm\,51.87$\\$[98.95,\ 227.75]$} \\
GAIN & \makecell[c]{$33.49\,\pm\,2.65$\\$[30.20,\ 36.78]$} & \makecell[c]{$37.40\,\pm\,2.29$\\$[34.56,\ 40.24]$} & \makecell[c]{$48.69\,\pm\,1.07$\\$[47.36,\ 50.01]$} & \makecell[c]{$124.06\,\pm\,14.27$\\$[106.34,\ 141.77]$} \\
GRIN & \makecell[c]{$\mathbf{27.24}\,\pm\,2.71$\\$[23.87,\ 30.61]$} & \makecell[c]{$\mathbf{28.33}\,\pm\,2.03$\\$[25.80,\ 30.85]$} & \makecell[c]{$\mathbf{29.47}\,\pm\,1.05$\\$[28.17,\ 30.77]$} & \makecell[c]{$107.69\,\pm\,24.63$\\$[77.11,\ 138.27]$} \\
\bottomrule
\end{tabular}
\endgroup
\end{table}

\clearpage

\begin{table}[htbp]
\centering
\caption{{MRE for sensor A023, October--December 2022 (dimensionless ratio). Each cell gives mean $\pm$ sample SD, followed by the 95\% confidence interval; $n=5$ seed runs. Bold means identify the lowest unrounded mean within each column, not statistical significance.}}
\label{tab:supp_p01_45}
\begingroup
\footnotesize
\setlength{\tabcolsep}{3pt}
\renewcommand{\arraystretch}{1.15}
\begin{tabular}{lcccc}
\toprule
Method & Block $B=1$ & Block $B=2$ & Block $B=4$ & Block $B=8$ \\
\midrule
Linear & \makecell[c]{$0.208\,\pm\,0.122$\\$[0.057,\ 0.359]$} & \makecell[c]{$0.215\,\pm\,0.163$\\$[0.012,\ 0.417]$} & \makecell[c]{$0.173\,\pm\,0.061$\\$[0.097,\ 0.249]$} & \makecell[c]{$0.131\,\pm\,0.108$\\$[-0.003,\ 0.266]$} \\
Median & \makecell[c]{$2.227\,\pm\,1.950$\\$[-0.194,\ 4.649]$} & \makecell[c]{$1.613\,\pm\,1.437$\\$[-0.171,\ 3.397]$} & \makecell[c]{$2.182\,\pm\,2.185$\\$[-0.531,\ 4.895]$} & \makecell[c]{$0.882\,\pm\,1.174$\\$[-0.576,\ 2.341]$} \\
FFill + BFill & \makecell[c]{$0.212\,\pm\,0.176$\\$[-0.006,\ 0.430]$} & \makecell[c]{$0.195\,\pm\,0.146$\\$[0.013,\ 0.376]$} & \makecell[c]{$0.205\,\pm\,0.099$\\$[0.082,\ 0.328]$} & \makecell[c]{$0.239\,\pm\,0.245$\\$[-0.066,\ 0.543]$} \\
Rolling mean & \makecell[c]{$0.290\,\pm\,0.210$\\$[0.029,\ 0.551]$} & \makecell[c]{$0.172\,\pm\,0.167$\\$[-0.035,\ 0.379]$} & \makecell[c]{$0.268\,\pm\,0.108$\\$[0.134,\ 0.403]$} & \makecell[c]{$0.146\,\pm\,0.122$\\$[-0.005,\ 0.298]$} \\
KNN (scaled) & \makecell[c]{$0.139\,\pm\,0.157$\\$[-0.056,\ 0.334]$} & \makecell[c]{$0.148\,\pm\,0.100$\\$[0.024,\ 0.272]$} & \makecell[c]{$0.152\,\pm\,0.097$\\$[0.031,\ 0.273]$} & \makecell[c]{$0.110\,\pm\,0.073$\\$[0.020,\ 0.201]$} \\
KNN (unscaled) & \makecell[c]{$0.202\,\pm\,0.104$\\$[0.073,\ 0.332]$} & \makecell[c]{$0.196\,\pm\,0.125$\\$[0.040,\ 0.351]$} & \makecell[c]{$0.172\,\pm\,0.147$\\$[-0.010,\ 0.355]$} & \makecell[c]{$0.159\,\pm\,0.130$\\$[-0.003,\ 0.321]$} \\
Historical & \makecell[c]{$0.305\,\pm\,0.404$\\$[-0.196,\ 0.807]$} & \makecell[c]{$0.519\,\pm\,0.492$\\$[-0.093,\ 1.130]$} & \makecell[c]{$0.293\,\pm\,0.179$\\$[0.071,\ 0.515]$} & \makecell[c]{$0.287\,\pm\,0.368$\\$[-0.170,\ 0.743]$} \\
Graph Laplacian & \makecell[c]{$0.446\,\pm\,0.463$\\$[-0.128,\ 1.021]$} & \makecell[c]{$0.321\,\pm\,0.284$\\$[-0.031,\ 0.673]$} & \makecell[c]{$0.419\,\pm\,0.489$\\$[-0.187,\ 1.026]$} & \makecell[c]{$0.195\,\pm\,0.231$\\$[-0.092,\ 0.482]$} \\
Graph Temporal & \makecell[c]{$0.408\,\pm\,0.363$\\$[-0.042,\ 0.859]$} & \makecell[c]{$0.297\,\pm\,0.227$\\$[0.014,\ 0.579]$} & \makecell[c]{$0.381\,\pm\,0.431$\\$[-0.154,\ 0.916]$} & \makecell[c]{$0.174\,\pm\,0.186$\\$[-0.058,\ 0.405]$} \\
Sage-Voronoi & \makecell[c]{$0.201\,\pm\,0.201$\\$[-0.049,\ 0.450]$} & \makecell[c]{$0.174\,\pm\,0.149$\\$[-0.010,\ 0.359]$} & \makecell[c]{$0.182\,\pm\,0.100$\\$[0.058,\ 0.306]$} & \makecell[c]{$0.109\,\pm\,0.065$\\$[0.028,\ 0.190]$} \\
GAIN & \makecell[c]{$0.175\,\pm\,0.196$\\$[-0.069,\ 0.418]$} & \makecell[c]{$0.166\,\pm\,0.118$\\$[0.019,\ 0.313]$} & \makecell[c]{$0.143\,\pm\,0.072$\\$[0.053,\ 0.233]$} & \makecell[c]{$0.125\,\pm\,0.059$\\$[0.052,\ 0.199]$} \\
GRIN & \makecell[c]{$\mathbf{0.137}\,\pm\,0.129$\\$[-0.023,\ 0.298]$} & \makecell[c]{$\mathbf{0.144}\,\pm\,0.134$\\$[-0.022,\ 0.310]$} & \makecell[c]{$\mathbf{0.114}\,\pm\,0.049$\\$[0.053,\ 0.174]$} & \makecell[c]{$\mathbf{0.102}\,\pm\,0.057$\\$[0.032,\ 0.173]$} \\
\bottomrule
\end{tabular}
\endgroup
\end{table}

\begin{table}[htbp]
\centering
\caption{{MRE for sensor A023, October--December 2022 (dimensionless ratio). Each cell gives mean $\pm$ sample SD, followed by the 95\% confidence interval; $n=5$ seed runs. Bold means identify the lowest unrounded mean within each column, not statistical significance.}}
\label{tab:supp_p01_46}
\begingroup
\footnotesize
\setlength{\tabcolsep}{3pt}
\renewcommand{\arraystretch}{1.15}
\begin{tabular}{lccc}
\toprule
Method & Block $B=16$ & Block $B=32$ & Block $B=64$ \\
\midrule
Linear & \makecell[c]{$0.121\,\pm\,0.013$\\$[0.104,\ 0.137]$} & \makecell[c]{$0.462\,\pm\,0.541$\\$[-0.210,\ 1.134]$} & \makecell[c]{$0.517\,\pm\,0.539$\\$[-0.152,\ 1.186]$} \\
Median & \makecell[c]{$0.436\,\pm\,0.523$\\$[-0.214,\ 1.085]$} & \makecell[c]{$1.889\,\pm\,2.543$\\$[-1.268,\ 5.047]$} & \makecell[c]{$1.200\,\pm\,1.379$\\$[-0.512,\ 2.912]$} \\
FFill + BFill & \makecell[c]{$0.250\,\pm\,0.298$\\$[-0.121,\ 0.620]$} & \makecell[c]{$0.988\,\pm\,0.820$\\$[-0.031,\ 2.006]$} & \makecell[c]{$0.795\,\pm\,0.831$\\$[-0.237,\ 1.826]$} \\
Rolling mean & \makecell[c]{$0.140\,\pm\,0.106$\\$[0.008,\ 0.272]$} & \makecell[c]{$0.505\,\pm\,0.535$\\$[-0.159,\ 1.169]$} & \makecell[c]{$0.506\,\pm\,0.503$\\$[-0.119,\ 1.131]$} \\
KNN (scaled) & \makecell[c]{$0.080\,\pm\,0.026$\\$[0.048,\ 0.112]$} & \makecell[c]{$0.162\,\pm\,0.090$\\$[0.050,\ 0.274]$} & \makecell[c]{$0.162\,\pm\,0.098$\\$[0.041,\ 0.284]$} \\
KNN (unscaled) & \makecell[c]{$0.103\,\pm\,0.033$\\$[0.062,\ 0.143]$} & \makecell[c]{$0.191\,\pm\,0.093$\\$[0.075,\ 0.307]$} & \makecell[c]{$0.201\,\pm\,0.128$\\$[0.043,\ 0.360]$} \\
Historical & \makecell[c]{$0.087\,\pm\,0.037$\\$[0.041,\ 0.133]$} & \makecell[c]{$0.408\,\pm\,0.176$\\$[0.189,\ 0.626]$} & \makecell[c]{$0.522\,\pm\,0.652$\\$[-0.288,\ 1.332]$} \\
Graph Laplacian & \makecell[c]{$0.151\,\pm\,0.083$\\$[0.048,\ 0.254]$} & \makecell[c]{$0.456\,\pm\,0.564$\\$[-0.244,\ 1.157]$} & \makecell[c]{$0.324\,\pm\,0.291$\\$[-0.037,\ 0.685]$} \\
Graph Temporal & \makecell[c]{$0.146\,\pm\,0.067$\\$[0.062,\ 0.230]$} & \makecell[c]{$0.446\,\pm\,0.546$\\$[-0.233,\ 1.124]$} & \makecell[c]{$0.307\,\pm\,0.260$\\$[-0.015,\ 0.630]$} \\
Sage-Voronoi & \makecell[c]{$0.168\,\pm\,0.131$\\$[0.006,\ 0.330]$} & \makecell[c]{$0.281\,\pm\,0.290$\\$[-0.079,\ 0.640]$} & \makecell[c]{$0.264\,\pm\,0.170$\\$[0.053,\ 0.475]$} \\
GAIN & \makecell[c]{$0.100\,\pm\,0.040$\\$[0.051,\ 0.150]$} & \makecell[c]{$0.166\,\pm\,0.063$\\$[0.088,\ 0.244]$} & \makecell[c]{$\mathbf{0.133}\,\pm\,0.062$\\$[0.057,\ 0.210]$} \\
GRIN & \makecell[c]{$\mathbf{0.074}\,\pm\,0.028$\\$[0.040,\ 0.109]$} & \makecell[c]{$\mathbf{0.134}\,\pm\,0.056$\\$[0.064,\ 0.203]$} & \makecell[c]{$0.140\,\pm\,0.058$\\$[0.068,\ 0.212]$} \\
\bottomrule
\end{tabular}
\endgroup
\end{table}

\clearpage

\begin{table}[htbp]
\centering
\caption{{MRE for sensor A023, October--December 2022 (dimensionless ratio). Each cell gives mean $\pm$ sample SD, followed by the 95\% confidence interval; $n=5$ seed runs. Bold means identify the lowest unrounded mean within each column, not statistical significance.}}
\label{tab:supp_p01_47}
\begingroup
\footnotesize
\setlength{\tabcolsep}{3pt}
\renewcommand{\arraystretch}{1.15}
\begin{tabular}{lcc}
\toprule
Method & Block $B=128$ & Block $B=144$ \\
\midrule
Linear & \makecell[c]{$1.073\,\pm\,0.596$\\$[0.333,\ 1.814]$} & \makecell[c]{$1.489\,\pm\,0.753$\\$[0.554,\ 2.424]$} \\
Median & \makecell[c]{$0.901\,\pm\,0.395$\\$[0.411,\ 1.392]$} & \makecell[c]{$1.508\,\pm\,0.339$\\$[1.087,\ 1.930]$} \\
FFill + BFill & \makecell[c]{$1.080\,\pm\,0.533$\\$[0.418,\ 1.743]$} & \makecell[c]{$1.428\,\pm\,0.736$\\$[0.514,\ 2.342]$} \\
Rolling mean & \makecell[c]{$1.072\,\pm\,0.597$\\$[0.332,\ 1.813]$} & \makecell[c]{$1.488\,\pm\,0.750$\\$[0.556,\ 2.419]$} \\
KNN (scaled) & \makecell[c]{$0.107\,\pm\,0.022$\\$[0.079,\ 0.134]$} & \makecell[c]{$0.168\,\pm\,0.029$\\$[0.133,\ 0.204]$} \\
KNN (unscaled) & \makecell[c]{$0.140\,\pm\,0.032$\\$[0.100,\ 0.180]$} & \makecell[c]{$0.199\,\pm\,0.041$\\$[0.148,\ 0.250]$} \\
Historical & \makecell[c]{$0.143\,\pm\,0.062$\\$[0.066,\ 0.220]$} & \makecell[c]{$0.364\,\pm\,0.279$\\$[0.018,\ 0.710]$} \\
Graph Laplacian & \makecell[c]{$0.243\,\pm\,0.085$\\$[0.137,\ 0.349]$} & \makecell[c]{$0.384\,\pm\,0.101$\\$[0.259,\ 0.510]$} \\
Graph Temporal & \makecell[c]{$0.232\,\pm\,0.083$\\$[0.130,\ 0.335]$} & \makecell[c]{$0.358\,\pm\,0.086$\\$[0.251,\ 0.465]$} \\
Sage-Voronoi & \makecell[c]{$0.166\,\pm\,0.040$\\$[0.116,\ 0.216]$} & \makecell[c]{$0.238\,\pm\,0.073$\\$[0.147,\ 0.328]$} \\
GAIN & \makecell[c]{$0.106\,\pm\,0.023$\\$[0.077,\ 0.134]$} & \makecell[c]{$0.155\,\pm\,0.023$\\$[0.126,\ 0.183]$} \\
GRIN & \makecell[c]{$\mathbf{0.092}\,\pm\,0.013$\\$[0.076,\ 0.109]$} & \makecell[c]{$\mathbf{0.145}\,\pm\,0.017$\\$[0.123,\ 0.166]$} \\
\bottomrule
\end{tabular}
\endgroup
\end{table}

\begin{table}[htbp]
\centering
\caption{{MRE for sensor A023, October--December 2022 (dimensionless ratio). Each cell gives mean $\pm$ sample SD, followed by the 95\% confidence interval; $n=5$ seed runs. Bold means identify the lowest unrounded mean within each column, not statistical significance.}}
\label{tab:supp_p01_48}
\begingroup
\footnotesize
\setlength{\tabcolsep}{3pt}
\renewcommand{\arraystretch}{1.15}
\begin{tabular}{lcccc}
\toprule
Method & MCAR 10\% & MCAR 30\% & MCAR 50\% & Synchronous $B=144$ \\
\midrule
Linear & \makecell[c]{$0.131\,\pm\,0.012$\\$[0.117,\ 0.145]$} & \makecell[c]{$0.136\,\pm\,0.005$\\$[0.130,\ 0.142]$} & \makecell[c]{$0.135\,\pm\,0.003$\\$[0.131,\ 0.140]$} & \makecell[c]{$0.895\,\pm\,0.518$\\$[0.251,\ 1.538]$} \\
Median & \makecell[c]{$1.254\,\pm\,0.034$\\$[1.212,\ 1.296]$} & \makecell[c]{$1.283\,\pm\,0.041$\\$[1.232,\ 1.335]$} & \makecell[c]{$1.232\,\pm\,0.027$\\$[1.198,\ 1.266]$} & \makecell[c]{$1.054\,\pm\,0.426$\\$[0.525,\ 1.582]$} \\
FFill + BFill & \makecell[c]{$0.160\,\pm\,0.015$\\$[0.142,\ 0.178]$} & \makecell[c]{$0.161\,\pm\,0.005$\\$[0.155,\ 0.167]$} & \makecell[c]{$0.168\,\pm\,0.006$\\$[0.160,\ 0.176]$} & \makecell[c]{$0.890\,\pm\,0.492$\\$[0.278,\ 1.501]$} \\
Rolling mean & \makecell[c]{$0.158\,\pm\,0.014$\\$[0.140,\ 0.175]$} & \makecell[c]{$0.160\,\pm\,0.006$\\$[0.152,\ 0.167]$} & \makecell[c]{$0.163\,\pm\,0.005$\\$[0.157,\ 0.169]$} & \makecell[c]{$0.896\,\pm\,0.517$\\$[0.254,\ 1.538]$} \\
KNN (scaled) & \makecell[c]{$0.146\,\pm\,0.009$\\$[0.135,\ 0.158]$} & \makecell[c]{$0.146\,\pm\,0.003$\\$[0.142,\ 0.150]$} & \makecell[c]{$0.140\,\pm\,0.003$\\$[0.136,\ 0.144]$} & \makecell[c]{$0.895\,\pm\,0.354$\\$[0.456,\ 1.335]$} \\
KNN (unscaled) & \makecell[c]{$0.180\,\pm\,0.014$\\$[0.163,\ 0.198]$} & \makecell[c]{$0.174\,\pm\,0.003$\\$[0.171,\ 0.177]$} & \makecell[c]{$0.163\,\pm\,0.004$\\$[0.158,\ 0.167]$} & \makecell[c]{$0.895\,\pm\,0.354$\\$[0.456,\ 1.335]$} \\
Historical & \makecell[c]{$0.324\,\pm\,0.041$\\$[0.273,\ 0.375]$} & \makecell[c]{$0.327\,\pm\,0.016$\\$[0.307,\ 0.347]$} & \makecell[c]{$0.317\,\pm\,0.009$\\$[0.306,\ 0.328]$} & \makecell[c]{$\mathbf{0.374}\,\pm\,0.212$\\$[0.112,\ 0.637]$} \\
Graph Laplacian & \makecell[c]{$0.319\,\pm\,0.009$\\$[0.308,\ 0.331]$} & \makecell[c]{$0.328\,\pm\,0.010$\\$[0.316,\ 0.340]$} & \makecell[c]{$0.321\,\pm\,0.010$\\$[0.308,\ 0.334]$} & \makecell[c]{$0.895\,\pm\,0.354$\\$[0.456,\ 1.335]$} \\
Graph Temporal & \makecell[c]{$0.266\,\pm\,0.014$\\$[0.249,\ 0.283]$} & \makecell[c]{$0.280\,\pm\,0.007$\\$[0.272,\ 0.288]$} & \makecell[c]{$0.282\,\pm\,0.008$\\$[0.272,\ 0.293]$} & \makecell[c]{$0.890\,\pm\,0.474$\\$[0.301,\ 1.478]$} \\
Sage-Voronoi & \makecell[c]{$0.134\,\pm\,0.016$\\$[0.114,\ 0.154]$} & \makecell[c]{$0.138\,\pm\,0.006$\\$[0.131,\ 0.145]$} & \makecell[c]{$0.145\,\pm\,0.005$\\$[0.140,\ 0.151]$} & \makecell[c]{$0.690\,\pm\,0.369$\\$[0.231,\ 1.149]$} \\
GAIN & \makecell[c]{$0.139\,\pm\,0.014$\\$[0.122,\ 0.157]$} & \makecell[c]{$0.190\,\pm\,0.004$\\$[0.186,\ 0.195]$} & \makecell[c]{$0.313\,\pm\,0.012$\\$[0.299,\ 0.327]$} & \makecell[c]{$0.992\,\pm\,0.403$\\$[0.492,\ 1.493]$} \\
GRIN & \makecell[c]{$\mathbf{0.111}\,\pm\,0.013$\\$[0.095,\ 0.127]$} & \makecell[c]{$\mathbf{0.115}\,\pm\,0.004$\\$[0.110,\ 0.121]$} & \makecell[c]{$\mathbf{0.116}\,\pm\,0.005$\\$[0.109,\ 0.122]$} & \makecell[c]{$0.842\,\pm\,0.481$\\$[0.244,\ 1.439]$} \\
\bottomrule
\end{tabular}
\endgroup
\end{table}

\clearpage

\begin{table}[htbp]
\centering
\caption{{MAE for sensor A027, January 2023 (vehicles per ten-minute interval). Each cell gives mean $\pm$ sample SD, followed by the 95\% confidence interval; $n=5$ seed runs. Bold means identify the lowest unrounded mean within each column, not statistical significance.}}
\label{tab:supp_p01_49}
\begingroup
\footnotesize
\setlength{\tabcolsep}{3pt}
\renewcommand{\arraystretch}{1.15}
\begin{tabular}{lcccc}
\toprule
Method & Block $B=1$ & Block $B=2$ & Block $B=4$ & Block $B=8$ \\
\midrule
Linear & \makecell[c]{$36.90\,\pm\,29.76$\\$[-0.05,\ 73.85]$} & \makecell[c]{$\mathbf{24.20}\,\pm\,6.26$\\$[16.43,\ 31.97]$} & \makecell[c]{$26.34\,\pm\,9.33$\\$[14.76,\ 37.92]$} & \makecell[c]{$28.88\,\pm\,13.73$\\$[11.83,\ 45.92]$} \\
Median & \makecell[c]{$108.40\,\pm\,55.79$\\$[39.12,\ 177.68]$} & \makecell[c]{$176.70\,\pm\,89.13$\\$[66.03,\ 287.37]$} & \makecell[c]{$143.60\,\pm\,68.41$\\$[58.66,\ 228.54]$} & \makecell[c]{$138.40\,\pm\,70.16$\\$[51.28,\ 225.52]$} \\
FFill + BFill & \makecell[c]{$36.80\,\pm\,45.89$\\$[-20.18,\ 93.78]$} & \makecell[c]{$34.10\,\pm\,24.30$\\$[3.93,\ 64.27]$} & \makecell[c]{$48.40\,\pm\,16.42$\\$[28.01,\ 68.79]$} & \makecell[c]{$65.80\,\pm\,44.55$\\$[10.49,\ 121.11]$} \\
Rolling mean & \makecell[c]{$29.00\,\pm\,40.14$\\$[-20.84,\ 78.84]$} & \makecell[c]{$48.40\,\pm\,17.82$\\$[26.27,\ 70.53]$} & \makecell[c]{$59.62\,\pm\,20.49$\\$[34.17,\ 85.06]$} & \makecell[c]{$44.60\,\pm\,30.92$\\$[6.21,\ 82.99]$} \\
KNN (scaled) & \makecell[c]{$149.44\,\pm\,81.67$\\$[48.04,\ 250.84]$} & \makecell[c]{$236.62\,\pm\,113.13$\\$[96.15,\ 377.09]$} & \makecell[c]{$182.10\,\pm\,127.10$\\$[24.28,\ 339.92]$} & \makecell[c]{$165.99\,\pm\,91.72$\\$[52.10,\ 279.87]$} \\
KNN (unscaled) & \makecell[c]{$40.08\,\pm\,39.37$\\$[-8.80,\ 88.96]$} & \makecell[c]{$32.08\,\pm\,25.35$\\$[0.60,\ 63.56]$} & \makecell[c]{$34.36\,\pm\,11.18$\\$[20.48,\ 48.24]$} & \makecell[c]{$31.85\,\pm\,12.99$\\$[15.72,\ 47.98]$} \\
Historical & \makecell[c]{$31.48\,\pm\,26.95$\\$[-1.98,\ 64.94]$} & \makecell[c]{$139.17\,\pm\,78.89$\\$[41.21,\ 237.12]$} & \makecell[c]{$98.45\,\pm\,60.37$\\$[23.49,\ 173.41]$} & \makecell[c]{$60.59\,\pm\,32.58$\\$[20.13,\ 101.05]$} \\
Graph Laplacian & \makecell[c]{$85.13\,\pm\,58.04$\\$[13.07,\ 157.19]$} & \makecell[c]{$85.38\,\pm\,63.78$\\$[6.19,\ 164.57]$} & \makecell[c]{$81.41\,\pm\,36.48$\\$[36.11,\ 126.70]$} & \makecell[c]{$64.88\,\pm\,36.06$\\$[20.10,\ 109.66]$} \\
Graph Temporal & \makecell[c]{$61.79\,\pm\,38.00$\\$[14.60,\ 108.98]$} & \makecell[c]{$133.32\,\pm\,137.35$\\$[-37.23,\ 303.86]$} & \makecell[c]{$54.54\,\pm\,33.47$\\$[12.98,\ 96.09]$} & \makecell[c]{$68.84\,\pm\,50.37$\\$[6.30,\ 131.39]$} \\
Sage-Voronoi & \makecell[c]{$26.51\,\pm\,32.63$\\$[-14.00,\ 67.02]$} & \makecell[c]{$27.90\,\pm\,14.74$\\$[9.59,\ 46.20]$} & \makecell[c]{$38.77\,\pm\,21.77$\\$[11.73,\ 65.80]$} & \makecell[c]{$34.61\,\pm\,15.93$\\$[14.83,\ 54.40]$} \\
GAIN & \makecell[c]{$39.88\,\pm\,23.44$\\$[10.77,\ 68.99]$} & \makecell[c]{$47.66\,\pm\,33.55$\\$[6.01,\ 89.31]$} & \makecell[c]{$31.25\,\pm\,11.94$\\$[16.43,\ 46.07]$} & \makecell[c]{$35.18\,\pm\,14.65$\\$[16.98,\ 53.37]$} \\
GRIN & \makecell[c]{$\mathbf{19.32}\,\pm\,18.21$\\$[-3.29,\ 41.94]$} & \makecell[c]{$28.89\,\pm\,19.55$\\$[4.62,\ 53.16]$} & \makecell[c]{$\mathbf{25.91}\,\pm\,13.07$\\$[9.68,\ 42.14]$} & \makecell[c]{$\mathbf{27.49}\,\pm\,6.48$\\$[19.45,\ 35.53]$} \\
\bottomrule
\end{tabular}
\endgroup
\end{table}

\begin{table}[htbp]
\centering
\caption{{MAE for sensor A027, January 2023 (vehicles per ten-minute interval). Each cell gives mean $\pm$ sample SD, followed by the 95\% confidence interval; $n=5$ seed runs. Bold means identify the lowest unrounded mean within each column, not statistical significance.}}
\label{tab:supp_p01_50}
\begingroup
\footnotesize
\setlength{\tabcolsep}{3pt}
\renewcommand{\arraystretch}{1.15}
\begin{tabular}{lccc}
\toprule
Method & Block $B=16$ & Block $B=32$ & Block $B=64$ \\
\midrule
Linear & \makecell[c]{$35.96\,\pm\,23.03$\\$[7.37,\ 64.55]$} & \makecell[c]{$93.15\,\pm\,57.76$\\$[21.43,\ 164.87]$} & \makecell[c]{$80.00\,\pm\,23.17$\\$[51.22,\ 108.77]$} \\
Median & \makecell[c]{$130.84\,\pm\,55.12$\\$[62.39,\ 199.28]$} & \makecell[c]{$128.14\,\pm\,49.14$\\$[67.12,\ 189.15]$} & \makecell[c]{$146.73\,\pm\,19.25$\\$[122.83,\ 170.63]$} \\
FFill + BFill & \makecell[c]{$54.59\,\pm\,30.60$\\$[16.60,\ 92.58]$} & \makecell[c]{$191.38\,\pm\,113.75$\\$[50.14,\ 332.61]$} & \makecell[c]{$213.74\,\pm\,71.91$\\$[124.45,\ 303.03]$} \\
Rolling mean & \makecell[c]{$35.28\,\pm\,23.17$\\$[6.51,\ 64.05]$} & \makecell[c]{$104.27\,\pm\,65.50$\\$[22.94,\ 185.60]$} & \makecell[c]{$83.67\,\pm\,25.37$\\$[52.16,\ 115.17]$} \\
KNN (scaled) & \makecell[c]{$101.38\,\pm\,70.62$\\$[13.69,\ 189.07]$} & \makecell[c]{$127.07\,\pm\,60.81$\\$[51.56,\ 202.58]$} & \makecell[c]{$162.68\,\pm\,48.07$\\$[102.99,\ 222.37]$} \\
KNN (unscaled) & \makecell[c]{$30.11\,\pm\,15.72$\\$[10.60,\ 49.63]$} & \makecell[c]{$34.86\,\pm\,13.37$\\$[18.26,\ 51.46]$} & \makecell[c]{$35.89\,\pm\,10.92$\\$[22.33,\ 49.45]$} \\
Historical & \makecell[c]{$29.43\,\pm\,6.49$\\$[21.37,\ 37.49]$} & \makecell[c]{$93.26\,\pm\,41.60$\\$[41.61,\ 144.91]$} & \makecell[c]{$57.38\,\pm\,14.98$\\$[38.78,\ 75.98]$} \\
Graph Laplacian & \makecell[c]{$69.86\,\pm\,25.20$\\$[38.56,\ 101.15]$} & \makecell[c]{$78.16\,\pm\,8.74$\\$[67.31,\ 89.02]$} & \makecell[c]{$92.70\,\pm\,16.67$\\$[72.00,\ 113.39]$} \\
Graph Temporal & \makecell[c]{$95.79\,\pm\,78.06$\\$[-1.14,\ 192.71]$} & \makecell[c]{$124.26\,\pm\,97.35$\\$[3.38,\ 245.13]$} & \makecell[c]{$105.24\,\pm\,43.06$\\$[51.78,\ 158.71]$} \\
Sage-Voronoi & \makecell[c]{$44.51\,\pm\,24.97$\\$[13.50,\ 75.52]$} & \makecell[c]{$58.25\,\pm\,14.81$\\$[39.87,\ 76.64]$} & \makecell[c]{$52.44\,\pm\,11.70$\\$[37.91,\ 66.97]$} \\
GAIN & \makecell[c]{$\mathbf{23.79}\,\pm\,11.00$\\$[10.14,\ 37.45]$} & \makecell[c]{$\mathbf{24.07}\,\pm\,9.11$\\$[12.76,\ 35.38]$} & \makecell[c]{$\mathbf{24.12}\,\pm\,3.70$\\$[19.53,\ 28.71]$} \\
GRIN & \makecell[c]{$33.98\,\pm\,29.59$\\$[-2.76,\ 70.72]$} & \makecell[c]{$55.19\,\pm\,50.14$\\$[-7.07,\ 117.45]$} & \makecell[c]{$44.54\,\pm\,26.34$\\$[11.84,\ 77.25]$} \\
\bottomrule
\end{tabular}
\endgroup
\end{table}

\clearpage

\begin{table}[htbp]
\centering
\caption{{MAE for sensor A027, January 2023 (vehicles per ten-minute interval). Each cell gives mean $\pm$ sample SD, followed by the 95\% confidence interval; $n=5$ seed runs. Bold means identify the lowest unrounded mean within each column, not statistical significance.}}
\label{tab:supp_p01_51}
\begingroup
\footnotesize
\setlength{\tabcolsep}{3pt}
\renewcommand{\arraystretch}{1.15}
\begin{tabular}{lcc}
\toprule
Method & Block $B=128$ & Block $B=144$ \\
\midrule
Linear & \makecell[c]{$208.26\,\pm\,78.89$\\$[110.31,\ 306.22]$} & \makecell[c]{$203.13\,\pm\,49.62$\\$[141.51,\ 264.74]$} \\
Median & \makecell[c]{$147.52\,\pm\,10.27$\\$[134.77,\ 160.27]$} & \makecell[c]{$148.14\,\pm\,14.22$\\$[130.49,\ 165.79]$} \\
FFill + BFill & \makecell[c]{$227.00\,\pm\,56.32$\\$[157.06,\ 296.93]$} & \makecell[c]{$204.96\,\pm\,58.51$\\$[132.31,\ 277.61]$} \\
Rolling mean & \makecell[c]{$208.30\,\pm\,78.59$\\$[110.72,\ 305.89]$} & \makecell[c]{$203.08\,\pm\,49.92$\\$[141.09,\ 265.06]$} \\
KNN (scaled) & \makecell[c]{$164.57\,\pm\,15.29$\\$[145.59,\ 183.55]$} & \makecell[c]{$148.46\,\pm\,31.91$\\$[108.84,\ 188.08]$} \\
KNN (unscaled) & \makecell[c]{$37.44\,\pm\,3.25$\\$[33.40,\ 41.47]$} & \makecell[c]{$26.87\,\pm\,4.77$\\$[20.95,\ 32.80]$} \\
Historical & \makecell[c]{$52.21\,\pm\,6.89$\\$[43.65,\ 60.77]$} & \makecell[c]{$57.47\,\pm\,4.69$\\$[51.65,\ 63.29]$} \\
Graph Laplacian & \makecell[c]{$98.97\,\pm\,13.68$\\$[81.98,\ 115.96]$} & \makecell[c]{$82.90\,\pm\,14.34$\\$[65.10,\ 100.70]$} \\
Graph Temporal & \makecell[c]{$119.39\,\pm\,60.94$\\$[43.72,\ 195.06]$} & \makecell[c]{$90.10\,\pm\,17.24$\\$[68.69,\ 111.51]$} \\
Sage-Voronoi & \makecell[c]{$66.60\,\pm\,10.10$\\$[54.07,\ 79.14]$} & \makecell[c]{$51.60\,\pm\,7.48$\\$[42.31,\ 60.89]$} \\
GAIN & \makecell[c]{$\mathbf{27.03}\,\pm\,1.71$\\$[24.90,\ 29.15]$} & \makecell[c]{$\mathbf{24.24}\,\pm\,2.28$\\$[21.41,\ 27.07]$} \\
GRIN & \makecell[c]{$46.45\,\pm\,32.45$\\$[6.16,\ 86.74]$} & \makecell[c]{$36.72\,\pm\,13.75$\\$[19.64,\ 53.80]$} \\
\bottomrule
\end{tabular}
\endgroup
\end{table}

\begin{table}[htbp]
\centering
\caption{{MAE for sensor A027, January 2023 (vehicles per ten-minute interval). Each cell gives mean $\pm$ sample SD, followed by the 95\% confidence interval; $n=5$ seed runs. Bold means identify the lowest unrounded mean within each column, not statistical significance.}}
\label{tab:supp_p01_52}
\begingroup
\footnotesize
\setlength{\tabcolsep}{3pt}
\renewcommand{\arraystretch}{1.15}
\begin{tabular}{lcccc}
\toprule
Method & MCAR 10\% & MCAR 30\% & MCAR 50\% & Synchronous $B=144$ \\
\midrule
Linear & \makecell[c]{$22.26\,\pm\,1.78$\\$[20.06,\ 24.47]$} & \makecell[c]{$23.17\,\pm\,0.81$\\$[22.16,\ 24.17]$} & \makecell[c]{$\mathbf{22.98}\,\pm\,0.82$\\$[21.96,\ 23.99]$} & \makecell[c]{$154.13\,\pm\,16.06$\\$[134.19,\ 174.07]$} \\
Median & \makecell[c]{$146.12\,\pm\,6.35$\\$[138.24,\ 154.00]$} & \makecell[c]{$145.01\,\pm\,1.76$\\$[142.82,\ 147.19]$} & \makecell[c]{$146.42\,\pm\,1.09$\\$[145.06,\ 147.77]$} & \makecell[c]{$147.67\,\pm\,16.56$\\$[127.10,\ 168.23]$} \\
FFill + BFill & \makecell[c]{$27.47\,\pm\,3.40$\\$[23.26,\ 31.69]$} & \makecell[c]{$29.66\,\pm\,1.01$\\$[28.40,\ 30.91]$} & \makecell[c]{$32.05\,\pm\,2.25$\\$[29.25,\ 34.85]$} & \makecell[c]{$163.17\,\pm\,25.71$\\$[131.25,\ 195.09]$} \\
Rolling mean & \makecell[c]{$30.40\,\pm\,2.05$\\$[27.86,\ 32.94]$} & \makecell[c]{$30.89\,\pm\,0.49$\\$[30.29,\ 31.50]$} & \makecell[c]{$31.00\,\pm\,1.08$\\$[29.66,\ 32.34]$} & \makecell[c]{$154.86\,\pm\,16.23$\\$[134.71,\ 175.02]$} \\
KNN (scaled) & \makecell[c]{$136.86\,\pm\,2.86$\\$[133.30,\ 140.41]$} & \makecell[c]{$130.16\,\pm\,5.39$\\$[123.47,\ 136.85]$} & \makecell[c]{$111.84\,\pm\,3.09$\\$[108.01,\ 115.68]$} & \makecell[c]{$150.19\,\pm\,19.41$\\$[126.08,\ 174.29]$} \\
KNN (unscaled) & \makecell[c]{$30.40\,\pm\,2.89$\\$[26.81,\ 33.98]$} & \makecell[c]{$33.19\,\pm\,0.98$\\$[31.97,\ 34.40]$} & \makecell[c]{$34.29\,\pm\,0.93$\\$[33.13,\ 35.44]$} & \makecell[c]{$150.19\,\pm\,19.41$\\$[126.08,\ 174.29]$} \\
Historical & \makecell[c]{$57.92\,\pm\,4.11$\\$[52.82,\ 63.02]$} & \makecell[c]{$57.88\,\pm\,1.80$\\$[55.64,\ 60.11]$} & \makecell[c]{$56.35\,\pm\,0.99$\\$[55.11,\ 57.58]$} & \makecell[c]{$\mathbf{57.54}\,\pm\,3.63$\\$[53.04,\ 62.05]$} \\
Graph Laplacian & \makecell[c]{$80.28\,\pm\,3.80$\\$[75.57,\ 85.00]$} & \makecell[c]{$78.64\,\pm\,1.50$\\$[76.78,\ 80.51]$} & \makecell[c]{$76.90\,\pm\,2.03$\\$[74.37,\ 79.42]$} & \makecell[c]{$150.19\,\pm\,19.41$\\$[126.08,\ 174.29]$} \\
Graph Temporal & \makecell[c]{$80.22\,\pm\,9.91$\\$[67.93,\ 92.52]$} & \makecell[c]{$87.26\,\pm\,6.19$\\$[79.57,\ 94.95]$} & \makecell[c]{$94.08\,\pm\,6.15$\\$[86.44,\ 101.71]$} & \makecell[c]{$198.23\,\pm\,63.63$\\$[119.22,\ 277.24]$} \\
Sage-Voronoi & \makecell[c]{$25.56\,\pm\,2.23$\\$[22.79,\ 28.32]$} & \makecell[c]{$27.57\,\pm\,1.09$\\$[26.22,\ 28.92]$} & \makecell[c]{$29.96\,\pm\,1.74$\\$[27.80,\ 32.13]$} & \makecell[c]{$156.16\,\pm\,37.34$\\$[109.80,\ 202.53]$} \\
GAIN & \makecell[c]{$26.06\,\pm\,2.17$\\$[23.36,\ 28.75]$} & \makecell[c]{$34.48\,\pm\,0.92$\\$[33.34,\ 35.62]$} & \makecell[c]{$48.33\,\pm\,2.32$\\$[45.45,\ 51.21]$} & \makecell[c]{$143.79\,\pm\,7.86$\\$[134.03,\ 153.55]$} \\
GRIN & \makecell[c]{$\mathbf{20.05}\,\pm\,1.62$\\$[18.03,\ 22.07]$} & \makecell[c]{$\mathbf{22.18}\,\pm\,1.09$\\$[20.83,\ 23.53]$} & \makecell[c]{$24.42\,\pm\,3.10$\\$[20.57,\ 28.27]$} & \makecell[c]{$154.93\,\pm\,23.91$\\$[125.23,\ 184.62]$} \\
\bottomrule
\end{tabular}
\endgroup
\end{table}

\clearpage

\begin{table}[htbp]
\centering
\caption{{RMSE for sensor A027, January 2023 (vehicles per ten-minute interval). Each cell gives mean $\pm$ sample SD, followed by the 95\% confidence interval; $n=5$ seed runs. Bold means identify the lowest unrounded mean within each column, not statistical significance.}}
\label{tab:supp_p01_53}
\begingroup
\footnotesize
\setlength{\tabcolsep}{3pt}
\renewcommand{\arraystretch}{1.15}
\begin{tabular}{lcccc}
\toprule
Method & Block $B=1$ & Block $B=2$ & Block $B=4$ & Block $B=8$ \\
\midrule
Linear & \makecell[c]{$36.90\,\pm\,29.76$\\$[-0.05,\ 73.85]$} & \makecell[c]{$\mathbf{26.81}\,\pm\,7.80$\\$[17.13,\ 36.49]$} & \makecell[c]{$\mathbf{30.35}\,\pm\,10.35$\\$[17.50,\ 43.20]$} & \makecell[c]{$33.76\,\pm\,14.90$\\$[15.25,\ 52.26]$} \\
Median & \makecell[c]{$108.40\,\pm\,55.79$\\$[39.12,\ 177.68]$} & \makecell[c]{$177.07\,\pm\,89.33$\\$[66.15,\ 287.99]$} & \makecell[c]{$147.83\,\pm\,66.93$\\$[64.73,\ 230.93]$} & \makecell[c]{$145.78\,\pm\,65.12$\\$[64.92,\ 226.64]$} \\
FFill + BFill & \makecell[c]{$36.80\,\pm\,45.89$\\$[-20.18,\ 93.78]$} & \makecell[c]{$36.41\,\pm\,24.25$\\$[6.29,\ 66.52]$} & \makecell[c]{$57.32\,\pm\,18.06$\\$[34.90,\ 79.74]$} & \makecell[c]{$73.94\,\pm\,49.31$\\$[12.71,\ 135.17]$} \\
Rolling mean & \makecell[c]{$29.00\,\pm\,40.14$\\$[-20.84,\ 78.84]$} & \makecell[c]{$50.05\,\pm\,17.92$\\$[27.80,\ 72.30]$} & \makecell[c]{$65.22\,\pm\,21.38$\\$[38.67,\ 91.77]$} & \makecell[c]{$50.29\,\pm\,32.16$\\$[10.36,\ 90.21]$} \\
KNN (scaled) & \makecell[c]{$149.44\,\pm\,81.67$\\$[48.04,\ 250.84]$} & \makecell[c]{$238.51\,\pm\,111.10$\\$[100.56,\ 376.45]$} & \makecell[c]{$186.54\,\pm\,124.52$\\$[31.93,\ 341.15]$} & \makecell[c]{$175.87\,\pm\,88.61$\\$[65.85,\ 285.89]$} \\
KNN (unscaled) & \makecell[c]{$40.08\,\pm\,39.37$\\$[-8.80,\ 88.96]$} & \makecell[c]{$32.72\,\pm\,25.16$\\$[1.48,\ 63.96]$} & \makecell[c]{$44.19\,\pm\,11.70$\\$[29.66,\ 58.72]$} & \makecell[c]{$38.05\,\pm\,15.70$\\$[18.56,\ 57.54]$} \\
Historical & \makecell[c]{$31.48\,\pm\,26.95$\\$[-1.98,\ 64.94]$} & \makecell[c]{$139.64\,\pm\,78.92$\\$[41.64,\ 237.64]$} & \makecell[c]{$105.50\,\pm\,54.52$\\$[37.81,\ 173.19]$} & \makecell[c]{$68.77\,\pm\,32.81$\\$[28.03,\ 109.52]$} \\
Graph Laplacian & \makecell[c]{$85.13\,\pm\,58.04$\\$[13.07,\ 157.19]$} & \makecell[c]{$88.03\,\pm\,64.52$\\$[7.91,\ 168.14]$} & \makecell[c]{$94.59\,\pm\,39.90$\\$[45.05,\ 144.12]$} & \makecell[c]{$70.70\,\pm\,38.46$\\$[22.95,\ 118.46]$} \\
Graph Temporal & \makecell[c]{$61.79\,\pm\,38.00$\\$[14.60,\ 108.98]$} & \makecell[c]{$136.31\,\pm\,135.63$\\$[-32.10,\ 304.72]$} & \makecell[c]{$59.81\,\pm\,33.19$\\$[18.60,\ 101.02]$} & \makecell[c]{$77.94\,\pm\,53.55$\\$[11.45,\ 144.42]$} \\
Sage-Voronoi & \makecell[c]{$26.51\,\pm\,32.63$\\$[-14.00,\ 67.02]$} & \makecell[c]{$29.16\,\pm\,15.59$\\$[9.81,\ 48.51]$} & \makecell[c]{$44.12\,\pm\,24.40$\\$[13.83,\ 74.42]$} & \makecell[c]{$40.45\,\pm\,17.37$\\$[18.88,\ 62.01]$} \\
GAIN & \makecell[c]{$39.88\,\pm\,23.44$\\$[10.77,\ 68.99]$} & \makecell[c]{$49.33\,\pm\,32.75$\\$[8.67,\ 89.99]$} & \makecell[c]{$38.16\,\pm\,14.79$\\$[19.79,\ 56.52]$} & \makecell[c]{$42.94\,\pm\,17.25$\\$[21.53,\ 64.35]$} \\
GRIN & \makecell[c]{$\mathbf{19.32}\,\pm\,18.21$\\$[-3.29,\ 41.94]$} & \makecell[c]{$31.27\,\pm\,18.72$\\$[8.03,\ 54.51]$} & \makecell[c]{$30.38\,\pm\,14.02$\\$[12.97,\ 47.78]$} & \makecell[c]{$\mathbf{32.38}\,\pm\,6.99$\\$[23.70,\ 41.05]$} \\
\bottomrule
\end{tabular}
\endgroup
\end{table}

\begin{table}[htbp]
\centering
\caption{{RMSE for sensor A027, January 2023 (vehicles per ten-minute interval). Each cell gives mean $\pm$ sample SD, followed by the 95\% confidence interval; $n=5$ seed runs. Bold means identify the lowest unrounded mean within each column, not statistical significance.}}
\label{tab:supp_p01_54}
\begingroup
\footnotesize
\setlength{\tabcolsep}{3pt}
\renewcommand{\arraystretch}{1.15}
\begin{tabular}{lccc}
\toprule
Method & Block $B=16$ & Block $B=32$ & Block $B=64$ \\
\midrule
Linear & \makecell[c]{$44.33\,\pm\,28.56$\\$[8.87,\ 79.79]$} & \makecell[c]{$106.68\,\pm\,66.68$\\$[23.89,\ 189.47]$} & \makecell[c]{$98.93\,\pm\,32.01$\\$[59.19,\ 138.68]$} \\
Median & \makecell[c]{$138.83\,\pm\,51.70$\\$[74.64,\ 203.02]$} & \makecell[c]{$140.45\,\pm\,45.84$\\$[83.53,\ 197.37]$} & \makecell[c]{$159.83\,\pm\,17.61$\\$[137.97,\ 181.69]$} \\
FFill + BFill & \makecell[c]{$65.15\,\pm\,36.74$\\$[19.52,\ 110.77]$} & \makecell[c]{$208.16\,\pm\,118.66$\\$[60.83,\ 355.49]$} & \makecell[c]{$244.09\,\pm\,69.76$\\$[157.47,\ 330.70]$} \\
Rolling mean & \makecell[c]{$42.79\,\pm\,30.11$\\$[5.40,\ 80.18]$} & \makecell[c]{$117.86\,\pm\,73.44$\\$[26.67,\ 209.06]$} & \makecell[c]{$102.64\,\pm\,35.11$\\$[59.04,\ 146.24]$} \\
KNN (scaled) & \makecell[c]{$112.99\,\pm\,73.71$\\$[21.47,\ 204.51]$} & \makecell[c]{$147.28\,\pm\,67.81$\\$[63.08,\ 231.47]$} & \makecell[c]{$185.23\,\pm\,39.78$\\$[135.84,\ 234.62]$} \\
KNN (unscaled) & \makecell[c]{$39.77\,\pm\,21.72$\\$[12.80,\ 66.73]$} & \makecell[c]{$43.08\,\pm\,14.83$\\$[24.66,\ 61.50]$} & \makecell[c]{$44.31\,\pm\,11.76$\\$[29.70,\ 58.91]$} \\
Historical & \makecell[c]{$38.15\,\pm\,15.44$\\$[18.98,\ 57.31]$} & \makecell[c]{$106.55\,\pm\,47.68$\\$[47.35,\ 165.75]$} & \makecell[c]{$71.34\,\pm\,19.44$\\$[47.20,\ 95.49]$} \\
Graph Laplacian & \makecell[c]{$76.52\,\pm\,28.71$\\$[40.88,\ 112.17]$} & \makecell[c]{$91.04\,\pm\,8.16$\\$[80.91,\ 101.16]$} & \makecell[c]{$104.18\,\pm\,17.59$\\$[82.35,\ 126.02]$} \\
Graph Temporal & \makecell[c]{$110.03\,\pm\,101.54$\\$[-16.05,\ 236.11]$} & \makecell[c]{$162.83\,\pm\,131.88$\\$[-0.92,\ 326.58]$} & \makecell[c]{$140.75\,\pm\,79.19$\\$[42.42,\ 239.08]$} \\
Sage-Voronoi & \makecell[c]{$53.60\,\pm\,27.03$\\$[20.03,\ 87.17]$} & \makecell[c]{$67.72\,\pm\,13.05$\\$[51.51,\ 83.92]$} & \makecell[c]{$62.23\,\pm\,15.29$\\$[43.24,\ 81.22]$} \\
GAIN & \makecell[c]{$\mathbf{31.70}\,\pm\,18.05$\\$[9.29,\ 54.11]$} & \makecell[c]{$\mathbf{29.64}\,\pm\,11.28$\\$[15.64,\ 43.65]$} & \makecell[c]{$\mathbf{30.82}\,\pm\,5.15$\\$[24.42,\ 37.22]$} \\
GRIN & \makecell[c]{$38.85\,\pm\,31.49$\\$[-0.26,\ 77.95]$} & \makecell[c]{$63.20\,\pm\,52.13$\\$[-1.52,\ 127.92]$} & \makecell[c]{$61.72\,\pm\,45.33$\\$[5.44,\ 118.01]$} \\
\bottomrule
\end{tabular}
\endgroup
\end{table}

\clearpage

\begin{table}[htbp]
\centering
\caption{{RMSE for sensor A027, January 2023 (vehicles per ten-minute interval). Each cell gives mean $\pm$ sample SD, followed by the 95\% confidence interval; $n=5$ seed runs. Bold means identify the lowest unrounded mean within each column, not statistical significance.}}
\label{tab:supp_p01_55}
\begingroup
\footnotesize
\setlength{\tabcolsep}{3pt}
\renewcommand{\arraystretch}{1.15}
\begin{tabular}{lcc}
\toprule
Method & Block $B=128$ & Block $B=144$ \\
\midrule
Linear & \makecell[c]{$249.50\,\pm\,86.08$\\$[142.62,\ 356.39]$} & \makecell[c]{$238.85\,\pm\,57.05$\\$[168.01,\ 309.69]$} \\
Median & \makecell[c]{$162.86\,\pm\,10.56$\\$[149.75,\ 175.97]$} & \makecell[c]{$164.31\,\pm\,10.95$\\$[150.71,\ 177.91]$} \\
FFill + BFill & \makecell[c]{$272.99\,\pm\,45.27$\\$[216.78,\ 329.21]$} & \makecell[c]{$238.79\,\pm\,68.29$\\$[154.00,\ 323.58]$} \\
Rolling mean & \makecell[c]{$249.55\,\pm\,85.92$\\$[142.86,\ 356.23]$} & \makecell[c]{$238.84\,\pm\,57.34$\\$[167.64,\ 310.04]$} \\
KNN (scaled) & \makecell[c]{$188.23\,\pm\,15.25$\\$[169.29,\ 207.17]$} & \makecell[c]{$176.00\,\pm\,37.61$\\$[129.30,\ 222.69]$} \\
KNN (unscaled) & \makecell[c]{$48.20\,\pm\,4.53$\\$[42.58,\ 53.83]$} & \makecell[c]{$35.23\,\pm\,6.02$\\$[27.75,\ 42.71]$} \\
Historical & \makecell[c]{$68.27\,\pm\,9.04$\\$[57.04,\ 79.50]$} & \makecell[c]{$77.21\,\pm\,7.57$\\$[67.82,\ 86.61]$} \\
Graph Laplacian & \makecell[c]{$109.77\,\pm\,13.50$\\$[93.00,\ 126.54]$} & \makecell[c]{$96.55\,\pm\,15.54$\\$[77.26,\ 115.84]$} \\
Graph Temporal & \makecell[c]{$151.15\,\pm\,82.68$\\$[48.49,\ 253.80]$} & \makecell[c]{$120.52\,\pm\,32.93$\\$[79.63,\ 161.41]$} \\
Sage-Voronoi & \makecell[c]{$80.44\,\pm\,12.80$\\$[64.55,\ 96.33]$} & \makecell[c]{$64.33\,\pm\,8.10$\\$[54.27,\ 74.39]$} \\
GAIN & \makecell[c]{$\mathbf{34.70}\,\pm\,2.78$\\$[31.25,\ 38.15]$} & \makecell[c]{$\mathbf{32.38}\,\pm\,2.82$\\$[28.87,\ 35.88]$} \\
GRIN & \makecell[c]{$61.33\,\pm\,46.56$\\$[3.52,\ 119.14]$} & \makecell[c]{$48.23\,\pm\,21.74$\\$[21.23,\ 75.22]$} \\
\bottomrule
\end{tabular}
\endgroup
\end{table}

\begin{table}[htbp]
\centering
\caption{{RMSE for sensor A027, January 2023 (vehicles per ten-minute interval). Each cell gives mean $\pm$ sample SD, followed by the 95\% confidence interval; $n=5$ seed runs. Bold means identify the lowest unrounded mean within each column, not statistical significance.}}
\label{tab:supp_p01_56}
\begingroup
\footnotesize
\setlength{\tabcolsep}{3pt}
\renewcommand{\arraystretch}{1.15}
\begin{tabular}{lcccc}
\toprule
Method & MCAR 10\% & MCAR 30\% & MCAR 50\% & Synchronous $B=144$ \\
\midrule
Linear & \makecell[c]{$29.08\,\pm\,2.47$\\$[26.01,\ 32.14]$} & \makecell[c]{$31.39\,\pm\,0.78$\\$[30.42,\ 32.37]$} & \makecell[c]{$\mathbf{31.21}\,\pm\,1.40$\\$[29.47,\ 32.96]$} & \makecell[c]{$196.01\,\pm\,29.13$\\$[159.84,\ 232.18]$} \\
Median & \makecell[c]{$163.15\,\pm\,5.12$\\$[156.79,\ 169.51]$} & \makecell[c]{$161.73\,\pm\,1.00$\\$[160.49,\ 162.97]$} & \makecell[c]{$162.69\,\pm\,1.02$\\$[161.43,\ 163.95]$} & \makecell[c]{$163.71\,\pm\,12.85$\\$[147.75,\ 179.66]$} \\
FFill + BFill & \makecell[c]{$36.05\,\pm\,3.87$\\$[31.24,\ 40.85]$} & \makecell[c]{$40.26\,\pm\,1.86$\\$[37.95,\ 42.58]$} & \makecell[c]{$43.70\,\pm\,3.99$\\$[38.75,\ 48.66]$} & \makecell[c]{$208.73\,\pm\,37.68$\\$[161.94,\ 255.52]$} \\
Rolling mean & \makecell[c]{$40.39\,\pm\,2.64$\\$[37.11,\ 43.67]$} & \makecell[c]{$41.60\,\pm\,0.83$\\$[40.57,\ 42.63]$} & \makecell[c]{$41.48\,\pm\,1.17$\\$[40.03,\ 42.94]$} & \makecell[c]{$196.58\,\pm\,28.88$\\$[160.72,\ 232.45]$} \\
KNN (scaled) & \makecell[c]{$170.94\,\pm\,2.63$\\$[167.68,\ 174.21]$} & \makecell[c]{$164.83\,\pm\,5.36$\\$[158.18,\ 171.49]$} & \makecell[c]{$150.16\,\pm\,2.87$\\$[146.60,\ 153.72]$} & \makecell[c]{$165.25\,\pm\,17.24$\\$[143.84,\ 186.65]$} \\
KNN (unscaled) & \makecell[c]{$40.43\,\pm\,2.95$\\$[36.77,\ 44.09]$} & \makecell[c]{$44.69\,\pm\,1.66$\\$[42.63,\ 46.75]$} & \makecell[c]{$48.59\,\pm\,1.43$\\$[46.81,\ 50.37]$} & \makecell[c]{$165.25\,\pm\,17.24$\\$[143.84,\ 186.65]$} \\
Historical & \makecell[c]{$78.80\,\pm\,4.61$\\$[73.08,\ 84.52]$} & \makecell[c]{$77.07\,\pm\,3.02$\\$[73.31,\ 80.82]$} & \makecell[c]{$75.25\,\pm\,1.32$\\$[73.62,\ 76.89]$} & \makecell[c]{$\mathbf{75.51}\,\pm\,8.10$\\$[65.45,\ 85.57]$} \\
Graph Laplacian & \makecell[c]{$97.07\,\pm\,3.14$\\$[93.17,\ 100.96]$} & \makecell[c]{$98.69\,\pm\,2.96$\\$[95.02,\ 102.36]$} & \makecell[c]{$101.39\,\pm\,2.62$\\$[98.14,\ 104.64]$} & \makecell[c]{$165.25\,\pm\,17.24$\\$[143.84,\ 186.65]$} \\
Graph Temporal & \makecell[c]{$123.35\,\pm\,15.36$\\$[104.28,\ 142.41]$} & \makecell[c]{$133.59\,\pm\,14.87$\\$[115.12,\ 152.06]$} & \makecell[c]{$137.97\,\pm\,9.71$\\$[125.91,\ 150.03]$} & \makecell[c]{$243.72\,\pm\,72.89$\\$[153.22,\ 334.22]$} \\
Sage-Voronoi & \makecell[c]{$33.40\,\pm\,2.85$\\$[29.86,\ 36.94]$} & \makecell[c]{$36.82\,\pm\,2.65$\\$[33.53,\ 40.11]$} & \makecell[c]{$39.33\,\pm\,2.91$\\$[35.71,\ 42.94]$} & \makecell[c]{$199.33\,\pm\,48.74$\\$[138.81,\ 259.84]$} \\
GAIN & \makecell[c]{$36.28\,\pm\,3.30$\\$[32.18,\ 40.38]$} & \makecell[c]{$44.56\,\pm\,1.66$\\$[42.50,\ 46.62]$} & \makecell[c]{$58.88\,\pm\,3.19$\\$[54.91,\ 62.84]$} & \makecell[c]{$177.46\,\pm\,10.70$\\$[164.17,\ 190.75]$} \\
GRIN & \makecell[c]{$\mathbf{26.99}\,\pm\,2.26$\\$[24.19,\ 29.79]$} & \makecell[c]{$\mathbf{29.80}\,\pm\,1.68$\\$[27.72,\ 31.89]$} & \makecell[c]{$33.49\,\pm\,5.37$\\$[26.82,\ 40.17]$} & \makecell[c]{$175.21\,\pm\,22.86$\\$[146.82,\ 203.59]$} \\
\bottomrule
\end{tabular}
\endgroup
\end{table}

\clearpage

\begin{table}[htbp]
\centering
\caption{{MRE for sensor A027, January 2023 (dimensionless ratio). Each cell gives mean $\pm$ sample SD, followed by the 95\% confidence interval; $n=5$ seed runs. Bold means identify the lowest unrounded mean within each column, not statistical significance.}}
\label{tab:supp_p01_57}
\begingroup
\footnotesize
\setlength{\tabcolsep}{3pt}
\renewcommand{\arraystretch}{1.15}
\begin{tabular}{lcccc}
\toprule
Method & Block $B=1$ & Block $B=2$ & Block $B=4$ & Block $B=8$ \\
\midrule
Linear & \makecell[c]{$0.170\,\pm\,0.146$\\$[-0.011,\ 0.352]$} & \makecell[c]{$\mathbf{0.059}\,\pm\,0.017$\\$[0.037,\ 0.080]$} & \makecell[c]{$\mathbf{0.097}\,\pm\,0.051$\\$[0.034,\ 0.161]$} & \makecell[c]{$\mathbf{0.104}\,\pm\,0.035$\\$[0.061,\ 0.147]$} \\
Median & \makecell[c]{$0.656\,\pm\,0.894$\\$[-0.454,\ 1.766]$} & \makecell[c]{$0.393\,\pm\,0.123$\\$[0.241,\ 0.546]$} & \makecell[c]{$0.527\,\pm\,0.386$\\$[0.047,\ 1.007]$} & \makecell[c]{$0.998\,\pm\,1.555$\\$[-0.933,\ 2.929]$} \\
FFill + BFill & \makecell[c]{$0.125\,\pm\,0.137$\\$[-0.046,\ 0.295]$} & \makecell[c]{$0.080\,\pm\,0.062$\\$[0.004,\ 0.157]$} & \makecell[c]{$0.204\,\pm\,0.204$\\$[-0.050,\ 0.457]$} & \makecell[c]{$0.234\,\pm\,0.107$\\$[0.101,\ 0.366]$} \\
Rolling mean & \makecell[c]{$0.106\,\pm\,0.130$\\$[-0.055,\ 0.267]$} & \makecell[c]{$0.113\,\pm\,0.031$\\$[0.074,\ 0.152]$} & \makecell[c]{$0.257\,\pm\,0.235$\\$[-0.034,\ 0.549]$} & \makecell[c]{$0.161\,\pm\,0.078$\\$[0.064,\ 0.258]$} \\
KNN (scaled) & \makecell[c]{$0.495\,\pm\,0.099$\\$[0.371,\ 0.618]$} & \makecell[c]{$0.537\,\pm\,0.166$\\$[0.332,\ 0.743]$} & \makecell[c]{$0.471\,\pm\,0.209$\\$[0.211,\ 0.731]$} & \makecell[c]{$0.649\,\pm\,0.381$\\$[0.175,\ 1.122]$} \\
KNN (unscaled) & \makecell[c]{$0.162\,\pm\,0.135$\\$[-0.006,\ 0.329]$} & \makecell[c]{$0.070\,\pm\,0.040$\\$[0.020,\ 0.120]$} & \makecell[c]{$0.130\,\pm\,0.070$\\$[0.043,\ 0.217]$} & \makecell[c]{$0.122\,\pm\,0.052$\\$[0.057,\ 0.187]$} \\
Historical & \makecell[c]{$0.095\,\pm\,0.065$\\$[0.014,\ 0.176]$} & \makecell[c]{$0.306\,\pm\,0.124$\\$[0.152,\ 0.461]$} & \makecell[c]{$0.285\,\pm\,0.096$\\$[0.166,\ 0.404]$} & \makecell[c]{$0.236\,\pm\,0.137$\\$[0.067,\ 0.406]$} \\
Graph Laplacian & \makecell[c]{$0.367\,\pm\,0.267$\\$[0.036,\ 0.698]$} & \makecell[c]{$0.189\,\pm\,0.112$\\$[0.049,\ 0.328]$} & \makecell[c]{$0.280\,\pm\,0.096$\\$[0.161,\ 0.399]$} & \makecell[c]{$0.434\,\pm\,0.597$\\$[-0.307,\ 1.174]$} \\
Graph Temporal & \makecell[c]{$0.215\,\pm\,0.089$\\$[0.105,\ 0.325]$} & \makecell[c]{$0.300\,\pm\,0.322$\\$[-0.100,\ 0.699]$} & \makecell[c]{$0.168\,\pm\,0.046$\\$[0.112,\ 0.225]$} & \makecell[c]{$0.528\,\pm\,0.826$\\$[-0.498,\ 1.554]$} \\
Sage-Voronoi & \makecell[c]{$0.112\,\pm\,0.114$\\$[-0.030,\ 0.254]$} & \makecell[c]{$0.064\,\pm\,0.030$\\$[0.027,\ 0.101]$} & \makecell[c]{$0.139\,\pm\,0.096$\\$[0.019,\ 0.258]$} & \makecell[c]{$0.136\,\pm\,0.059$\\$[0.063,\ 0.209]$} \\
GAIN & \makecell[c]{$0.175\,\pm\,0.122$\\$[0.024,\ 0.326]$} & \makecell[c]{$0.110\,\pm\,0.068$\\$[0.026,\ 0.195]$} & \makecell[c]{$0.109\,\pm\,0.056$\\$[0.040,\ 0.179]$} & \makecell[c]{$0.126\,\pm\,0.039$\\$[0.077,\ 0.174]$} \\
GRIN & \makecell[c]{$\mathbf{0.082}\,\pm\,0.069$\\$[-0.003,\ 0.168]$} & \makecell[c]{$0.067\,\pm\,0.046$\\$[0.010,\ 0.125]$} & \makecell[c]{$0.098\,\pm\,0.056$\\$[0.029,\ 0.168]$} & \makecell[c]{$0.162\,\pm\,0.181$\\$[-0.063,\ 0.387]$} \\
\bottomrule
\end{tabular}
\endgroup
\end{table}

\begin{table}[htbp]
\centering
\caption{{MRE for sensor A027, January 2023 (dimensionless ratio). Each cell gives mean $\pm$ sample SD, followed by the 95\% confidence interval; $n=5$ seed runs. Bold means identify the lowest unrounded mean within each column, not statistical significance.}}
\label{tab:supp_p01_58}
\begingroup
\footnotesize
\setlength{\tabcolsep}{3pt}
\renewcommand{\arraystretch}{1.15}
\begin{tabular}{lccc}
\toprule
Method & Block $B=16$ & Block $B=32$ & Block $B=64$ \\
\midrule
Linear & \makecell[c]{$0.183\,\pm\,0.048$\\$[0.123,\ 0.243]$} & \makecell[c]{$0.453\,\pm\,0.350$\\$[0.018,\ 0.888]$} & \makecell[c]{$0.965\,\pm\,1.430$\\$[-0.811,\ 2.740]$} \\
Median & \makecell[c]{$2.556\,\pm\,4.359$\\$[-2.857,\ 7.969]$} & \makecell[c]{$1.176\,\pm\,1.698$\\$[-0.932,\ 3.285]$} & \makecell[c]{$2.158\,\pm\,2.686$\\$[-1.178,\ 5.493]$} \\
FFill + BFill & \makecell[c]{$0.431\,\pm\,0.245$\\$[0.126,\ 0.736]$} & \makecell[c]{$0.694\,\pm\,0.230$\\$[0.409,\ 0.979]$} & \makecell[c]{$1.894\,\pm\,1.499$\\$[0.032,\ 3.755]$} \\
Rolling mean & \makecell[c]{$0.176\,\pm\,0.034$\\$[0.135,\ 0.218]$} & \makecell[c]{$0.466\,\pm\,0.259$\\$[0.144,\ 0.787]$} & \makecell[c]{$0.916\,\pm\,1.253$\\$[-0.640,\ 2.471]$} \\
KNN (scaled) & \makecell[c]{$0.896\,\pm\,1.064$\\$[-0.425,\ 2.217]$} & \makecell[c]{$0.620\,\pm\,0.372$\\$[0.157,\ 1.082]$} & \makecell[c]{$1.307\,\pm\,1.498$\\$[-0.552,\ 3.167]$} \\
KNN (unscaled) & \makecell[c]{$0.196\,\pm\,0.102$\\$[0.069,\ 0.323]$} & \makecell[c]{$0.209\,\pm\,0.104$\\$[0.079,\ 0.338]$} & \makecell[c]{$0.453\,\pm\,0.668$\\$[-0.377,\ 1.283]$} \\
Historical & \makecell[c]{$0.374\,\pm\,0.492$\\$[-0.237,\ 0.984]$} & \makecell[c]{$0.744\,\pm\,1.096$\\$[-0.617,\ 2.105]$} & \makecell[c]{$0.338\,\pm\,0.261$\\$[0.013,\ 0.662]$} \\
Graph Laplacian & \makecell[c]{$0.799\,\pm\,1.021$\\$[-0.469,\ 2.067]$} & \makecell[c]{$0.679\,\pm\,0.847$\\$[-0.374,\ 1.731]$} & \makecell[c]{$1.023\,\pm\,1.017$\\$[-0.241,\ 2.286]$} \\
Graph Temporal & \makecell[c]{$0.849\,\pm\,0.881$\\$[-0.245,\ 1.944]$} & \makecell[c]{$1.673\,\pm\,2.905$\\$[-1.934,\ 5.281]$} & \makecell[c]{$1.249\,\pm\,1.191$\\$[-0.230,\ 2.727]$} \\
Sage-Voronoi & \makecell[c]{$0.412\,\pm\,0.283$\\$[0.061,\ 0.764]$} & \makecell[c]{$0.532\,\pm\,0.694$\\$[-0.330,\ 1.394]$} & \makecell[c]{$0.840\,\pm\,0.867$\\$[-0.237,\ 1.917]$} \\
GAIN & \makecell[c]{$\mathbf{0.159}\,\pm\,0.103$\\$[0.031,\ 0.288]$} & \makecell[c]{$\mathbf{0.126}\,\pm\,0.061$\\$[0.050,\ 0.201]$} & \makecell[c]{$\mathbf{0.166}\,\pm\,0.143$\\$[-0.012,\ 0.344]$} \\
GRIN & \makecell[c]{$0.248\,\pm\,0.209$\\$[-0.012,\ 0.508]$} & \makecell[c]{$0.772\,\pm\,1.414$\\$[-0.983,\ 2.528]$} & \makecell[c]{$0.520\,\pm\,0.679$\\$[-0.322,\ 1.363]$} \\
\bottomrule
\end{tabular}
\endgroup
\end{table}

\clearpage

\begin{table}[htbp]
\centering
\caption{{MRE for sensor A027, January 2023 (dimensionless ratio). Each cell gives mean $\pm$ sample SD, followed by the 95\% confidence interval; $n=5$ seed runs. Bold means identify the lowest unrounded mean within each column, not statistical significance.}}
\label{tab:supp_p01_59}
\begingroup
\footnotesize
\setlength{\tabcolsep}{3pt}
\renewcommand{\arraystretch}{1.15}
\begin{tabular}{lcc}
\toprule
Method & Block $B=128$ & Block $B=144$ \\
\midrule
Linear & \makecell[c]{$3.428\,\pm\,2.936$\\$[-0.217,\ 7.074]$} & \makecell[c]{$1.222\,\pm\,0.666$\\$[0.395,\ 2.050]$} \\
Median & \makecell[c]{$2.604\,\pm\,1.339$\\$[0.941,\ 4.266]$} & \makecell[c]{$2.741\,\pm\,0.142$\\$[2.565,\ 2.917]$} \\
FFill + BFill & \makecell[c]{$4.137\,\pm\,3.071$\\$[0.324,\ 7.950]$} & \makecell[c]{$1.291\,\pm\,0.617$\\$[0.524,\ 2.058]$} \\
Rolling mean & \makecell[c]{$3.429\,\pm\,2.931$\\$[-0.210,\ 7.069]$} & \makecell[c]{$1.221\,\pm\,0.668$\\$[0.392,\ 2.050]$} \\
KNN (scaled) & \makecell[c]{$1.605\,\pm\,0.766$\\$[0.653,\ 2.556]$} & \makecell[c]{$1.146\,\pm\,0.091$\\$[1.034,\ 1.259]$} \\
KNN (unscaled) & \makecell[c]{$0.527\,\pm\,0.334$\\$[0.112,\ 0.943]$} & \makecell[c]{$0.280\,\pm\,0.089$\\$[0.169,\ 0.391]$} \\
Historical & \makecell[c]{$0.356\,\pm\,0.126$\\$[0.199,\ 0.513]$} & \makecell[c]{$0.479\,\pm\,0.249$\\$[0.169,\ 0.789]$} \\
Graph Laplacian & \makecell[c]{$1.282\,\pm\,0.657$\\$[0.467,\ 2.098]$} & \makecell[c]{$0.995\,\pm\,0.102$\\$[0.868,\ 1.121]$} \\
Graph Temporal & \makecell[c]{$1.631\,\pm\,1.811$\\$[-0.617,\ 3.880]$} & \makecell[c]{$1.054\,\pm\,0.438$\\$[0.510,\ 1.597]$} \\
Sage-Voronoi & \makecell[c]{$1.364\,\pm\,0.963$\\$[0.169,\ 2.560]$} & \makecell[c]{$0.734\,\pm\,0.354$\\$[0.295,\ 1.173]$} \\
GAIN & \makecell[c]{$\mathbf{0.278}\,\pm\,0.139$\\$[0.106,\ 0.450]$} & \makecell[c]{$\mathbf{0.204}\,\pm\,0.037$\\$[0.159,\ 0.250]$} \\
GRIN & \makecell[c]{$0.930\,\pm\,1.487$\\$[-0.916,\ 2.776]$} & \makecell[c]{$0.420\,\pm\,0.311$\\$[0.034,\ 0.806]$} \\
\bottomrule
\end{tabular}
\endgroup
\end{table}

\begin{table}[htbp]
\centering
\caption{{MRE for sensor A027, January 2023 (dimensionless ratio). Each cell gives mean $\pm$ sample SD, followed by the 95\% confidence interval; $n=5$ seed runs. Bold means identify the lowest unrounded mean within each column, not statistical significance.}}
\label{tab:supp_p01_60}
\begingroup
\footnotesize
\setlength{\tabcolsep}{3pt}
\renewcommand{\arraystretch}{1.15}
\begin{tabular}{lcccc}
\toprule
Method & MCAR 10\% & MCAR 30\% & MCAR 50\% & Synchronous $B=144$ \\
\midrule
Linear & \makecell[c]{$\mathbf{0.170}\,\pm\,0.023$\\$[0.141,\ 0.199]$} & \makecell[c]{$\mathbf{0.173}\,\pm\,0.015$\\$[0.154,\ 0.192]$} & \makecell[c]{$\mathbf{0.167}\,\pm\,0.019$\\$[0.143,\ 0.190]$} & \makecell[c]{$4.371\,\pm\,1.179$\\$[2.908,\ 5.835]$} \\
Median & \makecell[c]{$2.888\,\pm\,0.675$\\$[2.050,\ 3.726]$} & \makecell[c]{$2.666\,\pm\,0.149$\\$[2.481,\ 2.850]$} & \makecell[c]{$2.679\,\pm\,0.225$\\$[2.400,\ 2.959]$} & \makecell[c]{$3.030\,\pm\,0.245$\\$[2.726,\ 3.335]$} \\
FFill + BFill & \makecell[c]{$0.198\,\pm\,0.039$\\$[0.149,\ 0.246]$} & \makecell[c]{$0.210\,\pm\,0.025$\\$[0.179,\ 0.240]$} & \makecell[c]{$0.230\,\pm\,0.036$\\$[0.185,\ 0.274]$} & \makecell[c]{$4.084\,\pm\,2.135$\\$[1.433,\ 6.735]$} \\
Rolling mean & \makecell[c]{$0.211\,\pm\,0.044$\\$[0.156,\ 0.266]$} & \makecell[c]{$0.217\,\pm\,0.009$\\$[0.206,\ 0.228]$} & \makecell[c]{$0.224\,\pm\,0.030$\\$[0.187,\ 0.262]$} & \makecell[c]{$4.373\,\pm\,1.183$\\$[2.904,\ 5.842]$} \\
KNN (scaled) & \makecell[c]{$1.139\,\pm\,0.200$\\$[0.891,\ 1.388]$} & \makecell[c]{$1.082\,\pm\,0.124$\\$[0.928,\ 1.236]$} & \makecell[c]{$1.037\,\pm\,0.083$\\$[0.934,\ 1.140]$} & \makecell[c]{$2.866\,\pm\,0.236$\\$[2.573,\ 3.158]$} \\
KNN (unscaled) & \makecell[c]{$0.371\,\pm\,0.091$\\$[0.258,\ 0.485]$} & \makecell[c]{$0.382\,\pm\,0.024$\\$[0.353,\ 0.412]$} & \makecell[c]{$0.450\,\pm\,0.038$\\$[0.403,\ 0.497]$} & \makecell[c]{$2.866\,\pm\,0.236$\\$[2.573,\ 3.158]$} \\
Historical & \makecell[c]{$0.472\,\pm\,0.082$\\$[0.370,\ 0.574]$} & \makecell[c]{$0.466\,\pm\,0.013$\\$[0.449,\ 0.482]$} & \makecell[c]{$0.452\,\pm\,0.044$\\$[0.397,\ 0.507]$} & \makecell[c]{$\mathbf{0.496}\,\pm\,0.189$\\$[0.261,\ 0.731]$} \\
Graph Laplacian & \makecell[c]{$1.126\,\pm\,0.182$\\$[0.899,\ 1.352]$} & \makecell[c]{$1.028\,\pm\,0.071$\\$[0.940,\ 1.115]$} & \makecell[c]{$0.995\,\pm\,0.084$\\$[0.890,\ 1.099]$} & \makecell[c]{$2.866\,\pm\,0.236$\\$[2.573,\ 3.158]$} \\
Graph Temporal & \makecell[c]{$1.165\,\pm\,0.263$\\$[0.838,\ 1.491]$} & \makecell[c]{$1.260\,\pm\,0.199$\\$[1.012,\ 1.507]$} & \makecell[c]{$1.238\,\pm\,0.159$\\$[1.040,\ 1.436]$} & \makecell[c]{$5.431\,\pm\,1.843$\\$[3.142,\ 7.720]$} \\
Sage-Voronoi & \makecell[c]{$0.222\,\pm\,0.049$\\$[0.161,\ 0.282]$} & \makecell[c]{$0.238\,\pm\,0.022$\\$[0.211,\ 0.265]$} & \makecell[c]{$0.273\,\pm\,0.056$\\$[0.203,\ 0.343]$} & \makecell[c]{$3.460\,\pm\,1.981$\\$[1.000,\ 5.920]$} \\
GAIN & \makecell[c]{$0.216\,\pm\,0.022$\\$[0.188,\ 0.244]$} & \makecell[c]{$0.400\,\pm\,0.061$\\$[0.325,\ 0.476]$} & \makecell[c]{$0.766\,\pm\,0.060$\\$[0.691,\ 0.841]$} & \makecell[c]{$4.017\,\pm\,0.291$\\$[3.655,\ 4.378]$} \\
GRIN & \makecell[c]{$0.181\,\pm\,0.049$\\$[0.120,\ 0.241]$} & \makecell[c]{$0.222\,\pm\,0.025$\\$[0.190,\ 0.253]$} & \makecell[c]{$0.267\,\pm\,0.087$\\$[0.159,\ 0.375]$} & \makecell[c]{$2.257\,\pm\,0.332$\\$[1.844,\ 2.669]$} \\
\bottomrule
\end{tabular}
\endgroup
\end{table}

\clearpage

\begin{table}[htbp]
\centering
\caption{{MAE for sensor A027, October--December 2022 (vehicles per ten-minute interval). Each cell gives mean $\pm$ sample SD, followed by the 95\% confidence interval; $n=5$ seed runs. Bold means identify the lowest unrounded mean within each column, not statistical significance.}}
\label{tab:supp_p01_61}
\begingroup
\footnotesize
\setlength{\tabcolsep}{3pt}
\renewcommand{\arraystretch}{1.15}
\begin{tabular}{lcccc}
\toprule
Method & Block $B=1$ & Block $B=2$ & Block $B=4$ & Block $B=8$ \\
\midrule
Linear & \makecell[c]{$29.80\,\pm\,21.88$\\$[2.64,\ 56.96]$} & \makecell[c]{$12.87\,\pm\,4.27$\\$[7.56,\ 18.17]$} & \makecell[c]{$27.03\,\pm\,9.60$\\$[15.11,\ 38.95]$} & \makecell[c]{$25.30\,\pm\,11.29$\\$[11.29,\ 39.32]$} \\
Median & \makecell[c]{$134.30\,\pm\,86.02$\\$[27.49,\ 241.11]$} & \makecell[c]{$175.20\,\pm\,53.00$\\$[109.40,\ 241.00]$} & \makecell[c]{$119.50\,\pm\,115.99$\\$[-24.51,\ 263.51]$} & \makecell[c]{$131.32\,\pm\,64.61$\\$[51.10,\ 211.55]$} \\
FFill + BFill & \makecell[c]{$31.20\,\pm\,29.99$\\$[-6.04,\ 68.44]$} & \makecell[c]{$20.50\,\pm\,19.04$\\$[-3.14,\ 44.14]$} & \makecell[c]{$32.25\,\pm\,9.68$\\$[20.23,\ 44.27]$} & \makecell[c]{$37.67\,\pm\,17.21$\\$[16.30,\ 59.05]$} \\
Rolling mean & \makecell[c]{$\mathbf{15.20}\,\pm\,21.39$\\$[-11.36,\ 41.76]$} & \makecell[c]{$37.72\,\pm\,31.22$\\$[-1.04,\ 76.48]$} & \makecell[c]{$36.00\,\pm\,13.37$\\$[19.40,\ 52.60]$} & \makecell[c]{$30.87\,\pm\,9.92$\\$[18.55,\ 43.19]$} \\
KNN (scaled) & \makecell[c]{$55.58\,\pm\,30.94$\\$[17.16,\ 94.00]$} & \makecell[c]{$15.56\,\pm\,10.43$\\$[2.60,\ 28.52]$} & \makecell[c]{$20.36\,\pm\,8.96$\\$[9.24,\ 31.49]$} & \makecell[c]{$22.97\,\pm\,10.30$\\$[10.19,\ 35.75]$} \\
KNN (unscaled) & \makecell[c]{$38.92\,\pm\,30.67$\\$[0.84,\ 77.00]$} & \makecell[c]{$16.58\,\pm\,14.40$\\$[-1.30,\ 34.46]$} & \makecell[c]{$29.36\,\pm\,10.46$\\$[16.37,\ 42.35]$} & \makecell[c]{$23.86\,\pm\,9.86$\\$[11.61,\ 36.11]$} \\
Historical & \makecell[c]{$40.26\,\pm\,26.94$\\$[6.81,\ 73.71]$} & \makecell[c]{$24.96\,\pm\,10.44$\\$[12.00,\ 37.92]$} & \makecell[c]{$31.31\,\pm\,32.59$\\$[-9.15,\ 71.77]$} & \makecell[c]{$34.18\,\pm\,18.57$\\$[11.13,\ 57.24]$} \\
Graph Laplacian & \makecell[c]{$41.96\,\pm\,32.67$\\$[1.39,\ 82.53]$} & \makecell[c]{$24.20\,\pm\,18.50$\\$[1.23,\ 47.17]$} & \makecell[c]{$45.10\,\pm\,12.15$\\$[30.01,\ 60.19]$} & \makecell[c]{$35.69\,\pm\,14.41$\\$[17.81,\ 53.58]$} \\
Graph Temporal & \makecell[c]{$49.44\,\pm\,29.39$\\$[12.95,\ 85.94]$} & \makecell[c]{$23.55\,\pm\,17.11$\\$[2.30,\ 44.80]$} & \makecell[c]{$44.02\,\pm\,13.07$\\$[27.79,\ 60.26]$} & \makecell[c]{$56.58\,\pm\,14.65$\\$[38.39,\ 74.77]$} \\
Sage-Voronoi & \makecell[c]{$19.32\,\pm\,27.22$\\$[-14.48,\ 53.12]$} & \makecell[c]{$22.90\,\pm\,16.33$\\$[2.63,\ 43.18]$} & \makecell[c]{$16.87\,\pm\,8.95$\\$[5.76,\ 27.99]$} & \makecell[c]{$25.74\,\pm\,10.21$\\$[13.07,\ 38.42]$} \\
GAIN & \makecell[c]{$40.31\,\pm\,40.83$\\$[-10.38,\ 91.01]$} & \makecell[c]{$20.06\,\pm\,12.08$\\$[5.06,\ 35.06]$} & \makecell[c]{$29.14\,\pm\,18.95$\\$[5.61,\ 52.67]$} & \makecell[c]{$23.22\,\pm\,6.60$\\$[15.03,\ 31.41]$} \\
GRIN & \makecell[c]{$25.17\,\pm\,18.90$\\$[1.70,\ 48.63]$} & \makecell[c]{$\mathbf{8.72}\,\pm\,4.77$\\$[2.80,\ 14.65]$} & \makecell[c]{$\mathbf{13.91}\,\pm\,7.79$\\$[4.24,\ 23.58]$} & \makecell[c]{$\mathbf{21.70}\,\pm\,8.43$\\$[11.23,\ 32.17]$} \\
\bottomrule
\end{tabular}
\endgroup
\end{table}

\begin{table}[htbp]
\centering
\caption{{MAE for sensor A027, October--December 2022 (vehicles per ten-minute interval). Each cell gives mean $\pm$ sample SD, followed by the 95\% confidence interval; $n=5$ seed runs. Bold means identify the lowest unrounded mean within each column, not statistical significance.}}
\label{tab:supp_p01_62}
\begingroup
\footnotesize
\setlength{\tabcolsep}{3pt}
\renewcommand{\arraystretch}{1.15}
\begin{tabular}{lccc}
\toprule
Method & Block $B=16$ & Block $B=32$ & Block $B=64$ \\
\midrule
Linear & \makecell[c]{$27.11\,\pm\,10.93$\\$[13.53,\ 40.69]$} & \makecell[c]{$34.28\,\pm\,16.69$\\$[13.56,\ 55.00]$} & \makecell[c]{$85.91\,\pm\,60.05$\\$[11.35,\ 160.47]$} \\
Median & \makecell[c]{$167.45\,\pm\,61.40$\\$[91.21,\ 243.69]$} & \makecell[c]{$200.53\,\pm\,56.11$\\$[130.87,\ 270.20]$} & \makecell[c]{$133.16\,\pm\,68.01$\\$[48.72,\ 217.60]$} \\
FFill + BFill & \makecell[c]{$60.08\,\pm\,32.36$\\$[19.90,\ 100.25]$} & \makecell[c]{$91.77\,\pm\,46.52$\\$[34.01,\ 149.53]$} & \makecell[c]{$159.95\,\pm\,40.23$\\$[110.00,\ 209.90]$} \\
Rolling mean & \makecell[c]{$27.56\,\pm\,11.32$\\$[13.50,\ 41.62]$} & \makecell[c]{$40.70\,\pm\,12.90$\\$[24.68,\ 56.72]$} & \makecell[c]{$85.80\,\pm\,57.42$\\$[14.51,\ 157.10]$} \\
KNN (scaled) & \makecell[c]{$19.07\,\pm\,3.27$\\$[15.01,\ 23.14]$} & \makecell[c]{$24.23\,\pm\,15.08$\\$[5.51,\ 42.95]$} & \makecell[c]{$26.10\,\pm\,4.26$\\$[20.82,\ 31.39]$} \\
KNN (unscaled) & \makecell[c]{$21.36\,\pm\,4.76$\\$[15.45,\ 27.26]$} & \makecell[c]{$25.83\,\pm\,9.43$\\$[14.12,\ 37.54]$} & \makecell[c]{$32.38\,\pm\,2.81$\\$[28.90,\ 35.87]$} \\
Historical & \makecell[c]{$76.20\,\pm\,102.75$\\$[-51.38,\ 203.79]$} & \makecell[c]{$28.78\,\pm\,12.27$\\$[13.54,\ 44.01]$} & \makecell[c]{$48.99\,\pm\,25.68$\\$[17.10,\ 80.87]$} \\
Graph Laplacian & \makecell[c]{$29.18\,\pm\,13.21$\\$[12.78,\ 45.59]$} & \makecell[c]{$42.52\,\pm\,20.18$\\$[17.46,\ 67.57]$} & \makecell[c]{$44.44\,\pm\,10.65$\\$[31.22,\ 57.66]$} \\
Graph Temporal & \makecell[c]{$38.14\,\pm\,17.20$\\$[16.78,\ 59.49]$} & \makecell[c]{$42.71\,\pm\,10.31$\\$[29.91,\ 55.51]$} & \makecell[c]{$52.33\,\pm\,4.72$\\$[46.47,\ 58.19]$} \\
Sage-Voronoi & \makecell[c]{$22.66\,\pm\,4.92$\\$[16.56,\ 28.77]$} & \makecell[c]{$22.88\,\pm\,5.89$\\$[15.56,\ 30.20]$} & \makecell[c]{$40.36\,\pm\,13.85$\\$[23.16,\ 57.55]$} \\
GAIN & \makecell[c]{$30.35\,\pm\,12.52$\\$[14.80,\ 45.89]$} & \makecell[c]{$22.87\,\pm\,7.12$\\$[14.03,\ 31.71]$} & \makecell[c]{$30.65\,\pm\,14.02$\\$[13.24,\ 48.06]$} \\
GRIN & \makecell[c]{$\mathbf{15.97}\,\pm\,4.61$\\$[10.26,\ 21.69]$} & \makecell[c]{$\mathbf{14.81}\,\pm\,6.30$\\$[6.99,\ 22.63]$} & \makecell[c]{$\mathbf{25.31}\,\pm\,7.76$\\$[15.67,\ 34.95]$} \\
\bottomrule
\end{tabular}
\endgroup
\end{table}

\clearpage

\begin{table}[htbp]
\centering
\caption{{MAE for sensor A027, October--December 2022 (vehicles per ten-minute interval). Each cell gives mean $\pm$ sample SD, followed by the 95\% confidence interval; $n=5$ seed runs. Bold means identify the lowest unrounded mean within each column, not statistical significance.}}
\label{tab:supp_p01_63}
\begingroup
\footnotesize
\setlength{\tabcolsep}{3pt}
\renewcommand{\arraystretch}{1.15}
\begin{tabular}{lcc}
\toprule
Method & Block $B=128$ & Block $B=144$ \\
\midrule
Linear & \makecell[c]{$121.38\,\pm\,14.72$\\$[103.10,\ 139.66]$} & \makecell[c]{$139.12\,\pm\,10.89$\\$[125.60,\ 152.64]$} \\
Median & \makecell[c]{$134.35\,\pm\,10.90$\\$[120.82,\ 147.88]$} & \makecell[c]{$143.14\,\pm\,8.25$\\$[132.91,\ 153.38]$} \\
FFill + BFill & \makecell[c]{$134.49\,\pm\,16.46$\\$[114.05,\ 154.93]$} & \makecell[c]{$137.51\,\pm\,17.52$\\$[115.75,\ 159.26]$} \\
Rolling mean & \makecell[c]{$120.90\,\pm\,14.67$\\$[102.69,\ 139.11]$} & \makecell[c]{$139.23\,\pm\,10.82$\\$[125.79,\ 152.67]$} \\
KNN (scaled) & \makecell[c]{$\mathbf{25.84}\,\pm\,2.60$\\$[22.61,\ 29.07]$} & \makecell[c]{$\mathbf{24.54}\,\pm\,3.10$\\$[20.69,\ 28.39]$} \\
KNN (unscaled) & \makecell[c]{$30.27\,\pm\,3.81$\\$[25.54,\ 34.99]$} & \makecell[c]{$29.66\,\pm\,3.25$\\$[25.62,\ 33.69]$} \\
Historical & \makecell[c]{$48.95\,\pm\,9.61$\\$[37.01,\ 60.89]$} & \makecell[c]{$54.64\,\pm\,27.10$\\$[20.99,\ 88.29]$} \\
Graph Laplacian & \makecell[c]{$39.45\,\pm\,4.04$\\$[34.42,\ 44.47]$} & \makecell[c]{$40.72\,\pm\,7.43$\\$[31.49,\ 49.95]$} \\
Graph Temporal & \makecell[c]{$54.07\,\pm\,12.61$\\$[38.42,\ 69.73]$} & \makecell[c]{$54.26\,\pm\,12.06$\\$[39.28,\ 69.23]$} \\
Sage-Voronoi & \makecell[c]{$39.45\,\pm\,11.89$\\$[24.68,\ 54.21]$} & \makecell[c]{$35.46\,\pm\,5.17$\\$[29.04,\ 41.88]$} \\
GAIN & \makecell[c]{$27.56\,\pm\,4.89$\\$[21.48,\ 33.64]$} & \makecell[c]{$27.02\,\pm\,2.63$\\$[23.75,\ 30.29]$} \\
GRIN & \makecell[c]{$31.64\,\pm\,8.88$\\$[20.62,\ 42.66]$} & \makecell[c]{$31.17\,\pm\,18.38$\\$[8.35,\ 54.00]$} \\
\bottomrule
\end{tabular}
\endgroup
\end{table}

\begin{table}[htbp]
\centering
\caption{{MAE for sensor A027, October--December 2022 (vehicles per ten-minute interval). Each cell gives mean $\pm$ sample SD, followed by the 95\% confidence interval; $n=5$ seed runs. Bold means identify the lowest unrounded mean within each column, not statistical significance.}}
\label{tab:supp_p01_64}
\begingroup
\footnotesize
\setlength{\tabcolsep}{3pt}
\renewcommand{\arraystretch}{1.15}
\begin{tabular}{lcccc}
\toprule
Method & MCAR 10\% & MCAR 30\% & MCAR 50\% & Synchronous $B=144$ \\
\midrule
Linear & \makecell[c]{$22.89\,\pm\,1.11$\\$[21.51,\ 24.27]$} & \makecell[c]{$23.10\,\pm\,0.31$\\$[22.72,\ 23.48]$} & \makecell[c]{$23.69\,\pm\,0.51$\\$[23.06,\ 24.32]$} & \makecell[c]{$146.03\,\pm\,25.85$\\$[113.93,\ 178.13]$} \\
Median & \makecell[c]{$135.48\,\pm\,1.51$\\$[133.61,\ 137.36]$} & \makecell[c]{$137.23\,\pm\,1.66$\\$[135.16,\ 139.29]$} & \makecell[c]{$137.24\,\pm\,2.41$\\$[134.25,\ 140.23]$} & \makecell[c]{$136.19\,\pm\,10.56$\\$[123.08,\ 149.30]$} \\
FFill + BFill & \makecell[c]{$27.77\,\pm\,1.45$\\$[25.96,\ 29.57]$} & \makecell[c]{$28.87\,\pm\,1.22$\\$[27.36,\ 30.38]$} & \makecell[c]{$31.16\,\pm\,0.69$\\$[30.30,\ 32.02]$} & \makecell[c]{$143.93\,\pm\,24.08$\\$[114.03,\ 173.83]$} \\
Rolling mean & \makecell[c]{$28.71\,\pm\,1.42$\\$[26.95,\ 30.47]$} & \makecell[c]{$28.99\,\pm\,0.85$\\$[27.94,\ 30.05]$} & \makecell[c]{$30.18\,\pm\,0.71$\\$[29.30,\ 31.05]$} & \makecell[c]{$146.04\,\pm\,25.40$\\$[114.51,\ 177.58]$} \\
KNN (scaled) & \makecell[c]{$24.70\,\pm\,1.43$\\$[22.93,\ 26.47]$} & \makecell[c]{$24.87\,\pm\,0.75$\\$[23.94,\ 25.80]$} & \makecell[c]{$25.11\,\pm\,0.44$\\$[24.57,\ 25.65]$} & \makecell[c]{$129.34\,\pm\,16.24$\\$[109.18,\ 149.51]$} \\
KNN (unscaled) & \makecell[c]{$28.40\,\pm\,1.27$\\$[26.83,\ 29.97]$} & \makecell[c]{$27.80\,\pm\,0.90$\\$[26.68,\ 28.92]$} & \makecell[c]{$27.61\,\pm\,0.19$\\$[27.37,\ 27.84]$} & \makecell[c]{$129.34\,\pm\,16.24$\\$[109.18,\ 149.51]$} \\
Historical & \makecell[c]{$54.04\,\pm\,1.94$\\$[51.62,\ 56.45]$} & \makecell[c]{$50.60\,\pm\,1.12$\\$[49.21,\ 51.99]$} & \makecell[c]{$53.32\,\pm\,1.53$\\$[51.42,\ 55.22]$} & \makecell[c]{$\mathbf{61.57}\,\pm\,21.84$\\$[34.45,\ 88.69]$} \\
Graph Laplacian & \makecell[c]{$37.80\,\pm\,5.59$\\$[30.87,\ 44.74]$} & \makecell[c]{$37.82\,\pm\,3.26$\\$[33.77,\ 41.87]$} & \makecell[c]{$37.49\,\pm\,3.16$\\$[33.57,\ 41.42]$} & \makecell[c]{$129.34\,\pm\,16.24$\\$[109.18,\ 149.51]$} \\
Graph Temporal & \makecell[c]{$39.19\,\pm\,1.64$\\$[37.16,\ 41.23]$} & \makecell[c]{$40.38\,\pm\,0.63$\\$[39.60,\ 41.15]$} & \makecell[c]{$42.39\,\pm\,0.66$\\$[41.57,\ 43.21]$} & \makecell[c]{$145.32\,\pm\,27.82$\\$[110.78,\ 179.87]$} \\
Sage-Voronoi & \makecell[c]{$23.12\,\pm\,1.40$\\$[21.38,\ 24.86]$} & \makecell[c]{$23.06\,\pm\,0.47$\\$[22.48,\ 23.65]$} & \makecell[c]{$24.40\,\pm\,0.36$\\$[23.95,\ 24.85]$} & \makecell[c]{$162.10\,\pm\,58.68$\\$[89.23,\ 234.96]$} \\
GAIN & \makecell[c]{$25.68\,\pm\,2.05$\\$[23.14,\ 28.22]$} & \makecell[c]{$29.73\,\pm\,2.51$\\$[26.61,\ 32.85]$} & \makecell[c]{$47.18\,\pm\,4.89$\\$[41.11,\ 53.25]$} & \makecell[c]{$179.17\,\pm\,29.69$\\$[142.31,\ 216.04]$} \\
GRIN & \makecell[c]{$\mathbf{17.49}\,\pm\,1.38$\\$[15.78,\ 19.20]$} & \makecell[c]{$\mathbf{18.72}\,\pm\,0.49$\\$[18.11,\ 19.33]$} & \makecell[c]{$\mathbf{21.26}\,\pm\,0.27$\\$[20.93,\ 21.59]$} & \makecell[c]{$115.51\,\pm\,19.32$\\$[91.53,\ 139.50]$} \\
\bottomrule
\end{tabular}
\endgroup
\end{table}

\clearpage

\begin{table}[htbp]
\centering
\caption{{RMSE for sensor A027, October--December 2022 (vehicles per ten-minute interval). Each cell gives mean $\pm$ sample SD, followed by the 95\% confidence interval; $n=5$ seed runs. Bold means identify the lowest unrounded mean within each column, not statistical significance.}}
\label{tab:supp_p01_65}
\begingroup
\footnotesize
\setlength{\tabcolsep}{3pt}
\renewcommand{\arraystretch}{1.15}
\begin{tabular}{lcccc}
\toprule
Method & Block $B=1$ & Block $B=2$ & Block $B=4$ & Block $B=8$ \\
\midrule
Linear & \makecell[c]{$29.80\,\pm\,21.88$\\$[2.64,\ 56.96]$} & \makecell[c]{$15.53\,\pm\,6.11$\\$[7.94,\ 23.11]$} & \makecell[c]{$30.19\,\pm\,12.69$\\$[14.43,\ 45.95]$} & \makecell[c]{$29.59\,\pm\,13.67$\\$[12.62,\ 46.56]$} \\
Median & \makecell[c]{$134.30\,\pm\,86.02$\\$[27.49,\ 241.11]$} & \makecell[c]{$175.36\,\pm\,53.15$\\$[109.36,\ 241.36]$} & \makecell[c]{$121.71\,\pm\,114.60$\\$[-20.58,\ 264.01]$} & \makecell[c]{$134.84\,\pm\,62.93$\\$[56.70,\ 212.99]$} \\
FFill + BFill & \makecell[c]{$31.20\,\pm\,29.99$\\$[-6.04,\ 68.44]$} & \makecell[c]{$22.54\,\pm\,18.40$\\$[-0.31,\ 45.39]$} & \makecell[c]{$36.64\,\pm\,9.94$\\$[24.30,\ 48.98]$} & \makecell[c]{$43.39\,\pm\,16.93$\\$[22.37,\ 64.41]$} \\
Rolling mean & \makecell[c]{$\mathbf{15.20}\,\pm\,21.39$\\$[-11.36,\ 41.76]$} & \makecell[c]{$38.96\,\pm\,30.56$\\$[1.02,\ 76.90]$} & \makecell[c]{$40.22\,\pm\,13.05$\\$[24.01,\ 56.43]$} & \makecell[c]{$36.76\,\pm\,11.71$\\$[22.22,\ 51.31]$} \\
KNN (scaled) & \makecell[c]{$55.58\,\pm\,30.94$\\$[17.16,\ 94.00]$} & \makecell[c]{$17.58\,\pm\,11.69$\\$[3.06,\ 32.09]$} & \makecell[c]{$24.82\,\pm\,12.15$\\$[9.73,\ 39.90]$} & \makecell[c]{$\mathbf{27.83}\,\pm\,12.68$\\$[12.08,\ 43.58]$} \\
KNN (unscaled) & \makecell[c]{$38.92\,\pm\,30.67$\\$[0.84,\ 77.00]$} & \makecell[c]{$19.17\,\pm\,15.35$\\$[0.11,\ 38.24]$} & \makecell[c]{$35.78\,\pm\,13.72$\\$[18.75,\ 52.82]$} & \makecell[c]{$30.88\,\pm\,14.58$\\$[12.77,\ 48.99]$} \\
Historical & \makecell[c]{$40.26\,\pm\,26.94$\\$[6.81,\ 73.71]$} & \makecell[c]{$26.93\,\pm\,10.80$\\$[13.52,\ 40.33]$} & \makecell[c]{$35.18\,\pm\,33.41$\\$[-6.31,\ 76.67]$} & \makecell[c]{$38.48\,\pm\,18.21$\\$[15.87,\ 61.10]$} \\
Graph Laplacian & \makecell[c]{$41.96\,\pm\,32.67$\\$[1.39,\ 82.53]$} & \makecell[c]{$25.67\,\pm\,19.51$\\$[1.45,\ 49.90]$} & \makecell[c]{$46.05\,\pm\,12.14$\\$[30.98,\ 61.12]$} & \makecell[c]{$41.61\,\pm\,16.27$\\$[21.41,\ 61.82]$} \\
Graph Temporal & \makecell[c]{$49.44\,\pm\,29.39$\\$[12.95,\ 85.94]$} & \makecell[c]{$24.36\,\pm\,17.73$\\$[2.36,\ 46.37]$} & \makecell[c]{$44.86\,\pm\,12.70$\\$[29.08,\ 60.63]$} & \makecell[c]{$62.86\,\pm\,13.81$\\$[45.72,\ 80.01]$} \\
Sage-Voronoi & \makecell[c]{$19.32\,\pm\,27.22$\\$[-14.48,\ 53.12]$} & \makecell[c]{$24.92\,\pm\,15.55$\\$[5.61,\ 44.23]$} & \makecell[c]{$20.32\,\pm\,9.88$\\$[8.05,\ 32.58]$} & \makecell[c]{$30.90\,\pm\,12.30$\\$[15.62,\ 46.17]$} \\
GAIN & \makecell[c]{$40.31\,\pm\,40.83$\\$[-10.38,\ 91.01]$} & \makecell[c]{$21.90\,\pm\,12.54$\\$[6.33,\ 37.48]$} & \makecell[c]{$33.34\,\pm\,19.07$\\$[9.66,\ 57.03]$} & \makecell[c]{$29.01\,\pm\,9.51$\\$[17.20,\ 40.82]$} \\
GRIN & \makecell[c]{$25.17\,\pm\,18.90$\\$[1.70,\ 48.63]$} & \makecell[c]{$\mathbf{9.37}\,\pm\,4.60$\\$[3.66,\ 15.08]$} & \makecell[c]{$\mathbf{16.55}\,\pm\,9.31$\\$[5.00,\ 28.11]$} & \makecell[c]{$28.39\,\pm\,11.46$\\$[14.17,\ 42.62]$} \\
\bottomrule
\end{tabular}
\endgroup
\end{table}

\begin{table}[htbp]
\centering
\caption{{RMSE for sensor A027, October--December 2022 (vehicles per ten-minute interval). Each cell gives mean $\pm$ sample SD, followed by the 95\% confidence interval; $n=5$ seed runs. Bold means identify the lowest unrounded mean within each column, not statistical significance.}}
\label{tab:supp_p01_66}
\begingroup
\footnotesize
\setlength{\tabcolsep}{3pt}
\renewcommand{\arraystretch}{1.15}
\begin{tabular}{lccc}
\toprule
Method & Block $B=16$ & Block $B=32$ & Block $B=64$ \\
\midrule
Linear & \makecell[c]{$31.50\,\pm\,11.51$\\$[17.21,\ 45.78]$} & \makecell[c]{$39.14\,\pm\,16.19$\\$[19.03,\ 59.25]$} & \makecell[c]{$101.71\,\pm\,65.40$\\$[20.50,\ 182.92]$} \\
Median & \makecell[c]{$171.98\,\pm\,60.32$\\$[97.09,\ 246.88]$} & \makecell[c]{$206.51\,\pm\,54.52$\\$[138.82,\ 274.21]$} & \makecell[c]{$152.69\,\pm\,71.75$\\$[63.60,\ 241.78]$} \\
FFill + BFill & \makecell[c]{$69.89\,\pm\,31.80$\\$[30.40,\ 109.37]$} & \makecell[c]{$101.01\,\pm\,50.73$\\$[38.03,\ 164.00]$} & \makecell[c]{$190.86\,\pm\,52.48$\\$[125.69,\ 256.02]$} \\
Rolling mean & \makecell[c]{$33.64\,\pm\,12.07$\\$[18.65,\ 48.62]$} & \makecell[c]{$45.92\,\pm\,12.70$\\$[30.15,\ 61.69]$} & \makecell[c]{$101.07\,\pm\,63.38$\\$[22.38,\ 179.77]$} \\
KNN (scaled) & \makecell[c]{$23.75\,\pm\,4.27$\\$[18.44,\ 29.05]$} & \makecell[c]{$31.66\,\pm\,19.51$\\$[7.43,\ 55.89]$} & \makecell[c]{$40.47\,\pm\,10.78$\\$[27.09,\ 53.85]$} \\
KNN (unscaled) & \makecell[c]{$25.91\,\pm\,5.63$\\$[18.92,\ 32.90]$} & \makecell[c]{$32.45\,\pm\,11.47$\\$[18.21,\ 46.68]$} & \makecell[c]{$46.96\,\pm\,8.78$\\$[36.06,\ 57.86]$} \\
Historical & \makecell[c]{$80.46\,\pm\,101.78$\\$[-45.92,\ 206.84]$} & \makecell[c]{$33.98\,\pm\,14.11$\\$[16.46,\ 51.50]$} & \makecell[c]{$62.21\,\pm\,25.44$\\$[30.62,\ 93.80]$} \\
Graph Laplacian & \makecell[c]{$33.74\,\pm\,13.26$\\$[17.27,\ 50.21]$} & \makecell[c]{$47.04\,\pm\,18.76$\\$[23.75,\ 70.33]$} & \makecell[c]{$54.51\,\pm\,7.96$\\$[44.63,\ 64.40]$} \\
Graph Temporal & \makecell[c]{$42.74\,\pm\,16.75$\\$[21.95,\ 63.54]$} & \makecell[c]{$47.39\,\pm\,9.30$\\$[35.85,\ 58.93]$} & \makecell[c]{$62.02\,\pm\,5.40$\\$[55.31,\ 68.72]$} \\
Sage-Voronoi & \makecell[c]{$27.88\,\pm\,6.69$\\$[19.57,\ 36.19]$} & \makecell[c]{$27.32\,\pm\,6.88$\\$[18.78,\ 35.87]$} & \makecell[c]{$56.32\,\pm\,13.21$\\$[39.92,\ 72.72]$} \\
GAIN & \makecell[c]{$34.91\,\pm\,12.39$\\$[19.52,\ 50.30]$} & \makecell[c]{$28.61\,\pm\,8.15$\\$[18.49,\ 38.72]$} & \makecell[c]{$42.30\,\pm\,17.67$\\$[20.36,\ 64.24]$} \\
GRIN & \makecell[c]{$\mathbf{19.66}\,\pm\,6.33$\\$[11.81,\ 27.52]$} & \makecell[c]{$\mathbf{19.06}\,\pm\,7.20$\\$[10.12,\ 28.00]$} & \makecell[c]{$\mathbf{38.77}\,\pm\,13.92$\\$[21.48,\ 56.05]$} \\
\bottomrule
\end{tabular}
\endgroup
\end{table}

\clearpage

\begin{table}[htbp]
\centering
\caption{{RMSE for sensor A027, October--December 2022 (vehicles per ten-minute interval). Each cell gives mean $\pm$ sample SD, followed by the 95\% confidence interval; $n=5$ seed runs. Bold means identify the lowest unrounded mean within each column, not statistical significance.}}
\label{tab:supp_p01_67}
\begingroup
\footnotesize
\setlength{\tabcolsep}{3pt}
\renewcommand{\arraystretch}{1.15}
\begin{tabular}{lcc}
\toprule
Method & Block $B=128$ & Block $B=144$ \\
\midrule
Linear & \makecell[c]{$142.53\,\pm\,23.56$\\$[113.27,\ 171.79]$} & \makecell[c]{$171.31\,\pm\,20.85$\\$[145.42,\ 197.19]$} \\
Median & \makecell[c]{$163.70\,\pm\,3.34$\\$[159.55,\ 167.85]$} & \makecell[c]{$169.09\,\pm\,7.75$\\$[159.47,\ 178.71]$} \\
FFill + BFill & \makecell[c]{$151.59\,\pm\,12.99$\\$[135.46,\ 167.73]$} & \makecell[c]{$169.46\,\pm\,28.38$\\$[134.22,\ 204.70]$} \\
Rolling mean & \makecell[c]{$141.93\,\pm\,23.48$\\$[112.78,\ 171.08]$} & \makecell[c]{$171.28\,\pm\,20.88$\\$[145.36,\ 197.21]$} \\
KNN (scaled) & \makecell[c]{$\mathbf{36.51}\,\pm\,8.05$\\$[26.52,\ 46.50]$} & \makecell[c]{$\mathbf{32.90}\,\pm\,4.06$\\$[27.86,\ 37.93]$} \\
KNN (unscaled) & \makecell[c]{$41.39\,\pm\,8.28$\\$[31.12,\ 51.67]$} & \makecell[c]{$38.84\,\pm\,3.45$\\$[34.56,\ 43.12]$} \\
Historical & \makecell[c]{$66.53\,\pm\,14.66$\\$[48.33,\ 84.73]$} & \makecell[c]{$72.35\,\pm\,33.95$\\$[30.20,\ 114.50]$} \\
Graph Laplacian & \makecell[c]{$49.03\,\pm\,4.82$\\$[43.05,\ 55.01]$} & \makecell[c]{$48.05\,\pm\,7.33$\\$[38.95,\ 57.15]$} \\
Graph Temporal & \makecell[c]{$66.83\,\pm\,16.92$\\$[45.82,\ 87.84]$} & \makecell[c]{$63.69\,\pm\,17.60$\\$[41.83,\ 85.55]$} \\
Sage-Voronoi & \makecell[c]{$50.17\,\pm\,12.55$\\$[34.59,\ 65.75]$} & \makecell[c]{$46.89\,\pm\,7.96$\\$[37.01,\ 56.78]$} \\
GAIN & \makecell[c]{$38.28\,\pm\,9.22$\\$[26.83,\ 49.73]$} & \makecell[c]{$35.20\,\pm\,3.09$\\$[31.37,\ 39.04]$} \\
GRIN & \makecell[c]{$44.95\,\pm\,11.31$\\$[30.91,\ 58.98]$} & \makecell[c]{$42.09\,\pm\,23.23$\\$[13.24,\ 70.94]$} \\
\bottomrule
\end{tabular}
\endgroup
\end{table}

\begin{table}[htbp]
\centering
\caption{{RMSE for sensor A027, October--December 2022 (vehicles per ten-minute interval). Each cell gives mean $\pm$ sample SD, followed by the 95\% confidence interval; $n=5$ seed runs. Bold means identify the lowest unrounded mean within each column, not statistical significance.}}
\label{tab:supp_p01_68}
\begingroup
\footnotesize
\setlength{\tabcolsep}{3pt}
\renewcommand{\arraystretch}{1.15}
\begin{tabular}{lcccc}
\toprule
Method & MCAR 10\% & MCAR 30\% & MCAR 50\% & Synchronous $B=144$ \\
\midrule
Linear & \makecell[c]{$31.82\,\pm\,2.36$\\$[28.89,\ 34.75]$} & \makecell[c]{$32.69\,\pm\,0.90$\\$[31.57,\ 33.81]$} & \makecell[c]{$33.72\,\pm\,0.77$\\$[32.77,\ 34.68]$} & \makecell[c]{$179.75\,\pm\,37.36$\\$[133.35,\ 226.14]$} \\
Median & \makecell[c]{$162.30\,\pm\,1.99$\\$[159.82,\ 164.77]$} & \makecell[c]{$163.87\,\pm\,1.33$\\$[162.22,\ 165.52]$} & \makecell[c]{$164.06\,\pm\,2.19$\\$[161.34,\ 166.78]$} & \makecell[c]{$162.68\,\pm\,11.16$\\$[148.83,\ 176.54]$} \\
FFill + BFill & \makecell[c]{$39.73\,\pm\,2.55$\\$[36.57,\ 42.89]$} & \makecell[c]{$41.27\,\pm\,2.26$\\$[38.46,\ 44.08]$} & \makecell[c]{$44.55\,\pm\,1.32$\\$[42.91,\ 46.20]$} & \makecell[c]{$177.11\,\pm\,36.14$\\$[132.24,\ 221.99]$} \\
Rolling mean & \makecell[c]{$39.35\,\pm\,2.26$\\$[36.53,\ 42.16]$} & \makecell[c]{$40.26\,\pm\,1.73$\\$[38.11,\ 42.42]$} & \makecell[c]{$42.02\,\pm\,0.83$\\$[40.99,\ 43.06]$} & \makecell[c]{$179.70\,\pm\,37.07$\\$[133.67,\ 225.72]$} \\
KNN (scaled) & \makecell[c]{$34.49\,\pm\,3.48$\\$[30.17,\ 38.81]$} & \makecell[c]{$35.21\,\pm\,1.43$\\$[33.44,\ 36.98]$} & \makecell[c]{$36.14\,\pm\,0.91$\\$[35.01,\ 37.26]$} & \makecell[c]{$147.05\,\pm\,13.97$\\$[129.70,\ 164.39]$} \\
KNN (unscaled) & \makecell[c]{$39.04\,\pm\,2.86$\\$[35.49,\ 42.59]$} & \makecell[c]{$38.56\,\pm\,1.28$\\$[36.97,\ 40.15]$} & \makecell[c]{$38.82\,\pm\,0.58$\\$[38.10,\ 39.53]$} & \makecell[c]{$147.05\,\pm\,13.97$\\$[129.70,\ 164.39]$} \\
Historical & \makecell[c]{$78.04\,\pm\,2.61$\\$[74.80,\ 81.27]$} & \makecell[c]{$72.87\,\pm\,1.77$\\$[70.67,\ 75.07]$} & \makecell[c]{$76.57\,\pm\,2.10$\\$[73.96,\ 79.18]$} & \makecell[c]{$\mathbf{83.70}\,\pm\,27.63$\\$[49.40,\ 118.01]$} \\
Graph Laplacian & \makecell[c]{$46.84\,\pm\,6.55$\\$[38.70,\ 54.97]$} & \makecell[c]{$47.36\,\pm\,4.01$\\$[42.38,\ 52.34]$} & \makecell[c]{$47.73\,\pm\,3.87$\\$[42.93,\ 52.54]$} & \makecell[c]{$147.05\,\pm\,13.97$\\$[129.70,\ 164.39]$} \\
Graph Temporal & \makecell[c]{$47.97\,\pm\,2.51$\\$[44.85,\ 51.09]$} & \makecell[c]{$49.90\,\pm\,0.55$\\$[49.22,\ 50.59]$} & \makecell[c]{$52.74\,\pm\,0.73$\\$[51.83,\ 53.65]$} & \makecell[c]{$173.23\,\pm\,40.34$\\$[123.13,\ 223.32]$} \\
Sage-Voronoi & \makecell[c]{$32.68\,\pm\,2.83$\\$[29.16,\ 36.20]$} & \makecell[c]{$32.38\,\pm\,1.28$\\$[30.79,\ 33.97]$} & \makecell[c]{$34.44\,\pm\,0.53$\\$[33.77,\ 35.10]$} & \makecell[c]{$204.61\,\pm\,64.41$\\$[124.63,\ 284.59]$} \\
GAIN & \makecell[c]{$36.13\,\pm\,3.60$\\$[31.66,\ 40.59]$} & \makecell[c]{$39.07\,\pm\,2.77$\\$[35.62,\ 42.51]$} & \makecell[c]{$56.56\,\pm\,5.48$\\$[49.76,\ 63.36]$} & \makecell[c]{$220.70\,\pm\,19.73$\\$[196.20,\ 245.19]$} \\
GRIN & \makecell[c]{$\mathbf{25.92}\,\pm\,3.67$\\$[21.37,\ 30.47]$} & \makecell[c]{$\mathbf{28.10}\,\pm\,1.71$\\$[25.98,\ 30.22]$} & \makecell[c]{$\mathbf{31.98}\,\pm\,0.69$\\$[31.13,\ 32.83]$} & \makecell[c]{$134.51\,\pm\,24.83$\\$[103.68,\ 165.35]$} \\
\bottomrule
\end{tabular}
\endgroup
\end{table}

\clearpage

\begin{table}[htbp]
\centering
\caption{{MRE for sensor A027, October--December 2022 (dimensionless ratio). Each cell gives mean $\pm$ sample SD, followed by the 95\% confidence interval; $n=5$ seed runs. Bold means identify the lowest unrounded mean within each column, not statistical significance.}}
\label{tab:supp_p01_69}
\begingroup
\footnotesize
\setlength{\tabcolsep}{3pt}
\renewcommand{\arraystretch}{1.15}
\begin{tabular}{lcccc}
\toprule
Method & Block $B=1$ & Block $B=2$ & Block $B=4$ & Block $B=8$ \\
\midrule
Linear & \makecell[c]{$0.112\,\pm\,0.081$\\$[0.012,\ 0.212]$} & \makecell[c]{$0.127\,\pm\,0.117$\\$[-0.018,\ 0.273]$} & \makecell[c]{$0.266\,\pm\,0.294$\\$[-0.099,\ 0.632]$} & \makecell[c]{$0.095\,\pm\,0.057$\\$[0.024,\ 0.166]$} \\
Median & \makecell[c]{$1.284\,\pm\,2.341$\\$[-1.622,\ 4.191]$} & \makecell[c]{$1.639\,\pm\,1.461$\\$[-0.175,\ 3.454]$} & \makecell[c]{$1.909\,\pm\,2.568$\\$[-1.280,\ 5.097]$} & \makecell[c]{$1.021\,\pm\,1.737$\\$[-1.135,\ 3.177]$} \\
FFill + BFill & \makecell[c]{$0.094\,\pm\,0.058$\\$[0.021,\ 0.166]$} & \makecell[c]{$0.184\,\pm\,0.162$\\$[-0.018,\ 0.385]$} & \makecell[c]{$0.235\,\pm\,0.152$\\$[0.047,\ 0.423]$} & \makecell[c]{$0.223\,\pm\,0.298$\\$[-0.146,\ 0.593]$} \\
Rolling mean & \makecell[c]{$0.063\,\pm\,0.072$\\$[-0.027,\ 0.152]$} & \makecell[c]{$0.315\,\pm\,0.254$\\$[-0.001,\ 0.630]$} & \makecell[c]{$0.248\,\pm\,0.169$\\$[0.038,\ 0.458]$} & \makecell[c]{$0.148\,\pm\,0.154$\\$[-0.044,\ 0.339]$} \\
KNN (scaled) & \makecell[c]{$0.197\,\pm\,0.101$\\$[0.072,\ 0.322]$} & \makecell[c]{$0.123\,\pm\,0.067$\\$[0.040,\ 0.207]$} & \makecell[c]{$0.142\,\pm\,0.107$\\$[0.009,\ 0.274]$} & \makecell[c]{$0.088\,\pm\,0.061$\\$[0.012,\ 0.163]$} \\
KNN (unscaled) & \makecell[c]{$0.100\,\pm\,0.066$\\$[0.019,\ 0.182]$} & \makecell[c]{$0.115\,\pm\,0.095$\\$[-0.003,\ 0.233]$} & \makecell[c]{$0.205\,\pm\,0.112$\\$[0.066,\ 0.344]$} & \makecell[c]{$0.089\,\pm\,0.060$\\$[0.014,\ 0.163]$} \\
Historical & \makecell[c]{$0.297\,\pm\,0.473$\\$[-0.291,\ 0.884]$} & \makecell[c]{$0.210\,\pm\,0.150$\\$[0.023,\ 0.396]$} & \makecell[c]{$0.163\,\pm\,0.111$\\$[0.025,\ 0.300]$} & \makecell[c]{$0.133\,\pm\,0.101$\\$[0.007,\ 0.258]$} \\
Graph Laplacian & \makecell[c]{$0.285\,\pm\,0.434$\\$[-0.254,\ 0.825]$} & \makecell[c]{$0.280\,\pm\,0.349$\\$[-0.154,\ 0.714]$} & \makecell[c]{$0.502\,\pm\,0.573$\\$[-0.209,\ 1.213]$} & \makecell[c]{$0.215\,\pm\,0.304$\\$[-0.163,\ 0.592]$} \\
Graph Temporal & \makecell[c]{$0.252\,\pm\,0.296$\\$[-0.116,\ 0.619]$} & \makecell[c]{$0.247\,\pm\,0.316$\\$[-0.145,\ 0.639]$} & \makecell[c]{$0.431\,\pm\,0.452$\\$[-0.130,\ 0.992]$} & \makecell[c]{$0.265\,\pm\,0.280$\\$[-0.082,\ 0.612]$} \\
Sage-Voronoi & \makecell[c]{$\mathbf{0.048}\,\pm\,0.051$\\$[-0.016,\ 0.111]$} & \makecell[c]{$0.182\,\pm\,0.123$\\$[0.030,\ 0.335]$} & \makecell[c]{$0.125\,\pm\,0.079$\\$[0.027,\ 0.223]$} & \makecell[c]{$0.106\,\pm\,0.085$\\$[0.000,\ 0.211]$} \\
GAIN & \makecell[c]{$0.106\,\pm\,0.089$\\$[-0.005,\ 0.217]$} & \makecell[c]{$0.167\,\pm\,0.110$\\$[0.031,\ 0.303]$} & \makecell[c]{$0.157\,\pm\,0.044$\\$[0.102,\ 0.212]$} & \makecell[c]{$0.087\,\pm\,0.058$\\$[0.016,\ 0.159]$} \\
GRIN & \makecell[c]{$0.099\,\pm\,0.078$\\$[0.002,\ 0.196]$} & \makecell[c]{$\mathbf{0.078}\,\pm\,0.058$\\$[0.006,\ 0.150]$} & \makecell[c]{$\mathbf{0.086}\,\pm\,0.035$\\$[0.042,\ 0.130]$} & \makecell[c]{$\mathbf{0.071}\,\pm\,0.028$\\$[0.037,\ 0.106]$} \\
\bottomrule
\end{tabular}
\endgroup
\end{table}

\begin{table}[htbp]
\centering
\caption{{MRE for sensor A027, October--December 2022 (dimensionless ratio). Each cell gives mean $\pm$ sample SD, followed by the 95\% confidence interval; $n=5$ seed runs. Bold means identify the lowest unrounded mean within each column, not statistical significance.}}
\label{tab:supp_p01_70}
\begingroup
\footnotesize
\setlength{\tabcolsep}{3pt}
\renewcommand{\arraystretch}{1.15}
\begin{tabular}{lccc}
\toprule
Method & Block $B=16$ & Block $B=32$ & Block $B=64$ \\
\midrule
Linear & \makecell[c]{$0.237\,\pm\,0.202$\\$[-0.014,\ 0.487]$} & \makecell[c]{$0.820\,\pm\,1.203$\\$[-0.674,\ 2.314]$} & \makecell[c]{$2.885\,\pm\,4.530$\\$[-2.740,\ 8.509]$} \\
Median & \makecell[c]{$2.057\,\pm\,2.006$\\$[-0.434,\ 4.548]$} & \makecell[c]{$4.214\,\pm\,4.022$\\$[-0.780,\ 9.208]$} & \makecell[c]{$4.248\,\pm\,4.683$\\$[-1.566,\ 10.063]$} \\
FFill + BFill & \makecell[c]{$0.568\,\pm\,0.446$\\$[0.015,\ 1.121]$} & \makecell[c]{$1.882\,\pm\,1.107$\\$[0.507,\ 3.257]$} & \makecell[c]{$2.763\,\pm\,3.125$\\$[-1.117,\ 6.644]$} \\
Rolling mean & \makecell[c]{$0.222\,\pm\,0.127$\\$[0.064,\ 0.379]$} & \makecell[c]{$0.885\,\pm\,1.174$\\$[-0.573,\ 2.343]$} & \makecell[c]{$2.793\,\pm\,4.454$\\$[-2.738,\ 8.323]$} \\
KNN (scaled) & \makecell[c]{$0.157\,\pm\,0.099$\\$[0.035,\ 0.280]$} & \makecell[c]{$0.246\,\pm\,0.155$\\$[0.053,\ 0.438]$} & \makecell[c]{$0.270\,\pm\,0.307$\\$[-0.111,\ 0.650]$} \\
KNN (unscaled) & \makecell[c]{$0.167\,\pm\,0.094$\\$[0.050,\ 0.284]$} & \makecell[c]{$0.312\,\pm\,0.205$\\$[0.058,\ 0.567]$} & \makecell[c]{$0.424\,\pm\,0.438$\\$[-0.120,\ 0.968]$} \\
Historical & \makecell[c]{$1.207\,\pm\,2.167$\\$[-1.483,\ 3.897]$} & \makecell[c]{$0.341\,\pm\,0.182$\\$[0.114,\ 0.567]$} & \makecell[c]{$0.445\,\pm\,0.338$\\$[0.025,\ 0.865]$} \\
Graph Laplacian & \makecell[c]{$0.407\,\pm\,0.415$\\$[-0.108,\ 0.923]$} & \makecell[c]{$1.053\,\pm\,1.263$\\$[-0.515,\ 2.621]$} & \makecell[c]{$1.133\,\pm\,1.278$\\$[-0.454,\ 2.720]$} \\
Graph Temporal & \makecell[c]{$0.407\,\pm\,0.372$\\$[-0.055,\ 0.870]$} & \makecell[c]{$0.940\,\pm\,1.008$\\$[-0.311,\ 2.192]$} & \makecell[c]{$1.101\,\pm\,1.172$\\$[-0.354,\ 2.556]$} \\
Sage-Voronoi & \makecell[c]{$0.224\,\pm\,0.190$\\$[-0.012,\ 0.459]$} & \makecell[c]{$0.337\,\pm\,0.161$\\$[0.137,\ 0.536]$} & \makecell[c]{$0.575\,\pm\,0.775$\\$[-0.387,\ 1.537]$} \\
GAIN & \makecell[c]{$0.228\,\pm\,0.140$\\$[0.054,\ 0.402]$} & \makecell[c]{$0.243\,\pm\,0.120$\\$[0.095,\ 0.392]$} & \makecell[c]{$0.229\,\pm\,0.129$\\$[0.068,\ 0.389]$} \\
GRIN & \makecell[c]{$\mathbf{0.118}\,\pm\,0.061$\\$[0.042,\ 0.193]$} & \makecell[c]{$\mathbf{0.151}\,\pm\,0.078$\\$[0.054,\ 0.247]$} & \makecell[c]{$\mathbf{0.191}\,\pm\,0.160$\\$[-0.008,\ 0.390]$} \\
\bottomrule
\end{tabular}
\endgroup
\end{table}

\clearpage

\begin{table}[htbp]
\centering
\caption{{MRE for sensor A027, October--December 2022 (dimensionless ratio). Each cell gives mean $\pm$ sample SD, followed by the 95\% confidence interval; $n=5$ seed runs. Bold means identify the lowest unrounded mean within each column, not statistical significance.}}
\label{tab:supp_p01_71}
\begingroup
\footnotesize
\setlength{\tabcolsep}{3pt}
\renewcommand{\arraystretch}{1.15}
\begin{tabular}{lcc}
\toprule
Method & Block $B=128$ & Block $B=144$ \\
\midrule
Linear & \makecell[c]{$2.251\,\pm\,1.435$\\$[0.470,\ 4.033]$} & \makecell[c]{$2.865\,\pm\,1.162$\\$[1.422,\ 4.308]$} \\
Median & \makecell[c]{$2.739\,\pm\,1.062$\\$[1.421,\ 4.057]$} & \makecell[c]{$3.042\,\pm\,0.759$\\$[2.099,\ 3.984]$} \\
FFill + BFill & \makecell[c]{$2.093\,\pm\,1.206$\\$[0.596,\ 3.590]$} & \makecell[c]{$2.903\,\pm\,1.126$\\$[1.505,\ 4.300]$} \\
Rolling mean & \makecell[c]{$2.237\,\pm\,1.430$\\$[0.461,\ 4.013]$} & \makecell[c]{$2.867\,\pm\,1.158$\\$[1.429,\ 4.305]$} \\
KNN (scaled) & \makecell[c]{$0.177\,\pm\,0.063$\\$[0.099,\ 0.255]$} & \makecell[c]{$\mathbf{0.171}\,\pm\,0.030$\\$[0.134,\ 0.208]$} \\
KNN (unscaled) & \makecell[c]{$0.230\,\pm\,0.077$\\$[0.135,\ 0.326]$} & \makecell[c]{$0.241\,\pm\,0.054$\\$[0.174,\ 0.308]$} \\
Historical & \makecell[c]{$0.347\,\pm\,0.126$\\$[0.190,\ 0.503]$} & \makecell[c]{$0.519\,\pm\,0.525$\\$[-0.133,\ 1.171]$} \\
Graph Laplacian & \makecell[c]{$0.670\,\pm\,0.332$\\$[0.259,\ 1.082]$} & \makecell[c]{$0.735\,\pm\,0.277$\\$[0.391,\ 1.079]$} \\
Graph Temporal & \makecell[c]{$0.830\,\pm\,0.542$\\$[0.157,\ 1.503]$} & \makecell[c]{$0.909\,\pm\,0.414$\\$[0.396,\ 1.423]$} \\
Sage-Voronoi & \makecell[c]{$0.463\,\pm\,0.221$\\$[0.189,\ 0.737]$} & \makecell[c]{$0.421\,\pm\,0.153$\\$[0.231,\ 0.611]$} \\
GAIN & \makecell[c]{$\mathbf{0.172}\,\pm\,0.037$\\$[0.127,\ 0.218]$} & \makecell[c]{$0.179\,\pm\,0.046$\\$[0.122,\ 0.237]$} \\
GRIN & \makecell[c]{$0.382\,\pm\,0.407$\\$[-0.124,\ 0.887]$} & \makecell[c]{$0.367\,\pm\,0.529$\\$[-0.289,\ 1.023]$} \\
\bottomrule
\end{tabular}
\endgroup
\end{table}

\begin{table}[htbp]
\centering
\caption{{MRE for sensor A027, October--December 2022 (dimensionless ratio). Each cell gives mean $\pm$ sample SD, followed by the 95\% confidence interval; $n=5$ seed runs. Bold means identify the lowest unrounded mean within each column, not statistical significance.}}
\label{tab:supp_p01_72}
\begingroup
\footnotesize
\setlength{\tabcolsep}{3pt}
\renewcommand{\arraystretch}{1.15}
\begin{tabular}{lcccc}
\toprule
Method & MCAR 10\% & MCAR 30\% & MCAR 50\% & Synchronous $B=144$ \\
\midrule
Linear & \makecell[c]{$0.161\,\pm\,0.021$\\$[0.134,\ 0.187]$} & \makecell[c]{$0.165\,\pm\,0.007$\\$[0.157,\ 0.173]$} & \makecell[c]{$0.171\,\pm\,0.009$\\$[0.160,\ 0.182]$} & \makecell[c]{$1.814\,\pm\,1.486$\\$[-0.032,\ 3.659]$} \\
Median & \makecell[c]{$2.703\,\pm\,0.437$\\$[2.160,\ 3.245]$} & \makecell[c]{$2.919\,\pm\,0.165$\\$[2.714,\ 3.124]$} & \makecell[c]{$2.977\,\pm\,0.153$\\$[2.787,\ 3.167]$} & \makecell[c]{$2.184\,\pm\,0.934$\\$[1.024,\ 3.343]$} \\
FFill + BFill & \makecell[c]{$0.188\,\pm\,0.029$\\$[0.152,\ 0.224]$} & \makecell[c]{$0.200\,\pm\,0.010$\\$[0.188,\ 0.212]$} & \makecell[c]{$0.225\,\pm\,0.011$\\$[0.211,\ 0.238]$} & \makecell[c]{$1.741\,\pm\,1.460$\\$[-0.073,\ 3.554]$} \\
Rolling mean & \makecell[c]{$0.194\,\pm\,0.020$\\$[0.169,\ 0.219]$} & \makecell[c]{$0.209\,\pm\,0.008$\\$[0.199,\ 0.218]$} & \makecell[c]{$0.218\,\pm\,0.010$\\$[0.206,\ 0.230]$} & \makecell[c]{$1.813\,\pm\,1.483$\\$[-0.029,\ 3.655]$} \\
KNN (scaled) & \makecell[c]{$0.174\,\pm\,0.023$\\$[0.145,\ 0.202]$} & \makecell[c]{$0.179\,\pm\,0.017$\\$[0.158,\ 0.200]$} & \makecell[c]{$0.178\,\pm\,0.013$\\$[0.162,\ 0.195]$} & \makecell[c]{$1.846\,\pm\,0.799$\\$[0.854,\ 2.838]$} \\
KNN (unscaled) & \makecell[c]{$0.228\,\pm\,0.042$\\$[0.176,\ 0.280]$} & \makecell[c]{$0.222\,\pm\,0.022$\\$[0.195,\ 0.249]$} & \makecell[c]{$0.215\,\pm\,0.019$\\$[0.191,\ 0.239]$} & \makecell[c]{$1.846\,\pm\,0.799$\\$[0.854,\ 2.838]$} \\
Historical & \makecell[c]{$0.467\,\pm\,0.036$\\$[0.422,\ 0.513]$} & \makecell[c]{$0.433\,\pm\,0.022$\\$[0.406,\ 0.461]$} & \makecell[c]{$0.472\,\pm\,0.033$\\$[0.432,\ 0.513]$} & \makecell[c]{$\mathbf{0.559}\,\pm\,0.351$\\$[0.124,\ 0.995]$} \\
Graph Laplacian & \makecell[c]{$0.636\,\pm\,0.127$\\$[0.478,\ 0.794]$} & \makecell[c]{$0.667\,\pm\,0.071$\\$[0.578,\ 0.756]$} & \makecell[c]{$0.666\,\pm\,0.056$\\$[0.597,\ 0.735]$} & \makecell[c]{$1.846\,\pm\,0.799$\\$[0.854,\ 2.838]$} \\
Graph Temporal & \makecell[c]{$0.556\,\pm\,0.102$\\$[0.429,\ 0.682]$} & \makecell[c]{$0.640\,\pm\,0.048$\\$[0.580,\ 0.699]$} & \makecell[c]{$0.706\,\pm\,0.046$\\$[0.649,\ 0.763]$} & \makecell[c]{$1.739\,\pm\,1.228$\\$[0.214,\ 3.263]$} \\
Sage-Voronoi & \makecell[c]{$0.165\,\pm\,0.027$\\$[0.131,\ 0.198]$} & \makecell[c]{$0.173\,\pm\,0.009$\\$[0.162,\ 0.184]$} & \makecell[c]{$0.194\,\pm\,0.021$\\$[0.168,\ 0.220]$} & \makecell[c]{$1.196\,\pm\,0.847$\\$[0.144,\ 2.248]$} \\
GAIN & \makecell[c]{$0.165\,\pm\,0.016$\\$[0.145,\ 0.185]$} & \makecell[c]{$0.299\,\pm\,0.052$\\$[0.235,\ 0.363]$} & \makecell[c]{$0.693\,\pm\,0.074$\\$[0.602,\ 0.785]$} & \makecell[c]{$2.969\,\pm\,1.210$\\$[1.466,\ 4.472]$} \\
GRIN & \makecell[c]{$\mathbf{0.110}\,\pm\,0.015$\\$[0.092,\ 0.129]$} & \makecell[c]{$\mathbf{0.129}\,\pm\,0.012$\\$[0.114,\ 0.143]$} & \makecell[c]{$\mathbf{0.151}\,\pm\,0.007$\\$[0.143,\ 0.160]$} & \makecell[c]{$1.782\,\pm\,1.164$\\$[0.337,\ 3.227]$} \\
\bottomrule
\end{tabular}
\endgroup
\end{table}

\clearpage

\begin{table}[htbp]
\centering
\caption{{MAE for sensor A142, January 2023 (vehicles per ten-minute interval). Each cell gives mean $\pm$ sample SD, followed by the 95\% confidence interval; $n=5$ seed runs. Bold means identify the lowest unrounded mean within each column, not statistical significance.}}
\label{tab:supp_p01_73}
\begingroup
\footnotesize
\setlength{\tabcolsep}{3pt}
\renewcommand{\arraystretch}{1.15}
\begin{tabular}{lcccc}
\toprule
Method & Block $B=1$ & Block $B=2$ & Block $B=4$ & Block $B=8$ \\
\midrule
Linear & \makecell[c]{$70.30\,\pm\,73.46$\\$[-20.91,\ 161.51]$} & \makecell[c]{$22.53\,\pm\,11.87$\\$[7.79,\ 37.27]$} & \makecell[c]{$19.66\,\pm\,4.74$\\$[13.77,\ 25.55]$} & \makecell[c]{$48.67\,\pm\,71.58$\\$[-40.21,\ 137.54]$} \\
Median & \makecell[c]{$110.00\,\pm\,83.74$\\$[6.03,\ 213.97]$} & \makecell[c]{$79.00\,\pm\,58.32$\\$[6.59,\ 151.41]$} & \makecell[c]{$46.60\,\pm\,15.48$\\$[27.37,\ 65.83]$} & \makecell[c]{$140.45\,\pm\,86.01$\\$[33.65,\ 247.25]$} \\
FFill + BFill & \makecell[c]{$58.40\,\pm\,69.34$\\$[-27.70,\ 144.50]$} & \makecell[c]{$32.80\,\pm\,22.94$\\$[4.31,\ 61.29]$} & \makecell[c]{$34.70\,\pm\,11.27$\\$[20.71,\ 48.69]$} & \makecell[c]{$56.00\,\pm\,56.64$\\$[-14.33,\ 126.33]$} \\
Rolling mean & \makecell[c]{$63.55\,\pm\,90.75$\\$[-49.13,\ 176.23]$} & \makecell[c]{$41.91\,\pm\,17.40$\\$[20.31,\ 63.51]$} & \makecell[c]{$29.53\,\pm\,5.98$\\$[22.10,\ 36.96]$} & \makecell[c]{$47.99\,\pm\,52.43$\\$[-17.11,\ 113.09]$} \\
KNN (scaled) & \makecell[c]{$194.58\,\pm\,85.19$\\$[88.81,\ 300.35]$} & \makecell[c]{$80.04\,\pm\,51.49$\\$[16.11,\ 143.97]$} & \makecell[c]{$105.28\,\pm\,46.68$\\$[47.32,\ 163.25]$} & \makecell[c]{$72.47\,\pm\,52.68$\\$[7.06,\ 137.89]$} \\
KNN (unscaled) & \makecell[c]{$70.44\,\pm\,81.68$\\$[-30.98,\ 171.86]$} & \makecell[c]{$27.50\,\pm\,10.91$\\$[13.95,\ 41.05]$} & \makecell[c]{$26.16\,\pm\,8.64$\\$[15.43,\ 36.89]$} & \makecell[c]{$34.23\,\pm\,8.42$\\$[23.77,\ 44.68]$} \\
Historical & \makecell[c]{$88.62\,\pm\,91.63$\\$[-25.15,\ 202.39]$} & \makecell[c]{$36.77\,\pm\,15.97$\\$[16.94,\ 56.60]$} & \makecell[c]{$29.26\,\pm\,14.84$\\$[10.84,\ 47.69]$} & \makecell[c]{$36.55\,\pm\,14.93$\\$[18.01,\ 55.10]$} \\
Graph Laplacian & \makecell[c]{$83.31\,\pm\,60.39$\\$[8.33,\ 158.30]$} & \makecell[c]{$33.24\,\pm\,18.17$\\$[10.68,\ 55.81]$} & \makecell[c]{$51.68\,\pm\,22.46$\\$[23.80,\ 79.56]$} & \makecell[c]{$56.06\,\pm\,25.68$\\$[24.17,\ 87.96]$} \\
Graph Temporal & \makecell[c]{$71.90\,\pm\,73.14$\\$[-18.92,\ 162.72]$} & \makecell[c]{$31.48\,\pm\,18.36$\\$[8.68,\ 54.27]$} & \makecell[c]{$30.78\,\pm\,14.22$\\$[13.13,\ 48.44]$} & \makecell[c]{$46.38\,\pm\,16.68$\\$[25.67,\ 67.09]$} \\
Sage-Voronoi & \makecell[c]{$\mathbf{55.59}\,\pm\,69.09$\\$[-30.19,\ 141.38]$} & \makecell[c]{$32.34\,\pm\,18.10$\\$[9.87,\ 54.81]$} & \makecell[c]{$23.43\,\pm\,5.55$\\$[16.54,\ 30.32]$} & \makecell[c]{$25.35\,\pm\,19.78$\\$[0.79,\ 49.92]$} \\
GAIN & \makecell[c]{$79.97\,\pm\,90.98$\\$[-33.00,\ 192.94]$} & \makecell[c]{$27.41\,\pm\,21.25$\\$[1.03,\ 53.80]$} & \makecell[c]{$37.60\,\pm\,19.13$\\$[13.84,\ 61.35]$} & \makecell[c]{$29.13\,\pm\,24.14$\\$[-0.85,\ 59.10]$} \\
GRIN & \makecell[c]{$66.24\,\pm\,87.01$\\$[-41.80,\ 174.28]$} & \makecell[c]{$\mathbf{16.24}\,\pm\,7.66$\\$[6.72,\ 25.75]$} & \makecell[c]{$\mathbf{16.24}\,\pm\,7.95$\\$[6.37,\ 26.10]$} & \makecell[c]{$\mathbf{16.28}\,\pm\,12.03$\\$[1.34,\ 31.22]$} \\
\bottomrule
\end{tabular}
\endgroup
\end{table}

\begin{table}[htbp]
\centering
\caption{{MAE for sensor A142, January 2023 (vehicles per ten-minute interval). Each cell gives mean $\pm$ sample SD, followed by the 95\% confidence interval; $n=5$ seed runs. Bold means identify the lowest unrounded mean within each column, not statistical significance.}}
\label{tab:supp_p01_74}
\begingroup
\footnotesize
\setlength{\tabcolsep}{3pt}
\renewcommand{\arraystretch}{1.15}
\begin{tabular}{lccc}
\toprule
Method & Block $B=16$ & Block $B=32$ & Block $B=64$ \\
\midrule
Linear & \makecell[c]{$31.35\,\pm\,4.84$\\$[25.34,\ 37.36]$} & \makecell[c]{$40.62\,\pm\,18.64$\\$[17.48,\ 63.76]$} & \makecell[c]{$94.26\,\pm\,59.34$\\$[20.58,\ 167.94]$} \\
Median & \makecell[c]{$43.09\,\pm\,17.11$\\$[21.85,\ 64.33]$} & \makecell[c]{$70.88\,\pm\,31.71$\\$[31.50,\ 110.26]$} & \makecell[c]{$85.19\,\pm\,36.67$\\$[39.66,\ 130.72]$} \\
FFill + BFill & \makecell[c]{$47.06\,\pm\,28.58$\\$[11.57,\ 82.55]$} & \makecell[c]{$75.43\,\pm\,48.10$\\$[15.71,\ 135.16]$} & \makecell[c]{$129.42\,\pm\,91.39$\\$[15.94,\ 242.90]$} \\
Rolling mean & \makecell[c]{$31.92\,\pm\,8.12$\\$[21.84,\ 42.01]$} & \makecell[c]{$42.43\,\pm\,16.06$\\$[22.49,\ 62.37]$} & \makecell[c]{$96.06\,\pm\,57.65$\\$[24.48,\ 167.64]$} \\
KNN (scaled) & \makecell[c]{$115.91\,\pm\,27.04$\\$[82.34,\ 149.48]$} & \makecell[c]{$116.31\,\pm\,36.01$\\$[71.59,\ 161.02]$} & \makecell[c]{$87.13\,\pm\,36.90$\\$[41.31,\ 132.95]$} \\
KNN (unscaled) & \makecell[c]{$35.52\,\pm\,12.19$\\$[20.38,\ 50.66]$} & \makecell[c]{$35.15\,\pm\,10.36$\\$[22.29,\ 48.01]$} & \makecell[c]{$33.45\,\pm\,5.36$\\$[26.79,\ 40.11]$} \\
Historical & \makecell[c]{$30.49\,\pm\,14.85$\\$[12.05,\ 48.93]$} & \makecell[c]{$35.50\,\pm\,4.82$\\$[29.52,\ 41.48]$} & \makecell[c]{$47.14\,\pm\,14.41$\\$[29.25,\ 65.03]$} \\
Graph Laplacian & \makecell[c]{$64.01\,\pm\,11.77$\\$[49.39,\ 78.62]$} & \makecell[c]{$62.07\,\pm\,22.05$\\$[34.69,\ 89.45]$} & \makecell[c]{$54.88\,\pm\,8.36$\\$[44.50,\ 65.27]$} \\
Graph Temporal & \makecell[c]{$60.48\,\pm\,12.71$\\$[44.69,\ 76.26]$} & \makecell[c]{$46.01\,\pm\,13.21$\\$[29.61,\ 62.41]$} & \makecell[c]{$45.73\,\pm\,4.06$\\$[40.69,\ 50.77]$} \\
Sage-Voronoi & \makecell[c]{$39.02\,\pm\,7.92$\\$[29.19,\ 48.85]$} & \makecell[c]{$45.94\,\pm\,11.89$\\$[31.17,\ 60.70]$} & \makecell[c]{$48.36\,\pm\,12.68$\\$[32.62,\ 64.10]$} \\
GAIN & \makecell[c]{$49.01\,\pm\,28.85$\\$[13.19,\ 84.83]$} & \makecell[c]{$42.58\,\pm\,15.69$\\$[23.09,\ 62.06]$} & \makecell[c]{$37.13\,\pm\,11.90$\\$[22.35,\ 51.91]$} \\
GRIN & \makecell[c]{$\mathbf{23.04}\,\pm\,7.51$\\$[13.71,\ 32.36]$} & \makecell[c]{$\mathbf{25.58}\,\pm\,7.88$\\$[15.79,\ 35.37]$} & \makecell[c]{$\mathbf{25.80}\,\pm\,11.30$\\$[11.77,\ 39.83]$} \\
\bottomrule
\end{tabular}
\endgroup
\end{table}

\clearpage

\begin{table}[htbp]
\centering
\caption{{MAE for sensor A142, January 2023 (vehicles per ten-minute interval). Each cell gives mean $\pm$ sample SD, followed by the 95\% confidence interval; $n=5$ seed runs. Bold means identify the lowest unrounded mean within each column, not statistical significance.}}
\label{tab:supp_p01_75}
\begingroup
\footnotesize
\setlength{\tabcolsep}{3pt}
\renewcommand{\arraystretch}{1.15}
\begin{tabular}{lcc}
\toprule
Method & Block $B=128$ & Block $B=144$ \\
\midrule
Linear & \makecell[c]{$97.69\,\pm\,21.64$\\$[70.82,\ 124.56]$} & \makecell[c]{$141.76\,\pm\,40.99$\\$[90.86,\ 192.66]$} \\
Median & \makecell[c]{$107.13\,\pm\,14.85$\\$[88.70,\ 125.56]$} & \makecell[c]{$102.05\,\pm\,10.44$\\$[89.09,\ 115.01]$} \\
FFill + BFill & \makecell[c]{$124.96\,\pm\,20.58$\\$[99.41,\ 150.52]$} & \makecell[c]{$141.00\,\pm\,33.91$\\$[98.89,\ 183.10]$} \\
Rolling mean & \makecell[c]{$97.63\,\pm\,21.27$\\$[71.22,\ 124.04]$} & \makecell[c]{$141.31\,\pm\,41.34$\\$[89.98,\ 192.64]$} \\
KNN (scaled) & \makecell[c]{$104.97\,\pm\,12.44$\\$[89.52,\ 120.41]$} & \makecell[c]{$107.61\,\pm\,14.59$\\$[89.49,\ 125.72]$} \\
KNN (unscaled) & \makecell[c]{$33.86\,\pm\,1.77$\\$[31.66,\ 36.06]$} & \makecell[c]{$32.27\,\pm\,2.15$\\$[29.60,\ 34.95]$} \\
Historical & \makecell[c]{$46.92\,\pm\,2.95$\\$[43.25,\ 50.58]$} & \makecell[c]{$44.05\,\pm\,4.26$\\$[38.76,\ 49.34]$} \\
Graph Laplacian & \makecell[c]{$58.59\,\pm\,6.42$\\$[50.63,\ 66.56]$} & \makecell[c]{$58.67\,\pm\,6.44$\\$[50.67,\ 66.67]$} \\
Graph Temporal & \makecell[c]{$48.09\,\pm\,6.02$\\$[40.61,\ 55.56]$} & \makecell[c]{$47.07\,\pm\,2.88$\\$[43.49,\ 50.66]$} \\
Sage-Voronoi & \makecell[c]{$53.73\,\pm\,10.16$\\$[41.11,\ 66.34]$} & \makecell[c]{$50.56\,\pm\,8.87$\\$[39.55,\ 61.57]$} \\
GAIN & \makecell[c]{$35.53\,\pm\,5.66$\\$[28.51,\ 42.56]$} & \makecell[c]{$34.66\,\pm\,4.04$\\$[29.65,\ 39.67]$} \\
GRIN & \makecell[c]{$\mathbf{31.02}\,\pm\,5.77$\\$[23.85,\ 38.18]$} & \makecell[c]{$\mathbf{26.74}\,\pm\,1.62$\\$[24.73,\ 28.75]$} \\
\bottomrule
\end{tabular}
\endgroup
\end{table}

\begin{table}[htbp]
\centering
\caption{{MAE for sensor A142, January 2023 (vehicles per ten-minute interval). Each cell gives mean $\pm$ sample SD, followed by the 95\% confidence interval; $n=5$ seed runs. Bold means identify the lowest unrounded mean within each column, not statistical significance.}}
\label{tab:supp_p01_76}
\begingroup
\footnotesize
\setlength{\tabcolsep}{3pt}
\renewcommand{\arraystretch}{1.15}
\begin{tabular}{lcccc}
\toprule
Method & MCAR 10\% & MCAR 30\% & MCAR 50\% & Synchronous $B=144$ \\
\midrule
Linear & \makecell[c]{$26.16\,\pm\,1.90$\\$[23.79,\ 28.52]$} & \makecell[c]{$25.20\,\pm\,0.83$\\$[24.17,\ 26.24]$} & \makecell[c]{$25.57\,\pm\,0.75$\\$[24.63,\ 26.50]$} & \makecell[c]{$112.72\,\pm\,27.12$\\$[79.05,\ 146.40]$} \\
Median & \makecell[c]{$102.15\,\pm\,7.43$\\$[92.93,\ 111.37]$} & \makecell[c]{$101.29\,\pm\,2.23$\\$[98.52,\ 104.06]$} & \makecell[c]{$100.11\,\pm\,1.76$\\$[97.93,\ 102.30]$} & \makecell[c]{$101.49\,\pm\,5.86$\\$[94.21,\ 108.77]$} \\
FFill + BFill & \makecell[c]{$32.36\,\pm\,2.83$\\$[28.84,\ 35.87]$} & \makecell[c]{$31.30\,\pm\,1.86$\\$[28.99,\ 33.62]$} & \makecell[c]{$33.45\,\pm\,1.18$\\$[31.99,\ 34.92]$} & \makecell[c]{$128.00\,\pm\,49.32$\\$[66.75,\ 189.24]$} \\
Rolling mean & \makecell[c]{$32.49\,\pm\,2.78$\\$[29.03,\ 35.94]$} & \makecell[c]{$31.49\,\pm\,2.33$\\$[28.60,\ 34.38]$} & \makecell[c]{$32.74\,\pm\,0.90$\\$[31.62,\ 33.85]$} & \makecell[c]{$113.40\,\pm\,27.78$\\$[78.90,\ 147.90]$} \\
KNN (scaled) & \makecell[c]{$101.22\,\pm\,3.16$\\$[97.30,\ 105.14]$} & \makecell[c]{$96.21\,\pm\,1.52$\\$[94.32,\ 98.09]$} & \makecell[c]{$85.42\,\pm\,2.15$\\$[82.75,\ 88.08]$} & \makecell[c]{$106.75\,\pm\,6.12$\\$[99.15,\ 114.36]$} \\
KNN (unscaled) & \makecell[c]{$34.24\,\pm\,1.78$\\$[32.03,\ 36.45]$} & \makecell[c]{$33.87\,\pm\,1.22$\\$[32.36,\ 35.38]$} & \makecell[c]{$35.04\,\pm\,0.33$\\$[34.63,\ 35.46]$} & \makecell[c]{$106.75\,\pm\,6.12$\\$[99.15,\ 114.36]$} \\
Historical & \makecell[c]{$44.24\,\pm\,2.82$\\$[40.73,\ 47.74]$} & \makecell[c]{$46.04\,\pm\,1.50$\\$[44.18,\ 47.91]$} & \makecell[c]{$46.07\,\pm\,1.03$\\$[44.79,\ 47.34]$} & \makecell[c]{$\mathbf{45.65}\,\pm\,15.50$\\$[26.41,\ 64.90]$} \\
Graph Laplacian & \makecell[c]{$58.09\,\pm\,3.54$\\$[53.69,\ 62.49]$} & \makecell[c]{$58.33\,\pm\,1.31$\\$[56.70,\ 59.95]$} & \makecell[c]{$60.93\,\pm\,0.88$\\$[59.83,\ 62.03]$} & \makecell[c]{$106.75\,\pm\,6.12$\\$[99.15,\ 114.36]$} \\
Graph Temporal & \makecell[c]{$41.10\,\pm\,2.74$\\$[37.69,\ 44.50]$} & \makecell[c]{$40.92\,\pm\,1.12$\\$[39.53,\ 42.31]$} & \makecell[c]{$43.13\,\pm\,0.89$\\$[42.03,\ 44.24]$} & \makecell[c]{$163.51\,\pm\,59.33$\\$[89.84,\ 237.19]$} \\
Sage-Voronoi & \makecell[c]{$28.49\,\pm\,1.59$\\$[26.52,\ 30.47]$} & \makecell[c]{$29.43\,\pm\,1.91$\\$[27.06,\ 31.81]$} & \makecell[c]{$31.53\,\pm\,1.29$\\$[29.93,\ 33.13]$} & \makecell[c]{$113.25\,\pm\,32.08$\\$[73.42,\ 153.09]$} \\
GAIN & \makecell[c]{$35.78\,\pm\,5.58$\\$[28.86,\ 42.71]$} & \makecell[c]{$38.00\,\pm\,3.38$\\$[33.81,\ 42.20]$} & \makecell[c]{$48.15\,\pm\,3.04$\\$[44.38,\ 51.93]$} & \makecell[c]{$133.28\,\pm\,33.92$\\$[91.17,\ 175.40]$} \\
GRIN & \makecell[c]{$\mathbf{20.62}\,\pm\,1.83$\\$[18.34,\ 22.90]$} & \makecell[c]{$\mathbf{20.02}\,\pm\,0.66$\\$[19.20,\ 20.84]$} & \makecell[c]{$\mathbf{22.45}\,\pm\,0.86$\\$[21.39,\ 23.52]$} & \makecell[c]{$109.96\,\pm\,8.48$\\$[99.43,\ 120.48]$} \\
\bottomrule
\end{tabular}
\endgroup
\end{table}

\clearpage

\begin{table}[htbp]
\centering
\caption{{RMSE for sensor A142, January 2023 (vehicles per ten-minute interval). Each cell gives mean $\pm$ sample SD, followed by the 95\% confidence interval; $n=5$ seed runs. Bold means identify the lowest unrounded mean within each column, not statistical significance.}}
\label{tab:supp_p01_77}
\begingroup
\footnotesize
\setlength{\tabcolsep}{3pt}
\renewcommand{\arraystretch}{1.15}
\begin{tabular}{lcccc}
\toprule
Method & Block $B=1$ & Block $B=2$ & Block $B=4$ & Block $B=8$ \\
\midrule
Linear & \makecell[c]{$70.30\,\pm\,73.46$\\$[-20.91,\ 161.51]$} & \makecell[c]{$26.50\,\pm\,14.64$\\$[8.33,\ 44.68]$} & \makecell[c]{$23.66\,\pm\,6.79$\\$[15.23,\ 32.10]$} & \makecell[c]{$53.69\,\pm\,72.03$\\$[-35.74,\ 143.13]$} \\
Median & \makecell[c]{$110.00\,\pm\,83.74$\\$[6.03,\ 213.97]$} & \makecell[c]{$83.23\,\pm\,57.47$\\$[11.87,\ 154.59]$} & \makecell[c]{$49.98\,\pm\,15.89$\\$[30.25,\ 69.71]$} & \makecell[c]{$144.04\,\pm\,82.58$\\$[41.50,\ 246.58]$} \\
FFill + BFill & \makecell[c]{$58.40\,\pm\,69.34$\\$[-27.70,\ 144.50]$} & \makecell[c]{$38.51\,\pm\,27.81$\\$[3.97,\ 73.04]$} & \makecell[c]{$37.80\,\pm\,12.95$\\$[21.72,\ 53.88]$} & \makecell[c]{$60.45\,\pm\,55.84$\\$[-8.89,\ 129.79]$} \\
Rolling mean & \makecell[c]{$63.55\,\pm\,90.75$\\$[-49.13,\ 176.23]$} & \makecell[c]{$46.64\,\pm\,21.20$\\$[20.32,\ 72.96]$} & \makecell[c]{$35.03\,\pm\,6.64$\\$[26.79,\ 43.27]$} & \makecell[c]{$55.00\,\pm\,57.36$\\$[-16.23,\ 126.23]$} \\
KNN (scaled) & \makecell[c]{$194.58\,\pm\,85.19$\\$[88.81,\ 300.35]$} & \makecell[c]{$81.14\,\pm\,51.65$\\$[17.01,\ 145.27]$} & \makecell[c]{$106.93\,\pm\,45.77$\\$[50.10,\ 163.75]$} & \makecell[c]{$77.59\,\pm\,51.76$\\$[13.32,\ 141.86]$} \\
KNN (unscaled) & \makecell[c]{$70.44\,\pm\,81.68$\\$[-30.98,\ 171.86]$} & \makecell[c]{$30.85\,\pm\,13.54$\\$[14.04,\ 47.66]$} & \makecell[c]{$29.29\,\pm\,10.36$\\$[16.43,\ 42.15]$} & \makecell[c]{$40.50\,\pm\,8.73$\\$[29.67,\ 51.34]$} \\
Historical & \makecell[c]{$88.62\,\pm\,91.63$\\$[-25.15,\ 202.39]$} & \makecell[c]{$39.85\,\pm\,17.25$\\$[18.43,\ 61.28]$} & \makecell[c]{$34.34\,\pm\,14.83$\\$[15.93,\ 52.75]$} & \makecell[c]{$40.06\,\pm\,14.52$\\$[22.02,\ 58.09]$} \\
Graph Laplacian & \makecell[c]{$83.31\,\pm\,60.39$\\$[8.33,\ 158.30]$} & \makecell[c]{$36.93\,\pm\,19.12$\\$[13.20,\ 60.67]$} & \makecell[c]{$53.70\,\pm\,22.31$\\$[26.00,\ 81.41]$} & \makecell[c]{$59.59\,\pm\,25.87$\\$[27.47,\ 91.71]$} \\
Graph Temporal & \makecell[c]{$71.90\,\pm\,73.14$\\$[-18.92,\ 162.72]$} & \makecell[c]{$33.29\,\pm\,20.01$\\$[8.44,\ 58.13]$} & \makecell[c]{$36.45\,\pm\,13.16$\\$[20.11,\ 52.80]$} & \makecell[c]{$49.48\,\pm\,18.42$\\$[26.60,\ 72.35]$} \\
Sage-Voronoi & \makecell[c]{$\mathbf{55.59}\,\pm\,69.09$\\$[-30.19,\ 141.38]$} & \makecell[c]{$36.25\,\pm\,23.07$\\$[7.60,\ 64.89]$} & \makecell[c]{$27.90\,\pm\,5.38$\\$[21.22,\ 34.58]$} & \makecell[c]{$29.41\,\pm\,20.95$\\$[3.39,\ 55.42]$} \\
GAIN & \makecell[c]{$79.97\,\pm\,90.98$\\$[-33.00,\ 192.94]$} & \makecell[c]{$31.98\,\pm\,26.02$\\$[-0.33,\ 64.29]$} & \makecell[c]{$39.85\,\pm\,19.79$\\$[15.27,\ 64.43]$} & \makecell[c]{$33.53\,\pm\,26.76$\\$[0.31,\ 66.75]$} \\
GRIN & \makecell[c]{$66.24\,\pm\,87.01$\\$[-41.80,\ 174.28]$} & \makecell[c]{$\mathbf{19.41}\,\pm\,10.93$\\$[5.83,\ 32.98]$} & \makecell[c]{$\mathbf{18.00}\,\pm\,8.78$\\$[7.09,\ 28.91]$} & \makecell[c]{$\mathbf{20.01}\,\pm\,14.62$\\$[1.86,\ 38.17]$} \\
\bottomrule
\end{tabular}
\endgroup
\end{table}

\begin{table}[htbp]
\centering
\caption{{RMSE for sensor A142, January 2023 (vehicles per ten-minute interval). Each cell gives mean $\pm$ sample SD, followed by the 95\% confidence interval; $n=5$ seed runs. Bold means identify the lowest unrounded mean within each column, not statistical significance.}}
\label{tab:supp_p01_78}
\begingroup
\footnotesize
\setlength{\tabcolsep}{3pt}
\renewcommand{\arraystretch}{1.15}
\begin{tabular}{lccc}
\toprule
Method & Block $B=16$ & Block $B=32$ & Block $B=64$ \\
\midrule
Linear & \makecell[c]{$41.05\,\pm\,11.00$\\$[27.40,\ 54.71]$} & \makecell[c]{$54.47\,\pm\,21.84$\\$[27.36,\ 81.59]$} & \makecell[c]{$115.31\,\pm\,70.28$\\$[28.05,\ 202.58]$} \\
Median & \makecell[c]{$56.04\,\pm\,18.61$\\$[32.93,\ 79.15]$} & \makecell[c]{$87.52\,\pm\,35.79$\\$[43.09,\ 131.96]$} & \makecell[c]{$103.31\,\pm\,43.25$\\$[49.61,\ 157.02]$} \\
FFill + BFill & \makecell[c]{$57.24\,\pm\,29.98$\\$[20.02,\ 94.46]$} & \makecell[c]{$91.04\,\pm\,54.15$\\$[23.80,\ 158.27]$} & \makecell[c]{$147.16\,\pm\,101.00$\\$[21.75,\ 272.57]$} \\
Rolling mean & \makecell[c]{$41.11\,\pm\,13.95$\\$[23.80,\ 58.43]$} & \makecell[c]{$55.55\,\pm\,19.73$\\$[31.05,\ 80.05]$} & \makecell[c]{$117.70\,\pm\,69.61$\\$[31.27,\ 204.13]$} \\
KNN (scaled) & \makecell[c]{$124.04\,\pm\,25.95$\\$[91.82,\ 156.25]$} & \makecell[c]{$127.02\,\pm\,35.05$\\$[83.49,\ 170.54]$} & \makecell[c]{$101.97\,\pm\,34.29$\\$[59.39,\ 144.54]$} \\
KNN (unscaled) & \makecell[c]{$46.69\,\pm\,16.12$\\$[26.67,\ 66.70]$} & \makecell[c]{$45.25\,\pm\,15.25$\\$[26.32,\ 64.18]$} & \makecell[c]{$43.50\,\pm\,7.51$\\$[34.18,\ 52.82]$} \\
Historical & \makecell[c]{$38.99\,\pm\,17.32$\\$[17.49,\ 60.50]$} & \makecell[c]{$48.31\,\pm\,10.51$\\$[35.26,\ 61.36]$} & \makecell[c]{$57.78\,\pm\,16.05$\\$[37.85,\ 77.71]$} \\
Graph Laplacian & \makecell[c]{$70.25\,\pm\,11.21$\\$[56.33,\ 84.17]$} & \makecell[c]{$69.01\,\pm\,21.20$\\$[42.69,\ 95.33]$} & \makecell[c]{$62.89\,\pm\,8.83$\\$[51.93,\ 73.85]$} \\
Graph Temporal & \makecell[c]{$67.83\,\pm\,13.76$\\$[50.74,\ 84.92]$} & \makecell[c]{$54.47\,\pm\,15.00$\\$[35.84,\ 73.10]$} & \makecell[c]{$53.84\,\pm\,6.39$\\$[45.90,\ 61.78]$} \\
Sage-Voronoi & \makecell[c]{$48.46\,\pm\,7.75$\\$[38.83,\ 58.08]$} & \makecell[c]{$56.52\,\pm\,10.01$\\$[44.10,\ 68.94]$} & \makecell[c]{$61.04\,\pm\,15.73$\\$[41.51,\ 80.57]$} \\
GAIN & \makecell[c]{$60.17\,\pm\,29.77$\\$[23.20,\ 97.13]$} & \makecell[c]{$54.82\,\pm\,21.49$\\$[28.14,\ 81.50]$} & \makecell[c]{$49.12\,\pm\,14.92$\\$[30.60,\ 67.65]$} \\
GRIN & \makecell[c]{$\mathbf{32.45}\,\pm\,12.55$\\$[16.87,\ 48.04]$} & \makecell[c]{$\mathbf{36.66}\,\pm\,13.91$\\$[19.38,\ 53.93]$} & \makecell[c]{$\mathbf{32.81}\,\pm\,13.18$\\$[16.44,\ 49.18]$} \\
\bottomrule
\end{tabular}
\endgroup
\end{table}

\clearpage

\begin{table}[htbp]
\centering
\caption{{RMSE for sensor A142, January 2023 (vehicles per ten-minute interval). Each cell gives mean $\pm$ sample SD, followed by the 95\% confidence interval; $n=5$ seed runs. Bold means identify the lowest unrounded mean within each column, not statistical significance.}}
\label{tab:supp_p01_79}
\begingroup
\footnotesize
\setlength{\tabcolsep}{3pt}
\renewcommand{\arraystretch}{1.15}
\begin{tabular}{lcc}
\toprule
Method & Block $B=128$ & Block $B=144$ \\
\midrule
Linear & \makecell[c]{$126.94\,\pm\,27.54$\\$[92.74,\ 161.13]$} & \makecell[c]{$170.20\,\pm\,39.36$\\$[121.33,\ 219.06]$} \\
Median & \makecell[c]{$132.34\,\pm\,13.40$\\$[115.69,\ 148.98]$} & \makecell[c]{$127.30\,\pm\,10.03$\\$[114.84,\ 139.75]$} \\
FFill + BFill & \makecell[c]{$153.28\,\pm\,13.65$\\$[136.33,\ 170.23]$} & \makecell[c]{$165.80\,\pm\,36.06$\\$[121.03,\ 210.57]$} \\
Rolling mean & \makecell[c]{$126.87\,\pm\,27.16$\\$[93.15,\ 160.59]$} & \makecell[c]{$169.79\,\pm\,39.62$\\$[120.59,\ 218.99]$} \\
KNN (scaled) & \makecell[c]{$125.38\,\pm\,14.13$\\$[107.84,\ 142.92]$} & \makecell[c]{$127.19\,\pm\,17.25$\\$[105.78,\ 148.61]$} \\
KNN (unscaled) & \makecell[c]{$48.52\,\pm\,4.81$\\$[42.54,\ 54.50]$} & \makecell[c]{$46.30\,\pm\,3.83$\\$[41.54,\ 51.05]$} \\
Historical & \makecell[c]{$64.40\,\pm\,5.86$\\$[57.13,\ 71.68]$} & \makecell[c]{$60.80\,\pm\,9.19$\\$[49.40,\ 72.21]$} \\
Graph Laplacian & \makecell[c]{$69.24\,\pm\,7.23$\\$[60.26,\ 78.22]$} & \makecell[c]{$68.26\,\pm\,7.22$\\$[59.30,\ 77.23]$} \\
Graph Temporal & \makecell[c]{$60.06\,\pm\,6.08$\\$[52.51,\ 67.61]$} & \makecell[c]{$57.14\,\pm\,4.08$\\$[52.07,\ 62.20]$} \\
Sage-Voronoi & \makecell[c]{$70.12\,\pm\,12.87$\\$[54.14,\ 86.09]$} & \makecell[c]{$65.28\,\pm\,10.49$\\$[52.26,\ 78.30]$} \\
GAIN & \makecell[c]{$51.37\,\pm\,7.73$\\$[41.78,\ 60.97]$} & \makecell[c]{$49.42\,\pm\,5.01$\\$[43.20,\ 55.65]$} \\
GRIN & \makecell[c]{$\mathbf{45.92}\,\pm\,6.66$\\$[37.65,\ 54.19]$} & \makecell[c]{$\mathbf{39.79}\,\pm\,2.56$\\$[36.60,\ 42.97]$} \\
\bottomrule
\end{tabular}
\endgroup
\end{table}

\begin{table}[htbp]
\centering
\caption{{RMSE for sensor A142, January 2023 (vehicles per ten-minute interval). Each cell gives mean $\pm$ sample SD, followed by the 95\% confidence interval; $n=5$ seed runs. Bold means identify the lowest unrounded mean within each column, not statistical significance.}}
\label{tab:supp_p01_80}
\begingroup
\footnotesize
\setlength{\tabcolsep}{3pt}
\renewcommand{\arraystretch}{1.15}
\begin{tabular}{lcccc}
\toprule
Method & MCAR 10\% & MCAR 30\% & MCAR 50\% & Synchronous $B=144$ \\
\midrule
Linear & \makecell[c]{$38.10\,\pm\,3.47$\\$[33.80,\ 42.41]$} & \makecell[c]{$36.85\,\pm\,2.14$\\$[34.19,\ 39.51]$} & \makecell[c]{$36.61\,\pm\,1.43$\\$[34.84,\ 38.39]$} & \makecell[c]{$148.87\,\pm\,24.94$\\$[117.91,\ 179.84]$} \\
Median & \makecell[c]{$127.58\,\pm\,7.89$\\$[117.78,\ 137.38]$} & \makecell[c]{$125.95\,\pm\,2.30$\\$[123.10,\ 128.80]$} & \makecell[c]{$124.81\,\pm\,1.79$\\$[122.59,\ 127.04]$} & \makecell[c]{$125.68\,\pm\,7.68$\\$[116.14,\ 135.22]$} \\
FFill + BFill & \makecell[c]{$46.18\,\pm\,6.02$\\$[38.71,\ 53.65]$} & \makecell[c]{$45.71\,\pm\,2.97$\\$[42.02,\ 49.39]$} & \makecell[c]{$47.48\,\pm\,2.13$\\$[44.84,\ 50.13]$} & \makecell[c]{$164.29\,\pm\,47.52$\\$[105.28,\ 223.29]$} \\
Rolling mean & \makecell[c]{$44.50\,\pm\,4.19$\\$[39.30,\ 49.70]$} & \makecell[c]{$44.44\,\pm\,3.19$\\$[40.48,\ 48.41]$} & \makecell[c]{$46.05\,\pm\,1.70$\\$[43.94,\ 48.17]$} & \makecell[c]{$149.33\,\pm\,25.42$\\$[117.76,\ 180.90]$} \\
KNN (scaled) & \makecell[c]{$119.50\,\pm\,4.82$\\$[113.51,\ 125.48]$} & \makecell[c]{$117.77\,\pm\,1.75$\\$[115.60,\ 119.94]$} & \makecell[c]{$110.49\,\pm\,2.71$\\$[107.12,\ 113.86]$} & \makecell[c]{$124.29\,\pm\,4.61$\\$[118.57,\ 130.01]$} \\
KNN (unscaled) & \makecell[c]{$46.46\,\pm\,3.45$\\$[42.18,\ 50.74]$} & \makecell[c]{$46.06\,\pm\,2.38$\\$[43.10,\ 49.01]$} & \makecell[c]{$47.62\,\pm\,0.88$\\$[46.53,\ 48.71]$} & \makecell[c]{$124.29\,\pm\,4.61$\\$[118.57,\ 130.01]$} \\
Historical & \makecell[c]{$59.96\,\pm\,4.51$\\$[54.36,\ 65.57]$} & \makecell[c]{$63.58\,\pm\,2.72$\\$[60.20,\ 66.96]$} & \makecell[c]{$63.76\,\pm\,1.61$\\$[61.75,\ 65.76]$} & \makecell[c]{$\mathbf{60.69}\,\pm\,20.52$\\$[35.21,\ 86.17]$} \\
Graph Laplacian & \makecell[c]{$67.31\,\pm\,2.80$\\$[63.84,\ 70.79]$} & \makecell[c]{$69.13\,\pm\,1.57$\\$[67.18,\ 71.08]$} & \makecell[c]{$71.82\,\pm\,0.76$\\$[70.87,\ 72.77]$} & \makecell[c]{$124.29\,\pm\,4.61$\\$[118.57,\ 130.01]$} \\
Graph Temporal & \makecell[c]{$50.85\,\pm\,2.36$\\$[47.92,\ 53.77]$} & \makecell[c]{$51.35\,\pm\,2.13$\\$[48.71,\ 53.99]$} & \makecell[c]{$53.43\,\pm\,0.68$\\$[52.59,\ 54.28]$} & \makecell[c]{$199.46\,\pm\,61.02$\\$[123.69,\ 275.22]$} \\
Sage-Voronoi & \makecell[c]{$40.19\,\pm\,3.05$\\$[36.41,\ 43.98]$} & \makecell[c]{$41.54\,\pm\,2.65$\\$[38.25,\ 44.83]$} & \makecell[c]{$44.19\,\pm\,2.39$\\$[41.22,\ 47.15]$} & \makecell[c]{$149.97\,\pm\,34.76$\\$[106.80,\ 193.13]$} \\
GAIN & \makecell[c]{$49.25\,\pm\,6.49$\\$[41.18,\ 57.31]$} & \makecell[c]{$48.96\,\pm\,3.35$\\$[44.80,\ 53.11]$} & \makecell[c]{$59.37\,\pm\,3.80$\\$[54.65,\ 64.09]$} & \makecell[c]{$172.83\,\pm\,25.32$\\$[141.39,\ 204.27]$} \\
GRIN & \makecell[c]{$\mathbf{32.01}\,\pm\,3.79$\\$[27.31,\ 36.71]$} & \makecell[c]{$\mathbf{30.81}\,\pm\,1.86$\\$[28.50,\ 33.12]$} & \makecell[c]{$\mathbf{32.96}\,\pm\,1.31$\\$[31.33,\ 34.59]$} & \makecell[c]{$127.85\,\pm\,5.55$\\$[120.95,\ 134.74]$} \\
\bottomrule
\end{tabular}
\endgroup
\end{table}

\clearpage

\begin{table}[htbp]
\centering
\caption{{MRE for sensor A142, January 2023 (dimensionless ratio). Each cell gives mean $\pm$ sample SD, followed by the 95\% confidence interval; $n=5$ seed runs. Bold means identify the lowest unrounded mean within each column, not statistical significance.}}
\label{tab:supp_p01_81}
\begingroup
\footnotesize
\setlength{\tabcolsep}{3pt}
\renewcommand{\arraystretch}{1.15}
\begin{tabular}{lcccc}
\toprule
Method & Block $B=1$ & Block $B=2$ & Block $B=4$ & Block $B=8$ \\
\midrule
Linear & \makecell[c]{$0.172\,\pm\,0.146$\\$[-0.010,\ 0.353]$} & \makecell[c]{$0.117\,\pm\,0.094$\\$[-0.000,\ 0.233]$} & \makecell[c]{$0.075\,\pm\,0.025$\\$[0.043,\ 0.107]$} & \makecell[c]{$0.276\,\pm\,0.198$\\$[0.030,\ 0.523]$} \\
Median & \makecell[c]{$0.266\,\pm\,0.160$\\$[0.067,\ 0.466]$} & \makecell[c]{$0.643\,\pm\,0.798$\\$[-0.348,\ 1.635]$} & \makecell[c]{$0.176\,\pm\,0.094$\\$[0.059,\ 0.292]$} & \makecell[c]{$2.901\,\pm\,3.260$\\$[-1.147,\ 6.948]$} \\
FFill + BFill & \makecell[c]{$0.140\,\pm\,0.140$\\$[-0.034,\ 0.313]$} & \makecell[c]{$0.182\,\pm\,0.164$\\$[-0.022,\ 0.386]$} & \makecell[c]{$0.131\,\pm\,0.059$\\$[0.057,\ 0.205]$} & \makecell[c]{$0.449\,\pm\,0.312$\\$[0.061,\ 0.837]$} \\
Rolling mean & \makecell[c]{$0.151\,\pm\,0.186$\\$[-0.080,\ 0.381]$} & \makecell[c]{$0.242\,\pm\,0.199$\\$[-0.006,\ 0.489]$} & \makecell[c]{$0.116\,\pm\,0.056$\\$[0.047,\ 0.185]$} & \makecell[c]{$0.335\,\pm\,0.178$\\$[0.114,\ 0.556]$} \\
KNN (scaled) & \makecell[c]{$0.503\,\pm\,0.117$\\$[0.358,\ 0.648]$} & \makecell[c]{$0.342\,\pm\,0.118$\\$[0.195,\ 0.489]$} & \makecell[c]{$0.355\,\pm\,0.122$\\$[0.203,\ 0.506]$} & \makecell[c]{$0.694\,\pm\,0.483$\\$[0.094,\ 1.294]$} \\
KNN (unscaled) & \makecell[c]{$0.172\,\pm\,0.159$\\$[-0.025,\ 0.370]$} & \makecell[c]{$0.162\,\pm\,0.138$\\$[-0.010,\ 0.333]$} & \makecell[c]{$0.097\,\pm\,0.034$\\$[0.055,\ 0.140]$} & \makecell[c]{$0.481\,\pm\,0.432$\\$[-0.056,\ 1.018]$} \\
Historical & \makecell[c]{$0.207\,\pm\,0.181$\\$[-0.017,\ 0.432]$} & \makecell[c]{$0.189\,\pm\,0.090$\\$[0.076,\ 0.301]$} & \makecell[c]{$0.116\,\pm\,0.096$\\$[-0.003,\ 0.235]$} & \makecell[c]{$0.525\,\pm\,0.450$\\$[-0.035,\ 1.084]$} \\
Graph Laplacian & \makecell[c]{$0.218\,\pm\,0.116$\\$[0.074,\ 0.363]$} & \makecell[c]{$0.196\,\pm\,0.158$\\$[0.000,\ 0.392]$} & \makecell[c]{$0.178\,\pm\,0.057$\\$[0.108,\ 0.249]$} & \makecell[c]{$0.848\,\pm\,0.845$\\$[-0.202,\ 1.897]$} \\
Graph Temporal & \makecell[c]{$0.182\,\pm\,0.143$\\$[0.004,\ 0.360]$} & \makecell[c]{$0.179\,\pm\,0.138$\\$[0.007,\ 0.351]$} & \makecell[c]{$0.109\,\pm\,0.033$\\$[0.068,\ 0.151]$} & \makecell[c]{$0.715\,\pm\,0.720$\\$[-0.179,\ 1.609]$} \\
Sage-Voronoi & \makecell[c]{$\mathbf{0.132}\,\pm\,0.135$\\$[-0.037,\ 0.300]$} & \makecell[c]{$0.185\,\pm\,0.151$\\$[-0.003,\ 0.373]$} & \makecell[c]{$0.089\,\pm\,0.029$\\$[0.053,\ 0.125]$} & \makecell[c]{$0.277\,\pm\,0.217$\\$[0.007,\ 0.547]$} \\
GAIN & \makecell[c]{$0.202\,\pm\,0.182$\\$[-0.025,\ 0.428]$} & \makecell[c]{$0.136\,\pm\,0.142$\\$[-0.040,\ 0.313]$} & \makecell[c]{$0.129\,\pm\,0.050$\\$[0.067,\ 0.191]$} & \makecell[c]{$0.216\,\pm\,0.069$\\$[0.130,\ 0.302]$} \\
GRIN & \makecell[c]{$0.157\,\pm\,0.171$\\$[-0.056,\ 0.369]$} & \makecell[c]{$\mathbf{0.081}\,\pm\,0.054$\\$[0.014,\ 0.148]$} & \makecell[c]{$\mathbf{0.063}\,\pm\,0.036$\\$[0.018,\ 0.108]$} & \makecell[c]{$\mathbf{0.133}\,\pm\,0.051$\\$[0.070,\ 0.196]$} \\
\bottomrule
\end{tabular}
\endgroup
\end{table}

\begin{table}[htbp]
\centering
\caption{{MRE for sensor A142, January 2023 (dimensionless ratio). Each cell gives mean $\pm$ sample SD, followed by the 95\% confidence interval; $n=5$ seed runs. Bold means identify the lowest unrounded mean within each column, not statistical significance.}}
\label{tab:supp_p01_82}
\begingroup
\footnotesize
\setlength{\tabcolsep}{3pt}
\renewcommand{\arraystretch}{1.15}
\begin{tabular}{lccc}
\toprule
Method & Block $B=16$ & Block $B=32$ & Block $B=64$ \\
\midrule
Linear & \makecell[c]{$0.104\,\pm\,0.018$\\$[0.081,\ 0.127]$} & \makecell[c]{$0.176\,\pm\,0.113$\\$[0.036,\ 0.316]$} & \makecell[c]{$1.037\,\pm\,1.750$\\$[-1.137,\ 3.210]$} \\
Median & \makecell[c]{$0.133\,\pm\,0.045$\\$[0.077,\ 0.189]$} & \makecell[c]{$0.570\,\pm\,0.901$\\$[-0.549,\ 1.688]$} & \makecell[c]{$1.137\,\pm\,1.329$\\$[-0.513,\ 2.787]$} \\
FFill + BFill & \makecell[c]{$0.162\,\pm\,0.103$\\$[0.034,\ 0.291]$} & \makecell[c]{$0.659\,\pm\,1.110$\\$[-0.719,\ 2.037]$} & \makecell[c]{$0.973\,\pm\,1.208$\\$[-0.526,\ 2.473]$} \\
Rolling mean & \makecell[c]{$0.107\,\pm\,0.029$\\$[0.071,\ 0.143]$} & \makecell[c]{$0.197\,\pm\,0.154$\\$[0.005,\ 0.388]$} & \makecell[c]{$1.023\,\pm\,1.733$\\$[-1.128,\ 3.175]$} \\
KNN (scaled) & \makecell[c]{$0.376\,\pm\,0.063$\\$[0.298,\ 0.454]$} & \makecell[c]{$0.425\,\pm\,0.045$\\$[0.369,\ 0.481]$} & \makecell[c]{$0.492\,\pm\,0.292$\\$[0.129,\ 0.855]$} \\
KNN (unscaled) & \makecell[c]{$0.115\,\pm\,0.036$\\$[0.070,\ 0.160]$} & \makecell[c]{$0.145\,\pm\,0.056$\\$[0.075,\ 0.215]$} & \makecell[c]{$0.232\,\pm\,0.145$\\$[0.053,\ 0.412]$} \\
Historical & \makecell[c]{$0.102\,\pm\,0.058$\\$[0.029,\ 0.174]$} & \makecell[c]{$0.196\,\pm\,0.204$\\$[-0.057,\ 0.449]$} & \makecell[c]{$0.303\,\pm\,0.204$\\$[0.050,\ 0.556]$} \\
Graph Laplacian & \makecell[c]{$0.215\,\pm\,0.037$\\$[0.169,\ 0.262]$} & \makecell[c]{$0.296\,\pm\,0.184$\\$[0.067,\ 0.524]$} & \makecell[c]{$0.422\,\pm\,0.316$\\$[0.030,\ 0.813]$} \\
Graph Temporal & \makecell[c]{$0.204\,\pm\,0.046$\\$[0.147,\ 0.262]$} & \makecell[c]{$0.227\,\pm\,0.161$\\$[0.028,\ 0.427]$} & \makecell[c]{$0.349\,\pm\,0.255$\\$[0.033,\ 0.666]$} \\
Sage-Voronoi & \makecell[c]{$0.130\,\pm\,0.031$\\$[0.093,\ 0.168]$} & \makecell[c]{$0.274\,\pm\,0.296$\\$[-0.094,\ 0.642]$} & \makecell[c]{$0.333\,\pm\,0.360$\\$[-0.114,\ 0.779]$} \\
GAIN & \makecell[c]{$0.156\,\pm\,0.088$\\$[0.046,\ 0.265]$} & \makecell[c]{$0.151\,\pm\,0.026$\\$[0.119,\ 0.183]$} & \makecell[c]{$0.173\,\pm\,0.058$\\$[0.100,\ 0.245]$} \\
GRIN & \makecell[c]{$\mathbf{0.073}\,\pm\,0.021$\\$[0.047,\ 0.099]$} & \makecell[c]{$\mathbf{0.094}\,\pm\,0.022$\\$[0.067,\ 0.121]$} & \makecell[c]{$\mathbf{0.124}\,\pm\,0.050$\\$[0.062,\ 0.187]$} \\
\bottomrule
\end{tabular}
\endgroup
\end{table}

\clearpage

\begin{table}[htbp]
\centering
\caption{{MRE for sensor A142, January 2023 (dimensionless ratio). Each cell gives mean $\pm$ sample SD, followed by the 95\% confidence interval; $n=5$ seed runs. Bold means identify the lowest unrounded mean within each column, not statistical significance.}}
\label{tab:supp_p01_83}
\begingroup
\footnotesize
\setlength{\tabcolsep}{3pt}
\renewcommand{\arraystretch}{1.15}
\begin{tabular}{lcc}
\toprule
Method & Block $B=128$ & Block $B=144$ \\
\midrule
Linear & \makecell[c]{$1.547\,\pm\,0.705$\\$[0.672,\ 2.423]$} & \makecell[c]{$1.272\,\pm\,0.802$\\$[0.277,\ 2.268]$} \\
Median & \makecell[c]{$1.631\,\pm\,0.464$\\$[1.055,\ 2.207]$} & \makecell[c]{$1.484\,\pm\,0.385$\\$[1.006,\ 1.962]$} \\
FFill + BFill & \makecell[c]{$1.574\,\pm\,0.615$\\$[0.811,\ 2.338]$} & \makecell[c]{$1.271\,\pm\,0.766$\\$[0.320,\ 2.222]$} \\
Rolling mean & \makecell[c]{$1.538\,\pm\,0.706$\\$[0.661,\ 2.415]$} & \makecell[c]{$1.270\,\pm\,0.800$\\$[0.276,\ 2.263]$} \\
KNN (scaled) & \makecell[c]{$0.626\,\pm\,0.124$\\$[0.472,\ 0.779]$} & \makecell[c]{$0.553\,\pm\,0.072$\\$[0.464,\ 0.642]$} \\
KNN (unscaled) & \makecell[c]{$0.255\,\pm\,0.043$\\$[0.202,\ 0.308]$} & \makecell[c]{$0.228\,\pm\,0.037$\\$[0.182,\ 0.275]$} \\
Historical & \makecell[c]{$0.373\,\pm\,0.050$\\$[0.310,\ 0.435]$} & \makecell[c]{$0.356\,\pm\,0.036$\\$[0.312,\ 0.401]$} \\
Graph Laplacian & \makecell[c]{$0.539\,\pm\,0.133$\\$[0.374,\ 0.705]$} & \makecell[c]{$0.508\,\pm\,0.118$\\$[0.361,\ 0.655]$} \\
Graph Temporal & \makecell[c]{$0.465\,\pm\,0.135$\\$[0.298,\ 0.632]$} & \makecell[c]{$0.430\,\pm\,0.105$\\$[0.300,\ 0.559]$} \\
Sage-Voronoi & \makecell[c]{$0.462\,\pm\,0.176$\\$[0.244,\ 0.681]$} & \makecell[c]{$0.379\,\pm\,0.151$\\$[0.192,\ 0.566]$} \\
GAIN & \makecell[c]{$0.196\,\pm\,0.040$\\$[0.147,\ 0.245]$} & \makecell[c]{$0.183\,\pm\,0.028$\\$[0.148,\ 0.218]$} \\
GRIN & \makecell[c]{$\mathbf{0.154}\,\pm\,0.022$\\$[0.127,\ 0.181]$} & \makecell[c]{$\mathbf{0.135}\,\pm\,0.030$\\$[0.098,\ 0.173]$} \\
\bottomrule
\end{tabular}
\endgroup
\end{table}

\begin{table}[htbp]
\centering
\caption{{MRE for sensor A142, January 2023 (dimensionless ratio). Each cell gives mean $\pm$ sample SD, followed by the 95\% confidence interval; $n=5$ seed runs. Bold means identify the lowest unrounded mean within each column, not statistical significance.}}
\label{tab:supp_p01_84}
\begingroup
\footnotesize
\setlength{\tabcolsep}{3pt}
\renewcommand{\arraystretch}{1.15}
\begin{tabular}{lcccc}
\toprule
Method & MCAR 10\% & MCAR 30\% & MCAR 50\% & Synchronous $B=144$ \\
\midrule
Linear & \makecell[c]{$0.184\,\pm\,0.010$\\$[0.171,\ 0.197]$} & \makecell[c]{$0.147\,\pm\,0.008$\\$[0.137,\ 0.157]$} & \makecell[c]{$0.148\,\pm\,0.009$\\$[0.136,\ 0.160]$} & \makecell[c]{$1.693\,\pm\,0.432$\\$[1.156,\ 2.230]$} \\
Median & \makecell[c]{$1.676\,\pm\,0.311$\\$[1.290,\ 2.062]$} & \makecell[c]{$1.389\,\pm\,0.099$\\$[1.266,\ 1.512]$} & \makecell[c]{$1.334\,\pm\,0.075$\\$[1.240,\ 1.428]$} & \makecell[c]{$1.398\,\pm\,0.380$\\$[0.927,\ 1.869]$} \\
FFill + BFill & \makecell[c]{$0.209\,\pm\,0.020$\\$[0.184,\ 0.234]$} & \makecell[c]{$0.181\,\pm\,0.017$\\$[0.160,\ 0.202]$} & \makecell[c]{$0.190\,\pm\,0.009$\\$[0.179,\ 0.200]$} & \makecell[c]{$1.629\,\pm\,0.805$\\$[0.629,\ 2.629]$} \\
Rolling mean & \makecell[c]{$0.213\,\pm\,0.013$\\$[0.197,\ 0.229]$} & \makecell[c]{$0.186\,\pm\,0.018$\\$[0.163,\ 0.208]$} & \makecell[c]{$0.185\,\pm\,0.007$\\$[0.176,\ 0.193]$} & \makecell[c]{$1.694\,\pm\,0.433$\\$[1.157,\ 2.232]$} \\
KNN (scaled) & \makecell[c]{$0.703\,\pm\,0.089$\\$[0.593,\ 0.814]$} & \makecell[c]{$0.610\,\pm\,0.048$\\$[0.550,\ 0.669]$} & \makecell[c]{$0.566\,\pm\,0.043$\\$[0.513,\ 0.620]$} & \makecell[c]{$1.278\,\pm\,0.319$\\$[0.882,\ 1.674]$} \\
KNN (unscaled) & \makecell[c]{$0.304\,\pm\,0.049$\\$[0.244,\ 0.365]$} & \makecell[c]{$0.259\,\pm\,0.024$\\$[0.229,\ 0.290]$} & \makecell[c]{$0.272\,\pm\,0.021$\\$[0.246,\ 0.298]$} & \makecell[c]{$1.278\,\pm\,0.319$\\$[0.882,\ 1.674]$} \\
Historical & \makecell[c]{$0.403\,\pm\,0.055$\\$[0.335,\ 0.472]$} & \makecell[c]{$0.363\,\pm\,0.023$\\$[0.334,\ 0.391]$} & \makecell[c]{$0.350\,\pm\,0.020$\\$[0.325,\ 0.375]$} & \makecell[c]{$\mathbf{0.365}\,\pm\,0.244$\\$[0.062,\ 0.668]$} \\
Graph Laplacian & \makecell[c]{$0.571\,\pm\,0.098$\\$[0.450,\ 0.692]$} & \makecell[c]{$0.492\,\pm\,0.033$\\$[0.452,\ 0.533]$} & \makecell[c]{$0.494\,\pm\,0.026$\\$[0.462,\ 0.526]$} & \makecell[c]{$1.278\,\pm\,0.319$\\$[0.882,\ 1.674]$} \\
Graph Temporal & \makecell[c]{$0.400\,\pm\,0.066$\\$[0.318,\ 0.482]$} & \makecell[c]{$0.341\,\pm\,0.022$\\$[0.313,\ 0.368]$} & \makecell[c]{$0.340\,\pm\,0.022$\\$[0.312,\ 0.367]$} & \makecell[c]{$2.199\,\pm\,0.645$\\$[1.398,\ 3.000]$} \\
Sage-Voronoi & \makecell[c]{$0.190\,\pm\,0.025$\\$[0.160,\ 0.221]$} & \makecell[c]{$0.176\,\pm\,0.017$\\$[0.155,\ 0.196]$} & \makecell[c]{$0.182\,\pm\,0.005$\\$[0.176,\ 0.187]$} & \makecell[c]{$1.380\,\pm\,0.566$\\$[0.678,\ 2.083]$} \\
GAIN & \makecell[c]{$0.226\,\pm\,0.031$\\$[0.188,\ 0.265]$} & \makecell[c]{$0.293\,\pm\,0.039$\\$[0.244,\ 0.342]$} & \makecell[c]{$0.497\,\pm\,0.066$\\$[0.415,\ 0.579]$} & \makecell[c]{$2.041\,\pm\,0.619$\\$[1.273,\ 2.809]$} \\
GRIN & \makecell[c]{$\mathbf{0.114}\,\pm\,0.012$\\$[0.099,\ 0.129]$} & \makecell[c]{$\mathbf{0.106}\,\pm\,0.009$\\$[0.095,\ 0.117]$} & \makecell[c]{$\mathbf{0.126}\,\pm\,0.012$\\$[0.111,\ 0.140]$} & \makecell[c]{$1.223\,\pm\,0.415$\\$[0.708,\ 1.739]$} \\
\bottomrule
\end{tabular}
\endgroup
\end{table}

\clearpage

\begin{table}[htbp]
\centering
\caption{{MAE for sensor A142, October--December 2022 (vehicles per ten-minute interval). Each cell gives mean $\pm$ sample SD, followed by the 95\% confidence interval; $n=5$ seed runs. Bold means identify the lowest unrounded mean within each column, not statistical significance.}}
\label{tab:supp_p01_85}
\begingroup
\footnotesize
\setlength{\tabcolsep}{3pt}
\renewcommand{\arraystretch}{1.15}
\begin{tabular}{lcccc}
\toprule
Method & Block $B=1$ & Block $B=2$ & Block $B=4$ & Block $B=8$ \\
\midrule
Linear & \makecell[c]{$19.80\,\pm\,8.30$\\$[9.49,\ 30.11]$} & \makecell[c]{$21.87\,\pm\,10.95$\\$[8.27,\ 35.47]$} & \makecell[c]{$35.81\,\pm\,25.39$\\$[4.28,\ 67.34]$} & \makecell[c]{$\mathbf{16.81}\,\pm\,6.90$\\$[8.23,\ 25.38]$} \\
Median & \makecell[c]{$136.20\,\pm\,96.08$\\$[16.90,\ 255.50]$} & \makecell[c]{$148.20\,\pm\,61.76$\\$[71.52,\ 224.88]$} & \makecell[c]{$43.15\,\pm\,18.64$\\$[20.01,\ 66.29]$} & \makecell[c]{$93.20\,\pm\,92.21$\\$[-21.29,\ 207.69]$} \\
FFill + BFill & \makecell[c]{$20.40\,\pm\,14.47$\\$[2.44,\ 38.36]$} & \makecell[c]{$22.60\,\pm\,14.40$\\$[4.72,\ 40.48]$} & \makecell[c]{$42.85\,\pm\,16.43$\\$[22.44,\ 63.26]$} & \makecell[c]{$28.40\,\pm\,18.72$\\$[5.16,\ 51.64]$} \\
Rolling mean & \makecell[c]{$\mathbf{11.40}\,\pm\,6.76$\\$[3.01,\ 19.79]$} & \makecell[c]{$35.00\,\pm\,21.90$\\$[7.81,\ 62.19]$} & \makecell[c]{$57.09\,\pm\,24.92$\\$[26.14,\ 88.04]$} & \makecell[c]{$23.72\,\pm\,9.50$\\$[11.93,\ 35.52]$} \\
KNN (scaled) & \makecell[c]{$40.60\,\pm\,47.12$\\$[-17.91,\ 99.11]$} & \makecell[c]{$42.52\,\pm\,29.40$\\$[6.02,\ 79.02]$} & \makecell[c]{$\mathbf{24.73}\,\pm\,8.91$\\$[13.66,\ 35.79]$} & \makecell[c]{$30.91\,\pm\,21.45$\\$[4.27,\ 57.54]$} \\
KNN (unscaled) & \makecell[c]{$16.08\,\pm\,9.16$\\$[4.70,\ 27.46]$} & \makecell[c]{$43.64\,\pm\,28.17$\\$[8.67,\ 78.61]$} & \makecell[c]{$32.26\,\pm\,10.75$\\$[18.91,\ 45.61]$} & \makecell[c]{$35.92\,\pm\,19.84$\\$[11.29,\ 60.56]$} \\
Historical & \makecell[c]{$48.11\,\pm\,66.62$\\$[-34.61,\ 130.83]$} & \makecell[c]{$93.44\,\pm\,44.10$\\$[38.68,\ 148.20]$} & \makecell[c]{$26.18\,\pm\,9.35$\\$[14.57,\ 37.78]$} & \makecell[c]{$43.25\,\pm\,45.38$\\$[-13.10,\ 99.60]$} \\
Graph Laplacian & \makecell[c]{$42.74\,\pm\,25.57$\\$[10.98,\ 74.49]$} & \makecell[c]{$34.96\,\pm\,18.11$\\$[12.48,\ 57.44]$} & \makecell[c]{$61.51\,\pm\,25.10$\\$[30.35,\ 92.68]$} & \makecell[c]{$58.64\,\pm\,39.56$\\$[9.51,\ 107.76]$} \\
Graph Temporal & \makecell[c]{$20.10\,\pm\,18.73$\\$[-3.15,\ 43.36]$} & \makecell[c]{$28.19\,\pm\,18.79$\\$[4.86,\ 51.52]$} & \makecell[c]{$57.65\,\pm\,14.86$\\$[39.19,\ 76.10]$} & \makecell[c]{$55.31\,\pm\,45.31$\\$[-0.95,\ 111.58]$} \\
Sage-Voronoi & \makecell[c]{$14.21\,\pm\,9.79$\\$[2.06,\ 26.36]$} & \makecell[c]{$29.84\,\pm\,18.04$\\$[7.44,\ 52.24]$} & \makecell[c]{$53.97\,\pm\,21.86$\\$[26.83,\ 81.11]$} & \makecell[c]{$28.97\,\pm\,18.91$\\$[5.49,\ 52.45]$} \\
GAIN & \makecell[c]{$32.07\,\pm\,27.04$\\$[-1.51,\ 65.64]$} & \makecell[c]{$36.96\,\pm\,43.15$\\$[-16.63,\ 90.54]$} & \makecell[c]{$71.19\,\pm\,16.08$\\$[51.23,\ 91.15]$} & \makecell[c]{$59.59\,\pm\,53.91$\\$[-7.35,\ 126.53]$} \\
GRIN & \makecell[c]{$15.40\,\pm\,14.61$\\$[-2.75,\ 33.55]$} & \makecell[c]{$\mathbf{18.15}\,\pm\,9.67$\\$[6.14,\ 30.15]$} & \makecell[c]{$30.48\,\pm\,4.72$\\$[24.61,\ 36.34]$} & \makecell[c]{$26.98\,\pm\,21.02$\\$[0.88,\ 53.08]$} \\
\bottomrule
\end{tabular}
\endgroup
\end{table}

\begin{table}[htbp]
\centering
\caption{{MAE for sensor A142, October--December 2022 (vehicles per ten-minute interval). Each cell gives mean $\pm$ sample SD, followed by the 95\% confidence interval; $n=5$ seed runs. Bold means identify the lowest unrounded mean within each column, not statistical significance.}}
\label{tab:supp_p01_86}
\begingroup
\footnotesize
\setlength{\tabcolsep}{3pt}
\renewcommand{\arraystretch}{1.15}
\begin{tabular}{lccc}
\toprule
Method & Block $B=16$ & Block $B=32$ & Block $B=64$ \\
\midrule
Linear & \makecell[c]{$26.55\,\pm\,13.94$\\$[9.24,\ 43.85]$} & \makecell[c]{$55.85\,\pm\,29.34$\\$[19.41,\ 92.28]$} & \makecell[c]{$90.28\,\pm\,24.82$\\$[59.46,\ 121.10]$} \\
Median & \makecell[c]{$115.83\,\pm\,70.88$\\$[27.82,\ 203.83]$} & \makecell[c]{$99.17\,\pm\,76.49$\\$[4.20,\ 194.15]$} & \makecell[c]{$83.08\,\pm\,24.18$\\$[53.05,\ 113.11]$} \\
FFill + BFill & \makecell[c]{$49.55\,\pm\,21.07$\\$[23.39,\ 75.71]$} & \makecell[c]{$92.47\,\pm\,88.92$\\$[-17.93,\ 202.88]$} & \makecell[c]{$161.71\,\pm\,83.10$\\$[58.53,\ 264.89]$} \\
Rolling mean & \makecell[c]{$28.16\,\pm\,15.86$\\$[8.47,\ 47.85]$} & \makecell[c]{$58.67\,\pm\,37.25$\\$[12.42,\ 104.92]$} & \makecell[c]{$91.41\,\pm\,23.67$\\$[62.02,\ 120.81]$} \\
KNN (scaled) & \makecell[c]{$28.98\,\pm\,20.20$\\$[3.90,\ 54.07]$} & \makecell[c]{$30.48\,\pm\,13.61$\\$[13.58,\ 47.38]$} & \makecell[c]{$\mathbf{35.28}\,\pm\,8.44$\\$[24.80,\ 45.75]$} \\
KNN (unscaled) & \makecell[c]{$27.05\,\pm\,12.20$\\$[11.91,\ 42.20]$} & \makecell[c]{$34.34\,\pm\,11.26$\\$[20.36,\ 48.32]$} & \makecell[c]{$46.77\,\pm\,6.58$\\$[38.60,\ 54.94]$} \\
Historical & \makecell[c]{$68.36\,\pm\,56.42$\\$[-1.70,\ 138.41]$} & \makecell[c]{$35.87\,\pm\,13.68$\\$[18.88,\ 52.86]$} & \makecell[c]{$53.56\,\pm\,15.61$\\$[34.18,\ 72.94]$} \\
Graph Laplacian & \makecell[c]{$33.73\,\pm\,14.90$\\$[15.22,\ 52.23]$} & \makecell[c]{$51.12\,\pm\,7.94$\\$[41.26,\ 60.98]$} & \makecell[c]{$51.64\,\pm\,13.83$\\$[34.47,\ 68.81]$} \\
Graph Temporal & \makecell[c]{$30.62\,\pm\,14.43$\\$[12.69,\ 48.54]$} & \makecell[c]{$42.48\,\pm\,4.61$\\$[36.76,\ 48.21]$} & \makecell[c]{$46.20\,\pm\,13.24$\\$[29.76,\ 62.65]$} \\
Sage-Voronoi & \makecell[c]{$30.10\,\pm\,12.47$\\$[14.61,\ 45.58]$} & \makecell[c]{$42.31\,\pm\,15.14$\\$[23.51,\ 61.10]$} & \makecell[c]{$59.27\,\pm\,10.49$\\$[46.25,\ 72.30]$} \\
GAIN & \makecell[c]{$36.68\,\pm\,7.18$\\$[27.77,\ 45.60]$} & \makecell[c]{$56.56\,\pm\,33.41$\\$[15.08,\ 98.04]$} & \makecell[c]{$60.21\,\pm\,17.95$\\$[37.92,\ 82.50]$} \\
GRIN & \makecell[c]{$\mathbf{16.65}\,\pm\,8.47$\\$[6.13,\ 27.17]$} & \makecell[c]{$\mathbf{26.91}\,\pm\,11.56$\\$[12.55,\ 41.27]$} & \makecell[c]{$39.56\,\pm\,9.85$\\$[27.33,\ 51.80]$} \\
\bottomrule
\end{tabular}
\endgroup
\end{table}

\clearpage

\begin{table}[htbp]
\centering
\caption{{MAE for sensor A142, October--December 2022 (vehicles per ten-minute interval). Each cell gives mean $\pm$ sample SD, followed by the 95\% confidence interval; $n=5$ seed runs. Bold means identify the lowest unrounded mean within each column, not statistical significance.}}
\label{tab:supp_p01_87}
\begingroup
\footnotesize
\setlength{\tabcolsep}{3pt}
\renewcommand{\arraystretch}{1.15}
\begin{tabular}{lcc}
\toprule
Method & Block $B=128$ & Block $B=144$ \\
\midrule
Linear & \makecell[c]{$139.72\,\pm\,59.32$\\$[66.06,\ 213.37]$} & \makecell[c]{$136.05\,\pm\,47.73$\\$[76.79,\ 195.31]$} \\
Median & \makecell[c]{$93.65\,\pm\,20.25$\\$[68.51,\ 118.80]$} & \makecell[c]{$107.08\,\pm\,13.81$\\$[89.94,\ 124.22]$} \\
FFill + BFill & \makecell[c]{$185.98\,\pm\,34.20$\\$[143.52,\ 228.44]$} & \makecell[c]{$137.20\,\pm\,44.89$\\$[81.47,\ 192.94]$} \\
Rolling mean & \makecell[c]{$140.36\,\pm\,58.98$\\$[67.12,\ 213.59]$} & \makecell[c]{$136.10\,\pm\,47.52$\\$[77.09,\ 195.10]$} \\
KNN (scaled) & \makecell[c]{$35.26\,\pm\,9.30$\\$[23.71,\ 46.80]$} & \makecell[c]{$29.54\,\pm\,4.04$\\$[24.53,\ 34.55]$} \\
KNN (unscaled) & \makecell[c]{$35.83\,\pm\,4.88$\\$[29.77,\ 41.90]$} & \makecell[c]{$35.29\,\pm\,2.87$\\$[31.73,\ 38.85]$} \\
Historical & \makecell[c]{$43.86\,\pm\,17.02$\\$[22.73,\ 65.00]$} & \makecell[c]{$59.30\,\pm\,25.78$\\$[27.29,\ 91.31]$} \\
Graph Laplacian & \makecell[c]{$48.47\,\pm\,7.02$\\$[39.76,\ 57.19]$} & \makecell[c]{$48.08\,\pm\,6.96$\\$[39.43,\ 56.72]$} \\
Graph Temporal & \makecell[c]{$45.81\,\pm\,8.73$\\$[34.97,\ 56.66]$} & \makecell[c]{$42.72\,\pm\,6.96$\\$[34.09,\ 51.36]$} \\
Sage-Voronoi & \makecell[c]{$48.64\,\pm\,10.96$\\$[35.02,\ 62.25]$} & \makecell[c]{$44.95\,\pm\,9.12$\\$[33.63,\ 56.27]$} \\
GAIN & \makecell[c]{$42.62\,\pm\,10.15$\\$[30.01,\ 55.22]$} & \makecell[c]{$48.47\,\pm\,5.62$\\$[41.50,\ 55.45]$} \\
GRIN & \makecell[c]{$\mathbf{31.10}\,\pm\,7.82$\\$[21.39,\ 40.82]$} & \makecell[c]{$\mathbf{28.99}\,\pm\,5.56$\\$[22.09,\ 35.90]$} \\
\bottomrule
\end{tabular}
\endgroup
\end{table}

\begin{table}[htbp]
\centering
\caption{{MAE for sensor A142, October--December 2022 (vehicles per ten-minute interval). Each cell gives mean $\pm$ sample SD, followed by the 95\% confidence interval; $n=5$ seed runs. Bold means identify the lowest unrounded mean within each column, not statistical significance.}}
\label{tab:supp_p01_88}
\begingroup
\footnotesize
\setlength{\tabcolsep}{3pt}
\renewcommand{\arraystretch}{1.15}
\begin{tabular}{lcccc}
\toprule
Method & MCAR 10\% & MCAR 30\% & MCAR 50\% & Synchronous $B=144$ \\
\midrule
Linear & \makecell[c]{$26.34\,\pm\,1.09$\\$[24.99,\ 27.70]$} & \makecell[c]{$25.79\,\pm\,0.33$\\$[25.38,\ 26.20]$} & \makecell[c]{$26.68\,\pm\,0.78$\\$[25.71,\ 27.65]$} & \makecell[c]{$113.28\,\pm\,15.89$\\$[93.55,\ 133.00]$} \\
Median & \makecell[c]{$99.19\,\pm\,2.70$\\$[95.84,\ 102.55]$} & \makecell[c]{$97.47\,\pm\,0.99$\\$[96.23,\ 98.70]$} & \makecell[c]{$98.52\,\pm\,1.21$\\$[97.01,\ 100.03]$} & \makecell[c]{$100.44\,\pm\,16.20$\\$[80.33,\ 120.55]$} \\
FFill + BFill & \makecell[c]{$31.34\,\pm\,1.05$\\$[30.04,\ 32.65]$} & \makecell[c]{$31.21\,\pm\,0.23$\\$[30.93,\ 31.49]$} & \makecell[c]{$33.89\,\pm\,1.20$\\$[32.41,\ 35.38]$} & \makecell[c]{$121.23\,\pm\,20.69$\\$[95.55,\ 146.92]$} \\
Rolling mean & \makecell[c]{$30.40\,\pm\,0.63$\\$[29.61,\ 31.18]$} & \makecell[c]{$30.86\,\pm\,0.64$\\$[30.06,\ 31.65]$} & \makecell[c]{$32.22\,\pm\,0.94$\\$[31.05,\ 33.38]$} & \makecell[c]{$113.59\,\pm\,15.66$\\$[94.14,\ 133.03]$} \\
KNN (scaled) & \makecell[c]{$31.86\,\pm\,1.38$\\$[30.15,\ 33.56]$} & \makecell[c]{$30.09\,\pm\,0.73$\\$[29.19,\ 30.99]$} & \makecell[c]{$31.76\,\pm\,0.48$\\$[31.17,\ 32.36]$} & \makecell[c]{$94.70\,\pm\,13.95$\\$[77.37,\ 112.02]$} \\
KNN (unscaled) & \makecell[c]{$36.50\,\pm\,2.51$\\$[33.38,\ 39.62]$} & \makecell[c]{$34.79\,\pm\,0.76$\\$[33.85,\ 35.74]$} & \makecell[c]{$35.46\,\pm\,0.23$\\$[35.17,\ 35.75]$} & \makecell[c]{$94.70\,\pm\,13.95$\\$[77.37,\ 112.02]$} \\
Historical & \makecell[c]{$50.76\,\pm\,1.37$\\$[49.06,\ 52.45]$} & \makecell[c]{$49.08\,\pm\,1.39$\\$[47.35,\ 50.81]$} & \makecell[c]{$49.78\,\pm\,0.97$\\$[48.58,\ 50.99]$} & \makecell[c]{$\mathbf{55.84}\,\pm\,18.41$\\$[32.98,\ 78.70]$} \\
Graph Laplacian & \makecell[c]{$48.17\,\pm\,3.10$\\$[44.31,\ 52.02]$} & \makecell[c]{$47.49\,\pm\,1.75$\\$[45.32,\ 49.67]$} & \makecell[c]{$49.48\,\pm\,0.51$\\$[48.84,\ 50.11]$} & \makecell[c]{$94.70\,\pm\,13.95$\\$[77.37,\ 112.02]$} \\
Graph Temporal & \makecell[c]{$36.95\,\pm\,1.94$\\$[34.54,\ 39.37]$} & \makecell[c]{$37.41\,\pm\,0.92$\\$[36.26,\ 38.56]$} & \makecell[c]{$40.60\,\pm\,0.27$\\$[40.26,\ 40.94]$} & \makecell[c]{$110.49\,\pm\,20.37$\\$[85.20,\ 135.79]$} \\
Sage-Voronoi & \makecell[c]{$27.16\,\pm\,0.93$\\$[26.00,\ 28.32]$} & \makecell[c]{$27.98\,\pm\,0.80$\\$[26.99,\ 28.97]$} & \makecell[c]{$31.15\,\pm\,0.70$\\$[30.28,\ 32.02]$} & \makecell[c]{$132.20\,\pm\,36.42$\\$[86.98,\ 177.41]$} \\
GAIN & \makecell[c]{$40.57\,\pm\,4.84$\\$[34.56,\ 46.57]$} & \makecell[c]{$51.38\,\pm\,6.67$\\$[43.09,\ 59.66]$} & \makecell[c]{$120.88\,\pm\,13.36$\\$[104.28,\ 137.47]$} & \makecell[c]{$527.42\,\pm\,38.53$\\$[479.58,\ 575.26]$} \\
GRIN & \makecell[c]{$\mathbf{20.71}\,\pm\,0.92$\\$[19.57,\ 21.85]$} & \makecell[c]{$\mathbf{20.42}\,\pm\,0.50$\\$[19.80,\ 21.03]$} & \makecell[c]{$\mathbf{22.85}\,\pm\,0.38$\\$[22.37,\ 23.32]$} & \makecell[c]{$97.04\,\pm\,22.03$\\$[69.69,\ 124.39]$} \\
\bottomrule
\end{tabular}
\endgroup
\end{table}

\clearpage

\begin{table}[htbp]
\centering
\caption{{RMSE for sensor A142, October--December 2022 (vehicles per ten-minute interval). Each cell gives mean $\pm$ sample SD, followed by the 95\% confidence interval; $n=5$ seed runs. Bold means identify the lowest unrounded mean within each column, not statistical significance.}}
\label{tab:supp_p01_89}
\begingroup
\footnotesize
\setlength{\tabcolsep}{3pt}
\renewcommand{\arraystretch}{1.15}
\begin{tabular}{lcccc}
\toprule
Method & Block $B=1$ & Block $B=2$ & Block $B=4$ & Block $B=8$ \\
\midrule
Linear & \makecell[c]{$19.80\,\pm\,8.30$\\$[9.49,\ 30.11]$} & \makecell[c]{$22.50\,\pm\,11.16$\\$[8.63,\ 36.36]$} & \makecell[c]{$40.33\,\pm\,24.62$\\$[9.76,\ 70.89]$} & \makecell[c]{$\mathbf{24.00}\,\pm\,14.52$\\$[5.97,\ 42.03]$} \\
Median & \makecell[c]{$136.20\,\pm\,96.08$\\$[16.90,\ 255.50]$} & \makecell[c]{$149.06\,\pm\,61.32$\\$[72.93,\ 225.20]$} & \makecell[c]{$51.68\,\pm\,19.13$\\$[27.93,\ 75.43]$} & \makecell[c]{$97.68\,\pm\,90.35$\\$[-14.51,\ 209.87]$} \\
FFill + BFill & \makecell[c]{$20.40\,\pm\,14.47$\\$[2.44,\ 38.36]$} & \makecell[c]{$23.96\,\pm\,14.66$\\$[5.76,\ 42.16]$} & \makecell[c]{$51.41\,\pm\,14.49$\\$[33.43,\ 69.40]$} & \makecell[c]{$35.81\,\pm\,22.76$\\$[7.55,\ 64.08]$} \\
Rolling mean & \makecell[c]{$\mathbf{11.40}\,\pm\,6.76$\\$[3.01,\ 19.79]$} & \makecell[c]{$37.07\,\pm\,22.10$\\$[9.63,\ 64.50]$} & \makecell[c]{$66.70\,\pm\,18.27$\\$[44.00,\ 89.39]$} & \makecell[c]{$31.76\,\pm\,15.38$\\$[12.66,\ 50.86]$} \\
KNN (scaled) & \makecell[c]{$40.60\,\pm\,47.12$\\$[-17.91,\ 99.11]$} & \makecell[c]{$45.91\,\pm\,29.85$\\$[8.85,\ 82.97]$} & \makecell[c]{$\mathbf{32.27}\,\pm\,13.01$\\$[16.11,\ 48.42]$} & \makecell[c]{$38.47\,\pm\,23.66$\\$[9.09,\ 67.85]$} \\
KNN (unscaled) & \makecell[c]{$16.08\,\pm\,9.16$\\$[4.70,\ 27.46]$} & \makecell[c]{$44.64\,\pm\,29.01$\\$[8.63,\ 80.66]$} & \makecell[c]{$37.10\,\pm\,12.51$\\$[21.56,\ 52.63]$} & \makecell[c]{$43.73\,\pm\,23.84$\\$[14.13,\ 73.34]$} \\
Historical & \makecell[c]{$48.11\,\pm\,66.62$\\$[-34.61,\ 130.83]$} & \makecell[c]{$94.19\,\pm\,44.63$\\$[38.77,\ 149.61]$} & \makecell[c]{$33.10\,\pm\,13.16$\\$[16.77,\ 49.44]$} & \makecell[c]{$49.64\,\pm\,43.16$\\$[-3.95,\ 103.24]$} \\
Graph Laplacian & \makecell[c]{$42.74\,\pm\,25.57$\\$[10.98,\ 74.49]$} & \makecell[c]{$36.91\,\pm\,16.80$\\$[16.04,\ 57.78]$} & \makecell[c]{$69.58\,\pm\,19.28$\\$[45.65,\ 93.52]$} & \makecell[c]{$63.37\,\pm\,41.80$\\$[11.47,\ 115.26]$} \\
Graph Temporal & \makecell[c]{$20.10\,\pm\,18.73$\\$[-3.15,\ 43.36]$} & \makecell[c]{$30.03\,\pm\,17.67$\\$[8.09,\ 51.98]$} & \makecell[c]{$65.22\,\pm\,9.96$\\$[52.84,\ 77.59]$} & \makecell[c]{$59.99\,\pm\,47.56$\\$[0.94,\ 119.03]$} \\
Sage-Voronoi & \makecell[c]{$14.21\,\pm\,9.79$\\$[2.06,\ 26.36]$} & \makecell[c]{$32.00\,\pm\,18.45$\\$[9.09,\ 54.91]$} & \makecell[c]{$61.30\,\pm\,17.71$\\$[39.32,\ 83.29]$} & \makecell[c]{$35.52\,\pm\,25.65$\\$[3.68,\ 67.36]$} \\
GAIN & \makecell[c]{$32.07\,\pm\,27.04$\\$[-1.51,\ 65.64]$} & \makecell[c]{$39.00\,\pm\,42.19$\\$[-13.39,\ 91.39]$} & \makecell[c]{$76.39\,\pm\,15.66$\\$[56.95,\ 95.84]$} & \makecell[c]{$63.83\,\pm\,55.55$\\$[-5.14,\ 132.80]$} \\
GRIN & \makecell[c]{$15.40\,\pm\,14.61$\\$[-2.75,\ 33.55]$} & \makecell[c]{$\mathbf{19.91}\,\pm\,9.77$\\$[7.77,\ 32.04]$} & \makecell[c]{$36.93\,\pm\,7.07$\\$[28.15,\ 45.70]$} & \makecell[c]{$32.96\,\pm\,27.64$\\$[-1.36,\ 67.28]$} \\
\bottomrule
\end{tabular}
\endgroup
\end{table}

\begin{table}[htbp]
\centering
\caption{{RMSE for sensor A142, October--December 2022 (vehicles per ten-minute interval). Each cell gives mean $\pm$ sample SD, followed by the 95\% confidence interval; $n=5$ seed runs. Bold means identify the lowest unrounded mean within each column, not statistical significance.}}
\label{tab:supp_p01_90}
\begingroup
\footnotesize
\setlength{\tabcolsep}{3pt}
\renewcommand{\arraystretch}{1.15}
\begin{tabular}{lccc}
\toprule
Method & Block $B=16$ & Block $B=32$ & Block $B=64$ \\
\midrule
Linear & \makecell[c]{$36.08\,\pm\,19.47$\\$[11.90,\ 60.25]$} & \makecell[c]{$71.06\,\pm\,36.03$\\$[26.32,\ 115.79]$} & \makecell[c]{$117.85\,\pm\,35.00$\\$[74.39,\ 161.32]$} \\
Median & \makecell[c]{$122.57\,\pm\,68.81$\\$[37.13,\ 208.02]$} & \makecell[c]{$113.09\,\pm\,70.86$\\$[25.10,\ 201.07]$} & \makecell[c]{$108.92\,\pm\,31.48$\\$[69.83,\ 148.00]$} \\
FFill + BFill & \makecell[c]{$60.37\,\pm\,24.13$\\$[30.41,\ 90.33]$} & \makecell[c]{$107.35\,\pm\,98.49$\\$[-14.94,\ 229.64]$} & \makecell[c]{$187.31\,\pm\,90.96$\\$[74.36,\ 300.25]$} \\
Rolling mean & \makecell[c]{$36.08\,\pm\,20.82$\\$[10.22,\ 61.93]$} & \makecell[c]{$73.02\,\pm\,43.76$\\$[18.69,\ 127.36]$} & \makecell[c]{$120.18\,\pm\,35.50$\\$[76.10,\ 164.27]$} \\
KNN (scaled) & \makecell[c]{$34.94\,\pm\,21.75$\\$[7.93,\ 61.95]$} & \makecell[c]{$38.97\,\pm\,18.35$\\$[16.18,\ 61.76]$} & \makecell[c]{$\mathbf{54.06}\,\pm\,12.82$\\$[38.13,\ 69.98]$} \\
KNN (unscaled) & \makecell[c]{$33.23\,\pm\,12.78$\\$[17.36,\ 49.10]$} & \makecell[c]{$43.72\,\pm\,15.74$\\$[24.18,\ 63.26]$} & \makecell[c]{$68.12\,\pm\,14.29$\\$[50.37,\ 85.86]$} \\
Historical & \makecell[c]{$73.59\,\pm\,56.19$\\$[3.81,\ 143.36]$} & \makecell[c]{$46.82\,\pm\,19.66$\\$[22.41,\ 71.23]$} & \makecell[c]{$73.84\,\pm\,18.15$\\$[51.30,\ 96.37]$} \\
Graph Laplacian & \makecell[c]{$39.44\,\pm\,18.19$\\$[16.85,\ 62.03]$} & \makecell[c]{$59.99\,\pm\,7.15$\\$[51.12,\ 68.87]$} & \makecell[c]{$70.13\,\pm\,18.18$\\$[47.55,\ 92.70]$} \\
Graph Temporal & \makecell[c]{$37.75\,\pm\,18.55$\\$[14.71,\ 60.78]$} & \makecell[c]{$52.38\,\pm\,4.90$\\$[46.29,\ 58.46]$} & \makecell[c]{$66.35\,\pm\,17.78$\\$[44.27,\ 88.43]$} \\
Sage-Voronoi & \makecell[c]{$35.66\,\pm\,13.47$\\$[18.94,\ 52.38]$} & \makecell[c]{$51.72\,\pm\,18.68$\\$[28.53,\ 74.91]$} & \makecell[c]{$80.65\,\pm\,14.27$\\$[62.93,\ 98.37]$} \\
GAIN & \makecell[c]{$45.89\,\pm\,9.16$\\$[34.51,\ 57.26]$} & \makecell[c]{$66.97\,\pm\,35.07$\\$[23.42,\ 110.51]$} & \makecell[c]{$81.74\,\pm\,20.46$\\$[56.33,\ 107.15]$} \\
GRIN & \makecell[c]{$\mathbf{23.84}\,\pm\,14.07$\\$[6.37,\ 41.32]$} & \makecell[c]{$\mathbf{36.42}\,\pm\,17.41$\\$[14.81,\ 58.04]$} & \makecell[c]{$57.51\,\pm\,14.72$\\$[39.24,\ 75.79]$} \\
\bottomrule
\end{tabular}
\endgroup
\end{table}

\clearpage

\begin{table}[htbp]
\centering
\caption{{RMSE for sensor A142, October--December 2022 (vehicles per ten-minute interval). Each cell gives mean $\pm$ sample SD, followed by the 95\% confidence interval; $n=5$ seed runs. Bold means identify the lowest unrounded mean within each column, not statistical significance.}}
\label{tab:supp_p01_91}
\begingroup
\footnotesize
\setlength{\tabcolsep}{3pt}
\renewcommand{\arraystretch}{1.15}
\begin{tabular}{lcc}
\toprule
Method & Block $B=128$ & Block $B=144$ \\
\midrule
Linear & \makecell[c]{$162.21\,\pm\,60.84$\\$[86.67,\ 237.76]$} & \makecell[c]{$167.37\,\pm\,46.69$\\$[109.39,\ 225.34]$} \\
Median & \makecell[c]{$121.86\,\pm\,22.83$\\$[93.52,\ 150.21]$} & \makecell[c]{$133.75\,\pm\,11.56$\\$[119.40,\ 148.11]$} \\
FFill + BFill & \makecell[c]{$212.76\,\pm\,27.86$\\$[178.17,\ 247.36]$} & \makecell[c]{$168.77\,\pm\,43.29$\\$[115.01,\ 222.52]$} \\
Rolling mean & \makecell[c]{$162.93\,\pm\,60.58$\\$[87.72,\ 238.15]$} & \makecell[c]{$167.40\,\pm\,46.57$\\$[109.58,\ 225.22]$} \\
KNN (scaled) & \makecell[c]{$48.88\,\pm\,9.27$\\$[37.37,\ 60.39]$} & \makecell[c]{$42.83\,\pm\,6.85$\\$[34.32,\ 51.34]$} \\
KNN (unscaled) & \makecell[c]{$49.19\,\pm\,3.36$\\$[45.02,\ 53.36]$} & \makecell[c]{$50.33\,\pm\,6.75$\\$[41.94,\ 58.71]$} \\
Historical & \makecell[c]{$60.80\,\pm\,21.35$\\$[34.29,\ 87.32]$} & \makecell[c]{$79.53\,\pm\,32.43$\\$[39.26,\ 119.80]$} \\
Graph Laplacian & \makecell[c]{$61.10\,\pm\,7.72$\\$[51.52,\ 70.69]$} & \makecell[c]{$59.96\,\pm\,6.64$\\$[51.72,\ 68.20]$} \\
Graph Temporal & \makecell[c]{$60.07\,\pm\,10.20$\\$[47.41,\ 72.73]$} & \makecell[c]{$55.27\,\pm\,8.61$\\$[44.58,\ 65.96]$} \\
Sage-Voronoi & \makecell[c]{$63.16\,\pm\,11.75$\\$[48.58,\ 77.75]$} & \makecell[c]{$62.20\,\pm\,10.69$\\$[48.93,\ 75.48]$} \\
GAIN & \makecell[c]{$58.88\,\pm\,9.92$\\$[46.57,\ 71.20]$} & \makecell[c]{$66.12\,\pm\,7.06$\\$[57.36,\ 74.88]$} \\
GRIN & \makecell[c]{$\mathbf{44.63}\,\pm\,6.18$\\$[36.95,\ 52.30]$} & \makecell[c]{$\mathbf{42.60}\,\pm\,7.28$\\$[33.56,\ 51.63]$} \\
\bottomrule
\end{tabular}
\endgroup
\end{table}

\begin{table}[htbp]
\centering
\caption{{RMSE for sensor A142, October--December 2022 (vehicles per ten-minute interval). Each cell gives mean $\pm$ sample SD, followed by the 95\% confidence interval; $n=5$ seed runs. Bold means identify the lowest unrounded mean within each column, not statistical significance.}}
\label{tab:supp_p01_92}
\begingroup
\footnotesize
\setlength{\tabcolsep}{3pt}
\renewcommand{\arraystretch}{1.15}
\begin{tabular}{lcccc}
\toprule
Method & MCAR 10\% & MCAR 30\% & MCAR 50\% & Synchronous $B=144$ \\
\midrule
Linear & \makecell[c]{$39.72\,\pm\,2.81$\\$[36.24,\ 43.20]$} & \makecell[c]{$38.99\,\pm\,0.92$\\$[37.85,\ 40.13]$} & \makecell[c]{$40.70\,\pm\,1.27$\\$[39.13,\ 42.27]$} & \makecell[c]{$142.09\,\pm\,20.80$\\$[116.27,\ 167.91]$} \\
Median & \makecell[c]{$127.07\,\pm\,2.03$\\$[124.55,\ 129.60]$} & \makecell[c]{$125.09\,\pm\,1.10$\\$[123.73,\ 126.45]$} & \makecell[c]{$126.43\,\pm\,1.29$\\$[124.83,\ 128.03]$} & \makecell[c]{$126.16\,\pm\,15.77$\\$[106.57,\ 145.75]$} \\
FFill + BFill & \makecell[c]{$46.36\,\pm\,3.69$\\$[41.77,\ 50.94]$} & \makecell[c]{$46.94\,\pm\,1.64$\\$[44.91,\ 48.98]$} & \makecell[c]{$50.99\,\pm\,2.72$\\$[47.62,\ 54.37]$} & \makecell[c]{$149.78\,\pm\,24.93$\\$[118.83,\ 180.73]$} \\
Rolling mean & \makecell[c]{$43.41\,\pm\,1.80$\\$[41.18,\ 45.65]$} & \makecell[c]{$43.36\,\pm\,1.28$\\$[41.77,\ 44.96]$} & \makecell[c]{$46.85\,\pm\,1.92$\\$[44.46,\ 49.23]$} & \makecell[c]{$142.33\,\pm\,20.74$\\$[116.57,\ 168.09]$} \\
KNN (scaled) & \makecell[c]{$44.87\,\pm\,3.33$\\$[40.73,\ 49.00]$} & \makecell[c]{$41.93\,\pm\,1.73$\\$[39.79,\ 44.08]$} & \makecell[c]{$45.21\,\pm\,0.70$\\$[44.34,\ 46.08]$} & \makecell[c]{$111.97\,\pm\,15.02$\\$[93.33,\ 130.61]$} \\
KNN (unscaled) & \makecell[c]{$51.74\,\pm\,4.85$\\$[45.71,\ 57.76]$} & \makecell[c]{$48.80\,\pm\,1.45$\\$[46.99,\ 50.60]$} & \makecell[c]{$51.30\,\pm\,0.46$\\$[50.73,\ 51.88]$} & \makecell[c]{$111.97\,\pm\,15.02$\\$[93.33,\ 130.61]$} \\
Historical & \makecell[c]{$71.93\,\pm\,1.52$\\$[70.04,\ 73.82]$} & \makecell[c]{$69.08\,\pm\,1.83$\\$[66.81,\ 71.35]$} & \makecell[c]{$70.47\,\pm\,1.51$\\$[68.60,\ 72.34]$} & \makecell[c]{$\mathbf{74.81}\,\pm\,23.49$\\$[45.65,\ 103.98]$} \\
Graph Laplacian & \makecell[c]{$60.93\,\pm\,4.18$\\$[55.74,\ 66.12]$} & \makecell[c]{$59.21\,\pm\,1.63$\\$[57.18,\ 61.23]$} & \makecell[c]{$62.45\,\pm\,0.53$\\$[61.80,\ 63.11]$} & \makecell[c]{$111.97\,\pm\,15.02$\\$[93.33,\ 130.61]$} \\
Graph Temporal & \makecell[c]{$50.36\,\pm\,3.90$\\$[45.52,\ 55.19]$} & \makecell[c]{$49.27\,\pm\,1.14$\\$[47.86,\ 50.69]$} & \makecell[c]{$53.93\,\pm\,0.36$\\$[53.48,\ 54.37]$} & \makecell[c]{$133.50\,\pm\,26.57$\\$[100.51,\ 166.49]$} \\
Sage-Voronoi & \makecell[c]{$39.28\,\pm\,2.10$\\$[36.67,\ 41.89]$} & \makecell[c]{$38.74\,\pm\,1.35$\\$[37.06,\ 40.42]$} & \makecell[c]{$44.36\,\pm\,1.03$\\$[43.08,\ 45.65]$} & \makecell[c]{$161.96\,\pm\,39.10$\\$[113.41,\ 210.50]$} \\
GAIN & \makecell[c]{$56.18\,\pm\,5.98$\\$[48.75,\ 63.60]$} & \makecell[c]{$63.92\,\pm\,6.91$\\$[55.34,\ 72.50]$} & \makecell[c]{$133.53\,\pm\,13.01$\\$[117.38,\ 149.67]$} & \makecell[c]{$538.02\,\pm\,36.67$\\$[492.49,\ 583.54]$} \\
GRIN & \makecell[c]{$\mathbf{33.73}\,\pm\,2.91$\\$[30.12,\ 37.34]$} & \makecell[c]{$\mathbf{31.34}\,\pm\,1.26$\\$[29.77,\ 32.91]$} & \makecell[c]{$\mathbf{35.65}\,\pm\,0.74$\\$[34.73,\ 36.58]$} & \makecell[c]{$118.33\,\pm\,27.48$\\$[84.20,\ 152.45]$} \\
\bottomrule
\end{tabular}
\endgroup
\end{table}

\clearpage

\begin{table}[htbp]
\centering
\caption{{MRE for sensor A142, October--December 2022 (dimensionless ratio). Each cell gives mean $\pm$ sample SD, followed by the 95\% confidence interval; $n=5$ seed runs. Bold means identify the lowest unrounded mean within each column, not statistical significance.}}
\label{tab:supp_p01_93}
\begingroup
\footnotesize
\setlength{\tabcolsep}{3pt}
\renewcommand{\arraystretch}{1.15}
\begin{tabular}{lcccc}
\toprule
Method & Block $B=1$ & Block $B=2$ & Block $B=4$ & Block $B=8$ \\
\midrule
Linear & \makecell[c]{$0.157\,\pm\,0.072$\\$[0.067,\ 0.246]$} & \makecell[c]{$0.212\,\pm\,0.216$\\$[-0.056,\ 0.481]$} & \makecell[c]{$0.115\,\pm\,0.078$\\$[0.018,\ 0.211]$} & \makecell[c]{$0.154\,\pm\,0.204$\\$[-0.099,\ 0.407]$} \\
Median & \makecell[c]{$1.872\,\pm\,1.991$\\$[-0.600,\ 4.345]$} & \makecell[c]{$2.449\,\pm\,3.861$\\$[-2.345,\ 7.243]$} & \makecell[c]{$0.135\,\pm\,0.050$\\$[0.073,\ 0.197]$} & \makecell[c]{$2.214\,\pm\,4.503$\\$[-3.377,\ 7.805]$} \\
FFill + BFill & \makecell[c]{$0.139\,\pm\,0.121$\\$[-0.011,\ 0.289]$} & \makecell[c]{$0.179\,\pm\,0.099$\\$[0.056,\ 0.303]$} & \makecell[c]{$0.137\,\pm\,0.055$\\$[0.068,\ 0.205]$} & \makecell[c]{$0.170\,\pm\,0.123$\\$[0.017,\ 0.323]$} \\
Rolling mean & \makecell[c]{$\mathbf{0.080}\,\pm\,0.045$\\$[0.023,\ 0.136]$} & \makecell[c]{$0.270\,\pm\,0.136$\\$[0.101,\ 0.439]$} & \makecell[c]{$0.181\,\pm\,0.073$\\$[0.090,\ 0.271]$} & \makecell[c]{$0.163\,\pm\,0.138$\\$[-0.008,\ 0.334]$} \\
KNN (scaled) & \makecell[c]{$0.250\,\pm\,0.174$\\$[0.033,\ 0.466]$} & \makecell[c]{$0.290\,\pm\,0.148$\\$[0.106,\ 0.474]$} & \makecell[c]{$\mathbf{0.078}\,\pm\,0.038$\\$[0.031,\ 0.126]$} & \makecell[c]{$0.196\,\pm\,0.150$\\$[0.009,\ 0.382]$} \\
KNN (unscaled) & \makecell[c]{$0.150\,\pm\,0.112$\\$[0.011,\ 0.290]$} & \makecell[c]{$0.267\,\pm\,0.087$\\$[0.159,\ 0.374]$} & \makecell[c]{$0.104\,\pm\,0.038$\\$[0.056,\ 0.151]$} & \makecell[c]{$0.292\,\pm\,0.342$\\$[-0.132,\ 0.716]$} \\
Historical & \makecell[c]{$0.418\,\pm\,0.648$\\$[-0.386,\ 1.223]$} & \makecell[c]{$0.726\,\pm\,0.376$\\$[0.258,\ 1.193]$} & \makecell[c]{$0.081\,\pm\,0.027$\\$[0.047,\ 0.114]$} & \makecell[c]{$0.344\,\pm\,0.364$\\$[-0.108,\ 0.796]$} \\
Graph Laplacian & \makecell[c]{$0.348\,\pm\,0.282$\\$[-0.002,\ 0.697]$} & \makecell[c]{$0.618\,\pm\,1.041$\\$[-0.674,\ 1.910]$} & \makecell[c]{$0.191\,\pm\,0.089$\\$[0.081,\ 0.302]$} & \makecell[c]{$0.671\,\pm\,1.075$\\$[-0.663,\ 2.005]$} \\
Graph Temporal & \makecell[c]{$0.191\,\pm\,0.234$\\$[-0.099,\ 0.481]$} & \makecell[c]{$0.456\,\pm\,0.776$\\$[-0.508,\ 1.420]$} & \makecell[c]{$0.179\,\pm\,0.057$\\$[0.108,\ 0.250]$} & \makecell[c]{$0.562\,\pm\,0.845$\\$[-0.488,\ 1.611]$} \\
Sage-Voronoi & \makecell[c]{$0.131\,\pm\,0.100$\\$[0.007,\ 0.256]$} & \makecell[c]{$0.214\,\pm\,0.101$\\$[0.088,\ 0.340]$} & \makecell[c]{$0.167\,\pm\,0.068$\\$[0.083,\ 0.251]$} & \makecell[c]{$0.192\,\pm\,0.168$\\$[-0.017,\ 0.401]$} \\
GAIN & \makecell[c]{$0.297\,\pm\,0.257$\\$[-0.022,\ 0.615]$} & \makecell[c]{$0.181\,\pm\,0.080$\\$[0.081,\ 0.280]$} & \makecell[c]{$0.218\,\pm\,0.036$\\$[0.174,\ 0.262]$} & \makecell[c]{$0.254\,\pm\,0.131$\\$[0.092,\ 0.416]$} \\
GRIN & \makecell[c]{$0.104\,\pm\,0.132$\\$[-0.060,\ 0.267]$} & \makecell[c]{$\mathbf{0.139}\,\pm\,0.073$\\$[0.049,\ 0.229]$} & \makecell[c]{$0.093\,\pm\,0.018$\\$[0.071,\ 0.116]$} & \makecell[c]{$\mathbf{0.128}\,\pm\,0.054$\\$[0.062,\ 0.195]$} \\
\bottomrule
\end{tabular}
\endgroup
\end{table}

\begin{table}[htbp]
\centering
\caption{{MRE for sensor A142, October--December 2022 (dimensionless ratio). Each cell gives mean $\pm$ sample SD, followed by the 95\% confidence interval; $n=5$ seed runs. Bold means identify the lowest unrounded mean within each column, not statistical significance.}}
\label{tab:supp_p01_94}
\begingroup
\footnotesize
\setlength{\tabcolsep}{3pt}
\renewcommand{\arraystretch}{1.15}
\begin{tabular}{lccc}
\toprule
Method & Block $B=16$ & Block $B=32$ & Block $B=64$ \\
\midrule
Linear & \makecell[c]{$0.152\,\pm\,0.052$\\$[0.087,\ 0.217]$} & \makecell[c]{$0.388\,\pm\,0.488$\\$[-0.218,\ 0.994]$} & \makecell[c]{$0.368\,\pm\,0.088$\\$[0.258,\ 0.477]$} \\
Median & \makecell[c]{$1.320\,\pm\,1.554$\\$[-0.609,\ 3.250]$} & \makecell[c]{$1.335\,\pm\,2.195$\\$[-1.390,\ 4.061]$} & \makecell[c]{$0.815\,\pm\,0.576$\\$[0.101,\ 1.530]$} \\
FFill + BFill & \makecell[c]{$0.382\,\pm\,0.232$\\$[0.093,\ 0.670]$} & \makecell[c]{$0.545\,\pm\,0.576$\\$[-0.170,\ 1.260]$} & \makecell[c]{$0.588\,\pm\,0.248$\\$[0.280,\ 0.895]$} \\
Rolling mean & \makecell[c]{$0.171\,\pm\,0.057$\\$[0.100,\ 0.241]$} & \makecell[c]{$0.411\,\pm\,0.501$\\$[-0.211,\ 1.033]$} & \makecell[c]{$0.351\,\pm\,0.064$\\$[0.271,\ 0.431]$} \\
KNN (scaled) & \makecell[c]{$0.170\,\pm\,0.065$\\$[0.089,\ 0.250]$} & \makecell[c]{$0.136\,\pm\,0.049$\\$[0.075,\ 0.196]$} & \makecell[c]{$\mathbf{0.154}\,\pm\,0.032$\\$[0.114,\ 0.195]$} \\
KNN (unscaled) & \makecell[c]{$0.168\,\pm\,0.044$\\$[0.114,\ 0.223]$} & \makecell[c]{$0.177\,\pm\,0.114$\\$[0.035,\ 0.319]$} & \makecell[c]{$0.186\,\pm\,0.023$\\$[0.157,\ 0.214]$} \\
Historical & \makecell[c]{$0.458\,\pm\,0.446$\\$[-0.095,\ 1.012]$} & \makecell[c]{$0.180\,\pm\,0.132$\\$[0.016,\ 0.345]$} & \makecell[c]{$0.231\,\pm\,0.076$\\$[0.137,\ 0.325]$} \\
Graph Laplacian & \makecell[c]{$0.285\,\pm\,0.272$\\$[-0.053,\ 0.623]$} & \makecell[c]{$0.406\,\pm\,0.487$\\$[-0.199,\ 1.010]$} & \makecell[c]{$0.275\,\pm\,0.133$\\$[0.110,\ 0.440]$} \\
Graph Temporal & \makecell[c]{$0.257\,\pm\,0.257$\\$[-0.062,\ 0.576]$} & \makecell[c]{$0.329\,\pm\,0.390$\\$[-0.155,\ 0.813]$} & \makecell[c]{$0.235\,\pm\,0.102$\\$[0.109,\ 0.362]$} \\
Sage-Voronoi & \makecell[c]{$0.214\,\pm\,0.101$\\$[0.089,\ 0.339]$} & \makecell[c]{$0.226\,\pm\,0.157$\\$[0.032,\ 0.421]$} & \makecell[c]{$0.241\,\pm\,0.027$\\$[0.208,\ 0.275]$} \\
GAIN & \makecell[c]{$0.244\,\pm\,0.133$\\$[0.078,\ 0.410]$} & \makecell[c]{$0.210\,\pm\,0.070$\\$[0.123,\ 0.297]$} & \makecell[c]{$0.238\,\pm\,0.076$\\$[0.144,\ 0.333]$} \\
GRIN & \makecell[c]{$\mathbf{0.094}\,\pm\,0.025$\\$[0.064,\ 0.125]$} & \makecell[c]{$\mathbf{0.107}\,\pm\,0.020$\\$[0.082,\ 0.132]$} & \makecell[c]{$0.155\,\pm\,0.032$\\$[0.115,\ 0.195]$} \\
\bottomrule
\end{tabular}
\endgroup
\end{table}

\clearpage

\begin{table}[htbp]
\centering
\caption{{MRE for sensor A142, October--December 2022 (dimensionless ratio). Each cell gives mean $\pm$ sample SD, followed by the 95\% confidence interval; $n=5$ seed runs. Bold means identify the lowest unrounded mean within each column, not statistical significance.}}
\label{tab:supp_p01_95}
\begingroup
\footnotesize
\setlength{\tabcolsep}{3pt}
\renewcommand{\arraystretch}{1.15}
\begin{tabular}{lcc}
\toprule
Method & Block $B=128$ & Block $B=144$ \\
\midrule
Linear & \makecell[c]{$0.920\,\pm\,0.636$\\$[0.131,\ 1.710]$} & \makecell[c]{$1.005\,\pm\,0.579$\\$[0.286,\ 1.725]$} \\
Median & \makecell[c]{$1.261\,\pm\,0.534$\\$[0.598,\ 1.923]$} & \makecell[c]{$1.350\,\pm\,0.266$\\$[1.020,\ 1.681]$} \\
FFill + BFill & \makecell[c]{$1.388\,\pm\,0.936$\\$[0.226,\ 2.550]$} & \makecell[c]{$1.015\,\pm\,0.581$\\$[0.293,\ 1.737]$} \\
Rolling mean & \makecell[c]{$0.921\,\pm\,0.642$\\$[0.124,\ 1.718]$} & \makecell[c]{$1.005\,\pm\,0.580$\\$[0.285,\ 1.725]$} \\
KNN (scaled) & \makecell[c]{$0.173\,\pm\,0.040$\\$[0.124,\ 0.223]$} & \makecell[c]{$0.162\,\pm\,0.041$\\$[0.111,\ 0.214]$} \\
KNN (unscaled) & \makecell[c]{$0.190\,\pm\,0.039$\\$[0.143,\ 0.238]$} & \makecell[c]{$0.196\,\pm\,0.038$\\$[0.149,\ 0.243]$} \\
Historical & \makecell[c]{$0.280\,\pm\,0.200$\\$[0.031,\ 0.528]$} & \makecell[c]{$0.471\,\pm\,0.343$\\$[0.045,\ 0.897]$} \\
Graph Laplacian & \makecell[c]{$0.364\,\pm\,0.089$\\$[0.254,\ 0.474]$} & \makecell[c]{$0.371\,\pm\,0.073$\\$[0.281,\ 0.461]$} \\
Graph Temporal & \makecell[c]{$0.320\,\pm\,0.074$\\$[0.228,\ 0.411]$} & \makecell[c]{$0.314\,\pm\,0.044$\\$[0.259,\ 0.368]$} \\
Sage-Voronoi & \makecell[c]{$0.267\,\pm\,0.053$\\$[0.201,\ 0.332]$} & \makecell[c]{$0.260\,\pm\,0.063$\\$[0.182,\ 0.339]$} \\
GAIN & \makecell[c]{$0.194\,\pm\,0.062$\\$[0.117,\ 0.271]$} & \makecell[c]{$0.228\,\pm\,0.029$\\$[0.192,\ 0.265]$} \\
GRIN & \makecell[c]{$\mathbf{0.140}\,\pm\,0.034$\\$[0.098,\ 0.183]$} & \makecell[c]{$\mathbf{0.137}\,\pm\,0.027$\\$[0.103,\ 0.171]$} \\
\bottomrule
\end{tabular}
\endgroup
\end{table}

\begin{table}[htbp]
\centering
\caption{{MRE for sensor A142, October--December 2022 (dimensionless ratio). Each cell gives mean $\pm$ sample SD, followed by the 95\% confidence interval; $n=5$ seed runs. Bold means identify the lowest unrounded mean within each column, not statistical significance.}}
\label{tab:supp_p01_96}
\begingroup
\footnotesize
\setlength{\tabcolsep}{3pt}
\renewcommand{\arraystretch}{1.15}
\begin{tabular}{lcccc}
\toprule
Method & MCAR 10\% & MCAR 30\% & MCAR 50\% & Synchronous $B=144$ \\
\midrule
Linear & \makecell[c]{$0.140\,\pm\,0.001$\\$[0.138,\ 0.141]$} & \makecell[c]{$0.142\,\pm\,0.007$\\$[0.134,\ 0.151]$} & \makecell[c]{$0.145\,\pm\,0.006$\\$[0.138,\ 0.153]$} & \makecell[c]{$0.952\,\pm\,0.610$\\$[0.194,\ 1.710]$} \\
Median & \makecell[c]{$1.213\,\pm\,0.055$\\$[1.144,\ 1.281]$} & \makecell[c]{$1.204\,\pm\,0.022$\\$[1.177,\ 1.231]$} & \makecell[c]{$1.216\,\pm\,0.034$\\$[1.174,\ 1.259]$} & \makecell[c]{$1.174\,\pm\,0.502$\\$[0.551,\ 1.798]$} \\
FFill + BFill & \makecell[c]{$0.166\,\pm\,0.003$\\$[0.162,\ 0.170]$} & \makecell[c]{$0.175\,\pm\,0.011$\\$[0.161,\ 0.189]$} & \makecell[c]{$0.187\,\pm\,0.010$\\$[0.175,\ 0.199]$} & \makecell[c]{$0.996\,\pm\,0.625$\\$[0.220,\ 1.773]$} \\
Rolling mean & \makecell[c]{$0.170\,\pm\,0.006$\\$[0.162,\ 0.178]$} & \makecell[c]{$0.176\,\pm\,0.007$\\$[0.167,\ 0.185]$} & \makecell[c]{$0.180\,\pm\,0.007$\\$[0.172,\ 0.188]$} & \makecell[c]{$0.954\,\pm\,0.611$\\$[0.195,\ 1.713]$} \\
KNN (scaled) & \makecell[c]{$0.166\,\pm\,0.003$\\$[0.162,\ 0.169]$} & \makecell[c]{$0.158\,\pm\,0.003$\\$[0.154,\ 0.162]$} & \makecell[c]{$0.162\,\pm\,0.003$\\$[0.158,\ 0.167]$} & \makecell[c]{$0.998\,\pm\,0.430$\\$[0.465,\ 1.531]$} \\
KNN (unscaled) & \makecell[c]{$0.191\,\pm\,0.007$\\$[0.182,\ 0.200]$} & \makecell[c]{$0.186\,\pm\,0.006$\\$[0.179,\ 0.194]$} & \makecell[c]{$0.182\,\pm\,0.003$\\$[0.178,\ 0.185]$} & \makecell[c]{$0.998\,\pm\,0.430$\\$[0.465,\ 1.531]$} \\
Historical & \makecell[c]{$0.340\,\pm\,0.024$\\$[0.310,\ 0.370]$} & \makecell[c]{$0.337\,\pm\,0.017$\\$[0.315,\ 0.358]$} & \makecell[c]{$0.335\,\pm\,0.015$\\$[0.317,\ 0.354]$} & \makecell[c]{$\mathbf{0.425}\,\pm\,0.267$\\$[0.094,\ 0.756]$} \\
Graph Laplacian & \makecell[c]{$0.344\,\pm\,0.029$\\$[0.308,\ 0.380]$} & \makecell[c]{$0.347\,\pm\,0.023$\\$[0.318,\ 0.375]$} & \makecell[c]{$0.360\,\pm\,0.018$\\$[0.338,\ 0.382]$} & \makecell[c]{$0.998\,\pm\,0.430$\\$[0.465,\ 1.531]$} \\
Graph Temporal & \makecell[c]{$0.246\,\pm\,0.009$\\$[0.235,\ 0.258]$} & \makecell[c]{$0.254\,\pm\,0.006$\\$[0.247,\ 0.262]$} & \makecell[c]{$0.271\,\pm\,0.006$\\$[0.263,\ 0.279]$} & \makecell[c]{$0.968\,\pm\,0.508$\\$[0.337,\ 1.598]$} \\
Sage-Voronoi & \makecell[c]{$0.151\,\pm\,0.009$\\$[0.140,\ 0.162]$} & \makecell[c]{$0.161\,\pm\,0.007$\\$[0.152,\ 0.170]$} & \makecell[c]{$0.178\,\pm\,0.009$\\$[0.167,\ 0.189]$} & \makecell[c]{$0.816\,\pm\,0.393$\\$[0.328,\ 1.304]$} \\
GAIN & \makecell[c]{$0.192\,\pm\,0.019$\\$[0.168,\ 0.215]$} & \makecell[c]{$0.303\,\pm\,0.043$\\$[0.250,\ 0.356]$} & \makecell[c]{$0.783\,\pm\,0.075$\\$[0.690,\ 0.877]$} & \makecell[c]{$4.495\,\pm\,1.381$\\$[2.780,\ 6.209]$} \\
GRIN & \makecell[c]{$\mathbf{0.099}\,\pm\,0.004$\\$[0.094,\ 0.103]$} & \makecell[c]{$\mathbf{0.101}\,\pm\,0.003$\\$[0.097,\ 0.105]$} & \makecell[c]{$\mathbf{0.114}\,\pm\,0.004$\\$[0.109,\ 0.119]$} & \makecell[c]{$0.975\,\pm\,0.424$\\$[0.449,\ 1.501]$} \\
\bottomrule
\end{tabular}
\endgroup
\end{table}

\clearpage

\clearpage

\endgroup

\clearpage
\renewcommand{\theHtable}{A\arabic{table}}
% P03 supplementary fragment; keep separate from main preview.

\clearpage

\begingroup

\hypersetup{linkcolor=black}

\renewcommand{\thetable}{A\arabic{table}}

\setcounter{table}{0}

\section*{{Additional analyses: supporting tables}}

{Tables A1--A3 cover persistent stream loss; A4--A9 graph controls; A10--A12 common-$k$ KNN scaling; A13--A24 count/degree strata for $B=144$ and MCAR 30\%; A25--A30 daily sensor-count sums; A31 plain Sage-Voronoi forecasting. Every metric cell gives mean, sample SD and an unadjusted run-level 95\% Student-$t$ CI from five seeds. MAE/RMSE use vehicles per 10-minute interval, except daily sums, which use aggregated sensor counts; MRE is a dimensionless ratio on nonzero targets. Full auxiliary summaries for all completed scenarios are in \texttt{additional\_analyses\_p03.csv}; per-seed training records and rounding differences are supplied separately. No method is omitted from the primary-model comparisons. Low/high count thresholds are per-sensor training quartiles; degree groups use the road graph. Daily sums score affected test days only and can include partial boundary days. Negative CI endpoints, if present, are retained as calculated and do not imply negative errors.}

\begin{table}[htbp]
\centering
\caption{{MAE under persistent loss of sensor streams. Cells give mean $\pm$ sample SD and [95\% Student-$t$ CI] across five seed runs.}}
\label{tab:p03_A1}
\begingroup
\footnotesize
\setlength{\tabcolsep}{3pt}
\renewcommand{\arraystretch}{1.15}
\begin{tabular}{lcccc}
\toprule
Method / variant & \makecell{Jan\\25\% streams} & \makecell{Jan\\50\% streams} & \makecell{Oct--Dec\\25\% streams} & \makecell{Oct--Dec\\50\% streams} \\
\midrule
Linear & \makecell[c]{$114.98\,\pm\,17.12$\\$[93.72,\ 136.23]$} & \makecell[c]{$115.77\,\pm\,3.20$\\$[111.79,\ 119.74]$} & \makecell[c]{$107.72\,\pm\,17.43$\\$[86.07,\ 129.37]$} & \makecell[c]{$111.34\,\pm\,4.14$\\$[106.20,\ 116.48]$} \\
Median & \makecell[c]{$114.98\,\pm\,17.12$\\$[93.72,\ 136.23]$} & \makecell[c]{$115.77\,\pm\,3.20$\\$[111.79,\ 119.74]$} & \makecell[c]{$107.72\,\pm\,17.43$\\$[86.07,\ 129.37]$} & \makecell[c]{$111.34\,\pm\,4.14$\\$[106.20,\ 116.48]$} \\
FFill + BFill & \makecell[c]{$114.98\,\pm\,17.12$\\$[93.72,\ 136.23]$} & \makecell[c]{$115.77\,\pm\,3.20$\\$[111.79,\ 119.74]$} & \makecell[c]{$107.72\,\pm\,17.43$\\$[86.07,\ 129.37]$} & \makecell[c]{$111.34\,\pm\,4.14$\\$[106.20,\ 116.48]$} \\
Rolling mean & \makecell[c]{$114.98\,\pm\,17.12$\\$[93.72,\ 136.23]$} & \makecell[c]{$115.77\,\pm\,3.20$\\$[111.79,\ 119.74]$} & \makecell[c]{$107.72\,\pm\,17.43$\\$[86.07,\ 129.37]$} & \makecell[c]{$111.34\,\pm\,4.14$\\$[106.20,\ 116.48]$} \\
KNN (scaled) & \makecell[c]{$101.43\,\pm\,41.09$\\$[50.41,\ 152.46]$} & \makecell[c]{$90.03\,\pm\,36.07$\\$[45.25,\ 134.82]$} & \makecell[c]{$37.83\,\pm\,17.57$\\$[16.01,\ 59.65]$} & \makecell[c]{$41.45\,\pm\,3.47$\\$[37.14,\ 45.76]$} \\
KNN (unscaled) & \makecell[c]{$40.73\,\pm\,17.45$\\$[19.06,\ 62.40]$} & \makecell[c]{$51.45\,\pm\,2.04$\\$[48.92,\ 53.99]$} & \makecell[c]{$41.30\,\pm\,16.31$\\$[21.04,\ 61.56]$} & \makecell[c]{$42.91\,\pm\,2.44$\\$[39.87,\ 45.94]$} \\
Historical & \makecell[c]{$56.80\,\pm\,18.22$\\$[34.18,\ 79.43]$} & \makecell[c]{$65.17\,\pm\,5.06$\\$[58.88,\ 71.45]$} & \makecell[c]{$57.14\,\pm\,17.00$\\$[36.03,\ 78.25]$} & \makecell[c]{$60.65\,\pm\,3.75$\\$[55.99,\ 65.31]$} \\
Graph Laplacian & \makecell[c]{$67.03\,\pm\,33.89$\\$[24.95,\ 109.11]$} & \makecell[c]{$78.51\,\pm\,18.38$\\$[55.68,\ 101.34]$} & \makecell[c]{$44.54\,\pm\,21.79$\\$[17.48,\ 71.60]$} & \makecell[c]{$51.53\,\pm\,2.52$\\$[48.40,\ 54.66]$} \\
Graph Temporal & \makecell[c]{$83.67\,\pm\,36.20$\\$[38.72,\ 128.61]$} & \makecell[c]{$80.69\,\pm\,21.61$\\$[53.85,\ 107.52]$} & \makecell[c]{$46.23\,\pm\,20.41$\\$[20.89,\ 71.56]$} & \makecell[c]{$52.23\,\pm\,2.69$\\$[48.88,\ 55.57]$} \\
Sage-Voronoi & \makecell[c]{$60.22\,\pm\,29.52$\\$[23.57,\ 96.88]$} & \makecell[c]{$74.12\,\pm\,7.48$\\$[64.83,\ 83.41]$} & \makecell[c]{$48.27\,\pm\,21.91$\\$[21.06,\ 75.48]$} & \makecell[c]{$58.79\,\pm\,6.46$\\$[50.77,\ 66.80]$} \\
GAIN & \makecell[c]{$37.90\,\pm\,16.02$\\$[18.00,\ 57.80]$} & \makecell[c]{$58.94\,\pm\,3.49$\\$[54.61,\ 63.27]$} & \makecell[c]{$39.57\,\pm\,17.19$\\$[18.23,\ 60.91]$} & \makecell[c]{$69.30\,\pm\,2.84$\\$[65.78,\ 72.82]$} \\
GRIN & \makecell[c]{$40.24\,\pm\,15.27$\\$[21.28,\ 59.20]$} & \makecell[c]{$51.42\,\pm\,5.08$\\$[45.12,\ 57.72]$} & \makecell[c]{$36.23\,\pm\,16.78$\\$[15.39,\ 57.07]$} & \makecell[c]{$44.14\,\pm\,3.54$\\$[39.74,\ 48.54]$} \\
\bottomrule
\end{tabular}
\endgroup
\end{table}

\clearpage

\begin{table}[htbp]
\centering
\caption{{RMSE under persistent loss of sensor streams. Cells give mean $\pm$ sample SD and [95\% Student-$t$ CI] across five seed runs.}}
\label{tab:p03_A2}
\begingroup
\footnotesize
\setlength{\tabcolsep}{3pt}
\renewcommand{\arraystretch}{1.15}
\begin{tabular}{lcccc}
\toprule
Method / variant & \makecell{Jan\\25\% streams} & \makecell{Jan\\50\% streams} & \makecell{Oct--Dec\\25\% streams} & \makecell{Oct--Dec\\50\% streams} \\
\midrule
Linear & \makecell[c]{$192.51\,\pm\,27.32$\\$[158.59,\ 226.43]$} & \makecell[c]{$220.44\,\pm\,35.02$\\$[176.96,\ 263.92]$} & \makecell[c]{$191.38\,\pm\,43.45$\\$[137.43,\ 245.34]$} & \makecell[c]{$208.06\,\pm\,11.78$\\$[193.44,\ 222.69]$} \\
Median & \makecell[c]{$192.51\,\pm\,27.32$\\$[158.59,\ 226.43]$} & \makecell[c]{$220.44\,\pm\,35.02$\\$[176.96,\ 263.92]$} & \makecell[c]{$191.38\,\pm\,43.45$\\$[137.43,\ 245.34]$} & \makecell[c]{$208.06\,\pm\,11.78$\\$[193.44,\ 222.69]$} \\
FFill + BFill & \makecell[c]{$192.51\,\pm\,27.32$\\$[158.59,\ 226.43]$} & \makecell[c]{$220.44\,\pm\,35.02$\\$[176.96,\ 263.92]$} & \makecell[c]{$191.38\,\pm\,43.45$\\$[137.43,\ 245.34]$} & \makecell[c]{$208.06\,\pm\,11.78$\\$[193.44,\ 222.69]$} \\
Rolling mean & \makecell[c]{$192.51\,\pm\,27.32$\\$[158.59,\ 226.43]$} & \makecell[c]{$220.44\,\pm\,35.02$\\$[176.96,\ 263.92]$} & \makecell[c]{$191.38\,\pm\,43.45$\\$[137.43,\ 245.34]$} & \makecell[c]{$208.06\,\pm\,11.78$\\$[193.44,\ 222.69]$} \\
KNN (scaled) & \makecell[c]{$194.01\,\pm\,62.69$\\$[116.18,\ 271.85]$} & \makecell[c]{$221.71\,\pm\,32.26$\\$[181.66,\ 261.76]$} & \makecell[c]{$108.50\,\pm\,64.91$\\$[27.91,\ 189.09]$} & \makecell[c]{$141.00\,\pm\,10.53$\\$[127.93,\ 154.07]$} \\
KNN (unscaled) & \makecell[c]{$99.29\,\pm\,55.13$\\$[30.84,\ 167.74]$} & \makecell[c]{$187.68\,\pm\,48.36$\\$[127.63,\ 247.73]$} & \makecell[c]{$114.17\,\pm\,63.94$\\$[34.77,\ 193.57]$} & \makecell[c]{$144.55\,\pm\,10.99$\\$[130.90,\ 158.20]$} \\
Historical & \makecell[c]{$126.63\,\pm\,52.72$\\$[61.17,\ 192.09]$} & \makecell[c]{$198.42\,\pm\,44.64$\\$[142.99,\ 253.84]$} & \makecell[c]{$130.95\,\pm\,48.10$\\$[71.23,\ 190.67]$} & \makecell[c]{$150.03\,\pm\,12.53$\\$[134.47,\ 165.59]$} \\
Graph Laplacian & \makecell[c]{$152.73\,\pm\,107.43$\\$[19.34,\ 286.12]$} & \makecell[c]{$260.95\,\pm\,44.05$\\$[206.26,\ 315.65]$} & \makecell[c]{$118.44\,\pm\,80.98$\\$[17.89,\ 218.99]$} & \makecell[c]{$172.44\,\pm\,10.38$\\$[159.55,\ 185.33]$} \\
Graph Temporal & \makecell[c]{$208.31\,\pm\,92.64$\\$[93.28,\ 323.34]$} & \makecell[c]{$278.77\,\pm\,70.13$\\$[191.70,\ 365.85]$} & \makecell[c]{$121.89\,\pm\,74.83$\\$[28.98,\ 214.80]$} & \makecell[c]{$171.18\,\pm\,10.15$\\$[158.57,\ 183.78]$} \\
Sage-Voronoi & \makecell[c]{$135.77\,\pm\,91.60$\\$[22.03,\ 249.50]$} & \makecell[c]{$230.23\,\pm\,38.12$\\$[182.90,\ 277.56]$} & \makecell[c]{$126.15\,\pm\,79.00$\\$[28.06,\ 224.24]$} & \makecell[c]{$181.16\,\pm\,12.05$\\$[166.19,\ 196.13]$} \\
GAIN & \makecell[c]{$93.99\,\pm\,53.08$\\$[28.09,\ 159.90]$} & \makecell[c]{$170.54\,\pm\,48.88$\\$[109.85,\ 231.23]$} & \makecell[c]{$110.57\,\pm\,69.27$\\$[24.56,\ 196.59]$} & \makecell[c]{$161.52\,\pm\,5.80$\\$[154.33,\ 168.72]$} \\
GRIN & \makecell[c]{$103.82\,\pm\,52.37$\\$[38.79,\ 168.85]$} & \makecell[c]{$177.87\,\pm\,43.94$\\$[123.31,\ 232.42]$} & \makecell[c]{$108.96\,\pm\,72.16$\\$[19.36,\ 198.57]$} & \makecell[c]{$147.91\,\pm\,17.83$\\$[125.77,\ 170.05]$} \\
\bottomrule
\end{tabular}
\endgroup
\end{table}

\clearpage

\begin{table}[htbp]
\centering
\caption{{MRE under persistent loss of sensor streams. Cells give mean $\pm$ sample SD and [95\% Student-$t$ CI] across five seed runs.}}
\label{tab:p03_A3}
\begingroup
\footnotesize
\setlength{\tabcolsep}{3pt}
\renewcommand{\arraystretch}{1.15}
\begin{tabular}{lcccc}
\toprule
Method / variant & \makecell{Jan\\25\% streams} & \makecell{Jan\\50\% streams} & \makecell{Oct--Dec\\25\% streams} & \makecell{Oct--Dec\\50\% streams} \\
\midrule
Linear & \makecell[c]{$2.295\,\pm\,0.119$\\$[2.147,\ 2.443]$} & \makecell[c]{$2.273\,\pm\,0.031$\\$[2.235,\ 2.311]$} & \makecell[c]{$2.403\,\pm\,0.056$\\$[2.334,\ 2.472]$} & \makecell[c]{$2.366\,\pm\,0.074$\\$[2.275,\ 2.458]$} \\
Median & \makecell[c]{$2.295\,\pm\,0.119$\\$[2.147,\ 2.443]$} & \makecell[c]{$2.273\,\pm\,0.031$\\$[2.235,\ 2.311]$} & \makecell[c]{$2.403\,\pm\,0.056$\\$[2.334,\ 2.472]$} & \makecell[c]{$2.366\,\pm\,0.074$\\$[2.275,\ 2.458]$} \\
FFill + BFill & \makecell[c]{$2.295\,\pm\,0.119$\\$[2.147,\ 2.443]$} & \makecell[c]{$2.273\,\pm\,0.031$\\$[2.235,\ 2.311]$} & \makecell[c]{$2.403\,\pm\,0.056$\\$[2.334,\ 2.472]$} & \makecell[c]{$2.366\,\pm\,0.074$\\$[2.275,\ 2.458]$} \\
Rolling mean & \makecell[c]{$2.295\,\pm\,0.119$\\$[2.147,\ 2.443]$} & \makecell[c]{$2.273\,\pm\,0.031$\\$[2.235,\ 2.311]$} & \makecell[c]{$2.403\,\pm\,0.056$\\$[2.334,\ 2.472]$} & \makecell[c]{$2.366\,\pm\,0.074$\\$[2.275,\ 2.458]$} \\
KNN (scaled) & \makecell[c]{$1.167\,\pm\,0.145$\\$[0.988,\ 1.347]$} & \makecell[c]{$0.969\,\pm\,0.582$\\$[0.246,\ 1.691]$} & \makecell[c]{$0.402\,\pm\,0.112$\\$[0.264,\ 0.541]$} & \makecell[c]{$0.388\,\pm\,0.091$\\$[0.275,\ 0.502]$} \\
KNN (unscaled) & \makecell[c]{$0.531\,\pm\,0.078$\\$[0.434,\ 0.627]$} & \makecell[c]{$0.635\,\pm\,0.303$\\$[0.259,\ 1.012]$} & \makecell[c]{$0.469\,\pm\,0.101$\\$[0.344,\ 0.594]$} & \makecell[c]{$0.411\,\pm\,0.099$\\$[0.287,\ 0.534]$} \\
Historical & \makecell[c]{$0.546\,\pm\,0.057$\\$[0.476,\ 0.617]$} & \makecell[c]{$0.563\,\pm\,0.018$\\$[0.541,\ 0.585]$} & \makecell[c]{$0.690\,\pm\,0.118$\\$[0.543,\ 0.836]$} & \makecell[c]{$0.681\,\pm\,0.097$\\$[0.560,\ 0.801]$} \\
Graph Laplacian & \makecell[c]{$1.038\,\pm\,0.234$\\$[0.747,\ 1.329]$} & \makecell[c]{$0.902\,\pm\,0.401$\\$[0.404,\ 1.399]$} & \makecell[c]{$0.518\,\pm\,0.056$\\$[0.449,\ 0.588]$} & \makecell[c]{$0.513\,\pm\,0.049$\\$[0.452,\ 0.573]$} \\
Graph Temporal & \makecell[c]{$1.204\,\pm\,0.417$\\$[0.686,\ 1.722]$} & \makecell[c]{$1.061\,\pm\,0.632$\\$[0.276,\ 1.846]$} & \makecell[c]{$0.562\,\pm\,0.082$\\$[0.460,\ 0.663]$} & \makecell[c]{$0.549\,\pm\,0.052$\\$[0.484,\ 0.613]$} \\
Sage-Voronoi & \makecell[c]{$0.908\,\pm\,0.277$\\$[0.564,\ 1.252]$} & \makecell[c]{$0.882\,\pm\,0.319$\\$[0.486,\ 1.278]$} & \makecell[c]{$0.424\,\pm\,0.052$\\$[0.358,\ 0.489]$} & \makecell[c]{$0.452\,\pm\,0.035$\\$[0.408,\ 0.496]$} \\
GAIN & \makecell[c]{$0.508\,\pm\,0.088$\\$[0.399,\ 0.617]$} & \makecell[c]{$0.965\,\pm\,0.129$\\$[0.804,\ 1.126]$} & \makecell[c]{$0.443\,\pm\,0.080$\\$[0.343,\ 0.543]$} & \makecell[c]{$1.066\,\pm\,0.098$\\$[0.944,\ 1.188]$} \\
GRIN & \makecell[c]{$0.514\,\pm\,0.103$\\$[0.386,\ 0.641]$} & \makecell[c]{$0.590\,\pm\,0.103$\\$[0.462,\ 0.719]$} & \makecell[c]{$0.348\,\pm\,0.050$\\$[0.286,\ 0.411]$} & \makecell[c]{$0.459\,\pm\,0.068$\\$[0.374,\ 0.544]$} \\
\bottomrule
\end{tabular}
\endgroup
\end{table}

\clearpage

\begin{table}[htbp]
\centering
\caption{{Graph Temporal sensitivity controls, January 2023, MAE. Cells give mean $\pm$ sample SD and [95\% Student-$t$ CI] across five seed runs.}}
\label{tab:p03_A4}
\begingroup
\footnotesize
\setlength{\tabcolsep}{3pt}
\renewcommand{\arraystretch}{1.15}
\begin{tabular}{lccc}
\toprule
Method / variant & $B=16$ & $B=144$ & MCAR 30\% \\
\midrule
Road, standardized & \makecell[c]{$61.16\,\pm\,1.87$\\$[58.83,\ 63.48]$} & \makecell[c]{$72.02\,\pm\,4.78$\\$[66.09,\ 77.96]$} & \makecell[c]{$53.79\,\pm\,0.45$\\$[53.23,\ 54.34]$} \\
Road, raw counts & \makecell[c]{$141.52\,\pm\,9.21$\\$[130.08,\ 152.96]$} & \makecell[c]{$147.89\,\pm\,4.16$\\$[142.72,\ 153.05]$} & \makecell[c]{$106.69\,\pm\,0.49$\\$[106.08,\ 107.29]$} \\
No spatial edges & \makecell[c]{$34.77\,\pm\,7.82$\\$[25.06,\ 44.48]$} & \makecell[c]{$133.05\,\pm\,5.42$\\$[126.32,\ 139.79]$} & \makecell[c]{$23.16\,\pm\,0.19$\\$[22.93,\ 23.40]$} \\
Training correlation & \makecell[c]{$38.89\,\pm\,2.04$\\$[36.36,\ 41.43]$} & \makecell[c]{$42.24\,\pm\,4.19$\\$[37.04,\ 47.44]$} & \makecell[c]{$36.44\,\pm\,0.38$\\$[35.96,\ 36.91]$} \\
\bottomrule
\end{tabular}
\endgroup
\end{table}

\clearpage

\begin{table}[htbp]
\centering
\caption{{Graph Temporal sensitivity controls, January 2023, RMSE. Cells give mean $\pm$ sample SD and [95\% Student-$t$ CI] across five seed runs.}}
\label{tab:p03_A5}
\begingroup
\footnotesize
\setlength{\tabcolsep}{3pt}
\renewcommand{\arraystretch}{1.15}
\begin{tabular}{lccc}
\toprule
Method / variant & $B=16$ & $B=144$ & MCAR 30\% \\
\midrule
Road, standardized & \makecell[c]{$182.17\,\pm\,20.91$\\$[156.20,\ 208.14]$} & \makecell[c]{$240.03\,\pm\,26.62$\\$[206.97,\ 273.08]$} & \makecell[c]{$173.76\,\pm\,1.50$\\$[171.90,\ 175.62]$} \\
Road, raw counts & \makecell[c]{$247.24\,\pm\,31.93$\\$[207.59,\ 286.89]$} & \makecell[c]{$271.82\,\pm\,19.22$\\$[247.96,\ 295.67]$} & \makecell[c]{$203.05\,\pm\,1.56$\\$[201.12,\ 204.99]$} \\
No spatial edges & \makecell[c]{$117.12\,\pm\,55.76$\\$[47.88,\ 186.36]$} & \makecell[c]{$263.30\,\pm\,20.72$\\$[237.57,\ 289.02]$} & \makecell[c]{$77.64\,\pm\,1.42$\\$[75.87,\ 79.41]$} \\
Training correlation & \makecell[c]{$127.12\,\pm\,34.78$\\$[83.93,\ 170.30]$} & \makecell[c]{$161.02\,\pm\,35.90$\\$[116.44,\ 205.59]$} & \makecell[c]{$142.71\,\pm\,3.39$\\$[138.49,\ 146.92]$} \\
\bottomrule
\end{tabular}
\endgroup
\end{table}

\clearpage

\begin{table}[htbp]
\centering
\caption{{Graph Temporal sensitivity controls, January 2023, MRE. Cells give mean $\pm$ sample SD and [95\% Student-$t$ CI] across five seed runs.}}
\label{tab:p03_A6}
\begingroup
\footnotesize
\setlength{\tabcolsep}{3pt}
\renewcommand{\arraystretch}{1.15}
\begin{tabular}{lccc}
\toprule
Method / variant & $B=16$ & $B=144$ & MCAR 30\% \\
\midrule
Road, standardized & \makecell[c]{$0.879\,\pm\,0.495$\\$[0.265,\ 1.493]$} & \makecell[c]{$1.050\,\pm\,0.158$\\$[0.853,\ 1.247]$} & \makecell[c]{$0.750\,\pm\,0.031$\\$[0.711,\ 0.788]$} \\
Road, raw counts & \makecell[c]{$6.589\,\pm\,1.182$\\$[5.121,\ 8.057]$} & \makecell[c]{$7.450\,\pm\,0.737$\\$[6.535,\ 8.365]$} & \makecell[c]{$4.953\,\pm\,0.073$\\$[4.863,\ 5.043]$} \\
No spatial edges & \makecell[c]{$0.374\,\pm\,0.060$\\$[0.300,\ 0.448]$} & \makecell[c]{$2.913\,\pm\,0.286$\\$[2.558,\ 3.268]$} & \makecell[c]{$0.262\,\pm\,0.003$\\$[0.258,\ 0.266]$} \\
Training correlation & \makecell[c]{$0.582\,\pm\,0.256$\\$[0.265,\ 0.900]$} & \makecell[c]{$0.597\,\pm\,0.064$\\$[0.517,\ 0.676]$} & \makecell[c]{$0.525\,\pm\,0.014$\\$[0.508,\ 0.543]$} \\
\bottomrule
\end{tabular}
\endgroup
\end{table}

\clearpage

\begin{table}[htbp]
\centering
\caption{{Graph Temporal sensitivity controls, October--December 2022, MAE. Cells give mean $\pm$ sample SD and [95\% Student-$t$ CI] across five seed runs.}}
\label{tab:p03_A7}
\begingroup
\footnotesize
\setlength{\tabcolsep}{3pt}
\renewcommand{\arraystretch}{1.15}
\begin{tabular}{lccc}
\toprule
Method / variant & $B=16$ & $B=144$ & MCAR 30\% \\
\midrule
Road, standardized & \makecell[c]{$48.14\,\pm\,5.03$\\$[41.90,\ 54.38]$} & \makecell[c]{$44.67\,\pm\,3.18$\\$[40.72,\ 48.62]$} & \makecell[c]{$35.37\,\pm\,0.16$\\$[35.17,\ 35.57]$} \\
Road, raw counts & \makecell[c]{$130.75\,\pm\,8.45$\\$[120.26,\ 141.24]$} & \makecell[c]{$133.32\,\pm\,6.49$\\$[125.26,\ 141.38]$} & \makecell[c]{$97.55\,\pm\,0.16$\\$[97.35,\ 97.76]$} \\
No spatial edges & \makecell[c]{$29.12\,\pm\,2.21$\\$[26.38,\ 31.86]$} & \makecell[c]{$115.76\,\pm\,6.37$\\$[107.85,\ 123.68]$} & \makecell[c]{$20.58\,\pm\,0.14$\\$[20.41,\ 20.75]$} \\
Training correlation & \makecell[c]{$41.82\,\pm\,4.61$\\$[36.09,\ 47.54]$} & \makecell[c]{$38.54\,\pm\,4.57$\\$[32.87,\ 44.21]$} & \makecell[c]{$30.75\,\pm\,0.22$\\$[30.48,\ 31.02]$} \\
\bottomrule
\end{tabular}
\endgroup
\end{table}

\clearpage

\begin{table}[htbp]
\centering
\caption{{Graph Temporal sensitivity controls, October--December 2022, RMSE. Cells give mean $\pm$ sample SD and [95\% Student-$t$ CI] across five seed runs.}}
\label{tab:p03_A8}
\begingroup
\footnotesize
\setlength{\tabcolsep}{3pt}
\renewcommand{\arraystretch}{1.15}
\begin{tabular}{lccc}
\toprule
Method / variant & $B=16$ & $B=144$ & MCAR 30\% \\
\midrule
Road, standardized & \makecell[c]{$177.38\,\pm\,55.64$\\$[108.29,\ 246.47]$} & \makecell[c]{$133.66\,\pm\,21.55$\\$[106.90,\ 160.42]$} & \makecell[c]{$104.23\,\pm\,1.62$\\$[102.22,\ 106.23]$} \\
Road, raw counts & \makecell[c]{$261.77\,\pm\,39.41$\\$[212.83,\ 310.71]$} & \makecell[c]{$242.93\,\pm\,21.96$\\$[215.66,\ 270.19]$} & \makecell[c]{$182.93\,\pm\,1.26$\\$[181.36,\ 184.50]$} \\
No spatial edges & \makecell[c]{$86.28\,\pm\,17.56$\\$[64.48,\ 108.07]$} & \makecell[c]{$208.23\,\pm\,21.84$\\$[181.11,\ 235.36]$} & \makecell[c]{$60.51\,\pm\,1.16$\\$[59.06,\ 61.95]$} \\
Training correlation & \makecell[c]{$155.65\,\pm\,49.06$\\$[94.73,\ 216.57]$} & \makecell[c]{$124.18\,\pm\,27.44$\\$[90.11,\ 158.25]$} & \makecell[c]{$84.63\,\pm\,1.81$\\$[82.38,\ 86.88]$} \\
\bottomrule
\end{tabular}
\endgroup
\end{table}

\clearpage

\begin{table}[htbp]
\centering
\caption{{Graph Temporal sensitivity controls, October--December 2022, MRE. Cells give mean $\pm$ sample SD and [95\% Student-$t$ CI] across five seed runs.}}
\label{tab:p03_A9}
\begingroup
\footnotesize
\setlength{\tabcolsep}{3pt}
\renewcommand{\arraystretch}{1.15}
\begin{tabular}{lccc}
\toprule
Method / variant & $B=16$ & $B=144$ & MCAR 30\% \\
\midrule
Road, standardized & \makecell[c]{$0.552\,\pm\,0.189$\\$[0.317,\ 0.787]$} & \makecell[c]{$0.620\,\pm\,0.089$\\$[0.509,\ 0.731]$} & \makecell[c]{$0.457\,\pm\,0.003$\\$[0.454,\ 0.461]$} \\
Road, raw counts & \makecell[c]{$4.886\,\pm\,0.677$\\$[4.044,\ 5.727]$} & \makecell[c]{$6.021\,\pm\,0.319$\\$[5.625,\ 6.417]$} & \makecell[c]{$4.240\,\pm\,0.048$\\$[4.180,\ 4.300]$} \\
No spatial edges & \makecell[c]{$0.371\,\pm\,0.052$\\$[0.306,\ 0.436]$} & \makecell[c]{$2.537\,\pm\,0.109$\\$[2.402,\ 2.673]$} & \makecell[c]{$0.251\,\pm\,0.001$\\$[0.249,\ 0.252]$} \\
Training correlation & \makecell[c]{$0.752\,\pm\,0.730$\\$[-0.154,\ 1.658]$} & \makecell[c]{$0.585\,\pm\,0.246$\\$[0.280,\ 0.890]$} & \makecell[c]{$0.414\,\pm\,0.005$\\$[0.407,\ 0.421]$} \\
\bottomrule
\end{tabular}
\endgroup
\end{table}

\clearpage

\begin{table}[htbp]
\centering
\caption{{KNN scaling control at common $k=10$, MAE. Cells give mean $\pm$ sample SD and [95\% Student-$t$ CI] across five seed runs.}}
\label{tab:p03_A10}
\begingroup
\footnotesize
\setlength{\tabcolsep}{3pt}
\renewcommand{\arraystretch}{1.15}
\begin{tabular}{lcccc}
\toprule
Method / variant & \makecell{Jan\\$B=144$} & \makecell{Jan\\MCAR 30\%} & \makecell{Oct--Dec\\$B=144$} & \makecell{Oct--Dec\\MCAR 30\%} \\
\midrule
KNN, standardized & \makecell[c]{$113.60\,\pm\,5.43$\\$[106.85,\ 120.34]$} & \makecell[c]{$104.12\,\pm\,1.46$\\$[102.30,\ 105.93]$} & \makecell[c]{$35.87\,\pm\,1.81$\\$[33.63,\ 38.11]$} & \makecell[c]{$36.78\,\pm\,0.19$\\$[36.55,\ 37.02]$} \\
KNN, unscaled & \makecell[c]{$41.52\,\pm\,3.35$\\$[37.37,\ 45.68]$} & \makecell[c]{$41.89\,\pm\,0.22$\\$[41.61,\ 42.16]$} & \makecell[c]{$38.87\,\pm\,1.84$\\$[36.58,\ 41.15]$} & \makecell[c]{$39.13\,\pm\,0.12$\\$[38.98,\ 39.29]$} \\
\bottomrule
\end{tabular}
\endgroup
\end{table}

\clearpage

\begin{table}[htbp]
\centering
\caption{{KNN scaling control at common $k=10$, RMSE. Cells give mean $\pm$ sample SD and [95\% Student-$t$ CI] across five seed runs.}}
\label{tab:p03_A11}
\begingroup
\footnotesize
\setlength{\tabcolsep}{3pt}
\renewcommand{\arraystretch}{1.15}
\begin{tabular}{lcccc}
\toprule
Method / variant & \makecell{Jan\\$B=144$} & \makecell{Jan\\MCAR 30\%} & \makecell{Oct--Dec\\$B=144$} & \makecell{Oct--Dec\\MCAR 30\%} \\
\midrule
KNN, standardized & \makecell[c]{$237.08\,\pm\,24.33$\\$[206.87,\ 267.29]$} & \makecell[c]{$227.19\,\pm\,1.50$\\$[225.33,\ 229.06]$} & \makecell[c]{$112.74\,\pm\,6.58$\\$[104.57,\ 120.92]$} & \makecell[c]{$123.63\,\pm\,1.28$\\$[122.05,\ 125.22]$} \\
KNN, unscaled & \makecell[c]{$162.90\,\pm\,31.08$\\$[124.31,\ 201.49]$} & \makecell[c]{$169.71\,\pm\,2.82$\\$[166.21,\ 173.22]$} & \makecell[c]{$114.80\,\pm\,6.69$\\$[106.50,\ 123.11]$} & \makecell[c]{$125.07\,\pm\,1.30$\\$[123.46,\ 126.69]$} \\
\bottomrule
\end{tabular}
\endgroup
\end{table}

\clearpage

\begin{table}[htbp]
\centering
\caption{{KNN scaling control at common $k=10$, MRE. Cells give mean $\pm$ sample SD and [95\% Student-$t$ CI] across five seed runs.}}
\label{tab:p03_A12}
\begingroup
\footnotesize
\setlength{\tabcolsep}{3pt}
\renewcommand{\arraystretch}{1.15}
\begin{tabular}{lcccc}
\toprule
Method / variant & \makecell{Jan\\$B=144$} & \makecell{Jan\\MCAR 30\%} & \makecell{Oct--Dec\\$B=144$} & \makecell{Oct--Dec\\MCAR 30\%} \\
\midrule
KNN, standardized & \makecell[c]{$1.051\,\pm\,0.068$\\$[0.966,\ 1.136]$} & \makecell[c]{$1.045\,\pm\,0.024$\\$[1.015,\ 1.074]$} & \makecell[c]{$0.491\,\pm\,0.125$\\$[0.335,\ 0.646]$} & \makecell[c]{$0.457\,\pm\,0.008$\\$[0.447,\ 0.467]$} \\
KNN, unscaled & \makecell[c]{$0.448\,\pm\,0.034$\\$[0.406,\ 0.491]$} & \makecell[c]{$0.451\,\pm\,0.012$\\$[0.437,\ 0.466]$} & \makecell[c]{$0.582\,\pm\,0.141$\\$[0.407,\ 0.757]$} & \makecell[c]{$0.512\,\pm\,0.009$\\$[0.500,\ 0.524]$} \\
\bottomrule
\end{tabular}
\endgroup
\end{table}

\clearpage

\begin{table}[htbp]
\centering
\caption{{January 2023, $B=144$, MAE by count and road-graph degree strata. Cells give mean $\pm$ sample SD and [95\% Student-$t$ CI] across five seed runs.}}
\label{tab:p03_A13}
\begingroup
\footnotesize
\setlength{\tabcolsep}{3pt}
\renewcommand{\arraystretch}{1.15}
\begin{tabular}{lcccc}
\toprule
Method / variant & Low count & High count & Degree 0--2 & Degree $\geq3$ \\
\midrule
Linear & \makecell[c]{$129.15\,\pm\,12.72$\\$[113.36,\ 144.94]$} & \makecell[c]{$134.22\,\pm\,13.09$\\$[117.97,\ 150.47]$} & \makecell[c]{$87.87\,\pm\,10.17$\\$[75.24,\ 100.51]$} & \makecell[c]{$162.12\,\pm\,5.22$\\$[155.64,\ 168.61]$} \\
Median & \makecell[c]{$108.45\,\pm\,2.55$\\$[105.29,\ 111.62]$} & \makecell[c]{$154.48\,\pm\,10.02$\\$[142.04,\ 166.92]$} & \makecell[c]{$75.89\,\pm\,2.79$\\$[72.42,\ 79.35]$} & \makecell[c]{$136.81\,\pm\,6.77$\\$[128.40,\ 145.21]$} \\
FFill + BFill & \makecell[c]{$128.71\,\pm\,11.64$\\$[114.26,\ 143.17]$} & \makecell[c]{$144.35\,\pm\,10.77$\\$[130.98,\ 157.73]$} & \makecell[c]{$88.10\,\pm\,8.27$\\$[77.84,\ 98.37]$} & \makecell[c]{$173.53\,\pm\,10.92$\\$[159.98,\ 187.08]$} \\
Rolling mean & \makecell[c]{$129.13\,\pm\,12.61$\\$[113.47,\ 144.78]$} & \makecell[c]{$134.33\,\pm\,13.03$\\$[118.15,\ 150.52]$} & \makecell[c]{$87.71\,\pm\,10.25$\\$[74.99,\ 100.44]$} & \makecell[c]{$162.42\,\pm\,5.31$\\$[155.82,\ 169.01]$} \\
KNN (scaled) & \makecell[c]{$39.93\,\pm\,3.90$\\$[35.09,\ 44.77]$} & \makecell[c]{$189.84\,\pm\,11.67$\\$[175.36,\ 204.33]$} & \makecell[c]{$74.00\,\pm\,4.83$\\$[68.01,\ 80.00]$} & \makecell[c]{$137.19\,\pm\,7.11$\\$[128.36,\ 146.02]$} \\
KNN (unscaled) & \makecell[c]{$28.31\,\pm\,3.85$\\$[23.53,\ 33.09]$} & \makecell[c]{$52.90\,\pm\,8.39$\\$[42.48,\ 63.32]$} & \makecell[c]{$33.89\,\pm\,3.01$\\$[30.16,\ 37.63]$} & \makecell[c]{$47.95\,\pm\,5.44$\\$[41.19,\ 54.72]$} \\
Historical & \makecell[c]{$33.57\,\pm\,3.67$\\$[29.01,\ 38.12]$} & \makecell[c]{$83.06\,\pm\,9.86$\\$[70.82,\ 95.31]$} & \makecell[c]{$45.56\,\pm\,2.76$\\$[42.13,\ 48.99]$} & \makecell[c]{$67.43\,\pm\,6.11$\\$[59.84,\ 75.02]$} \\
Graph Laplacian & \makecell[c]{$54.09\,\pm\,3.01$\\$[50.35,\ 57.83]$} & \makecell[c]{$66.93\,\pm\,9.68$\\$[54.91,\ 78.95]$} & \makecell[c]{$64.72\,\pm\,5.57$\\$[57.81,\ 71.63]$} & \makecell[c]{$67.45\,\pm\,5.30$\\$[60.87,\ 74.04]$} \\
Graph Temporal & \makecell[c]{$59.63\,\pm\,6.45$\\$[51.61,\ 67.64]$} & \makecell[c]{$77.39\,\pm\,8.52$\\$[66.81,\ 87.98]$} & \makecell[c]{$64.13\,\pm\,3.85$\\$[59.35,\ 68.91]$} & \makecell[c]{$76.73\,\pm\,5.54$\\$[69.85,\ 83.60]$} \\
Sage-Voronoi & \makecell[c]{$48.83\,\pm\,7.43$\\$[39.61,\ 58.06]$} & \makecell[c]{$58.60\,\pm\,8.71$\\$[47.79,\ 69.41]$} & \makecell[c]{$55.63\,\pm\,4.35$\\$[50.23,\ 61.04]$} & \makecell[c]{$62.91\,\pm\,3.86$\\$[58.12,\ 67.70]$} \\
GAIN & \makecell[c]{$18.12\,\pm\,1.66$\\$[16.07,\ 20.18]$} & \makecell[c]{$58.97\,\pm\,10.74$\\$[45.64,\ 72.31]$} & \makecell[c]{$29.65\,\pm\,2.64$\\$[26.36,\ 32.93]$} & \makecell[c]{$46.72\,\pm\,6.49$\\$[38.66,\ 54.79]$} \\
GRIN & \makecell[c]{$22.89\,\pm\,2.68$\\$[19.57,\ 26.22]$} & \makecell[c]{$56.50\,\pm\,9.24$\\$[45.03,\ 67.96]$} & \makecell[c]{$29.01\,\pm\,2.10$\\$[26.40,\ 31.62]$} & \makecell[c]{$46.99\,\pm\,6.66$\\$[38.71,\ 55.26]$} \\
\bottomrule
\end{tabular}
\endgroup
\end{table}

\clearpage

\begin{table}[htbp]
\centering
\caption{{January 2023, $B=144$, RMSE by count and road-graph degree strata. Cells give mean $\pm$ sample SD and [95\% Student-$t$ CI] across five seed runs.}}
\label{tab:p03_A14}
\begingroup
\footnotesize
\setlength{\tabcolsep}{3pt}
\renewcommand{\arraystretch}{1.15}
\begin{tabular}{lcccc}
\toprule
Method / variant & Low count & High count & Degree 0--2 & Degree $\geq3$ \\
\midrule
Linear & \makecell[c]{$249.86\,\pm\,31.56$\\$[210.67,\ 289.04]$} & \makecell[c]{$299.75\,\pm\,35.86$\\$[255.22,\ 344.28]$} & \makecell[c]{$197.58\,\pm\,57.78$\\$[125.84,\ 269.32]$} & \makecell[c]{$297.74\,\pm\,22.91$\\$[269.30,\ 326.18]$} \\
Median & \makecell[c]{$167.53\,\pm\,7.38$\\$[158.37,\ 176.70]$} & \makecell[c]{$301.48\,\pm\,42.66$\\$[248.52,\ 354.45]$} & \makecell[c]{$164.78\,\pm\,17.99$\\$[142.45,\ 187.11]$} & \makecell[c]{$244.87\,\pm\,36.05$\\$[200.10,\ 289.63]$} \\
FFill + BFill & \makecell[c]{$254.95\,\pm\,38.45$\\$[207.22,\ 302.69]$} & \makecell[c]{$333.95\,\pm\,51.05$\\$[270.55,\ 397.34]$} & \makecell[c]{$193.82\,\pm\,45.53$\\$[137.29,\ 250.36]$} & \makecell[c]{$342.28\,\pm\,52.07$\\$[277.62,\ 406.93]$} \\
Rolling mean & \makecell[c]{$249.84\,\pm\,31.53$\\$[210.70,\ 288.99]$} & \makecell[c]{$300.43\,\pm\,35.98$\\$[255.76,\ 345.10]$} & \makecell[c]{$197.03\,\pm\,57.72$\\$[125.35,\ 268.70]$} & \makecell[c]{$298.69\,\pm\,23.53$\\$[269.47,\ 327.91]$} \\
KNN (scaled) & \makecell[c]{$99.08\,\pm\,11.03$\\$[85.39,\ 112.77]$} & \makecell[c]{$345.68\,\pm\,36.60$\\$[300.23,\ 391.12]$} & \makecell[c]{$174.19\,\pm\,21.57$\\$[147.41,\ 200.98]$} & \makecell[c]{$266.91\,\pm\,33.74$\\$[225.03,\ 308.80]$} \\
KNN (unscaled) & \makecell[c]{$112.27\,\pm\,21.52$\\$[85.55,\ 138.98]$} & \makecell[c]{$223.36\,\pm\,53.90$\\$[156.44,\ 290.28]$} & \makecell[c]{$131.53\,\pm\,18.77$\\$[108.23,\ 154.84]$} & \makecell[c]{$178.47\,\pm\,46.19$\\$[121.12,\ 235.82]$} \\
Historical & \makecell[c]{$127.16\,\pm\,13.92$\\$[109.88,\ 144.44]$} & \makecell[c]{$242.32\,\pm\,52.00$\\$[177.75,\ 306.88]$} & \makecell[c]{$158.46\,\pm\,17.73$\\$[136.45,\ 180.48]$} & \makecell[c]{$189.71\,\pm\,45.65$\\$[133.03,\ 246.39]$} \\
Graph Laplacian & \makecell[c]{$189.36\,\pm\,17.88$\\$[167.16,\ 211.57]$} & \makecell[c]{$233.80\,\pm\,54.26$\\$[166.43,\ 301.17]$} & \makecell[c]{$252.93\,\pm\,40.08$\\$[203.16,\ 302.70]$} & \makecell[c]{$191.85\,\pm\,41.59$\\$[140.21,\ 243.49]$} \\
Graph Temporal & \makecell[c]{$212.06\,\pm\,22.79$\\$[183.76,\ 240.35]$} & \makecell[c]{$264.79\,\pm\,45.74$\\$[208.00,\ 321.59]$} & \makecell[c]{$254.63\,\pm\,20.27$\\$[229.47,\ 279.80]$} & \makecell[c]{$229.84\,\pm\,38.58$\\$[181.94,\ 277.74]$} \\
Sage-Voronoi & \makecell[c]{$165.69\,\pm\,15.07$\\$[146.97,\ 184.41]$} & \makecell[c]{$223.55\,\pm\,55.30$\\$[154.89,\ 292.21]$} & \makecell[c]{$215.71\,\pm\,32.91$\\$[174.85,\ 256.57]$} & \makecell[c]{$188.25\,\pm\,36.72$\\$[142.65,\ 233.85]$} \\
GAIN & \makecell[c]{$62.39\,\pm\,10.33$\\$[49.56,\ 75.22]$} & \makecell[c]{$239.55\,\pm\,55.44$\\$[170.71,\ 308.39]$} & \makecell[c]{$125.00\,\pm\,21.02$\\$[98.90,\ 151.10]$} & \makecell[c]{$179.85\,\pm\,46.47$\\$[122.14,\ 237.55]$} \\
GRIN & \makecell[c]{$80.64\,\pm\,15.56$\\$[61.32,\ 99.96]$} & \makecell[c]{$232.63\,\pm\,54.66$\\$[164.75,\ 300.50]$} & \makecell[c]{$117.38\,\pm\,17.03$\\$[96.24,\ 138.52]$} & \makecell[c]{$180.09\,\pm\,47.77$\\$[120.79,\ 239.40]$} \\
\bottomrule
\end{tabular}
\endgroup
\end{table}

\clearpage

\begin{table}[htbp]
\centering
\caption{{January 2023, $B=144$, MRE by count and road-graph degree strata. Cells give mean $\pm$ sample SD and [95\% Student-$t$ CI] across five seed runs.}}
\label{tab:p03_A15}
\begingroup
\footnotesize
\setlength{\tabcolsep}{3pt}
\renewcommand{\arraystretch}{1.15}
\begin{tabular}{lcccc}
\toprule
Method / variant & Low count & High count & Degree 0--2 & Degree $\geq3$ \\
\midrule
Linear & \makecell[c]{$10.958\,\pm\,1.427$\\$[9.186,\ 12.731]$} & \makecell[c]{$0.512\,\pm\,0.034$\\$[0.469,\ 0.554]$} & \makecell[c]{$3.224\,\pm\,0.447$\\$[2.668,\ 3.779]$} & \makecell[c]{$2.774\,\pm\,0.251$\\$[2.462,\ 3.085]$} \\
Median & \makecell[c]{$8.649\,\pm\,0.343$\\$[8.223,\ 9.076]$} & \makecell[c]{$0.567\,\pm\,0.005$\\$[0.561,\ 0.574]$} & \makecell[c]{$2.431\,\pm\,0.079$\\$[2.333,\ 2.530]$} & \makecell[c]{$2.178\,\pm\,0.071$\\$[2.090,\ 2.267]$} \\
FFill + BFill & \makecell[c]{$10.669\,\pm\,1.203$\\$[9.176,\ 12.163]$} & \makecell[c]{$0.551\,\pm\,0.046$\\$[0.493,\ 0.608]$} & \makecell[c]{$3.182\,\pm\,0.383$\\$[2.706,\ 3.658]$} & \makecell[c]{$2.769\,\pm\,0.185$\\$[2.539,\ 2.999]$} \\
Rolling mean & \makecell[c]{$10.952\,\pm\,1.424$\\$[9.184,\ 12.720]$} & \makecell[c]{$0.512\,\pm\,0.034$\\$[0.470,\ 0.554]$} & \makecell[c]{$3.221\,\pm\,0.446$\\$[2.666,\ 3.775]$} & \makecell[c]{$2.774\,\pm\,0.249$\\$[2.464,\ 3.084]$} \\
KNN (scaled) & \makecell[c]{$2.891\,\pm\,0.302$\\$[2.516,\ 3.265]$} & \makecell[c]{$0.660\,\pm\,0.013$\\$[0.644,\ 0.676]$} & \makecell[c]{$1.164\,\pm\,0.078$\\$[1.067,\ 1.262]$} & \makecell[c]{$0.987\,\pm\,0.068$\\$[0.902,\ 1.071]$} \\
KNN (unscaled) & \makecell[c]{$1.187\,\pm\,0.173$\\$[0.972,\ 1.402]$} & \makecell[c]{$0.216\,\pm\,0.005$\\$[0.211,\ 0.222]$} & \makecell[c]{$0.602\,\pm\,0.086$\\$[0.495,\ 0.709]$} & \makecell[c]{$0.376\,\pm\,0.020$\\$[0.351,\ 0.402]$} \\
Historical & \makecell[c]{$1.452\,\pm\,0.092$\\$[1.338,\ 1.567]$} & \makecell[c]{$0.288\,\pm\,0.005$\\$[0.282,\ 0.293]$} & \makecell[c]{$0.672\,\pm\,0.048$\\$[0.613,\ 0.732]$} & \makecell[c]{$0.484\,\pm\,0.021$\\$[0.459,\ 0.510]$} \\
Graph Laplacian & \makecell[c]{$2.874\,\pm\,0.292$\\$[2.511,\ 3.237]$} & \makecell[c]{$0.258\,\pm\,0.008$\\$[0.249,\ 0.268]$} & \makecell[c]{$1.265\,\pm\,0.132$\\$[1.101,\ 1.430]$} & \makecell[c]{$0.695\,\pm\,0.027$\\$[0.661,\ 0.728]$} \\
Graph Temporal & \makecell[c]{$3.406\,\pm\,0.608$\\$[2.651,\ 4.160]$} & \makecell[c]{$0.313\,\pm\,0.009$\\$[0.302,\ 0.325]$} & \makecell[c]{$1.661\,\pm\,0.336$\\$[1.244,\ 2.078]$} & \makecell[c]{$0.706\,\pm\,0.066$\\$[0.624,\ 0.789]$} \\
Sage-Voronoi & \makecell[c]{$2.420\,\pm\,0.587$\\$[1.691,\ 3.148]$} & \makecell[c]{$0.232\,\pm\,0.020$\\$[0.207,\ 0.256]$} & \makecell[c]{$1.053\,\pm\,0.178$\\$[0.833,\ 1.274]$} & \makecell[c]{$0.651\,\pm\,0.118$\\$[0.505,\ 0.797]$} \\
GAIN & \makecell[c]{$0.844\,\pm\,0.170$\\$[0.633,\ 1.055]$} & \makecell[c]{$0.218\,\pm\,0.008$\\$[0.208,\ 0.227]$} & \makecell[c]{$0.470\,\pm\,0.076$\\$[0.375,\ 0.565]$} & \makecell[c]{$0.293\,\pm\,0.025$\\$[0.262,\ 0.324]$} \\
GRIN & \makecell[c]{$1.219\,\pm\,0.133$\\$[1.053,\ 1.384]$} & \makecell[c]{$0.217\,\pm\,0.014$\\$[0.200,\ 0.235]$} & \makecell[c]{$0.624\,\pm\,0.071$\\$[0.537,\ 0.712]$} & \makecell[c]{$0.337\,\pm\,0.031$\\$[0.298,\ 0.376]$} \\
\bottomrule
\end{tabular}
\endgroup
\end{table}

\clearpage

\begin{table}[htbp]
\centering
\caption{{January 2023, MCAR 30\%, MAE by count and road-graph degree strata. Cells give mean $\pm$ sample SD and [95\% Student-$t$ CI] across five seed runs.}}
\label{tab:p03_A16}
\begingroup
\footnotesize
\setlength{\tabcolsep}{3pt}
\renewcommand{\arraystretch}{1.15}
\begin{tabular}{lcccc}
\toprule
Method / variant & Low count & High count & Degree 0--2 & Degree $\geq3$ \\
\midrule
Linear & \makecell[c]{$9.19\,\pm\,0.28$\\$[8.85,\ 9.54]$} & \makecell[c]{$31.67\,\pm\,0.42$\\$[31.15,\ 32.19]$} & \makecell[c]{$17.40\,\pm\,0.36$\\$[16.94,\ 17.85]$} & \makecell[c]{$27.07\,\pm\,0.16$\\$[26.88,\ 27.27]$} \\
Median & \makecell[c]{$109.30\,\pm\,1.35$\\$[107.62,\ 110.97]$} & \makecell[c]{$150.21\,\pm\,0.65$\\$[149.40,\ 151.02]$} & \makecell[c]{$72.42\,\pm\,0.37$\\$[71.96,\ 72.87]$} & \makecell[c]{$135.85\,\pm\,0.97$\\$[134.65,\ 137.06]$} \\
FFill + BFill & \makecell[c]{$11.46\,\pm\,0.55$\\$[10.77,\ 12.15]$} & \makecell[c]{$39.06\,\pm\,0.84$\\$[38.01,\ 40.11]$} & \makecell[c]{$21.86\,\pm\,0.43$\\$[21.33,\ 22.39]$} & \makecell[c]{$34.51\,\pm\,0.34$\\$[34.09,\ 34.93]$} \\
Rolling mean & \makecell[c]{$12.51\,\pm\,0.40$\\$[12.01,\ 13.01]$} & \makecell[c]{$38.77\,\pm\,0.59$\\$[38.03,\ 39.50]$} & \makecell[c]{$22.13\,\pm\,0.50$\\$[21.51,\ 22.75]$} & \makecell[c]{$36.04\,\pm\,0.21$\\$[35.78,\ 36.29]$} \\
KNN (scaled) & \makecell[c]{$43.70\,\pm\,1.26$\\$[42.13,\ 45.26]$} & \makecell[c]{$166.96\,\pm\,2.70$\\$[163.62,\ 170.31]$} & \makecell[c]{$67.10\,\pm\,0.96$\\$[65.91,\ 68.29]$} & \makecell[c]{$126.25\,\pm\,2.09$\\$[123.65,\ 128.85]$} \\
KNN (unscaled) & \makecell[c]{$30.87\,\pm\,1.09$\\$[29.51,\ 32.23]$} & \makecell[c]{$51.44\,\pm\,0.75$\\$[50.51,\ 52.38]$} & \makecell[c]{$33.39\,\pm\,0.78$\\$[32.42,\ 34.36]$} & \makecell[c]{$48.95\,\pm\,0.46$\\$[48.38,\ 49.52]$} \\
Historical & \makecell[c]{$35.37\,\pm\,1.03$\\$[34.09,\ 36.65]$} & \makecell[c]{$80.61\,\pm\,0.51$\\$[79.97,\ 81.25]$} & \makecell[c]{$44.53\,\pm\,0.84$\\$[43.49,\ 45.58]$} & \makecell[c]{$68.21\,\pm\,0.58$\\$[67.49,\ 68.92]$} \\
Graph Laplacian & \makecell[c]{$52.49\,\pm\,1.75$\\$[50.31,\ 54.66]$} & \makecell[c]{$65.34\,\pm\,1.05$\\$[64.03,\ 66.64]$} & \makecell[c]{$63.00\,\pm\,0.66$\\$[62.17,\ 63.82]$} & \makecell[c]{$68.46\,\pm\,0.61$\\$[67.71,\ 69.22]$} \\
Graph Temporal & \makecell[c]{$42.49\,\pm\,1.01$\\$[41.24,\ 43.75]$} & \makecell[c]{$58.77\,\pm\,1.32$\\$[57.13,\ 60.41]$} & \makecell[c]{$39.67\,\pm\,0.71$\\$[38.79,\ 40.56]$} & \makecell[c]{$62.23\,\pm\,0.90$\\$[61.11,\ 63.34]$} \\
Sage-Voronoi & \makecell[c]{$15.23\,\pm\,0.95$\\$[14.05,\ 16.42]$} & \makecell[c]{$38.64\,\pm\,1.18$\\$[37.17,\ 40.11]$} & \makecell[c]{$25.07\,\pm\,1.13$\\$[23.68,\ 26.47]$} & \makecell[c]{$34.67\,\pm\,0.59$\\$[33.93,\ 35.40]$} \\
GAIN & \makecell[c]{$27.79\,\pm\,1.39$\\$[26.07,\ 29.52]$} & \makecell[c]{$56.98\,\pm\,1.40$\\$[55.23,\ 58.72]$} & \makecell[c]{$29.14\,\pm\,0.91$\\$[28.01,\ 30.27]$} & \makecell[c]{$50.78\,\pm\,1.11$\\$[49.39,\ 52.16]$} \\
GRIN & \makecell[c]{$9.21\,\pm\,0.51$\\$[8.58,\ 9.85]$} & \makecell[c]{$35.68\,\pm\,1.10$\\$[34.31,\ 37.05]$} & \makecell[c]{$15.98\,\pm\,0.46$\\$[15.41,\ 16.54]$} & \makecell[c]{$28.49\,\pm\,0.48$\\$[27.89,\ 29.08]$} \\
\bottomrule
\end{tabular}
\endgroup
\end{table}

\clearpage

\begin{table}[htbp]
\centering
\caption{{January 2023, MCAR 30\%, RMSE by count and road-graph degree strata. Cells give mean $\pm$ sample SD and [95\% Student-$t$ CI] across five seed runs.}}
\label{tab:p03_A17}
\begingroup
\footnotesize
\setlength{\tabcolsep}{3pt}
\renewcommand{\arraystretch}{1.15}
\begin{tabular}{lcccc}
\toprule
Method / variant & Low count & High count & Degree 0--2 & Degree $\geq3$ \\
\midrule
Linear & \makecell[c]{$38.99\,\pm\,2.07$\\$[36.43,\ 41.56]$} & \makecell[c]{$107.26\,\pm\,1.94$\\$[104.85,\ 109.68]$} & \makecell[c]{$69.40\,\pm\,3.34$\\$[65.25,\ 73.54]$} & \makecell[c]{$82.68\,\pm\,0.59$\\$[81.95,\ 83.41]$} \\
Median & \makecell[c]{$171.10\,\pm\,2.49$\\$[168.01,\ 174.20]$} & \makecell[c]{$298.90\,\pm\,3.56$\\$[294.48,\ 303.32]$} & \makecell[c]{$145.43\,\pm\,3.96$\\$[140.52,\ 150.35]$} & \makecell[c]{$250.39\,\pm\,2.72$\\$[247.02,\ 253.77]$} \\
FFill + BFill & \makecell[c]{$56.98\,\pm\,4.08$\\$[51.92,\ 62.04]$} & \makecell[c]{$133.74\,\pm\,6.35$\\$[125.86,\ 141.62]$} & \makecell[c]{$89.43\,\pm\,4.27$\\$[84.13,\ 94.73]$} & \makecell[c]{$107.34\,\pm\,4.47$\\$[101.79,\ 112.89]$} \\
Rolling mean & \makecell[c]{$59.62\,\pm\,3.18$\\$[55.67,\ 63.57]$} & \makecell[c]{$135.08\,\pm\,5.60$\\$[128.13,\ 142.03]$} & \makecell[c]{$88.82\,\pm\,6.35$\\$[80.93,\ 96.70]$} & \makecell[c]{$109.88\,\pm\,2.71$\\$[106.52,\ 113.24]$} \\
KNN (scaled) & \makecell[c]{$112.19\,\pm\,2.57$\\$[109.00,\ 115.39]$} & \makecell[c]{$327.87\,\pm\,2.83$\\$[324.36,\ 331.39]$} & \makecell[c]{$155.48\,\pm\,3.48$\\$[151.16,\ 159.81]$} & \makecell[c]{$260.79\,\pm\,3.17$\\$[256.86,\ 264.73]$} \\
KNN (unscaled) & \makecell[c]{$123.26\,\pm\,3.50$\\$[118.91,\ 127.61]$} & \makecell[c]{$233.18\,\pm\,5.58$\\$[226.25,\ 240.11]$} & \makecell[c]{$129.22\,\pm\,5.19$\\$[122.76,\ 135.67]$} & \makecell[c]{$191.52\,\pm\,3.70$\\$[186.92,\ 196.12]$} \\
Historical & \makecell[c]{$135.18\,\pm\,4.40$\\$[129.71,\ 140.64]$} & \makecell[c]{$248.44\,\pm\,4.93$\\$[242.32,\ 254.56]$} & \makecell[c]{$151.15\,\pm\,5.58$\\$[144.22,\ 158.08]$} & \makecell[c]{$202.19\,\pm\,3.58$\\$[197.76,\ 206.63]$} \\
Graph Laplacian & \makecell[c]{$177.73\,\pm\,8.62$\\$[167.02,\ 188.43]$} & \makecell[c]{$241.53\,\pm\,4.99$\\$[235.34,\ 247.72]$} & \makecell[c]{$250.30\,\pm\,3.62$\\$[245.81,\ 254.79]$} & \makecell[c]{$204.18\,\pm\,3.62$\\$[199.69,\ 208.67]$} \\
Graph Temporal & \makecell[c]{$137.88\,\pm\,3.25$\\$[133.85,\ 141.91]$} & \makecell[c]{$210.72\,\pm\,4.29$\\$[205.39,\ 216.05]$} & \makecell[c]{$142.61\,\pm\,4.21$\\$[137.38,\ 147.83]$} & \makecell[c]{$189.93\,\pm\,3.26$\\$[185.88,\ 193.98]$} \\
Sage-Voronoi & \makecell[c]{$60.31\,\pm\,5.24$\\$[53.81,\ 66.82]$} & \makecell[c]{$157.20\,\pm\,14.03$\\$[139.77,\ 174.62]$} & \makecell[c]{$96.36\,\pm\,5.04$\\$[90.10,\ 102.62]$} & \makecell[c]{$125.27\,\pm\,9.83$\\$[113.07,\ 137.48]$} \\
GAIN & \makecell[c]{$82.17\,\pm\,3.76$\\$[77.51,\ 86.84]$} & \makecell[c]{$238.49\,\pm\,5.62$\\$[231.51,\ 245.47]$} & \makecell[c]{$101.60\,\pm\,4.55$\\$[95.95,\ 107.26]$} & \makecell[c]{$190.63\,\pm\,4.21$\\$[185.41,\ 195.85]$} \\
GRIN & \makecell[c]{$32.39\,\pm\,4.67$\\$[26.59,\ 38.20]$} & \makecell[c]{$160.30\,\pm\,12.17$\\$[145.18,\ 175.41]$} & \makecell[c]{$66.49\,\pm\,3.98$\\$[61.55,\ 71.43]$} & \makecell[c]{$124.34\,\pm\,8.79$\\$[113.43,\ 135.25]$} \\
\bottomrule
\end{tabular}
\endgroup
\end{table}

\clearpage

\begin{table}[htbp]
\centering
\caption{{January 2023, MCAR 30\%, MRE by count and road-graph degree strata. Cells give mean $\pm$ sample SD and [95\% Student-$t$ CI] across five seed runs.}}
\label{tab:p03_A18}
\begingroup
\footnotesize
\setlength{\tabcolsep}{3pt}
\renewcommand{\arraystretch}{1.15}
\begin{tabular}{lcccc}
\toprule
Method / variant & Low count & High count & Degree 0--2 & Degree $\geq3$ \\
\midrule
Linear & \makecell[c]{$0.552\,\pm\,0.009$\\$[0.541,\ 0.562]$} & \makecell[c]{$0.178\,\pm\,0.002$\\$[0.175,\ 0.181]$} & \makecell[c]{$0.308\,\pm\,0.004$\\$[0.303,\ 0.313]$} & \makecell[c]{$0.242\,\pm\,0.005$\\$[0.235,\ 0.248]$} \\
Median & \makecell[c]{$8.810\,\pm\,0.107$\\$[8.677,\ 8.943]$} & \makecell[c]{$0.563\,\pm\,0.001$\\$[0.562,\ 0.565]$} & \makecell[c]{$2.341\,\pm\,0.045$\\$[2.286,\ 2.397]$} & \makecell[c]{$2.228\,\pm\,0.016$\\$[2.207,\ 2.248]$} \\
FFill + BFill & \makecell[c]{$0.647\,\pm\,0.017$\\$[0.626,\ 0.668]$} & \makecell[c]{$0.213\,\pm\,0.004$\\$[0.209,\ 0.217]$} & \makecell[c]{$0.362\,\pm\,0.003$\\$[0.357,\ 0.366]$} & \makecell[c]{$0.294\,\pm\,0.007$\\$[0.285,\ 0.303]$} \\
Rolling mean & \makecell[c]{$0.655\,\pm\,0.013$\\$[0.638,\ 0.671]$} & \makecell[c]{$0.194\,\pm\,0.002$\\$[0.191,\ 0.197]$} & \makecell[c]{$0.343\,\pm\,0.002$\\$[0.340,\ 0.346]$} & \makecell[c]{$0.301\,\pm\,0.006$\\$[0.293,\ 0.308]$} \\
KNN (scaled) & \makecell[c]{$3.049\,\pm\,0.088$\\$[2.940,\ 3.158]$} & \makecell[c]{$0.596\,\pm\,0.008$\\$[0.586,\ 0.606]$} & \makecell[c]{$1.137\,\pm\,0.036$\\$[1.091,\ 1.182]$} & \makecell[c]{$0.993\,\pm\,0.020$\\$[0.967,\ 1.018]$} \\
KNN (unscaled) & \makecell[c]{$1.251\,\pm\,0.064$\\$[1.172,\ 1.330]$} & \makecell[c]{$0.211\,\pm\,0.002$\\$[0.208,\ 0.214]$} & \makecell[c]{$0.583\,\pm\,0.023$\\$[0.554,\ 0.612]$} & \makecell[c]{$0.404\,\pm\,0.012$\\$[0.390,\ 0.419]$} \\
Historical & \makecell[c]{$1.414\,\pm\,0.025$\\$[1.383,\ 1.445]$} & \makecell[c]{$0.284\,\pm\,0.001$\\$[0.283,\ 0.285]$} & \makecell[c]{$0.625\,\pm\,0.011$\\$[0.611,\ 0.639]$} & \makecell[c]{$0.490\,\pm\,0.005$\\$[0.484,\ 0.496]$} \\
Graph Laplacian & \makecell[c]{$2.813\,\pm\,0.081$\\$[2.712,\ 2.914]$} & \makecell[c]{$0.256\,\pm\,0.000$\\$[0.255,\ 0.256]$} & \makecell[c]{$1.181\,\pm\,0.045$\\$[1.126,\ 1.237]$} & \makecell[c]{$0.717\,\pm\,0.007$\\$[0.709,\ 0.725]$} \\
Graph Temporal & \makecell[c]{$2.356\,\pm\,0.153$\\$[2.165,\ 2.546]$} & \makecell[c]{$0.248\,\pm\,0.002$\\$[0.246,\ 0.251]$} & \makecell[c]{$1.074\,\pm\,0.097$\\$[0.954,\ 1.194]$} & \makecell[c]{$0.567\,\pm\,0.014$\\$[0.549,\ 0.585]$} \\
Sage-Voronoi & \makecell[c]{$0.698\,\pm\,0.054$\\$[0.631,\ 0.765]$} & \makecell[c]{$0.164\,\pm\,0.003$\\$[0.161,\ 0.167]$} & \makecell[c]{$0.363\,\pm\,0.021$\\$[0.336,\ 0.389]$} & \makecell[c]{$0.274\,\pm\,0.007$\\$[0.266,\ 0.282]$} \\
GAIN & \makecell[c]{$1.558\,\pm\,0.160$\\$[1.359,\ 1.757]$} & \makecell[c]{$0.218\,\pm\,0.002$\\$[0.215,\ 0.220]$} & \makecell[c]{$0.642\,\pm\,0.049$\\$[0.582,\ 0.703]$} & \makecell[c]{$0.463\,\pm\,0.035$\\$[0.420,\ 0.506]$} \\
GRIN & \makecell[c]{$0.552\,\pm\,0.037$\\$[0.506,\ 0.598]$} & \makecell[c]{$0.152\,\pm\,0.006$\\$[0.145,\ 0.159]$} & \makecell[c]{$0.312\,\pm\,0.019$\\$[0.288,\ 0.336]$} & \makecell[c]{$0.203\,\pm\,0.004$\\$[0.198,\ 0.208]$} \\
\bottomrule
\end{tabular}
\endgroup
\end{table}

\clearpage

\begin{table}[htbp]
\centering
\caption{{October--December 2022, $B=144$, MAE by count and road-graph degree strata. Cells give mean $\pm$ sample SD and [95\% Student-$t$ CI] across five seed runs.}}
\label{tab:p03_A19}
\begingroup
\footnotesize
\setlength{\tabcolsep}{3pt}
\renewcommand{\arraystretch}{1.15}
\begin{tabular}{lcccc}
\toprule
Method / variant & Low count & High count & Degree 0--2 & Degree $\geq3$ \\
\midrule
Linear & \makecell[c]{$105.65\,\pm\,9.76$\\$[93.53,\ 117.77]$} & \makecell[c]{$126.14\,\pm\,9.67$\\$[114.14,\ 138.14]$} & \makecell[c]{$82.83\,\pm\,12.29$\\$[67.57,\ 98.08]$} & \makecell[c]{$137.36\,\pm\,4.42$\\$[131.87,\ 142.85]$} \\
Median & \makecell[c]{$113.24\,\pm\,1.21$\\$[111.75,\ 114.74]$} & \makecell[c]{$150.25\,\pm\,11.22$\\$[136.31,\ 164.19]$} & \makecell[c]{$74.32\,\pm\,9.13$\\$[62.99,\ 85.65]$} & \makecell[c]{$119.92\,\pm\,1.19$\\$[118.45,\ 121.40]$} \\
FFill + BFill & \makecell[c]{$106.13\,\pm\,9.29$\\$[94.60,\ 117.67]$} & \makecell[c]{$130.70\,\pm\,14.74$\\$[112.40,\ 149.00]$} & \makecell[c]{$83.08\,\pm\,12.26$\\$[67.86,\ 98.31]$} & \makecell[c]{$140.07\,\pm\,6.20$\\$[132.38,\ 147.77]$} \\
Rolling mean & \makecell[c]{$105.69\,\pm\,9.74$\\$[93.60,\ 117.78]$} & \makecell[c]{$126.09\,\pm\,9.63$\\$[114.13,\ 138.06]$} & \makecell[c]{$82.77\,\pm\,12.28$\\$[67.53,\ 98.02]$} & \makecell[c]{$137.41\,\pm\,4.45$\\$[131.88,\ 142.94]$} \\
KNN (scaled) & \makecell[c]{$18.29\,\pm\,2.26$\\$[15.48,\ 21.10]$} & \makecell[c]{$63.58\,\pm\,7.54$\\$[54.22,\ 72.95]$} & \makecell[c]{$30.59\,\pm\,4.99$\\$[24.40,\ 36.78]$} & \makecell[c]{$38.99\,\pm\,0.95$\\$[37.82,\ 40.17]$} \\
KNN (unscaled) & \makecell[c]{$20.35\,\pm\,2.87$\\$[16.78,\ 23.91]$} & \makecell[c]{$64.38\,\pm\,7.50$\\$[55.07,\ 73.69]$} & \makecell[c]{$34.28\,\pm\,4.65$\\$[28.51,\ 40.05]$} & \makecell[c]{$43.11\,\pm\,2.02$\\$[40.60,\ 45.62]$} \\
Historical & \makecell[c]{$32.08\,\pm\,1.79$\\$[29.86,\ 34.30]$} & \makecell[c]{$74.55\,\pm\,10.73$\\$[61.22,\ 87.88]$} & \makecell[c]{$39.88\,\pm\,6.52$\\$[31.78,\ 47.97]$} & \makecell[c]{$59.64\,\pm\,2.27$\\$[56.82,\ 62.45]$} \\
Graph Laplacian & \makecell[c]{$20.20\,\pm\,1.72$\\$[18.06,\ 22.33]$} & \makecell[c]{$71.75\,\pm\,11.99$\\$[56.86,\ 86.65]$} & \makecell[c]{$39.81\,\pm\,8.28$\\$[29.53,\ 50.08]$} & \makecell[c]{$46.59\,\pm\,0.55$\\$[45.91,\ 47.27]$} \\
Graph Temporal & \makecell[c]{$19.27\,\pm\,2.03$\\$[16.75,\ 21.78]$} & \makecell[c]{$72.76\,\pm\,11.50$\\$[58.48,\ 87.04]$} & \makecell[c]{$40.55\,\pm\,9.26$\\$[29.05,\ 52.04]$} & \makecell[c]{$47.10\,\pm\,1.18$\\$[45.64,\ 48.57]$} \\
Sage-Voronoi & \makecell[c]{$15.20\,\pm\,2.02$\\$[12.70,\ 17.70]$} & \makecell[c]{$65.63\,\pm\,9.92$\\$[53.31,\ 77.95]$} & \makecell[c]{$39.58\,\pm\,6.36$\\$[31.68,\ 47.49]$} & \makecell[c]{$44.42\,\pm\,2.33$\\$[41.52,\ 47.31]$} \\
GAIN & \makecell[c]{$13.55\,\pm\,4.89$\\$[7.47,\ 19.62]$} & \makecell[c]{$76.18\,\pm\,11.78$\\$[61.55,\ 90.81]$} & \makecell[c]{$34.49\,\pm\,9.89$\\$[22.22,\ 46.77]$} & \makecell[c]{$42.25\,\pm\,1.02$\\$[40.99,\ 43.52]$} \\
GRIN & \makecell[c]{$9.59\,\pm\,1.62$\\$[7.57,\ 11.60]$} & \makecell[c]{$62.63\,\pm\,11.54$\\$[48.31,\ 76.96]$} & \makecell[c]{$28.80\,\pm\,8.29$\\$[18.51,\ 39.09]$} & \makecell[c]{$33.10\,\pm\,0.97$\\$[31.90,\ 34.30]$} \\
\bottomrule
\end{tabular}
\endgroup
\end{table}

\clearpage

\begin{table}[htbp]
\centering
\caption{{October--December 2022, $B=144$, RMSE by count and road-graph degree strata. Cells give mean $\pm$ sample SD and [95\% Student-$t$ CI] across five seed runs.}}
\label{tab:p03_A20}
\begingroup
\footnotesize
\setlength{\tabcolsep}{3pt}
\renewcommand{\arraystretch}{1.15}
\begin{tabular}{lcccc}
\toprule
Method / variant & Low count & High count & Degree 0--2 & Degree $\geq3$ \\
\midrule
Linear & \makecell[c]{$201.37\,\pm\,19.72$\\$[176.88,\ 225.86]$} & \makecell[c]{$247.61\,\pm\,26.46$\\$[214.75,\ 280.46]$} & \makecell[c]{$177.53\,\pm\,51.59$\\$[113.47,\ 241.58]$} & \makecell[c]{$227.74\,\pm\,9.88$\\$[215.47,\ 240.01]$} \\
Median & \makecell[c]{$172.09\,\pm\,2.29$\\$[169.25,\ 174.94]$} & \makecell[c]{$265.55\,\pm\,39.38$\\$[216.65,\ 314.44]$} & \makecell[c]{$166.32\,\pm\,45.07$\\$[110.36,\ 222.29]$} & \makecell[c]{$183.57\,\pm\,2.43$\\$[180.56,\ 186.58]$} \\
FFill + BFill & \makecell[c]{$200.29\,\pm\,20.43$\\$[174.92,\ 225.65]$} & \makecell[c]{$252.74\,\pm\,41.51$\\$[201.20,\ 304.28]$} & \makecell[c]{$175.55\,\pm\,51.39$\\$[111.74,\ 239.36]$} & \makecell[c]{$230.87\,\pm\,12.86$\\$[214.91,\ 246.84]$} \\
Rolling mean & \makecell[c]{$201.31\,\pm\,19.57$\\$[177.01,\ 225.61]$} & \makecell[c]{$247.57\,\pm\,26.45$\\$[214.73,\ 280.41]$} & \makecell[c]{$177.34\,\pm\,51.48$\\$[113.43,\ 241.26]$} & \makecell[c]{$227.76\,\pm\,9.80$\\$[215.60,\ 239.93]$} \\
KNN (scaled) & \makecell[c]{$78.86\,\pm\,9.73$\\$[66.78,\ 90.95]$} & \makecell[c]{$180.42\,\pm\,25.00$\\$[149.38,\ 211.47]$} & \makecell[c]{$117.55\,\pm\,17.43$\\$[95.90,\ 139.19]$} & \makecell[c]{$109.13\,\pm\,8.66$\\$[98.38,\ 119.88]$} \\
KNN (unscaled) & \makecell[c]{$90.05\,\pm\,15.64$\\$[70.64,\ 109.47]$} & \makecell[c]{$181.02\,\pm\,26.38$\\$[148.27,\ 213.77]$} & \makecell[c]{$124.32\,\pm\,17.05$\\$[103.14,\ 145.49]$} & \makecell[c]{$115.99\,\pm\,10.58$\\$[102.86,\ 129.12]$} \\
Historical & \makecell[c]{$91.73\,\pm\,9.89$\\$[79.45,\ 104.01]$} & \makecell[c]{$177.91\,\pm\,34.10$\\$[135.56,\ 220.25]$} & \makecell[c]{$120.76\,\pm\,27.52$\\$[86.59,\ 154.94]$} & \makecell[c]{$121.05\,\pm\,6.98$\\$[112.38,\ 129.72]$} \\
Graph Laplacian & \makecell[c]{$43.37\,\pm\,5.44$\\$[36.62,\ 50.12]$} & \makecell[c]{$218.53\,\pm\,54.97$\\$[150.28,\ 286.78]$} & \makecell[c]{$161.35\,\pm\,52.35$\\$[96.34,\ 226.36]$} & \makecell[c]{$112.57\,\pm\,8.94$\\$[101.48,\ 123.67]$} \\
Graph Temporal & \makecell[c]{$45.16\,\pm\,8.38$\\$[34.75,\ 55.57]$} & \makecell[c]{$216.62\,\pm\,52.66$\\$[151.23,\ 282.00]$} & \makecell[c]{$159.28\,\pm\,50.83$\\$[96.16,\ 222.40]$} & \makecell[c]{$112.78\,\pm\,10.46$\\$[99.80,\ 125.77]$} \\
Sage-Voronoi & \makecell[c]{$46.08\,\pm\,12.17$\\$[30.97,\ 61.18]$} & \makecell[c]{$201.44\,\pm\,47.25$\\$[142.77,\ 260.11]$} & \makecell[c]{$157.53\,\pm\,41.84$\\$[105.57,\ 209.48]$} & \makecell[c]{$106.81\,\pm\,10.25$\\$[94.09,\ 119.54]$} \\
GAIN & \makecell[c]{$42.53\,\pm\,29.89$\\$[5.42,\ 79.65]$} & \makecell[c]{$222.35\,\pm\,44.50$\\$[167.09,\ 277.61]$} & \makecell[c]{$135.57\,\pm\,58.92$\\$[62.41,\ 208.74]$} & \makecell[c]{$113.95\,\pm\,5.36$\\$[107.29,\ 120.61]$} \\
GRIN & \makecell[c]{$28.00\,\pm\,5.07$\\$[21.71,\ 34.30]$} & \makecell[c]{$186.91\,\pm\,40.22$\\$[136.96,\ 236.85]$} & \makecell[c]{$118.72\,\pm\,45.80$\\$[61.85,\ 175.59]$} & \makecell[c]{$90.23\,\pm\,11.46$\\$[76.01,\ 104.46]$} \\
\bottomrule
\end{tabular}
\endgroup
\end{table}

\clearpage

\begin{table}[htbp]
\centering
\caption{{October--December 2022, $B=144$, MRE by count and road-graph degree strata. Cells give mean $\pm$ sample SD and [95\% Student-$t$ CI] across five seed runs.}}
\label{tab:p03_A21}
\begingroup
\footnotesize
\setlength{\tabcolsep}{3pt}
\renewcommand{\arraystretch}{1.15}
\begin{tabular}{lcccc}
\toprule
Method / variant & Low count & High count & Degree 0--2 & Degree $\geq3$ \\
\midrule
Linear & \makecell[c]{$8.470\,\pm\,0.543$\\$[7.796,\ 9.144]$} & \makecell[c]{$0.455\,\pm\,0.021$\\$[0.429,\ 0.481]$} & \makecell[c]{$2.634\,\pm\,0.336$\\$[2.217,\ 3.051]$} & \makecell[c]{$2.508\,\pm\,0.122$\\$[2.356,\ 2.659]$} \\
Median & \makecell[c]{$8.770\,\pm\,0.301$\\$[8.397,\ 9.143]$} & \makecell[c]{$0.513\,\pm\,0.007$\\$[0.505,\ 0.521]$} & \makecell[c]{$2.428\,\pm\,0.100$\\$[2.304,\ 2.551]$} & \makecell[c]{$2.460\,\pm\,0.049$\\$[2.400,\ 2.521]$} \\
FFill + BFill & \makecell[c]{$8.232\,\pm\,0.532$\\$[7.571,\ 8.893]$} & \makecell[c]{$0.467\,\pm\,0.023$\\$[0.438,\ 0.496]$} & \makecell[c]{$2.557\,\pm\,0.309$\\$[2.173,\ 2.941]$} & \makecell[c]{$2.467\,\pm\,0.141$\\$[2.291,\ 2.642]$} \\
Rolling mean & \makecell[c]{$8.463\,\pm\,0.543$\\$[7.789,\ 9.138]$} & \makecell[c]{$0.455\,\pm\,0.021$\\$[0.429,\ 0.481]$} & \makecell[c]{$2.631\,\pm\,0.338$\\$[2.212,\ 3.050]$} & \makecell[c]{$2.507\,\pm\,0.122$\\$[2.356,\ 2.658]$} \\
KNN (scaled) & \makecell[c]{$1.391\,\pm\,0.541$\\$[0.719,\ 2.063]$} & \makecell[c]{$0.197\,\pm\,0.010$\\$[0.184,\ 0.210]$} & \makecell[c]{$0.542\,\pm\,0.159$\\$[0.345,\ 0.739]$} & \makecell[c]{$0.463\,\pm\,0.118$\\$[0.316,\ 0.610]$} \\
KNN (unscaled) & \makecell[c]{$1.677\,\pm\,0.710$\\$[0.794,\ 2.559]$} & \makecell[c]{$0.217\,\pm\,0.013$\\$[0.201,\ 0.233]$} & \makecell[c]{$0.619\,\pm\,0.149$\\$[0.434,\ 0.803]$} & \makecell[c]{$0.571\,\pm\,0.199$\\$[0.323,\ 0.818]$} \\
Historical & \makecell[c]{$2.071\,\pm\,0.223$\\$[1.794,\ 2.347]$} & \makecell[c]{$0.228\,\pm\,0.008$\\$[0.218,\ 0.238]$} & \makecell[c]{$0.735\,\pm\,0.156$\\$[0.542,\ 0.929]$} & \makecell[c]{$0.756\,\pm\,0.031$\\$[0.717,\ 0.795]$} \\
Graph Laplacian & \makecell[c]{$1.617\,\pm\,0.278$\\$[1.272,\ 1.962]$} & \makecell[c]{$0.217\,\pm\,0.013$\\$[0.201,\ 0.233]$} & \makecell[c]{$0.772\,\pm\,0.199$\\$[0.525,\ 1.019]$} & \makecell[c]{$0.491\,\pm\,0.009$\\$[0.479,\ 0.502]$} \\
Graph Temporal & \makecell[c]{$1.708\,\pm\,0.355$\\$[1.268,\ 2.149]$} & \makecell[c]{$0.230\,\pm\,0.010$\\$[0.217,\ 0.242]$} & \makecell[c]{$0.820\,\pm\,0.269$\\$[0.486,\ 1.154]$} & \makecell[c]{$0.514\,\pm\,0.013$\\$[0.498,\ 0.529]$} \\
Sage-Voronoi & \makecell[c]{$0.953\,\pm\,0.139$\\$[0.781,\ 1.125]$} & \makecell[c]{$0.213\,\pm\,0.012$\\$[0.198,\ 0.229]$} & \makecell[c]{$0.589\,\pm\,0.129$\\$[0.429,\ 0.749]$} & \makecell[c]{$0.361\,\pm\,0.025$\\$[0.329,\ 0.392]$} \\
GAIN & \makecell[c]{$1.128\,\pm\,0.804$\\$[0.129,\ 2.126]$} & \makecell[c]{$0.229\,\pm\,0.009$\\$[0.218,\ 0.240]$} & \makecell[c]{$0.508\,\pm\,0.170$\\$[0.297,\ 0.720]$} & \makecell[c]{$0.405\,\pm\,0.240$\\$[0.106,\ 0.703]$} \\
GRIN & \makecell[c]{$0.776\,\pm\,0.219$\\$[0.504,\ 1.049]$} & \makecell[c]{$0.205\,\pm\,0.018$\\$[0.183,\ 0.227]$} & \makecell[c]{$0.485\,\pm\,0.167$\\$[0.278,\ 0.692]$} & \makecell[c]{$0.260\,\pm\,0.009$\\$[0.249,\ 0.272]$} \\
\bottomrule
\end{tabular}
\endgroup
\end{table}

\clearpage

\begin{table}[htbp]
\centering
\caption{{October--December 2022, MCAR 30\%, MAE by count and road-graph degree strata. Cells give mean $\pm$ sample SD and [95\% Student-$t$ CI] across five seed runs.}}
\label{tab:p03_A22}
\begingroup
\footnotesize
\setlength{\tabcolsep}{3pt}
\renewcommand{\arraystretch}{1.15}
\begin{tabular}{lcccc}
\toprule
Method / variant & Low count & High count & Degree 0--2 & Degree $\geq3$ \\
\midrule
Linear & \makecell[c]{$7.95\,\pm\,0.15$\\$[7.76,\ 8.14]$} & \makecell[c]{$32.14\,\pm\,0.40$\\$[31.65,\ 32.63]$} & \makecell[c]{$17.93\,\pm\,0.13$\\$[17.76,\ 18.10]$} & \makecell[c]{$22.63\,\pm\,0.19$\\$[22.39,\ 22.87]$} \\
Median & \makecell[c]{$113.98\,\pm\,0.47$\\$[113.40,\ 114.56]$} & \makecell[c]{$149.75\,\pm\,0.91$\\$[148.63,\ 150.88]$} & \makecell[c]{$71.92\,\pm\,0.45$\\$[71.37,\ 72.48]$} & \makecell[c]{$119.65\,\pm\,0.60$\\$[118.91,\ 120.39]$} \\
FFill + BFill & \makecell[c]{$9.58\,\pm\,0.13$\\$[9.42,\ 9.73]$} & \makecell[c]{$39.06\,\pm\,0.63$\\$[38.27,\ 39.84]$} & \makecell[c]{$22.01\,\pm\,0.25$\\$[21.70,\ 22.32]$} & \makecell[c]{$28.14\,\pm\,0.22$\\$[27.87,\ 28.42]$} \\
Rolling mean & \makecell[c]{$10.72\,\pm\,0.21$\\$[10.46,\ 10.98]$} & \makecell[c]{$38.05\,\pm\,0.60$\\$[37.31,\ 38.79]$} & \makecell[c]{$21.73\,\pm\,0.18$\\$[21.50,\ 21.96]$} & \makecell[c]{$29.06\,\pm\,0.25$\\$[28.74,\ 29.37]$} \\
KNN (scaled) & \makecell[c]{$18.56\,\pm\,0.28$\\$[18.21,\ 18.90]$} & \makecell[c]{$66.98\,\pm\,0.96$\\$[65.79,\ 68.16]$} & \makecell[c]{$30.15\,\pm\,0.30$\\$[29.77,\ 30.52]$} & \makecell[c]{$40.71\,\pm\,0.44$\\$[40.16,\ 41.26]$} \\
KNN (unscaled) & \makecell[c]{$20.01\,\pm\,0.48$\\$[19.41,\ 20.62]$} & \makecell[c]{$68.11\,\pm\,0.84$\\$[67.07,\ 69.15]$} & \makecell[c]{$33.14\,\pm\,0.33$\\$[32.73,\ 33.55]$} & \makecell[c]{$44.04\,\pm\,0.35$\\$[43.61,\ 44.47]$} \\
Historical & \makecell[c]{$32.73\,\pm\,0.26$\\$[32.41,\ 33.06]$} & \makecell[c]{$76.66\,\pm\,0.95$\\$[75.48,\ 77.84]$} & \makecell[c]{$40.27\,\pm\,0.14$\\$[40.10,\ 40.45]$} & \makecell[c]{$63.00\,\pm\,0.45$\\$[62.44,\ 63.56]$} \\
Graph Laplacian & \makecell[c]{$19.77\,\pm\,0.42$\\$[19.25,\ 20.29]$} & \makecell[c]{$74.07\,\pm\,0.90$\\$[72.95,\ 75.18]$} & \makecell[c]{$38.90\,\pm\,0.46$\\$[38.32,\ 39.47]$} & \makecell[c]{$46.96\,\pm\,0.36$\\$[46.51,\ 47.41]$} \\
Graph Temporal & \makecell[c]{$14.98\,\pm\,0.10$\\$[14.86,\ 15.10]$} & \makecell[c]{$57.67\,\pm\,0.89$\\$[56.56,\ 58.77]$} & \makecell[c]{$26.79\,\pm\,0.20$\\$[26.55,\ 27.03]$} & \makecell[c]{$40.44\,\pm\,0.35$\\$[40.01,\ 40.87]$} \\
Sage-Voronoi & \makecell[c]{$8.99\,\pm\,0.26$\\$[8.67,\ 9.31]$} & \makecell[c]{$35.26\,\pm\,0.68$\\$[34.41,\ 36.10]$} & \makecell[c]{$20.07\,\pm\,0.27$\\$[19.74,\ 20.40]$} & \makecell[c]{$24.84\,\pm\,0.20$\\$[24.58,\ 25.09]$} \\
GAIN & \makecell[c]{$20.97\,\pm\,0.88$\\$[19.88,\ 22.05]$} & \makecell[c]{$71.29\,\pm\,1.15$\\$[69.86,\ 72.72]$} & \makecell[c]{$33.43\,\pm\,0.39$\\$[32.95,\ 33.90]$} & \makecell[c]{$46.40\,\pm\,0.49$\\$[45.79,\ 47.00]$} \\
GRIN & \makecell[c]{$6.15\,\pm\,0.08$\\$[6.05,\ 6.25]$} & \makecell[c]{$30.57\,\pm\,0.45$\\$[30.01,\ 31.14]$} & \makecell[c]{$15.98\,\pm\,0.08$\\$[15.88,\ 16.08]$} & \makecell[c]{$18.98\,\pm\,0.15$\\$[18.79,\ 19.16]$} \\
\bottomrule
\end{tabular}
\endgroup
\end{table}

\clearpage

\begin{table}[htbp]
\centering
\caption{{October--December 2022, MCAR 30\%, RMSE by count and road-graph degree strata. Cells give mean $\pm$ sample SD and [95\% Student-$t$ CI] across five seed runs.}}
\label{tab:p03_A23}
\begingroup
\footnotesize
\setlength{\tabcolsep}{3pt}
\renewcommand{\arraystretch}{1.15}
\begin{tabular}{lcccc}
\toprule
Method / variant & Low count & High count & Degree 0--2 & Degree $\geq3$ \\
\midrule
Linear & \makecell[c]{$34.25\,\pm\,1.89$\\$[31.91,\ 36.59]$} & \makecell[c]{$93.76\,\pm\,1.74$\\$[91.60,\ 95.92]$} & \makecell[c]{$68.46\,\pm\,0.87$\\$[67.38,\ 69.55]$} & \makecell[c]{$56.02\,\pm\,1.42$\\$[54.26,\ 57.78]$} \\
Median & \makecell[c]{$173.61\,\pm\,0.75$\\$[172.68,\ 174.55]$} & \makecell[c]{$282.24\,\pm\,2.06$\\$[279.68,\ 284.80]$} & \makecell[c]{$163.64\,\pm\,2.28$\\$[160.81,\ 166.47]$} & \makecell[c]{$193.95\,\pm\,2.15$\\$[191.28,\ 196.62]$} \\
FFill + BFill & \makecell[c]{$46.65\,\pm\,2.00$\\$[44.16,\ 49.13]$} & \makecell[c]{$115.36\,\pm\,3.12$\\$[111.49,\ 119.23]$} & \makecell[c]{$86.76\,\pm\,2.12$\\$[84.13,\ 89.39]$} & \makecell[c]{$71.86\,\pm\,1.65$\\$[69.81,\ 73.91]$} \\
Rolling mean & \makecell[c]{$52.34\,\pm\,2.53$\\$[49.20,\ 55.48]$} & \makecell[c]{$114.38\,\pm\,2.91$\\$[110.77,\ 117.99]$} & \makecell[c]{$83.97\,\pm\,1.74$\\$[81.81,\ 86.13]$} & \makecell[c]{$74.43\,\pm\,1.63$\\$[72.41,\ 76.45]$} \\
KNN (scaled) & \makecell[c]{$86.30\,\pm\,1.87$\\$[83.97,\ 88.62]$} & \makecell[c]{$206.56\,\pm\,3.51$\\$[202.20,\ 210.92]$} & \makecell[c]{$118.04\,\pm\,1.91$\\$[115.68,\ 120.41]$} & \makecell[c]{$126.80\,\pm\,2.97$\\$[123.11,\ 130.48]$} \\
KNN (unscaled) & \makecell[c]{$91.76\,\pm\,3.08$\\$[87.94,\ 95.58]$} & \makecell[c]{$209.12\,\pm\,2.78$\\$[205.67,\ 212.57]$} & \makecell[c]{$123.48\,\pm\,2.30$\\$[120.63,\ 126.33]$} & \makecell[c]{$131.30\,\pm\,2.63$\\$[128.03,\ 134.57]$} \\
Historical & \makecell[c]{$95.28\,\pm\,1.03$\\$[94.00,\ 96.56]$} & \makecell[c]{$203.19\,\pm\,3.35$\\$[199.02,\ 207.35]$} & \makecell[c]{$119.75\,\pm\,1.28$\\$[118.17,\ 121.34]$} & \makecell[c]{$139.93\,\pm\,2.54$\\$[136.78,\ 143.09]$} \\
Graph Laplacian & \makecell[c]{$41.23\,\pm\,0.86$\\$[40.16,\ 42.29]$} & \makecell[c]{$240.15\,\pm\,2.09$\\$[237.55,\ 242.75]$} & \makecell[c]{$164.32\,\pm\,2.82$\\$[160.82,\ 167.82]$} & \makecell[c]{$126.07\,\pm\,2.54$\\$[122.92,\ 129.22]$} \\
Graph Temporal & \makecell[c]{$34.95\,\pm\,0.78$\\$[33.98,\ 35.92]$} & \makecell[c]{$173.04\,\pm\,4.39$\\$[167.59,\ 178.49]$} & \makecell[c]{$100.09\,\pm\,1.56$\\$[98.16,\ 102.02]$} & \makecell[c]{$106.58\,\pm\,2.99$\\$[102.86,\ 110.30]$} \\
Sage-Voronoi & \makecell[c]{$37.12\,\pm\,1.96$\\$[34.69,\ 39.56]$} & \makecell[c]{$114.52\,\pm\,3.43$\\$[110.27,\ 118.78]$} & \makecell[c]{$79.57\,\pm\,1.77$\\$[77.37,\ 81.76]$} & \makecell[c]{$65.59\,\pm\,1.96$\\$[63.15,\ 68.03]$} \\
GAIN & \makecell[c]{$57.76\,\pm\,3.85$\\$[52.98,\ 62.54]$} & \makecell[c]{$224.94\,\pm\,3.36$\\$[220.77,\ 229.12]$} & \makecell[c]{$128.78\,\pm\,3.24$\\$[124.75,\ 132.81]$} & \makecell[c]{$126.19\,\pm\,2.43$\\$[123.17,\ 129.21]$} \\
GRIN & \makecell[c]{$27.84\,\pm\,1.55$\\$[25.92,\ 29.76]$} & \makecell[c]{$97.91\,\pm\,2.98$\\$[94.21,\ 101.60]$} & \makecell[c]{$68.01\,\pm\,0.73$\\$[67.10,\ 68.92]$} & \makecell[c]{$51.53\,\pm\,1.70$\\$[49.43,\ 53.64]$} \\
\bottomrule
\end{tabular}
\endgroup
\end{table}

\clearpage

\begin{table}[htbp]
\centering
\caption{{October--December 2022, MCAR 30\%, MRE by count and road-graph degree strata. Cells give mean $\pm$ sample SD and [95\% Student-$t$ CI] across five seed runs.}}
\label{tab:p03_A24}
\begingroup
\footnotesize
\setlength{\tabcolsep}{3pt}
\renewcommand{\arraystretch}{1.15}
\begin{tabular}{lcccc}
\toprule
Method / variant & Low count & High count & Degree 0--2 & Degree $\geq3$ \\
\midrule
Linear & \makecell[c]{$0.511\,\pm\,0.005$\\$[0.505,\ 0.517]$} & \makecell[c]{$0.169\,\pm\,0.001$\\$[0.167,\ 0.170]$} & \makecell[c]{$0.302\,\pm\,0.004$\\$[0.297,\ 0.307]$} & \makecell[c]{$0.229\,\pm\,0.003$\\$[0.225,\ 0.233]$} \\
Median & \makecell[c]{$8.597\,\pm\,0.036$\\$[8.553,\ 8.642]$} & \makecell[c]{$0.515\,\pm\,0.001$\\$[0.514,\ 0.516]$} & \makecell[c]{$2.452\,\pm\,0.002$\\$[2.450,\ 2.455]$} & \makecell[c]{$2.350\,\pm\,0.020$\\$[2.326,\ 2.374]$} \\
FFill + BFill & \makecell[c]{$0.589\,\pm\,0.007$\\$[0.580,\ 0.597]$} & \makecell[c]{$0.201\,\pm\,0.003$\\$[0.197,\ 0.204]$} & \makecell[c]{$0.352\,\pm\,0.005$\\$[0.347,\ 0.358]$} & \makecell[c]{$0.273\,\pm\,0.004$\\$[0.269,\ 0.278]$} \\
Rolling mean & \makecell[c]{$0.592\,\pm\,0.008$\\$[0.582,\ 0.602]$} & \makecell[c]{$0.179\,\pm\,0.002$\\$[0.177,\ 0.181]$} & \makecell[c]{$0.329\,\pm\,0.004$\\$[0.324,\ 0.333]$} & \makecell[c]{$0.276\,\pm\,0.004$\\$[0.271,\ 0.280]$} \\
KNN (scaled) & \makecell[c]{$1.255\,\pm\,0.033$\\$[1.214,\ 1.295]$} & \makecell[c]{$0.204\,\pm\,0.001$\\$[0.202,\ 0.205]$} & \makecell[c]{$0.481\,\pm\,0.010$\\$[0.468,\ 0.493]$} & \makecell[c]{$0.444\,\pm\,0.015$\\$[0.425,\ 0.463]$} \\
KNN (unscaled) & \makecell[c]{$1.404\,\pm\,0.058$\\$[1.332,\ 1.476]$} & \makecell[c]{$0.224\,\pm\,0.002$\\$[0.222,\ 0.226]$} & \makecell[c]{$0.548\,\pm\,0.012$\\$[0.533,\ 0.563]$} & \makecell[c]{$0.503\,\pm\,0.021$\\$[0.477,\ 0.529]$} \\
Historical & \makecell[c]{$1.904\,\pm\,0.019$\\$[1.880,\ 1.928]$} & \makecell[c]{$0.230\,\pm\,0.001$\\$[0.228,\ 0.231]$} & \makecell[c]{$0.707\,\pm\,0.008$\\$[0.697,\ 0.717]$} & \makecell[c]{$0.718\,\pm\,0.010$\\$[0.706,\ 0.730]$} \\
Graph Laplacian & \makecell[c]{$1.469\,\pm\,0.014$\\$[1.451,\ 1.486]$} & \makecell[c]{$0.223\,\pm\,0.003$\\$[0.219,\ 0.226]$} & \makecell[c]{$0.674\,\pm\,0.009$\\$[0.664,\ 0.685]$} & \makecell[c]{$0.481\,\pm\,0.004$\\$[0.476,\ 0.486]$} \\
Graph Temporal & \makecell[c]{$1.185\,\pm\,0.008$\\$[1.175,\ 1.196]$} & \makecell[c]{$0.194\,\pm\,0.001$\\$[0.193,\ 0.196]$} & \makecell[c]{$0.516\,\pm\,0.003$\\$[0.513,\ 0.520]$} & \makecell[c]{$0.425\,\pm\,0.004$\\$[0.421,\ 0.430]$} \\
Sage-Voronoi & \makecell[c]{$0.518\,\pm\,0.020$\\$[0.494,\ 0.542]$} & \makecell[c]{$0.146\,\pm\,0.002$\\$[0.144,\ 0.148]$} & \makecell[c]{$0.292\,\pm\,0.004$\\$[0.287,\ 0.297]$} & \makecell[c]{$0.235\,\pm\,0.006$\\$[0.228,\ 0.243]$} \\
GAIN & \makecell[c]{$1.489\,\pm\,0.073$\\$[1.399,\ 1.579]$} & \makecell[c]{$0.212\,\pm\,0.001$\\$[0.210,\ 0.214]$} & \makecell[c]{$0.616\,\pm\,0.017$\\$[0.595,\ 0.637]$} & \makecell[c]{$0.510\,\pm\,0.024$\\$[0.481,\ 0.539]$} \\
GRIN & \makecell[c]{$0.443\,\pm\,0.010$\\$[0.430,\ 0.456]$} & \makecell[c]{$0.142\,\pm\,0.009$\\$[0.131,\ 0.153]$} & \makecell[c]{$0.272\,\pm\,0.004$\\$[0.267,\ 0.277]$} & \makecell[c]{$0.179\,\pm\,0.003$\\$[0.176,\ 0.183]$} \\
\bottomrule
\end{tabular}
\endgroup
\end{table}

\clearpage

\begin{table}[htbp]
\centering
\caption{{January 2023, MAE of reconstructed daily sensor-count sums on affected test days. Cells give mean $\pm$ sample SD and [95\% Student-$t$ CI] across five seed runs.}}
\label{tab:p03_A25}
\begingroup
\footnotesize
\setlength{\tabcolsep}{3pt}
\renewcommand{\arraystretch}{1.15}
\begin{tabular}{lccc}
\toprule
Method / variant & $B=144$ & MCAR 30\% & Synchronous \\
\midrule
Linear & \makecell[c]{$72410.75\,\pm\,12109.05$\\$[57375.38,\ 87446.11]$} & \makecell[c]{$6150.56\,\pm\,2243.87$\\$[3364.43,\ 8936.69]$} & \makecell[c]{$980293.78\,\pm\,377426.58$\\$[511656.79,\ 1448930.77]$} \\
Median & \makecell[c]{$90035.00\,\pm\,16381.39$\\$[69694.82,\ 110375.18]$} & \makecell[c]{$231525.20\,\pm\,2903.38$\\$[227920.17,\ 235130.23]$} & \makecell[c]{$611806.00\,\pm\,163725.07$\\$[408514.47,\ 815097.53]$} \\
FFill + BFill & \makecell[c]{$85656.88\,\pm\,18957.11$\\$[62118.52,\ 109195.24]$} & \makecell[c]{$8610.08\,\pm\,1657.37$\\$[6552.18,\ 10667.98]$} & \makecell[c]{$972934.90\,\pm\,366138.17$\\$[518314.31,\ 1427555.49]$} \\
Rolling mean & \makecell[c]{$72274.58\,\pm\,12233.86$\\$[57084.23,\ 87464.93]$} & \makecell[c]{$9196.60\,\pm\,2676.48$\\$[5873.31,\ 12519.90]$} & \makecell[c]{$980797.96\,\pm\,372104.43$\\$[518769.29,\ 1442826.64]$} \\
KNN (scaled) & \makecell[c]{$190782.77\,\pm\,15735.16$\\$[171244.99,\ 210320.55]$} & \makecell[c]{$446441.30\,\pm\,10826.67$\\$[432998.21,\ 459884.38]$} & \makecell[c]{$589644.89\,\pm\,154911.41$\\$[397296.97,\ 781992.81]$} \\
KNN (unscaled) & \makecell[c]{$31607.26\,\pm\,9718.70$\\$[19539.90,\ 43674.62]$} & \makecell[c]{$15783.40\,\pm\,2244.63$\\$[12996.33,\ 18570.47]$} & \makecell[c]{$589644.89\,\pm\,154911.41$\\$[397296.97,\ 781992.81]$} \\
Historical & \makecell[c]{$73200.83\,\pm\,15691.08$\\$[53717.78,\ 92683.88]$} & \makecell[c]{$198564.57\,\pm\,2480.62$\\$[195484.47,\ 201644.66]$} & \makecell[c]{$451834.57\,\pm\,100964.69$\\$[326470.35,\ 577198.80]$} \\
Graph Laplacian & \makecell[c]{$85428.42\,\pm\,8061.82$\\$[75418.35,\ 95438.48]$} & \makecell[c]{$212077.53\,\pm\,5574.20$\\$[205156.24,\ 218998.82]$} & \makecell[c]{$589644.89\,\pm\,154911.41$\\$[397296.97,\ 781992.81]$} \\
Graph Temporal & \makecell[c]{$80436.01\,\pm\,15084.37$\\$[61706.29,\ 99165.74]$} & \makecell[c]{$122172.38\,\pm\,2556.79$\\$[118997.70,\ 125347.05]$} & \makecell[c]{$1453677.99\,\pm\,738136.06$\\$[537161.02,\ 2370194.96]$} \\
Sage-Voronoi & \makecell[c]{$54436.33\,\pm\,22194.49$\\$[26878.23,\ 81994.43]$} & \makecell[c]{$16735.64\,\pm\,5408.87$\\$[10019.65,\ 23451.63]$} & \makecell[c]{$945616.77\,\pm\,491318.39$\\$[335564.41,\ 1555669.12]$} \\
GAIN & \makecell[c]{$39654.53\,\pm\,7559.95$\\$[30267.62,\ 49041.45]$} & \makecell[c]{$44242.60\,\pm\,7305.22$\\$[35171.98,\ 53313.23]$} & \makecell[c]{$962114.41\,\pm\,465157.59$\\$[384544.98,\ 1539683.84]$} \\
GRIN & \makecell[c]{$31783.75\,\pm\,6919.57$\\$[23191.97,\ 40375.53]$} & \makecell[c]{$42622.32\,\pm\,7061.68$\\$[33854.09,\ 51390.56]$} & \makecell[c]{$839838.44\,\pm\,321189.96$\\$[441028.42,\ 1238648.45]$} \\
\bottomrule
\end{tabular}
\endgroup
\end{table}

\clearpage

\begin{table}[htbp]
\centering
\caption{{January 2023, RMSE of reconstructed daily sensor-count sums on affected test days. Cells give mean $\pm$ sample SD and [95\% Student-$t$ CI] across five seed runs.}}
\label{tab:p03_A26}
\begingroup
\footnotesize
\setlength{\tabcolsep}{3pt}
\renewcommand{\arraystretch}{1.15}
\begin{tabular}{lccc}
\toprule
Method / variant & $B=144$ & MCAR 30\% & Synchronous \\
\midrule
Linear & \makecell[c]{$93735.94\,\pm\,25103.74$\\$[62565.54,\ 124906.35]$} & \makecell[c]{$7565.38\,\pm\,2288.36$\\$[4724.00,\ 10406.75]$} & \makecell[c]{$1060912.23\,\pm\,344025.04$\\$[633748.73,\ 1488075.73]$} \\
Median & \makecell[c]{$109975.87\,\pm\,22153.24$\\$[82468.99,\ 137482.76]$} & \makecell[c]{$252456.67\,\pm\,3676.35$\\$[247891.87,\ 257021.46]$} & \makecell[c]{$691354.84\,\pm\,182989.29$\\$[464143.63,\ 918566.06]$} \\
FFill + BFill & \makecell[c]{$114645.91\,\pm\,39533.28$\\$[65558.86,\ 163732.96]$} & \makecell[c]{$10368.82\,\pm\,1634.50$\\$[8339.31,\ 12398.32]$} & \makecell[c]{$1113379.52\,\pm\,313875.24$\\$[723651.93,\ 1503107.11]$} \\
Rolling mean & \makecell[c]{$93662.62\,\pm\,25340.16$\\$[62198.66,\ 125126.59]$} & \makecell[c]{$10891.94\,\pm\,2575.31$\\$[7694.27,\ 14089.61]$} & \makecell[c]{$1062969.39\,\pm\,336871.97$\\$[644687.59,\ 1481251.19]$} \\
KNN (scaled) & \makecell[c]{$228627.08\,\pm\,20218.51$\\$[203522.48,\ 253731.68]$} & \makecell[c]{$485432.27\,\pm\,10598.40$\\$[472272.61,\ 498591.92]$} & \makecell[c]{$654996.70\,\pm\,160460.72$\\$[455758.40,\ 854235.01]$} \\
KNN (unscaled) & \makecell[c]{$45674.69\,\pm\,15208.28$\\$[26791.12,\ 64558.26]$} & \makecell[c]{$20444.83\,\pm\,3757.69$\\$[15779.04,\ 25110.62]$} & \makecell[c]{$654996.70\,\pm\,160460.72$\\$[455758.40,\ 854235.01]$} \\
Historical & \makecell[c]{$90330.33\,\pm\,22764.42$\\$[62064.57,\ 118596.08]$} & \makecell[c]{$216997.29\,\pm\,2943.84$\\$[213342.03,\ 220652.54]$} & \makecell[c]{$517418.71\,\pm\,166511.62$\\$[310667.23,\ 724170.19]$} \\
Graph Laplacian & \makecell[c]{$111593.30\,\pm\,12648.95$\\$[95887.56,\ 127299.05]$} & \makecell[c]{$221781.34\,\pm\,6144.90$\\$[214151.44,\ 229411.24]$} & \makecell[c]{$654996.70\,\pm\,160460.72$\\$[455758.40,\ 854235.01]$} \\
Graph Temporal & \makecell[c]{$106584.35\,\pm\,16174.92$\\$[86500.54,\ 126668.17]$} & \makecell[c]{$128068.48\,\pm\,2290.52$\\$[125224.42,\ 130912.54]$} & \makecell[c]{$1519204.64\,\pm\,727950.07$\\$[615335.24,\ 2423074.04]$} \\
Sage-Voronoi & \makecell[c]{$76745.87\,\pm\,27492.24$\\$[42609.74,\ 110881.99]$} & \makecell[c]{$19582.00\,\pm\,5656.03$\\$[12559.11,\ 26604.89]$} & \makecell[c]{$1077640.27\,\pm\,516068.67$\\$[436856.39,\ 1718424.16]$} \\
GAIN & \makecell[c]{$56597.96\,\pm\,11846.24$\\$[41888.92,\ 71307.01]$} & \makecell[c]{$51489.16\,\pm\,8047.30$\\$[41497.12,\ 61481.21]$} & \makecell[c]{$1056563.77\,\pm\,453209.88$\\$[493829.38,\ 1619298.16]$} \\
GRIN & \makecell[c]{$45683.62\,\pm\,10595.07$\\$[32528.10,\ 58839.14]$} & \makecell[c]{$46143.35\,\pm\,7370.48$\\$[36991.70,\ 55295.01]$} & \makecell[c]{$953917.01\,\pm\,384726.11$\\$[476216.45,\ 1431617.56]$} \\
\bottomrule
\end{tabular}
\endgroup
\end{table}

\clearpage

\begin{table}[htbp]
\centering
\caption{{January 2023, MRE of reconstructed daily sensor-count sums on affected test days. Cells give mean $\pm$ sample SD and [95\% Student-$t$ CI] across five seed runs.}}
\label{tab:p03_A27}
\begingroup
\footnotesize
\setlength{\tabcolsep}{3pt}
\renewcommand{\arraystretch}{1.15}
\begin{tabular}{lccc}
\toprule
Method / variant & $B=144$ & MCAR 30\% & Synchronous \\
\midrule
Linear & \makecell[c]{$0.019\,\pm\,0.004$\\$[0.013,\ 0.024]$} & \makecell[c]{$0.002\,\pm\,0.001$\\$[0.001,\ 0.003]$} & \makecell[c]{$0.225\,\pm\,0.076$\\$[0.131,\ 0.319]$} \\
Median & \makecell[c]{$0.022\,\pm\,0.004$\\$[0.017,\ 0.027]$} & \makecell[c]{$0.058\,\pm\,0.001$\\$[0.057,\ 0.059]$} & \makecell[c]{$0.140\,\pm\,0.033$\\$[0.098,\ 0.181]$} \\
FFill + BFill & \makecell[c]{$0.022\,\pm\,0.007$\\$[0.014,\ 0.031]$} & \makecell[c]{$0.003\,\pm\,0.001$\\$[0.002,\ 0.004]$} & \makecell[c]{$0.236\,\pm\,0.092$\\$[0.121,\ 0.350]$} \\
Rolling mean & \makecell[c]{$0.019\,\pm\,0.004$\\$[0.013,\ 0.024]$} & \makecell[c]{$0.003\,\pm\,0.001$\\$[0.002,\ 0.005]$} & \makecell[c]{$0.226\,\pm\,0.075$\\$[0.132,\ 0.319]$} \\
KNN (scaled) & \makecell[c]{$0.044\,\pm\,0.003$\\$[0.040,\ 0.048]$} & \makecell[c]{$0.108\,\pm\,0.003$\\$[0.104,\ 0.112]$} & \makecell[c]{$0.137\,\pm\,0.037$\\$[0.090,\ 0.183]$} \\
KNN (unscaled) & \makecell[c]{$0.008\,\pm\,0.003$\\$[0.004,\ 0.011]$} & \makecell[c]{$0.004\,\pm\,0.000$\\$[0.004,\ 0.005]$} & \makecell[c]{$0.137\,\pm\,0.037$\\$[0.090,\ 0.183]$} \\
Historical & \makecell[c]{$0.018\,\pm\,0.004$\\$[0.013,\ 0.023]$} & \makecell[c]{$0.055\,\pm\,0.000$\\$[0.054,\ 0.055]$} & \makecell[c]{$0.115\,\pm\,0.054$\\$[0.047,\ 0.182]$} \\
Graph Laplacian & \makecell[c]{$0.020\,\pm\,0.002$\\$[0.018,\ 0.023]$} & \makecell[c]{$0.054\,\pm\,0.001$\\$[0.052,\ 0.055]$} & \makecell[c]{$0.137\,\pm\,0.037$\\$[0.090,\ 0.183]$} \\
Graph Temporal & \makecell[c]{$0.019\,\pm\,0.004$\\$[0.014,\ 0.024]$} & \makecell[c]{$0.031\,\pm\,0.001$\\$[0.030,\ 0.032]$} & \makecell[c]{$0.322\,\pm\,0.153$\\$[0.132,\ 0.511]$} \\
Sage-Voronoi & \makecell[c]{$0.014\,\pm\,0.006$\\$[0.006,\ 0.022]$} & \makecell[c]{$0.006\,\pm\,0.002$\\$[0.004,\ 0.008]$} & \makecell[c]{$0.207\,\pm\,0.109$\\$[0.072,\ 0.343]$} \\
GAIN & \makecell[c]{$0.010\,\pm\,0.001$\\$[0.008,\ 0.011]$} & \makecell[c]{$0.011\,\pm\,0.002$\\$[0.009,\ 0.013]$} & \makecell[c]{$0.244\,\pm\,0.181$\\$[0.020,\ 0.469]$} \\
GRIN & \makecell[c]{$0.008\,\pm\,0.001$\\$[0.006,\ 0.009]$} & \makecell[c]{$0.011\,\pm\,0.002$\\$[0.009,\ 0.013]$} & \makecell[c]{$0.187\,\pm\,0.063$\\$[0.108,\ 0.265]$} \\
\bottomrule
\end{tabular}
\endgroup
\end{table}

\clearpage

\begin{table}[htbp]
\centering
\caption{{October--December 2022, MAE of reconstructed daily sensor-count sums on affected test days. Cells give mean $\pm$ sample SD and [95\% Student-$t$ CI] across five seed runs.}}
\label{tab:p03_A28}
\begingroup
\footnotesize
\setlength{\tabcolsep}{3pt}
\renewcommand{\arraystretch}{1.15}
\begin{tabular}{lccc}
\toprule
Method / variant & $B=144$ & MCAR 30\% & Synchronous \\
\midrule
Linear & \makecell[c]{$31034.47\,\pm\,8046.66$\\$[21043.22,\ 41025.72]$} & \makecell[c]{$4149.24\,\pm\,662.27$\\$[3326.93,\ 4971.56]$} & \makecell[c]{$928936.29\,\pm\,470488.66$\\$[344747.46,\ 1513125.12]$} \\
Median & \makecell[c]{$27838.31\,\pm\,4548.63$\\$[22190.44,\ 33486.17]$} & \makecell[c]{$193375.25\,\pm\,2431.43$\\$[190356.23,\ 196394.26]$} & \makecell[c]{$391730.10\,\pm\,291623.22$\\$[29632.05,\ 753828.15]$} \\
FFill + BFill & \makecell[c]{$32613.12\,\pm\,7468.78$\\$[23339.40,\ 41886.84]$} & \makecell[c]{$5690.29\,\pm\,1070.17$\\$[4361.50,\ 7019.08]$} & \makecell[c]{$964772.10\,\pm\,547238.63$\\$[285285.59,\ 1644258.61]$} \\
Rolling mean & \makecell[c]{$30995.58\,\pm\,7975.92$\\$[21092.16,\ 40898.99]$} & \makecell[c]{$5960.51\,\pm\,1020.42$\\$[4693.48,\ 7227.53]$} & \makecell[c]{$930489.38\,\pm\,470877.39$\\$[345817.88,\ 1515160.88]$} \\
KNN (scaled) & \makecell[c]{$12934.76\,\pm\,1929.70$\\$[10538.72,\ 15330.80]$} & \makecell[c]{$58304.09\,\pm\,1738.37$\\$[56145.61,\ 60462.56]$} & \makecell[c]{$395205.05\,\pm\,290778.10$\\$[34156.36,\ 756253.74]$} \\
KNN (unscaled) & \makecell[c]{$13671.03\,\pm\,2226.58$\\$[10906.36,\ 16435.70]$} & \makecell[c]{$33660.05\,\pm\,1392.31$\\$[31931.26,\ 35388.83]$} & \makecell[c]{$395205.05\,\pm\,290778.10$\\$[34156.36,\ 756253.74]$} \\
Historical & \makecell[c]{$25138.25\,\pm\,2588.35$\\$[21924.38,\ 28352.12]$} & \makecell[c]{$198701.36\,\pm\,1566.27$\\$[196756.58,\ 200646.14]$} & \makecell[c]{$337485.08\,\pm\,280791.67$\\$[-11163.83,\ 686133.98]$} \\
Graph Laplacian & \makecell[c]{$14703.42\,\pm\,3090.71$\\$[10865.79,\ 18541.05]$} & \makecell[c]{$27298.32\,\pm\,2212.11$\\$[24551.62,\ 30045.02]$} & \makecell[c]{$395205.05\,\pm\,290778.10$\\$[34156.36,\ 756253.74]$} \\
Graph Temporal & \makecell[c]{$15291.87\,\pm\,3150.52$\\$[11379.98,\ 19203.76]$} & \makecell[c]{$18969.44\,\pm\,2183.87$\\$[16257.82,\ 21681.07]$} & \makecell[c]{$877365.88\,\pm\,466034.82$\\$[298707.22,\ 1456024.53]$} \\
Sage-Voronoi & \makecell[c]{$14930.71\,\pm\,2077.24$\\$[12351.47,\ 17509.94]$} & \makecell[c]{$11986.24\,\pm\,2893.52$\\$[8393.46,\ 15579.02]$} & \makecell[c]{$1198269.84\,\pm\,494366.23$\\$[584433.08,\ 1812106.59]$} \\
GAIN & \makecell[c]{$19546.61\,\pm\,3036.68$\\$[15776.07,\ 23317.14]$} & \makecell[c]{$33139.96\,\pm\,1489.54$\\$[31290.46,\ 34989.47]$} & \makecell[c]{$1993002.23\,\pm\,395703.37$\\$[1501671.59,\ 2484332.86]$} \\
GRIN & \makecell[c]{$12752.65\,\pm\,3329.79$\\$[8618.17,\ 16887.14]$} & \makecell[c]{$20186.61\,\pm\,2476.21$\\$[17111.99,\ 23261.23]$} & \makecell[c]{$380694.76\,\pm\,116215.08$\\$[236394.68,\ 524994.84]$} \\
\bottomrule
\end{tabular}
\endgroup
\end{table}

\clearpage

\begin{table}[htbp]
\centering
\caption{{October--December 2022, RMSE of reconstructed daily sensor-count sums on affected test days. Cells give mean $\pm$ sample SD and [95\% Student-$t$ CI] across five seed runs.}}
\label{tab:p03_A29}
\begingroup
\footnotesize
\setlength{\tabcolsep}{3pt}
\renewcommand{\arraystretch}{1.15}
\begin{tabular}{lccc}
\toprule
Method / variant & $B=144$ & MCAR 30\% & Synchronous \\
\midrule
Linear & \makecell[c]{$42371.99\,\pm\,12266.90$\\$[27140.61,\ 57603.36]$} & \makecell[c]{$5220.02\,\pm\,758.17$\\$[4278.62,\ 6161.41]$} & \makecell[c]{$1151177.89\,\pm\,631814.50$\\$[366676.57,\ 1935679.21]$} \\
Median & \makecell[c]{$37182.80\,\pm\,5998.52$\\$[29734.65,\ 44630.94]$} & \makecell[c]{$228507.00\,\pm\,3142.43$\\$[224605.15,\ 232408.84]$} & \makecell[c]{$429435.01\,\pm\,330486.13$\\$[19082.28,\ 839787.74]$} \\
FFill + BFill & \makecell[c]{$45373.38\,\pm\,11657.47$\\$[30898.72,\ 59848.04]$} & \makecell[c]{$7169.11\,\pm\,907.52$\\$[6042.27,\ 8295.95]$} & \makecell[c]{$1197030.00\,\pm\,691932.87$\\$[337881.87,\ 2056178.14]$} \\
Rolling mean & \makecell[c]{$42324.26\,\pm\,12206.92$\\$[27167.37,\ 57481.16]$} & \makecell[c]{$7366.62\,\pm\,1233.95$\\$[5834.46,\ 8898.77]$} & \makecell[c]{$1152280.61\,\pm\,630178.26$\\$[369810.95,\ 1934750.27]$} \\
KNN (scaled) & \makecell[c]{$23804.14\,\pm\,2738.26$\\$[20404.14,\ 27204.13]$} & \makecell[c]{$70492.53\,\pm\,2513.60$\\$[67371.48,\ 73613.58]$} & \makecell[c]{$432017.13\,\pm\,329251.82$\\$[23197.00,\ 840837.25]$} \\
KNN (unscaled) & \makecell[c]{$24235.77\,\pm\,3901.40$\\$[19391.54,\ 29079.99]$} & \makecell[c]{$44152.82\,\pm\,3043.37$\\$[40373.98,\ 47931.66]$} & \makecell[c]{$432017.13\,\pm\,329251.82$\\$[23197.00,\ 840837.25]$} \\
Historical & \makecell[c]{$32296.63\,\pm\,3429.63$\\$[28038.18,\ 36555.08]$} & \makecell[c]{$231819.60\,\pm\,1790.85$\\$[229595.97,\ 234043.23]$} & \makecell[c]{$379035.11\,\pm\,273528.22$\\$[39404.97,\ 718665.25]$} \\
Graph Laplacian & \makecell[c]{$24835.53\,\pm\,6201.29$\\$[17135.62,\ 32535.45]$} & \makecell[c]{$35485.53\,\pm\,2436.84$\\$[32459.79,\ 38511.27]$} & \makecell[c]{$432017.13\,\pm\,329251.82$\\$[23197.00,\ 840837.25]$} \\
Graph Temporal & \makecell[c]{$25373.72\,\pm\,6073.65$\\$[17832.28,\ 32915.16]$} & \makecell[c]{$24406.78\,\pm\,3891.38$\\$[19575.00,\ 29238.57]$} & \makecell[c]{$1097689.10\,\pm\,612649.02$\\$[336984.87,\ 1858393.34]$} \\
Sage-Voronoi & \makecell[c]{$21867.12\,\pm\,4852.48$\\$[15841.98,\ 27892.27]$} & \makecell[c]{$14572.64\,\pm\,3140.81$\\$[10672.81,\ 18472.47]$} & \makecell[c]{$1514135.61\,\pm\,687397.74$\\$[660618.59,\ 2367652.63]$} \\
GAIN & \makecell[c]{$31231.74\,\pm\,5273.92$\\$[24683.30,\ 37780.18]$} & \makecell[c]{$44286.82\,\pm\,2021.76$\\$[41776.47,\ 46797.16]$} & \makecell[c]{$2305153.39\,\pm\,407108.01$\\$[1799662.03,\ 2810644.75]$} \\
GRIN & \makecell[c]{$21933.16\,\pm\,5011.53$\\$[15710.52,\ 28155.81]$} & \makecell[c]{$22050.23\,\pm\,2362.88$\\$[19116.33,\ 24984.13]$} & \makecell[c]{$442367.46\,\pm\,144468.27$\\$[262986.41,\ 621748.51]$} \\
\bottomrule
\end{tabular}
\endgroup
\end{table}

\clearpage

\begin{table}[htbp]
\centering
\caption{{October--December 2022, MRE of reconstructed daily sensor-count sums on affected test days. Cells give mean $\pm$ sample SD and [95\% Student-$t$ CI] across five seed runs.}}
\label{tab:p03_A30}
\begingroup
\footnotesize
\setlength{\tabcolsep}{3pt}
\renewcommand{\arraystretch}{1.15}
\begin{tabular}{lccc}
\toprule
Method / variant & $B=144$ & MCAR 30\% & Synchronous \\
\midrule
Linear & \makecell[c]{$0.008\,\pm\,0.002$\\$[0.005,\ 0.010]$} & \makecell[c]{$0.001\,\pm\,0.000$\\$[0.001,\ 0.001]$} & \makecell[c]{$0.259\,\pm\,0.117$\\$[0.113,\ 0.404]$} \\
Median & \makecell[c]{$0.008\,\pm\,0.001$\\$[0.006,\ 0.009]$} & \makecell[c]{$0.057\,\pm\,0.001$\\$[0.057,\ 0.058]$} & \makecell[c]{$0.124\,\pm\,0.123$\\$[-0.029,\ 0.277]$} \\
FFill + BFill & \makecell[c]{$0.008\,\pm\,0.002$\\$[0.006,\ 0.011]$} & \makecell[c]{$0.002\,\pm\,0.000$\\$[0.001,\ 0.002]$} & \makecell[c]{$0.267\,\pm\,0.135$\\$[0.099,\ 0.435]$} \\
Rolling mean & \makecell[c]{$0.008\,\pm\,0.002$\\$[0.005,\ 0.010]$} & \makecell[c]{$0.002\,\pm\,0.000$\\$[0.001,\ 0.002]$} & \makecell[c]{$0.259\,\pm\,0.117$\\$[0.113,\ 0.405]$} \\
KNN (scaled) & \makecell[c]{$0.003\,\pm\,0.000$\\$[0.003,\ 0.004]$} & \makecell[c]{$0.014\,\pm\,0.000$\\$[0.014,\ 0.015]$} & \makecell[c]{$0.125\,\pm\,0.124$\\$[-0.028,\ 0.279]$} \\
KNN (unscaled) & \makecell[c]{$0.004\,\pm\,0.001$\\$[0.003,\ 0.004]$} & \makecell[c]{$0.008\,\pm\,0.000$\\$[0.008,\ 0.009]$} & \makecell[c]{$0.125\,\pm\,0.124$\\$[-0.028,\ 0.279]$} \\
Historical & \makecell[c]{$0.007\,\pm\,0.001$\\$[0.006,\ 0.008]$} & \makecell[c]{$0.060\,\pm\,0.000$\\$[0.060,\ 0.061]$} & \makecell[c]{$0.110\,\pm\,0.128$\\$[-0.048,\ 0.269]$} \\
Graph Laplacian & \makecell[c]{$0.004\,\pm\,0.001$\\$[0.003,\ 0.005]$} & \makecell[c]{$0.007\,\pm\,0.001$\\$[0.006,\ 0.008]$} & \makecell[c]{$0.125\,\pm\,0.124$\\$[-0.028,\ 0.279]$} \\
Graph Temporal & \makecell[c]{$0.004\,\pm\,0.001$\\$[0.003,\ 0.005]$} & \makecell[c]{$0.005\,\pm\,0.001$\\$[0.004,\ 0.006]$} & \makecell[c]{$0.245\,\pm\,0.117$\\$[0.100,\ 0.390]$} \\
Sage-Voronoi & \makecell[c]{$0.004\,\pm\,0.001$\\$[0.003,\ 0.005]$} & \makecell[c]{$0.003\,\pm\,0.001$\\$[0.002,\ 0.004]$} & \makecell[c]{$0.320\,\pm\,0.089$\\$[0.209,\ 0.431]$} \\
GAIN & \makecell[c]{$0.005\,\pm\,0.001$\\$[0.004,\ 0.006]$} & \makecell[c]{$0.010\,\pm\,0.000$\\$[0.009,\ 0.010]$} & \makecell[c]{$0.594\,\pm\,0.283$\\$[0.242,\ 0.945]$} \\
GRIN & \makecell[c]{$0.003\,\pm\,0.001$\\$[0.002,\ 0.004]$} & \makecell[c]{$0.005\,\pm\,0.001$\\$[0.004,\ 0.006]$} & \makecell[c]{$0.110\,\pm\,0.055$\\$[0.041,\ 0.178]$} \\
\bottomrule
\end{tabular}
\endgroup
\end{table}

\clearpage

\begin{table}[htbp]
\centering
\caption{{Plain one-step (10-minute) Sage-Voronoi forecasting with observed past inputs. MAE and RMSE are in vehicles per 10-minute interval; MRE is dimensionless. Cells give mean $\pm$ sample SD and [95\% Student-$t$ CI] across five seed runs.}}
\label{tab:p03_A31}
\begingroup
\footnotesize
\setlength{\tabcolsep}{3pt}
\renewcommand{\arraystretch}{1.15}
\begin{tabular}{lccc}
\toprule
Method / variant & MAE & RMSE & MRE \\
\midrule
January 2023 & \makecell[c]{$28.24\,\pm\,0.37$\\$[27.78,\ 28.69]$} & \makecell[c]{$105.57\,\pm\,4.79$\\$[99.63,\ 111.51]$} & \makecell[c]{$0.288\,\pm\,0.009$\\$[0.277,\ 0.299]$} \\
October--December 2022 & \makecell[c]{$21.58\,\pm\,0.19$\\$[21.35,\ 21.81]$} & \makecell[c]{$66.06\,\pm\,0.47$\\$[65.48,\ 66.65]$} & \makecell[c]{$0.244\,\pm\,0.004$\\$[0.239,\ 0.250]$} \\
\bottomrule
\end{tabular}
\endgroup
\end{table}

\clearpage

\endgroup

\end{document}
